%&preformat-disser
\RequirePackage[l2tabu,orthodox]{nag} % Раскомментировав, можно в логе получать рекомендации относительно правильного использования пакетов и предупреждения об устаревших и нерекомендуемых пакетах
% Формат А4, 14pt (ГОСТ Р 7.0.11-2011, 5.3.6)
\documentclass[a4paper,14pt,oneside,openany]{memoir}

%%%%%%%%%%%%%%%%%%%%%%%%%%%%%%%%%%%%%%%%%%%%%%%%%%%%%%%%%%%%%%%%%%%%%%%%%%%%%%%%
%%%% Файл упрощённых настроек шаблона, общих для диссертации и автореферата %%%%
%%%%%%%%%%%%%%%%%%%%%%%%%%%%%%%%%%%%%%%%%%%%%%%%%%%%%%%%%%%%%%%%%%%%%%%%%%%%%%%%

%%% Режим черновика %%%
\makeatletter
\@ifundefined{c@draft}{
  \newcounter{draft}
  \setcounter{draft}{0}  % 0 --- чистовик (максимальное соблюдение ГОСТ)
                         % 1 --- черновик (отклонения от ГОСТ, но быстрая
                         %       сборка итоговых PDF)
}{}
\makeatother

%%% Пометки в тексте %%%
\makeatletter
\@ifundefined{c@showmarkup}{
  \newcounter{showmarkup}
  \setcounter{showmarkup}{0}  % 0 --- скрыть пометки
                              % 1 --- показывать пометки
}{}
\makeatother

%%% Использование в pdflatex шрифтов не по-умолчанию %%%
\makeatletter
\@ifundefined{c@usealtfont}{
  \newcounter{usealtfont}
  \setcounter{usealtfont}{1}    % 0 --- шрифты на базе Computer Modern
                                % 1 --- использовать пакет pscyr, при его
                                %       наличии
                                % 2 --- использовать пакет XCharter, при наличии
                                %       подходящей версии
}{}
\makeatother

%%% Использование в xelatex и lualatex семейств шрифтов %%%
\makeatletter
\@ifundefined{c@fontfamily}{
  \newcounter{fontfamily}
  \setcounter{fontfamily}{1}  % 0 --- CMU семейство. Используется как fallback;
                              % 1 --- Шрифты от MS (Times New Roman и компания)
                              % 2 --- Семейство Liberation
}{}
\makeatother

%%% Библиография %%%
\makeatletter
\@ifundefined{c@bibliosel}{
  \newcounter{bibliosel}
  \setcounter{bibliosel}{1}   % 0 --- встроенная реализация с загрузкой файла
                              %       через движок bibtex8;
                              % 1 --- реализация пакетом biblatex через движок
                              %       biber
}{}
\makeatother

%%% Вывод типов ссылок в библиографии %%%
\makeatletter
\@ifundefined{c@mediadisplay}{
  \newcounter{mediadisplay}
  \setcounter{mediadisplay}{1}   % 0 --- не делать ничего; надписи [Текст] и
                                 %       [Эл. ресурс] будут выводиться только в ссылках с
                                 %       заполненным полем `media`;
                                 % 1 --- автоматически добавлять надпись [Текст] к ссылкам с
                                 %       незаполненным полем `media`; таким образом, у всех
                                 %       источников будет указан тип, что соответствует
                                 %       требованиям ГОСТ
                                 % 2 --- автоматически удалять надписи [Текст], [Эл. Ресурс] и др.;
                                 %       не соответствует ГОСТ
                                 % 3 --- автоматически удалять надпись [Текст];
                                 %       не соответствует ГОСТ
                                 % 4 --- автоматически удалять надпись [Эл. Ресурс];
                                 %       не соответствует ГОСТ
}{}
\makeatother

%%% Предкомпиляция tikz рисунков для ускорения работы %%%
\makeatletter
\@ifundefined{c@imgprecompile}{
  \newcounter{imgprecompile}
  \setcounter{imgprecompile}{0}   % 0 --- без предкомпиляции;
                                  % 1 --- пользоваться предварительно
                                  %       скомпилированными pdf вместо генерации
                                  %       заново из tikz
}{}
\makeatother
            % общие настройки шаблона
%%% Проверка используемого TeX-движка %%%
\newif\ifxetexorluatex   % определяем новый условный оператор (http://tex.stackexchange.com/a/47579)
\ifxetex
    \xetexorluatextrue
\else
    \ifluatex
        \xetexorluatextrue
    \else
        \xetexorluatexfalse
    \fi
\fi

\newif\ifsynopsis           % Условие, проверяющее, что документ --- автореферат

\usepackage{etoolbox}[2015/08/02]   % Для продвинутой проверки разных условий
\providebool{presentation}

\usepackage{comment}    % Позволяет убирать блоки текста (добавляет
                        % окружение comment и команду \excludecomment)

%%% Поля и разметка страницы %%%
\usepackage{pdflscape}  % Для включения альбомных страниц
\usepackage{geometry}   % Для последующего задания полей

%%% Математические пакеты %%%
\usepackage{amsthm,amsmath,amscd}   % Математические дополнения от AMS
\usepackage{amsfonts,amssymb}       % Математические дополнения от AMS
\usepackage{mathtools}              % Добавляет окружение multlined
\usepackage{xfrac}                  % Красивые дроби
\usepackage[
    locale = DE,
    list-separator       = {;\,},
    list-final-separator = {;\,},
    list-pair-separator  = {;\,},
    list-units           = single,
    range-units          = single,
    range-phrase={\text{\ensuremath{-}}},
    % quotient-mode        = fraction, % красивые дроби могут не соответствовать ГОСТ
    fraction-function    = \sfrac,
    separate-uncertainty,
    ]{siunitx}[=v2]                 % Размерности SI
\sisetup{inter-unit-product = \ensuremath{{}\cdot{}}}

% Кириллица в нумерации subequations
% Для правильной работы требуется выполнение сразу после загрузки пакетов
\patchcmd{\subequations}{\def\theequation{\theparentequation\alph{equation}}}
{\def\theequation{\theparentequation\asbuk{equation}}}
{\typeout{subequations patched}}{\typeout{subequations not patched}}

%%%% Установки для размера шрифта 14 pt %%%%
%% Формирование переменных и констант для сравнения (один раз для всех подключаемых файлов)%%
%% должно располагаться до вызова пакета fontspec или polyglossia, потому что они сбивают его работу
\newlength{\curtextsize}
\newlength{\bigtextsize}
\setlength{\bigtextsize}{13.9pt}

\makeatletter
%\show\f@size    % неплохо для отслеживания, но вызывает стопорение процесса,
                 % если документ компилируется без команды  -interaction=nonstopmode
\setlength{\curtextsize}{\f@size pt}
\makeatother

%%% Кодировки и шрифты %%%
\ifxetexorluatex
    \ifpresentation
        \providecommand*\autodot{} % quick fix for polyglossia 1.50
    \fi
    \PassOptionsToPackage{no-math}{fontspec}    % https://tex.stackexchange.com/a/26295/104425
    \usepackage{polyglossia}[2014/05/21]        % Поддержка многоязычности
                                        % (fontspec подгружается автоматически)
\else
   %%% Решение проблемы копирования текста в буфер кракозябрами
    \ifnumequal{\value{usealtfont}}{0}{}{
        \input glyphtounicode.tex
        \input glyphtounicode-cmr.tex %from pdfx package
        \pdfgentounicode=1
    }
    \usepackage{cmap}   % Улучшенный поиск русских слов в полученном pdf-файле
    \ifnumequal{\value{usealtfont}}{2}{}{
        \defaulthyphenchar=127  % Если стоит до fontenc, то переносы
                                % не впишутся в выделяемый текст при
                                % копировании его в буфер обмена
    }
    \usepackage{textcomp}
    \usepackage[T1,T2A]{fontenc}                    % Поддержка русских букв
    \DeclareTextSymbolDefault{\CYRK}{T2A}
    \ifnumequal{\value{usealtfont}}{1}{% Используется pscyr, при наличии
        \IfFileExists{pscyr.sty}{\usepackage{pscyr}}{}  % Подключение pscyr
    }{}
    \usepackage[utf8]{inputenc}[2014/04/30]         % Кодировка utf8
    \usepackage[english, russian]{babel}[2014/03/24]% Языки: русский, английский
    \makeatletter\AtBeginDocument{\let\@elt\relax}\makeatother % babel 3.40 fix
    \ifnumequal{\value{usealtfont}}{2}{
        % http://dxdy.ru/post1238763.html#p1238763
        \usepackage[scaled=0.914]{XCharter}[2017/12/19] % Подключение русифицированных шрифтов XCharter
        \usepackage[charter, vvarbb, scaled=1.048]{newtxmath}[2017/12/14]
        \ifpresentation
        \else
            \setDisplayskipStretch{-0.078}
        \fi
    }{}
\fi

%%% Оформление абзацев %%%
\ifpresentation
\else
    \indentafterchapter     % Красная строка после заголовков типа chapter
    \usepackage{indentfirst}
\fi

%%% Цвета %%%
\ifpresentation
\else
    \usepackage[dvipsnames, table, hyperref]{xcolor} % Совместимо с tikz
\fi

%%% Таблицы %%%
\usepackage{longtable,ltcaption} % Длинные таблицы
\usepackage{multirow,makecell}   % Улучшенное форматирование таблиц
\usepackage{tabu, tabulary}      % таблицы с автоматически подбирающейся
                                 % шириной столбцов (tabu обязательно
                                 % до hyperref вызывать)
\makeatletter
%https://github.com/tabu-issues-for-future-maintainer/tabu/issues/26
\@ifpackagelater{longtable}{2020/02/07}{
\def\tabuendlongtrial{%
    \LT@echunk  \global\setbox\LT@gbox \hbox{\unhbox\LT@gbox}\kern\wd\LT@gbox
                \LT@get@widths
}%
}{}
\makeatother

\usepackage{threeparttable}      % автоматический подгон ширины подписи таблицы

%%% Общее форматирование
\usepackage{soulutf8}% Поддержка переносоустойчивых подчёркиваний и зачёркиваний
\usepackage{icomma}  % Запятая в десятичных дробях

%%% Оптимизация расстановки переносов и длины последней строки абзаца
\IfFileExists{impnattypo.sty}{% проверка установленности пакета impnattypo
    \ifluatex
        \ifnumequal{\value{draft}}{1}{% Черновик
            \usepackage[hyphenation, lastparline, nosingleletter, homeoarchy,
            rivers, draft]{impnattypo}
        }{% Чистовик
            \usepackage[hyphenation, lastparline, nosingleletter]{impnattypo}
        }
    \else
        \usepackage[hyphenation, lastparline]{impnattypo}
    \fi
}{}

%% Векторная графика

\usepackage{tikz}                   % Продвинутый пакет векторной графики
\usetikzlibrary{chains}             % Для примера tikz рисунка
\usetikzlibrary{shapes.geometric}   % Для примера tikz рисунка
\usetikzlibrary{shapes.symbols}     % Для примера tikz рисунка
\usetikzlibrary{arrows}             % Для примера tikz рисунка

%%% Гиперссылки %%%
\ifxetexorluatex
    \let\CYRDZE\relax
\fi
\usepackage{hyperref}[2012/11/06]

%%% Изображения %%%
\usepackage{graphicx}[2014/04/25]   % Подключаем пакет работы с графикой
\usepackage{caption}                % Подписи рисунков и таблиц
\usepackage{subcaption}             % Подписи подрисунков и подтаблиц
\usepackage{pdfpages}               % Добавление внешних pdf файлов

%%% Счётчики %%%
\usepackage{aliascnt}
\usepackage[figure,table]{totalcount}   % Счётчик рисунков и таблиц
\usepackage{totcount}   % Пакет создания счётчиков на основе последнего номера
                        % подсчитываемого элемента (может требовать дважды
                        % компилировать документ)
\usepackage{totpages}   % Счётчик страниц, совместимый с hyperref (ссылается
                        % на номер последней страницы). Желательно ставить
                        % последним пакетом в преамбуле

%%% Продвинутое управление групповыми ссылками (пока только формулами) %%%
\ifpresentation
\else
    \usepackage[russian]{cleveref} % cleveref имеет сложности со считыванием
    % языка из babel. Такое решение русификации вывода выбрано вместо
    % определения в documentclass из опасности что-то лишнее передать во все
    % остальные пакеты, включая библиографию.

    % Добавление возможности использования пробелов в \labelcref
    % https://tex.stackexchange.com/a/340502/104425
    \usepackage{kvsetkeys}
    \makeatletter
    \let\org@@cref\@cref
    \renewcommand*{\@cref}[2]{%
        \edef\process@me{%
            \noexpand\org@@cref{#1}{\zap@space#2 \@empty}%
        }\process@me
    }
    \makeatother
\fi

\usepackage{placeins} % для \FloatBarrier

\ifnumequal{\value{draft}}{1}{% Черновик
    \usepackage[firstpage]{draftwatermark}
    \SetWatermarkText{DRAFT}
    \SetWatermarkFontSize{14pt}
    \SetWatermarkScale{15}
    \SetWatermarkAngle{45}
}{}

%%% Цитата, не приводимая в автореферате:
% возможно, актуальна только для biblatex
%\newcommand{\citeinsynopsis}[1]{\ifsynopsis\else ~\cite{#1} \fi}

% если текущий процесс запущен библиотекой tikz-external, то прекомпиляция должна быть включена
\ifdefined\tikzexternalrealjob
    \setcounter{imgprecompile}{1}
\fi

\ifnumequal{\value{imgprecompile}}{1}{% Только если у нас включена предкомпиляция
    \usetikzlibrary{external}   % подключение возможности предкомпиляции
    \tikzexternalize[prefix=images/cache/,optimize command away=\includepdf] % activate! % здесь можно указать отдельную папку для скомпилированных файлов
    \ifxetex
        \tikzset{external/up to date check={diff}}
    \fi
}{}

\usepackage{amsmath} % без этого никуда -МП
\usepackage{adjustbox}         % Пакеты общие для диссертации и автореферата
\synopsisfalse                      % Этот документ --- не автореферат
\input{Dissertation/dispackages}    % Пакеты для диссертации
\input{Dissertation/userpackages}   % Пакеты для специфических пользовательских задач

%%%%%%%%%%%%%%%%%%%%%%%%%%%%%%%%%%%%%%%%%%%%%%%%%%%%%%%%%%%%%%%%%%%%%%%%%%%%%%%%
%%%% Файл упрощённых настроек шаблона, общих для диссертации и автореферата %%%%
%%%%%%%%%%%%%%%%%%%%%%%%%%%%%%%%%%%%%%%%%%%%%%%%%%%%%%%%%%%%%%%%%%%%%%%%%%%%%%%%

%%% Режим черновика %%%
\makeatletter
\@ifundefined{c@draft}{
  \newcounter{draft}
  \setcounter{draft}{0}  % 0 --- чистовик (максимальное соблюдение ГОСТ)
                         % 1 --- черновик (отклонения от ГОСТ, но быстрая
                         %       сборка итоговых PDF)
}{}
\makeatother

%%% Пометки в тексте %%%
\makeatletter
\@ifundefined{c@showmarkup}{
  \newcounter{showmarkup}
  \setcounter{showmarkup}{0}  % 0 --- скрыть пометки
                              % 1 --- показывать пометки
}{}
\makeatother

%%% Использование в pdflatex шрифтов не по-умолчанию %%%
\makeatletter
\@ifundefined{c@usealtfont}{
  \newcounter{usealtfont}
  \setcounter{usealtfont}{1}    % 0 --- шрифты на базе Computer Modern
                                % 1 --- использовать пакет pscyr, при его
                                %       наличии
                                % 2 --- использовать пакет XCharter, при наличии
                                %       подходящей версии
}{}
\makeatother

%%% Использование в xelatex и lualatex семейств шрифтов %%%
\makeatletter
\@ifundefined{c@fontfamily}{
  \newcounter{fontfamily}
  \setcounter{fontfamily}{1}  % 0 --- CMU семейство. Используется как fallback;
                              % 1 --- Шрифты от MS (Times New Roman и компания)
                              % 2 --- Семейство Liberation
}{}
\makeatother

%%% Библиография %%%
\makeatletter
\@ifundefined{c@bibliosel}{
  \newcounter{bibliosel}
  \setcounter{bibliosel}{1}   % 0 --- встроенная реализация с загрузкой файла
                              %       через движок bibtex8;
                              % 1 --- реализация пакетом biblatex через движок
                              %       biber
}{}
\makeatother

%%% Вывод типов ссылок в библиографии %%%
\makeatletter
\@ifundefined{c@mediadisplay}{
  \newcounter{mediadisplay}
  \setcounter{mediadisplay}{1}   % 0 --- не делать ничего; надписи [Текст] и
                                 %       [Эл. ресурс] будут выводиться только в ссылках с
                                 %       заполненным полем `media`;
                                 % 1 --- автоматически добавлять надпись [Текст] к ссылкам с
                                 %       незаполненным полем `media`; таким образом, у всех
                                 %       источников будет указан тип, что соответствует
                                 %       требованиям ГОСТ
                                 % 2 --- автоматически удалять надписи [Текст], [Эл. Ресурс] и др.;
                                 %       не соответствует ГОСТ
                                 % 3 --- автоматически удалять надпись [Текст];
                                 %       не соответствует ГОСТ
                                 % 4 --- автоматически удалять надпись [Эл. Ресурс];
                                 %       не соответствует ГОСТ
}{}
\makeatother

%%% Предкомпиляция tikz рисунков для ускорения работы %%%
\makeatletter
\@ifundefined{c@imgprecompile}{
  \newcounter{imgprecompile}
  \setcounter{imgprecompile}{0}   % 0 --- без предкомпиляции;
                                  % 1 --- пользоваться предварительно
                                  %       скомпилированными pdf вместо генерации
                                  %       заново из tikz
}{}
\makeatother
      % Упрощённые настройки шаблона

\input{common/newnames}         % Новые переменные, для всего проекта

\input{common/data}             % Основные сведения
\input{common/fonts}            % Определение шрифтов (частичное)
\input{common/styles}           % Стили общие для диссертации и автореферата
%%% Переопределение именований, если иначе не сработает %%%
%\gappto\captionsrussian{
%    \renewcommand{\chaptername}{Глава}
%    \renewcommand{\appendixname}{Приложение} % (ГОСТ Р 7.0.11-2011, 5.7)
%}

%%% Изображения %%%
\graphicspath{{images/}{Dissertation/images/}}         % Пути к изображениям

%%% Интервалы %%%
%% По ГОСТ Р 7.0.11-2011, пункту 5.3.6 требуется полуторный интервал
%% Реализация средствами класса (на основе setspace) ближе к типографской классике.
%% И правит сразу и в таблицах (если со звёздочкой)
%\DoubleSpacing*     % Двойной интервал
\OnehalfSpacing*    % Полуторный интервал
%\setSpacing{1.42}   % Полуторный интервал, подобный Ворду (возможно, стоит включать вместе с предыдущей строкой)

%%% Макет страницы %%%
% Выставляем значения полей (ГОСТ 7.0.11-2011, 5.3.7)
\geometry{a4paper, top=2cm, bottom=2cm, left=2.5cm, right=1cm, nofoot, nomarginpar} %, heightrounded, showframe
\setlength{\topskip}{0pt}   %размер дополнительного верхнего поля
\setlength{\footskip}{12.3pt} % снимет warning, согласно https://tex.stackexchange.com/a/334346

%%% Выравнивание и переносы %%%
%% http://tex.stackexchange.com/questions/241343/what-is-the-meaning-of-fussy-sloppy-emergencystretch-tolerance-hbadness
%% http://www.latex-community.org/forum/viewtopic.php?p=70342#p70342
\tolerance 1414
\hbadness 1414
\emergencystretch 1.5em % В случае проблем регулировать в первую очередь
\hfuzz 0.3pt
\vfuzz \hfuzz
%\raggedbottom
%\sloppy                 % Избавляемся от переполнений
\clubpenalty=10000      % Запрещаем разрыв страницы после первой строки абзаца
\widowpenalty=10000     % Запрещаем разрыв страницы после последней строки абзаца
\brokenpenalty=4991     % Ограничение на разрыв страницы, если строка заканчивается переносом

%%% Блок управления параметрами для выравнивания заголовков в тексте %%%
\newlength{\otstuplen}
\setlength{\otstuplen}{\theotstup\parindent}
\ifnumequal{\value{headingalign}}{0}{% выравнивание заголовков в тексте
    \newcommand{\hdngalign}{\centering}                % по центру
    \newcommand{\hdngaligni}{}% по центру
    \setlength{\otstuplen}{0pt}
}{%
    \newcommand{\hdngalign}{}                 % по левому краю
    \newcommand{\hdngaligni}{\hspace{\otstuplen}}      % по левому краю
} % В обоих случаях вроде бы без переноса, как и надо (ГОСТ Р 7.0.11-2011, 5.3.5)

%%% Оглавление %%%
\renewcommand{\cftchapterdotsep}{\cftdotsep}                % отбивка точками до номера страницы начала главы/раздела

%% Переносить слова в заголовке не допускается (ГОСТ Р 7.0.11-2011, 5.3.5). Заголовки в оглавлении должны точно повторять заголовки в тексте (ГОСТ Р 7.0.11-2011, 5.2.3). Прямого указания на запрет переносов в оглавлении нет, но по той же логике невнесения искажений в смысл, лучше в оглавлении не переносить:
\setrmarg{2.55em plus1fil}                             %To have the (sectional) titles in the ToC, etc., typeset ragged right with no hyphenation
\renewcommand{\cftchapterpagefont}{\normalfont}        % нежирные номера страниц у глав в оглавлении
\renewcommand{\cftchapterleader}{\cftdotfill{\cftchapterdotsep}}% нежирные точки до номеров страниц у глав в оглавлении
%\renewcommand{\cftchapterfont}{}                       % нежирные названия глав в оглавлении

\ifnumgreater{\value{headingdelim}}{0}{%
    \renewcommand\cftchapteraftersnum{.\space}       % добавляет точку с пробелом после номера раздела в оглавлении
}{}
\ifnumgreater{\value{headingdelim}}{1}{%
    \renewcommand\cftsectionaftersnum{.\space}       % добавляет точку с пробелом после номера подраздела в оглавлении
    \renewcommand\cftsubsectionaftersnum{.\space}    % добавляет точку с пробелом после номера подподраздела в оглавлении
    \renewcommand\cftsubsubsectionaftersnum{.\space} % добавляет точку с пробелом после номера подподподраздела в оглавлении
    \AfterEndPreamble{% без этого polyglossia сама всё переопределяет
        \setsecnumformat{\csname the#1\endcsname.\space}
    }
}{%
    \AfterEndPreamble{% без этого polyglossia сама всё переопределяет
        \setsecnumformat{\csname the#1\endcsname\quad}
    }
}

\renewcommand*{\cftappendixname}{\appendixname\space} % Слово Приложение в оглавлении

%%% Колонтитулы %%%
% Порядковый номер страницы печатают на середине верхнего поля страницы (ГОСТ Р 7.0.11-2011, 5.3.8)
\makeevenhead{plain}{}{\rmfamily\thepage}{}
\makeoddhead{plain}{}{\rmfamily\thepage}{}
\makeevenfoot{plain}{}{}{}
\makeoddfoot{plain}{}{}{}
\pagestyle{plain}

%%% добавить Стр. над номерами страниц в оглавлении
%%% http://tex.stackexchange.com/a/306950
\newif\ifendTOC

\newcommand*{\tocheader}{
\ifnumequal{\value{pgnum}}{1}{%
    \ifendTOC\else\hbox to \linewidth%
      {\noindent{}~\hfill{Стр.}}\par%
      \ifnumless{\value{page}}{3}{}{%
        \vspace{0.5\onelineskip}
      }
      \afterpage{\tocheader}
    \fi%
}{}%
}%

%%% Оформление заголовков глав, разделов, подразделов %%%
%% Работа должна быть выполнена ... размером шрифта 12-14 пунктов (ГОСТ Р 7.0.11-2011, 5.3.8). То есть не должно быть надписей шрифтом более 14. Так и поставим.
%% Эти установки будут давать одинаковый результат независимо от выбора базовым шрифтом 12 пт или 14 пт
\newcommand{\basegostsectionfont}{\fontsize{14pt}{16pt}\selectfont\bfseries}

\makechapterstyle{thesisgost}{%
    \chapterstyle{default}
    \setlength{\beforechapskip}{0pt}
    \setlength{\midchapskip}{0pt}
    \setlength{\afterchapskip}{\theintvl\curtextsize}
    \renewcommand*{\chapnamefont}{\basegostsectionfont}
    \renewcommand*{\chapnumfont}{\basegostsectionfont}
    \renewcommand*{\chaptitlefont}{\basegostsectionfont}
    \renewcommand*{\chapterheadstart}{}
    \ifnumgreater{\value{headingdelim}}{0}{%
        \renewcommand*{\afterchapternum}{.\space}   % добавляет точку с пробелом после номера раздела
    }{%
        \renewcommand*{\afterchapternum}{\quad}     % добавляет \quad после номера раздела
    }
    \renewcommand*{\printchapternum}{\hdngaligni\hdngalign\chapnumfont \thechapter}
    \renewcommand*{\printchaptername}{}
    \renewcommand*{\printchapternonum}{\hdngaligni\hdngalign}
}

\makeatletter
\makechapterstyle{thesisgostchapname}{%
    \chapterstyle{thesisgost}
    \renewcommand*{\printchapternum}{\chapnumfont \thechapter}
    \renewcommand*{\printchaptername}{\hdngaligni\hdngalign\chapnamefont \@chapapp} %
}
\makeatother

\chapterstyle{thesisgost}

\setsecheadstyle{\basegostsectionfont\hdngalign}
\setsecindent{\otstuplen}

\setsubsecheadstyle{\basegostsectionfont\hdngalign}
\setsubsecindent{\otstuplen}

\setsubsubsecheadstyle{\basegostsectionfont\hdngalign}
\setsubsubsecindent{\otstuplen}

\sethangfrom{\noindent #1} %все заголовки подразделов центрируются с учетом номера, как block

\ifnumequal{\value{chapstyle}}{1}{%
    \chapterstyle{thesisgostchapname}
    \renewcommand*{\cftchaptername}{\chaptername\space} % будет вписано слово Глава перед каждым номером раздела в оглавлении
}{}%

%%% Интервалы между заголовками
\setbeforesecskip{\theintvl\curtextsize}% Заголовки отделяют от текста сверху и снизу тремя интервалами (ГОСТ Р 7.0.11-2011, 5.3.5).
\setaftersecskip{\theintvl\curtextsize}
\setbeforesubsecskip{\theintvl\curtextsize}
\setaftersubsecskip{\theintvl\curtextsize}
\setbeforesubsubsecskip{\theintvl\curtextsize}
\setaftersubsubsecskip{\theintvl\curtextsize}

%%% Вертикальные интервалы глав (\chapter) в оглавлении как и у заголовков
% раскомментировать следующие 2
% \setlength{\cftbeforechapterskip}{0pt plus 0pt}   % ИЛИ эти 2 строки из учебника
% \renewcommand*{\insertchapterspace}{}
% или эту
% \renewcommand*{\cftbeforechapterskip}{0em}


%%% Блок дополнительного управления размерами заголовков
\ifnumequal{\value{headingsize}}{1}{% Пропорциональные заголовки и базовый шрифт 14 пт
    \renewcommand{\basegostsectionfont}{\large\bfseries}
    \renewcommand*{\chapnamefont}{\Large\bfseries}
    \renewcommand*{\chapnumfont}{\Large\bfseries}
    \renewcommand*{\chaptitlefont}{\Large\bfseries}
}{}

%%% Счётчики %%%

%% Упрощённые настройки шаблона диссертации: нумерация формул, таблиц, рисунков
\ifnumequal{\value{contnumeq}}{1}{%
    \counterwithout{equation}{chapter} % Убираем связанность номера формулы с номером главы/раздела
}{}
\ifnumequal{\value{contnumfig}}{1}{%
    \counterwithout{figure}{chapter}   % Убираем связанность номера рисунка с номером главы/раздела
}{}
\ifnumequal{\value{contnumtab}}{1}{%
    \counterwithout{table}{chapter}    % Убираем связанность номера таблицы с номером главы/раздела
}{}

\AfterEndPreamble{
%% регистрируем счётчики в системе totcounter
    \regtotcounter{totalcount@figure}
    \regtotcounter{totalcount@table}       % Если иным способом поставить в преамбуле то ошибка в числе таблиц
    \regtotcounter{TotPages}               % Если иным способом поставить в преамбуле то ошибка в числе страниц
    \newtotcounter{totalappendix}
    \newtotcounter{totalchapter}
}

%%% Дополнительно %%%

%% используется в длинных таблицах
\UseTblrLibrary{booktabs} 
\DefTblrTemplate{contfoot-text}{default}{}
\DefTblrTemplate{middlehead,lasthead}{default}{Продолжение таблицы \thetable}
\DefTblrTemplate{caption-sep}{default}{\enskip---\enskip}

\DefTblrTemplate{caption}{default}{
	\par
	\UseTblrTemplate{caption-tag}{default}
	\UseTblrTemplate{caption-sep}{default}
	\UseTblrTemplate{caption-text}{default}
	\par
}

\DefTblrTemplate{capcont}{default}{%
	\par
	\UseTblrTemplate {caption-tag}{default}%
	\UseTblrTemplate {caption-sep}{default}%
	\UseTblrTemplate {caption-text}{default}
	\UseTblrTemplate {conthead-text}{default}
	\par
}

\DefTblrTemplate { remark } { normal }
{
	\UseTblrAlign { remark }
	\UseTblrIndent { remark }
	\MapTblrRemarks
	{
		\hangindent = 2.5em
		\hangafter = 1
		\UseTblrHang { remark }
		\leavevmode
		\UseTblrTemplate { remark-tag } { default }
		\UseTblrTemplate { remark-sep } { default }
		\UseTblrTemplate { remark-text } { default }
		\par
	}
}
  % Стили для диссертации
\input{Dissertation/userstyles} % Стили для специфических пользовательских задач

%%% Библиография. Выбор движка для реализации %%%
% Здесь только проверка установленного ключа. Сама настройка выбора движка
% размещена в common/setup.tex
\ifnumequal{\value{bibliosel}}{0}{%
    \input{biblio/predefined}   % Встроенная реализация с загрузкой файла через движок bibtex8
}{
    \input{biblio/biblatex}     % Реализация пакетом biblatex через движок biber
}

% Вывести информацию о выбранных опциях в лог сборки
\typeout{Selected options:}
\typeout{Draft mode: \arabic{draft}}
\typeout{Font: \arabic{fontfamily}}
\typeout{AltFont: \arabic{usealtfont}}
\typeout{Bibliography backend: \arabic{bibliosel}}
\typeout{Precompile images: \arabic{imgprecompile}}
% Вывести информацию о версиях используемых библиотек в лог сборки
\listfiles

%%% Управление компиляцией отдельных частей диссертации %%%
% Необходимо сначала иметь полностью скомпилированный документ, чтобы все
% промежуточные файлы были в наличии
% Затем, для вывода отдельных частей можно воспользоваться командой \includeonly
% Ниже примеры использования команды:
%
%\includeonly{Dissertation/part2}
%\includeonly{Dissertation/contents,Dissertation/appendix,Dissertation/conclusion}
%
% Если все команды закомментированы, то документ будет выведен в PDF файл полностью

\begin{document}
%%% Переопределение именований типовых разделов
% https://tex.stackexchange.com/a/156050
\gappto\captionsrussian{\input{common/renames}\unskip} % for polyglossia and babel
\input{common/renames}

%%% Структура диссертации (ГОСТ Р 7.0.11-2011, 4)
\include{Dissertation/title}           % Титульный лист
\include{Dissertation/contents}        % Оглавление
\ifnumequal{\value{contnumfig}}{1}{}{\counterwithout{figure}{chapter}}
\ifnumequal{\value{contnumtab}}{1}{}{\counterwithout{table}{chapter}}
\include{Dissertation/introduction}    % Введение
\ifnumequal{\value{contnumfig}}{1}{\counterwithout{figure}{chapter}
}{\counterwithin{figure}{chapter}}
\ifnumequal{\value{contnumtab}}{1}{\counterwithout{table}{chapter}
}{\counterwithin{table}{chapter}}
\chapter{Web 1.0}\label{ch:ch1}

\section{Концептуальное назначение}\label{sec:ch1/sec1}

\[
--
\]

\section{Обобщенная структура}\label{sec:ch1/sec2}

\subsection{Анализ топологии крупных сегментов Веба с помощью модели Бродера “Bow-tie”}\label{subsec:ch1/sec2/sub1}

\paragraph{Введение.} В настоящее время все больше организаций, стремясь отразить свою деятельность в Вебе, создают отдельные веб-ресурсы (сайты, группы сайтов) и заботятся об улучшении их рейтингов в информационно-поисковых системах с целью накопления символического капитала, повышения прямых продаж. Важную роль в формировании рейтинга сайтов и локализированных веб-сегментов в целом играют поисковые машины.

С 2006 года основные поисковые машины (например, в мире -- Google, в России -- Яндекс) в своих методах ранжирования наряду с классическими факторами \cite{Kleinberg, BrinPage, Chakrabarti} (частота встречаемости слов, объем цитирования и авторитетность сайтов, частота обновления сайтов и др.) по-новому используют факторы, связанные с поведением пользователя на веб-ресурсах \cite{GuhaKunduBhadra,AntoniouPlegasTsakalidis,FeuerSavevAslam}. Так, действия пользователя, ранее приводившие к повышению рейтингов сайтов (например, частота захода пользователей), ныне могут приводить к “пессимизации” показателей, поскольку поисковые системы становятся более “социальными” и учитывают уже не только профиль пользователя и частоту захода на ресурсы, но также мотивацию посетителей и стратегии их поведения (частоту повторного захода, время нахождения на странице и сайте, логику и маршруты переходов, пользовательские интересы, тип потребляемого контента и многие иные факторы). При ранжировании современные поисковые машины в совокупности учитывают около 300-500 таких факторов.

С другой стороны, существует небольшое количество структурных характеристик Веб-сегментов \cite{ChoRoy}, которые существенно влияют практически на все эти факторы. Имеется значительное число исследований \cite{Kleinberg,ChoRoy,BroderKumarMaghoul,AguilloGranadinoOrtega,StuartThelwallHarries,Chakrabarti,Thelwall,Pechnikov,PechnikovNwohiri}, в которых строится модель фрагментов веба, описывающаяся сравнительно небольшим набором параметров.

Общим направлением исследований авторов является изучение связей между рейтингом сайта и его глобальными характеристиками. По-видимому, эти связи различны для сайтов разных категорий, например, естественнонаучной направленности и гуманитарных. Для выявления существенных различий между категориями сайтов, необходимо исследовать большое количество сайтов с последующей статистической обработкой материала. Авторы располагают разработанной ими аналитической системой для вебометрических исследований \cite{BlekanovSergeevMaksimov,BlekanovSergeevMartynenko}, способной выполнить эту работу. Ввиду чрезвычайно больших затрат машинных ресурсов на подобное исследование, представляет интерес машинный эксперимент, заключающийся в определении затрат на полное исследование типичных сайтов, описание которого и является темой данной статьи.

\paragraph{Основная часть.} В эксперименте ставилась задача построения топологии на основе модели Бродера “Bow-tie” \cite{BroderKumarMaghoul,Thelwall} для типичных крупных гуманитарно-ориентированных и естественнонаучно-ориентированных сегментов Веба с последующим вычислением их структурных характеристик и определением произведенных затрат ресурсов.

В качестве гуманитарно-ориентированного веб-сегмента был выбран сайт Факультета журналистики (JF) Санкт-Петербургского государственного университета (СПбГУ) -- www.jf.spbu.ru, естественнонаучно-ориентированного -- сайт факультета Прикладной математики -- процессов управления (APMATH) СПбГУ -- www.apmath.spbu.ru.

С помощью специализированной сборки аналитической системы для вебометрических исследований на основе ядра поискового робота, успешно апробированной в исследованиях \cite{BlekanovSergeev,MaksimovBlekanov} , были выявлены гиперссылочные структуры обоих сайтов, которые описываются в виде двух ориентированных веб-графов \(G_{apmath}\) и \(G_{jf}\), в которых вершины соответствуют страницам, а дуги – соединяющим эти страницы гиперссылкам. Для получения глобальных характеристик использовалась модель Бродера “Bow-tie” \cite{BroderKumarMaghoul}, которая во множестве вершин исследуемого графа выделяет четыре основных подмножества: 1) центральное ядро (компонента сильной связности SCC); 2) истоки (IN); 3) стоки (OUT); 4) “TENDRILS” и “TUBES”. Поиск сильно-связной компоненты выполнялся с помощью алгоритма Косарайю \cite{Sedgewick} (Kosaraju's algorithm), время выполнения которого для разреженных графов (а такими являются исследуемые веб-графы) пропорционально сумме всех их вершин и направленных ребер.

В эксперименте вычислялись следующие характеристики веб-сайтов:
\begin{itemize}
	\item размер центрального ядра \(S_{scc}\);
	\item размер стоков \(S_{out}\) и истоков \(S_{in}\);
	\item размер “TENDRILS” и “TUBES” \(S_{tubes}\);
	\item отношение \({\lvert S_{scc} \rvert + \lvert S_{in} \rvert}\) к \(\lvert S_{scc} \rvert + \lvert S_{out} \rvert\).
\end{itemize}

Кроме того, для каждой страницы центрального ядра был введен показатель ее связности с другими страницами из ядра:
\[
	\text{\textit{VC}}(u, v) = 
	\begin{cases} 
		1, & \text{\textit{if }} u \rightarrow v,\\
		0, & \text{\textit{else}}
	\end{cases},
	\forall u, v \in S_{scc}
\]

Этот показатель лег в основу усредненной меры связности каждой страницы ядра \(\text{\textit{MA\text{\textit{VC}}}}(u)\) и усредненной меры связности центрального ядра \(\text{\textit{MAC}}(S_{scc})\):
\[
	\text{\textit{MAVC}}(u) = \frac{1}{\lvert S_{scc} \rvert - 1} \sum_{\forall v \in S_{scc}, \; v \neq u}\text{\textit{VC}}(u, v),
\] 
\[
	\text{\textit{MAC}}(S_{scc}) = \frac{1}{|S_{scc}|} \sum_{\forall u \in S_{scc}}\text{\textit{MAVC}}(u).
\]

Приведенные меры были введены для оценки качества связанности каждой топологии, полученной в ходе поставленного эксперимента.

Для оценки трудоемкости эксперимента по каждому исследуемому сайту использовались следующие характеристики:
\begin{itemize}
	\item время (\(T_{total}^{js}\) и \(T_{total}^{apmath}\)), необходимое для выявления топологии и вычисления указанных обобщенных характеристик;
	\item объем памяти (\(\text{\textit{VMem}}_{js}\) и \(\text{\textit{VMem}}_{apmath}\)), требуемый для хранения обработанных данных;
	\item объем Веб трафика \(\text{\textit{Traffic}}_{js}\) и \(\text{\textit{Traffic}}_{apmath}\)).
\end{itemize}

\paragraph{Результаты эксперимента.} Эксперимент проводился на тестовой персональной электронно-вычислительной машине (ПЭВМ) с процессором \textit{Intel Core i7 CPU 950 @ 3.07 GHz x 8}, \textit{12 GB} объемом жесткого диска и операционной системой \textit{Ubuntu 13.10}. Загрузка данных с Интернета выполнялась посредством канала с пропускной способностью до \textit{50 Мбит/с}.

В ходе эксперимента было выявлено, что Веб-граф \(G_{apmath}\) содержит \textit{26 148} вершин и \textit{2 025 909} направленных ребер, а Веб-граф \(G_{jf}\) -- \textit{27 361} вершин и \textit{4 252 424} направленных ребра. На рис.~\cref{fig:webGraphs} приведены топологии, построенные для каждого полученного графа с помощью модели Бродера “Bow-tie”.

\begin{figure}[ht]
    \centerfloat{
	        \hfill
	        \subcaptionbox[List-of-Figures entry]{\label{fig:webGraphs-1}}{%
		            \includegraphics[width=0.5\linewidth]{webGraph(apmath)}}
	        \subcaptionbox{\label{fig:webGraphs-2}}{%
		            \includegraphics[width=0.5\linewidth]{webGraph(jf)}}
	        \hfill
	    }
    \caption{Топологии Веб-графов \(G_{apmath}\) и \(G_{jf}\), соответственно.}\label{fig:webGraphs}
\end{figure}

Значительную долю в топологии Веб-графа \(G_{apmath}\) среди всех ее компонент занимает SCC (\textit{65.56\%} страниц от общего их числа), \textit{8.8\%} занимают истоки (IN), - \textit{25.5\%} стоки (OUT), \textit{3.96\%} -- “TENDRILS” и “TUBES” (рис.~\cref{fig:webGraphs-1}). В то время как доли компонент топологии Веб-графа (рис.~\cref{fig:webGraphs}\subcaptionref*{fig:webGraphs-2}) приблизительно одинаковые (более \textit{46\%} от общего числа всех страниц) за исключением компоненты OUT (\textit{0.76\%}).

В табл.~\cref{tab:webGraphTable} приведены полученные численные значения структурных характеристик, которые были введены в эксперименте для анализа исследуемых сайтов.

\begin{table} [htbp]%
    \centering
    \caption{Структурные характеристики Веб-графов \(G_{apmath}\) и \(G_{jf}\).}%
    \label{tab:webGraphTable}% label всегда желательно идти после caption
    \renewcommand{\arraystretch}{1.5}%% Увеличение расстояния между рядами, для улучшения восприятия.
    \begin{SingleSpace}
    	\begin{tabulary}{\textwidth}{@{}>{\zz}L >{\zz}C >{\zz}C@{}} %Вертикальные полосы не используются принципиально, как и лишние горизонтальные (допускается по ГОСТ 2.105 пункт 4.4.5) % @{} позволяет прижиматься к краям
		            \toprule     %%% верхняя линейка
		            Характеристика & Топология сайта APMATH &Топология сайта JS \\
		            \midrule %%% тонкий разделитель. Отделяет названия столбцов. Обязателен по ГОСТ 2.105 пункт 4.4.5
		            Мощность множества ядра \(\lvert S_{scc} \rvert\)  & 17142 & 13893    \\
		            Мощность множества IN \(\lvert S_{in} \rvert\)       & 2302     & 13165    \\
		            Мощность множества OUT \(\lvert S_{out} \rvert\)       & 6669     & 206    \\
					Мощность множества \newline TENDRILS и TUBES  \(\lvert S_{tubes} \rvert\)      & 1035    & 12809     \\
					Отношение \(\frac{\lvert S_{scc} \rvert + \lvert S_{in} \rvert}{\lvert S_{scc} \rvert + \lvert S_{out} \rvert}\) & 0.82 & 1.92 \\
					Усредненная мера связности центрального ядра \(\text{\textit{MAC}}(S_{scc})\) & 0.00244 & 0.00193 \\
					Усредненная мера связности всего Веб-графа \textit{MAC(G)} & 0.00106 & 0.00098 \\
		            \bottomrule %%% нижняя линейка
		        \end{tabulary}%
	    \end{SingleSpace}
\end{table}

Затраты ресурсов, требуемых для полного анализа исследуемых в эксперименте сегментов Веба, приведены в табл.~\cref{tab:webGraphCostTable}.

\begin{table}[ht]%
	\caption{Характеристики, оценивающие трудоемкость поставленного эксперимента.}%
	\label{tab:webGraphCostTable}% label всегда желательно идти после caption
    \renewcommand{\arraystretch}{1.6}%% Увеличение расстояния между рядами, для улучшения восприятия.
    \def\tabularxcolumn#1{m{#1}}
    \begin{tabularx}{\textwidth}{@{}>{\raggedright}X>{\centering}m{3.5cm}  >{\centering\arraybackslash}m{3.5cm}@{}}% Вертикальные полосы не используются принципиально, как и лишние горизонтальные (допускается по ГОСТ 2.105 пункт 4.4.5) % @{} позволяет прижиматься к краям
			\toprule     %%% верхняя линейка
			Характеристика & Сайт APMATH & Сайт JS \\
			\midrule %%% тонкий разделитель. Отделяет названия столбцов. Обязателен по ГОСТ 2.105 пункт 4.4.5
			Время, необходимое для построения топологии и вычисления структурных характеристик, \textit{min.}  & \(T_{total}^{apmath} = 672\) &  \(T_{total}^{js} = 271\)  \\
			Среднее время загрузки одной страницы, \textit{min.} & 0.026 & 0.001 \\
			Объем памяти на жестком диске для хранения данных, \textit{MB} & \(\text{\textit{VMem}}_{apmath} = 154.743\) & \(\text{\textit{VMem}}_{js} = 318.463\) \\
			Минимальный объем \newline оперативной памяти \newline для обработки данных, \textit{GB} & \(\approx 1.5\) & \(\approx 1.5\) \\
			Средний объем памяти, требуемый  для полной обработки \newline одной веб-страницы, \textit{kB} & 6.06 & 11.92 \\
			Объем Веб трафика, \textit{GB} &  \(\text{\textit{Traffic}}_{apmath} = 9.66\) &  \(\text{\textit{Traffic}}_{js} = 3.1\) \\
			\bottomrule %%% нижняя линейка
	    \end{tabularx}%
\end{table}

\paragraph{Выводы.} Результаты эксперимента подтвердили существенные различия между гиперссылочной структурой представителя сайтов естественнонаучной направленности и структурой представителя гуманитарно-ориентированных сайтов. Так, например, в топологии первого сайта (рис.~\cref{fig:webGraphs-1}) большая доля всех веб-страниц сосредоточена в центральном ядре, которое имеет большую связность между своими элементами по сравнению с элементами ядра второго сайта (рис.~\cref{fig:webGraphs}\subcaptionref*{fig:webGraphs-2}). Кроме, того значительное различие компоненты OUT первой топологии в сравнении со второй объясняется наличием в ней большого числа полнотекстовых документов (в формате PDF, DOC, DOCX и др.). Однако, количество гиперссылок в Веб-графе \(G_{jf}\) и характеристик размеров компонент SCC, IN, TUBES/TENDRILS (Таб.~\cref{tab:webGraphTable}) говорят о хорошей коммуникабельности между компонентами топологии Веб-графа \(G_{jf}\).

Для получения статистически достоверных результатов необходимо провести экспериментальные исследования топологий большого числа сайтов для каждой выделенной категории. Опираясь на полученные значения характеристик (Таб.~\cref{tab:webGraphCostTable}), оценивающие трудоемкость поставленного в данной статье эксперимента, можно приблизительно оценить ресурсные затраты на полное исследование любого количества сайтов.

\subsection{Research of University Sites Internal Links Distribution}\label{subsec:ch1/sec2/sub2}

\paragraph{Intoduction.} The paper is concerned with the university sites structural properties study. An oriented graph is used as a generally accepted website structure representation. Nodes of the graph correspond to webpages, while arcs correspond to links between webpages. Consider that webpage does not have its own graph structure in terms of this graph.

Any site is a part of the World Wide Web, so we have to consider how to lay emphasis on a particular site.

First of all, technically speaking, the site has a domain name and actually only webpages corresponding to the domain can be referred to the organization website. Secondly, we consider subdomains to be external in relation to the website main domain. Thirdly, all links that are relevant to the main site webpages and subdomain webpages or other domain webpages are also considered to be external and are not part of the website.

University websites are usually both extremely large and significantly vary in size. A typical website has hundreds or thousands of pages, whereas there are also sites including 25, 10 thousand and 1 million pages (an approximate number of website pages can be estimated using Google's «site:» search term). Since the websites are so large, the task to determine structural, thematic, ergonomic characteristics of the site becomes truly laborious.

Knowledge about these characteristics will allow owners of large scientific and educational organizations’ websites to analyze the quality of their web resources presence in Web. Particularly, studies of structural characteristics will allow to monitor connections between web resources of organizations, to capture target groups of users and to evaluate site relevancy for them, to organize effective logistics of hyperlink structure \cite{StuartThelwallHarries,Thelwall,PechnikovNwohiri,PechnikovNwohiri2}. Ergonomic characteristics are estimating the adaptation of web resources of organizations to relevancy and ranking evaluation criteria of top search engines \cite{Huang,BodrunovaYakuninSmolin,LangvilleMeyer}. Enhancement of these characteristics will help stimulating user interest for every site and extend the time users spend on the sites which is the most important factor of relevancy in search engines rankings. However, estimating these characteristics for big sites is not an easy task in terms of time, computing power and system resources.

It seems natural to try to find an approximate solution to this issue, looking at a relatively small part of the site. Though, when you try to do this, you will get stuck at the point where you need to know what part of the site has already been looked through. To answer this question, you still need to look through the whole site to find out the total number of pages.

In our paper \cite{SergeevBlekanovMaksimov} we offered a method for website size (the number of webpages and links) estimation by its fraction (10\%). By collecting 11 university websites it has been shown that this method gives an estimate site webpages number with an average error of 40\%. The analysis was based on studies in \cite{BarabasiAlbert,BroderKumarMaghoul}, revealing that the pages are distributed according to the power law by the number of incoming links. We developed the idea that university sites are special in some way and that they have a different exponent. The preliminary analysis revealed that the changes of the exponent of the distribution law in our formulas \cite{SergeevBlekanovMaksimov} could have a significant influence on the site’s webpages and hyperlinks count estimation error. So, we performed an experiment using hyperlink structure of 97 universities from top 500 Webometrics  \cite{RankingWeb} rating. The information has been collected by means of authors' specialized web crawler \cite{BlekanovSergeevMartynenko}, adapted for webometric studies. Search engines like Google, Yahoo, etc. do not provide this kind of information.

\paragraph{Experiment.} \textit{Goal.} The experiment specifies the distribution law of university sites webpages by the number of incoming links. We assume that the distribution law is a power law. We need to find out the exponent for a power law.

\textit{Selecting university websites.} Since the websites of the world's leading universities are in great interests of the Webometrics rating, we decided to take first 500 sites into account. In order to reduce the huge amount of work, we have randomly chosen 100 university sites (from 1 to 500 of Webometrics rating), about 2-3 sites from every dozen.

Using web crawler designed by authors all pages and websites internal links of 97 universities have been collected and processed \cite{WebometricsLab}. We were unable to collect data from three sites due to their special layout. As a result of random selection universities were distributed by region as follows:
\begin{itemize}
	\item North America -- 42
	\item South America -- 2
	\item Europe -- 36
	\item Asia -- 13
	\item Australia and New Zealand -- 4
\end{itemize}

\textit{Experimental results.} Preliminary analysis displayed that university sites can be divided into several distinct groups by the number of webpages (Table~\cref{tab:sitesGroupedByNoOfPages}). This distinction caused some dubiety, whether or not it is possible to calculate the universal distribution exponent. 

\begin{table} [htbp]%
	\centering
	\caption{Sites grouped by the number of pages}%
	\label{tab:sitesGroupedByNoOfPages}% label всегда желательно идти после caption
	\renewcommand{\arraystretch}{1.5}%% Увеличение расстояния между рядами, для улучшения восприятия.
	\begin{SingleSpace}
		\begin{tabulary}{\textwidth}{@{}>{\zz}L >{\zz}C@{}} %Вертикальные полосы не используются принципиально, как и лишние горизонтальные (допускается по ГОСТ 2.105 пункт 4.4.5) % @{} позволяет прижиматься к краям
			\toprule     %%% верхняя линейка
			Pages count & Sites count  \\
			\midrule %%% тонкий разделитель. Отделяет названия столбцов. Обязателен по ГОСТ 2.105 пункт 4.4.5
			<1 300 & \textit{12} \\
			2 000-9 200 & \textit{10} \\
			10 000-25 000 & \textit{16} \\ 
			25 000-61 000 & \textit{12} \\
			61 000-100 000 &  \textit{7} \\
			100 000-400 000 & \textit{15} \\
			400 000-1 000 000  & \textit{9} \\ 
			>1 000 000  & \textit{16} \\
			\bottomrule %%% нижняя линейка
		\end{tabulary}%
	\end{SingleSpace}
\end{table}

We formed a plane for each university site, where horizontal coordinate is the number of incoming links, and vertical coordinate is the proportion of the page’s number. A set of points was applied to the plane. A power approximating curve was drawn, that helped to determine the exponent. We omitted some of the points while drawing the curve. These were points with the proportion of page count less than \(10^{-5}\). A
typical picture is shown in figure~\cref{fig:uniCollegeLondonCurve}. 

\begin{figure}[ht]
	\centerfloat{
		\includegraphics[scale=0.5]{uniCollegeLondonCurve}
	}
	\caption{Approximation curve for University College London site.}\label{fig:uniCollegeLondonCurve}
\end{figure}

We created a plane for each exponent of all university websites pages’ power distribution law where X axis represents the number of pages and Y axis represents exponent module. (Fig.~\cref{fig:sitesExponent}). 

\begin{figure}[ht]
	\centerfloat{
		\includegraphics[scale=0.4]{sitesExponent}
	}
	\caption{The exponent for different sites (X axis represents the number of pages of these sites). Black dot separates sites from first two groups from the rest of the sites.}\label{fig:sitesExponent}
\end{figure}

Figure 2 shows that sites from first and second groups have a vastly varying exponent that makes them stand out from other sites. This is most likely because their structure is influenced by two factors: management decisions regarding the structure and random spontaneous structure influences by active users. When a website is small (less than 10 thousand pages), management decisions influence is prevailing. And since every university has its own unique administrative unit, sites have different structural characteristics and different incoming link distributions. But when a website is big, the number of random structure influences becomes so large, that the role of management becomes less significant. 

\begin{figure}[ht]
	\centerfloat{
		\includegraphics[scale=0.4]{bigUniAveragedDistribution}
	}
	\caption{Averaged distribution for big university sites. Solid line represents power law distribution from A. Broder and R. Kumar (2000), Barabasi and Albert (1999) papers. Circles show variation areas of the averaged distribution values.}\label{fig:bigUniAveragedDistribution}
\end{figure}

\begin{figure}[ht]
	\centerfloat{
		\includegraphics[scale=0.4]{allUniAveragedDistribution}
	}
	\caption{Averaged distribution for all university websites. Solid line represents power law distribution from A. Broder and R. Kumar (2000), Barabasi and Albert (1999) papers.}\label{fig:allUniAveragedDistribution}
\end{figure}

We calculated two average exponents (Fig.~\cref{fig:bigUniAveragedDistribution},~\cref{fig:allUniAveragedDistribution}): \(k_\textit{uni\_all} = 1.8\) for all 97 sites and \(k_\textit{uni\_max} = 1.9\) just for big sites (75 sites with more than 10 thousand pages). These exponents differ from those obtained in \cite{BarabasiAlbert,BroderKumarMaghoul}. 

Figure~\cref{fig:bigUniAveragedDistribution} shows averaged distribution for all big sites. Two splashes (highlighted with circles) attract our attention: large deviations from the approximating curve. The first -- from 8 to 10 incoming links, the second -- around 69-71. The second can obviously be explained by the fact that the main page for most sites is here. The first is harder to explain. Perhaps, such a number of links is typical for main pages of different structural divisions of universities.

Figure~\cref{fig:allUniAveragedDistribution} shows that despite the fact there are many splashes due to considering small sized sites in terms of page count, the average exponent stays the same -- 1.8.

\paragraph{Conclusions and future work.} Our main results: 

\begin{itemize}
	\item All pages of 97 sites have been collected
	\item All internal incoming links for each page have been
	collected
	\item Exponent of power distribution for each university site
	has been calculated.
	\item The proximity of large (more than 10 thousand pages)
	sites distributions (\(k_\textit{uni\_max} = 1.9\)) and difference in
	distributions of small sites have been revealed (with
	the exponent between 0.0 and 5.0).
	\item The averaged approximating distribution curve (the
	exponent \(k_\textit{uni\_all} = 1.8\)) has been constructed and
	analyzed.
	\item Average exponent values have been calculated.
	Hyperlink structure of university sites was found to be
	different from \cite{BarabasiAlbert,BroderKumarMaghoul} in terms of the exponent.
\end{itemize}

The final results will be used to improve the algorithms for estimation of the size of any university site (page and hyperlink count) by studying just a part of the site. In the future, we plan to estimate structural, thematic and ergonomic characteristics of university websites by studying a randomly selected part. 

\section{Сбор данных}\label{sec:ch1/sec3}

\subsection{Построение тематико-ориентированных веб-краулеров с использованием обобщенного ядра}\label{subsec:ch1/sec3/sub1}

За последние несколько лет в области информационного Веб-поиска все чаще проводятся исследования \cite{Pechnikov,PechnikovLugovayaChuiko,PechnikovChirkovChuiko}, связанные с развивающимся научным направлением вебометрика (webometrics) \cite{PechnikovLugovayaChuiko}. В частности, методами вебометрики исследуются такие вопросы, как оценка присутствия информационных веб-ресурсов в Вебе, ранжирование веб-сайтов университетов и научных учреждений \cite{Pechnikov,PechnikovChirkovChuiko} или повышение вебометрического рейтинга, введенного испанской группой Cybermetrics Lab (http://www.webometrics.info), различных университетов мира.

К актуальным направлениям вебометрики относятся исследования по анализу и выявлению гиперссылочных структур различных сегментов Веб-пространства (например, академический сегмент Веба, корпоративный и др.). Для получения больших объемов информации о гиперссылках используются Веб-краулеры. Общей задачей таких инструментов является специализированный обход Веба с целью сбора информации или определения гиперссылочной структуры и полезности каких-либо информационных ресурсов \cite{BlekanovBondarenko1,BlekanovBondarenko2}. Имеющиеся в настоящее время Веб-краулеры можно разделить на две группы "--- универсальные и специализированные (тематико-ориентированные). Универсальные обычно избыточны, сильно нагружают сетевые ресурсы и проигрывают специализированным по таким показателям, как время обхода Веб-пространства, производительность обработки информации, а также возможность направленного поиска информационных источников в рамках определенной метрики значимости. В этом смысле преимущества специализированных систем очевидны, однако, их минусом являются значительные затраты на изготовление -- для каждой специализации нужен свой краулер.

Вместе с тем, значительное число функций всех краулеров совпадает \cite{ArasuChoGM,Bar-Ilan,ChoGM,NajorkHeydon}.

В данной статье предлагается модель специализированных Веб-краулеров, состоящих из обобщенного ядра и специализированных дополнений. При таком подходе затраты на изготовление специализированного Веб-краулера минимизируются. Предлагается архитектура обобщенного ядра поискового робота и ее реализация. Ставится эксперимент, в котором сравнивается качество работы Веб-краулеров на основе обобщенного ядра с качеством работы зарубежных аналогов, на примере Heritrix, OpenWebSpider и Methanol Web Crawler.

\paragraph{Архитектура обобщенного ядра.} Масштабность Веб-пространства порождают ряд проблем, влияющих на эффективность работы Веб-краулеров любых типов. Анализ известных публикаций \cite{ArasuChoGM,ChoGM,NajorkHeydon,RaghavanGM}, в которых затрагиваются проблемы эффективной работы поисковых роботов, показывает, что любой специализированный поисковый робот должен решать следующие основные проблемы:
\begin{itemize}
	\item \textit{построение специализированного алгоритма обхода Веб-пространства;}
	\item \textit{обновление информационных Веб-ресурсов в общей коллекции документов, собранной Веб-краулером;}
	\item \textit{учет разновидностей форматов Веб-ресурсов;}
	\item \textit{организация информационного поиска по скрытому Вебу} \cite{RaghavanGM}\textit{;}
	\item \textit{минимизация нагрузки на информационный источник;}
	\item \textit{распараллеливание и масштабирование процесса сбора информационных Веб-ресурсов;}
	\item \textit{минимизация нагрузки на каналы связи в Веб-пространстве.}
\end{itemize}

Для их решения была разработана и реализована архитектура обобщенного ядра (рис.~\cref{fig:kernelArchitecture}), которая может применяться для создания любых типов Веб-краулеров. В основе этой архитектуры лежит краулер-процесс. По структуре \textit{краулер-процесс} является классом, содержащим текущую информацию о состоянии поискового робота (текущий список проанализированных гиперссылок, содержимое Веб-станиц, кэш адресов Веб-страниц, список значимых источников информации), а также набор методов для осуществления поиска и обновления информации в индексе. Данные методы позволяют запускать и останавливать процесс поиска.

Под процессом поиска в рамках рассматриваемой архитектуры понимается итеративный процесс, останавливающийся при выполнении заданных условий (например, число пройденных итераций, количество полученных гиперссылок, сходимость результатов), которые зависят от настроек конфигурации. Каждая итерация данного процесса заключается в последовательном выполнении следующих основных модулей (рис.~\cref{fig:kernelArchitecture}).

\textit{Многопоточный загрузчик Веб-страниц.} Данный модуль исполняет роль загрузчика Веб-страниц с сервера, на котором они расположены в Веб-пространстве. Также он является менеджером безопасности, который контролирует количество потоков, выделяемое для загрузки всех информационных источников в рамках одной итерации, и не обрабатывает Веб-ресурсы, время отклика которых превышает заданное. Все параметры менеджера загрузок настраиваемые, по умолчанию число потоков равно \textit{10}, а время отклика -- \textit{2 секунды}.

\textit{Модуль извлечения гиперссылок (URLs extraction)}, который ищет дочерние элементы из заданного набора Веб-страниц. Поиск таких элементов осуществляется путем извлечения всех гиперссылок, находящихся в начальном множестве веб-страниц, и их добавления в очередь гиперссылок. Процесс обработки элементов очереди выполняется с помощью набора параллельных и синхронизированных между собой потоков, которые независимо извлекают ссылки и добавляют результаты обработки в общую коллекцию полученных дочерних Веб-страниц.

\textit{Модуль синтаксического анализа Веб-страниц (HTML parser)} -- извлекает текст из загруженных Веб-страниц и определяет его кодировку.

\textit{Нормализация гиперссылок (URLs normalization).} После того, как из начального множества Веб-страниц извлечены все гиперссылки, необходимо отфильтровать их адреса. За данную операцию в обобщенном ядре отвечает модуль фильтрации адресов для Веб-страниц. Основная задача такого модуля заключается в приведении адреса каждого Веб-ресурса к стандартизированному виду по заданным критериям \cite{LeeKimHong} (например, имя хоста и протокола каждого ресурса должны состоять из символов нижнего регистра или адреса всех страниц должны приводиться к одинаковой канонической форме, то есть, заканчиваться символом «/»). Другими словами, модуль фильтрации отсеивает Веб-страницы с некорректными и дублированными адресами, тем самым улучшая производительности работы поискового робота.

\textit{Модуль кэширования.} Операция извлечения содержимого Веб-страниц затрачивает ресурсы системы (интенсивность затрат ресурсов системы зависит от роста числа Веб-страниц), что значительным образом влияет на ее производительность. Поэтому была разработана система кэширования, предназначенная для ускорения процесса извлечения гиперссылок с Веб-страниц за счет повторного использования ранее извлеченного контента. Хранение контента продолжается до тех пор, пока не превышен допустимый порог памяти компьютера, используемой приложением. Данный порог приложение определяет динамически, исходя из имеющегося объема и требований по использованию ресурсов компьютера, а также производительности. При превышении порога система освобождает память путем удаления из кэша контента Веб-страниц, которые реже всего используются.

\textit{Коллекция найденных документов.} Данный модуль хранит информацию обо всех веб-ресурсах и их гиперссылочных структурах, полученных Веб-краулером на каждой итерации краулер-процесса, и предоставляет ее пользователю для дальнейших вебометрических исследований.

\begin{figure}[ht]
	\centerfloat{
		\includegraphics[scale=0.6]{kernelArchitecture}
	}
	\caption{Архитектура обобщенного ядра.}\label{fig:kernelArchitecture}
\end{figure}

Данная архитектура обобщенного ядра предоставляет конфигурируемые настройки интеграции и интерфейсы, использование которых существенно упрощают процесс и минимизируют время добавления в архитектуру новых модулей (рис.~\cref{fig:kernelModuleLink}), для создания потенциальных внешних модулей (новые алгоритм обхода, синтаксический и семантический анализаторы и др.).

Описанная реализация поискового робота является базовым инструментом для вебометрических исследований, поэтому на ее основе могут создаваться Веб-краулеры со специализированными алгоритмами обхода Веба.

\begin{figure}[ht]
	\centerfloat{
		\includegraphics[scale=0.6]{kernelModuleLink}
	}
	\caption{Связь обобщенного ядра с внешними модулями.}\label{fig:kernelModuleLink}
\end{figure}

\paragraph{Архитектурные особенности обобщенного ядра.} Перед экспериментом в данных исследованиях проводилась оценка архитектурных возможностей обобщенного ядра и зарубежных аналогов Веб-краулеров (Heritrix, OpenWebSpider, Methanol Web Crawler) по наличию следующей функциональности (табл.~\cref{tab:webCrawlerComparison}):
\begin{itemize}
	\item \textit{возможность многопоточной загрузки Веб-страниц;}
	\item \textit{масштабирования процесса сбора Веб-страниц (на программно-аппаратном уровне);}
	\item \textit{минимизация нагрузки на информационные Веб-ресурсы;}
	\item \textit{минимизации нагрузки на каналы связи;}
	\item \textit{гибкость архитектуры (возможность добавлять новые модули, алгоритмы обхода);}
	\item \textit{возможность специализированного обновления информационных Веб-ресурсов в индексе;}
	\item \textit{учет разновидностей форматов Веб-ресурсов;}
	\item \textit{обработка текста Веб-страниц;}
	\item \textit{приспособленность к российскому сегменту Веба;}
	\item \textit{приспособленность к англоязычному сегменту Веба.}
\end{itemize}

\begin{table} [htbp]%
	\centering
	\caption{Сравнение особенностей архитектуры Веб-краулеров.}%
	\label{tab:webCrawlerComparison}% label всегда желательно идти после caption
	\renewcommand{\arraystretch}{1.5}%% Увеличение расстояния между рядами, для улучшения восприятия.
	\begin{SingleSpace}
		\begin{tabulary}{\textwidth}{@{}>{\zz}L >{\zz}C >{\zz}C >{\zz}C >{\zz}C@{}}%Вертикальные полосы не используются принципиально, как и лишние горизонтальные (допускается по ГОСТ 2.105 пункт 4.4.5) % @{} позволяет прижиматься к краям
			\toprule     %%% верхняя линейка
			Функциональность &Heri-\linebreak trix & OpenWeb Spider & Methanol Web Crawler & Обобщенное ядро \\
			\midrule %%% тонкий разделитель. Отделяет названия столбцов. Обязателен по ГОСТ 2.105 пункт 4.4.5
			Многопоточная загрузка &+ &+ &+ &+ \\				
			Масштабирования процесса сбора Веб-ресурсов & -- & -- & -- &+ \\
			Минимизация нагрузки на Веб-ресурсы & + &+ &+ &+ \\
			Минимизации нагрузки на каналы связи &+ & -- & -- &+\\
			Гибкость архитектуры & --& -- & -- &+ \\
			Обновление индекса & -- & -- & -- &+ \\
			Учет разновидностей форматов Веб-ресурсов &+ &+ &+ &+ \\
			Обработка текста Веб-страниц &  -- & -- & -- &+ \\
			Адаптированность к рос. Вебу & -- & -- & -- &+ \\
			Адаптированность к англ. Вебу &+ &+ &+ &+ \\
			\bottomrule %%% нижняя линейка
		\end{tabulary}%
	\end{SingleSpace}
\end{table}

Оценка архитектурных возможностей показала, что обобщенное ядро и Heritrix имеют гибкие настройки (например, объем скачиваемой информации с одного источника, время доступа к серверу Веб-ресурса, количество итераций краулер-процесса и др.), которые при получении и обработке информационных Веб-ресурсов легко позволяют контролировать нагрузку на ресурсы и каналы связи (табл.~\cref{tab:webCrawlerComparison}). Веб-краулеры OpenWebSpider и Methanol Web Crawler плохо настраиваются (OpenWebSpider конфигурируется в течение \textit{34 минут}, а Methanol Web Crawler -- \textit{30 минут} (табл.~\cref{tab:crawlersEnglish},~\cref{tab:crawlersRussian})) и при обработке информации источников максимально нагружают каналы связи и используют ресурсы обрабатываемого источника. Вместе с этим зарубежные аналоги плохо адаптированы к российскому Вебу, что подтверждается далее в эксперименте.

\paragraph{Эксперимент.} В эксперименте ставилась задача сравнения эффективности работы четырех реализаций Веб-краулеров:
\begin{itemize}
	\item \textit{Веб-краулер на основе обобщенного ядра} с тремя дополнительными модулями: \textit{Module\_BDD}, \textit{Module\_MDD} и \textit{Module\_SDD}, ориентированные на загрузку большого количества Веб-страниц, среднего и малого, соответственно;
	\item \textit{Heritrix (http://crawler.archive.org)} \cite{MohrKimptonStack}, имеющий хорошую производительность при обработке большого количества Веб-страниц;
	\item \textit{Methanol Web Crawler (http://metha-sys.org)}, имеющий хорошую производительность при обработке малого количества с Веб-страниц;
	\item \textit{Open WebSpider (http://www.openwebspider.org)}, имеющий хорошую производительность при обработке небольшого количества с Веб-страниц.
\end{itemize}

Сравнивалась производительность работы и время конфигурирования Веб-краулеров на основе обобщенного ядра с зарубежными аналогами на примере загрузки \textit{2000}, \textit{300000} и \textit{500000} Веб-страниц.

Производительность поисковых роботов проверялась в российском и англоязычном сегментах Веба при одинаковом тестовом наборе данных (ЭВМ, ОС, начальное множество ссылок в очереди и др.). В качестве начальных параметров были заданы начальное множество Веб-страниц для каждого из сегмента (по \textit{10} Веб-страниц), и частота запуска, равная \textit{10}.

\paragraph{Результаты эксперимента.} В ходе сравнительного анализа производительности Веб-краулера на основе обобщенного ядра (с подключенными дополнительными модулями \textit{Module\_BDD}, \textit{Module\_MDD} и \textit{Module\_SDD}) с зарубежными аналогами (Heritrix, OpenWebSpider, Methanol Web Crawler) были получены следующие результаты (табл.~\cref{tab:crawlersEnglish},~\cref{tab:crawlersRussian}):

\begin{table} [htbp]%
	\centering
	\caption{Средняя производительность обработки информации Веб-краулерами в англоязычном сегменте Веб-пространства и время их настройки.}%
	\label{tab:crawlersEnglish}% label всегда желательно идти после caption
	\renewcommand{\arraystretch}{1.5}%% Увеличение расстояния между рядами, для улучшения восприятия.
	\begin{SingleSpace}
		\begin{tabulary}{\textwidth}{@{}>{\zz}L >{\zz}C >{\zz}C >{\zz}C >{\zz}C@{}}%Вертикальные полосы не используются принципиально, как и лишние горизонтальные (допускается по ГОСТ 2.105 пункт 4.4.5) % @{} позволяет прижиматься к краям
			\toprule     %%% верхняя линейка
			Веб-краулеры & Время загрузки 2000 страниц (мин.) & Время загрузки 30000 страниц (мин.) & Время загрузки 50000 страниц (мин.) & Время настройки Веб-краулера (мин.) \\
			\midrule %%% тонкий разделитель. Отделяет названия столбцов. Обязателен по ГОСТ 2.105 пункт 4.4.5
			Heritrix & 1.9 & 8.0 & 11.2 & 20.0 \\				
			OpenWebSpider & 3.0 & 8.4 & 13.0 & 34.0 \\
			Methanol & 1.9 & 8.6 & 14.1 & 30.0 \\			
			Обобщенное ядро с \textit{Module\_BDD} модулем & -- & -- & \textbf{9.1} & 4.0\\
			Обобщенное ядро с \textit{Module\_MDD} модулем & -- & \textbf{6.0} & -- & 4.0\\			
			Обобщенное ядро с \textit{Module\_SDD} модулем & \textbf{1.7} & -- & -- & 4.0\\		
			Обобщенное ядро с модулем анализа текста и ссылок & \textbf{2.0} & \textbf{8.5} & \textbf{11.5} & 4.0\\		
			\bottomrule %%% нижняя линейка
		\end{tabulary}%
	\end{SingleSpace}
\end{table}

\begin{table} [htbp]%
	\centering
	\caption{Средняя производительность обработки информации Веб-краулерами в российском сегменте Веб-пространства и время их настройки.}%
	\label{tab:crawlersRussian}% label всегда желательно идти после caption
	\renewcommand{\arraystretch}{1.5}%% Увеличение расстояния между рядами, для улучшения восприятия.
	\begin{SingleSpace}
		\begin{tabulary}{\textwidth}{@{}>{\zz}L >{\zz}C >{\zz}C >{\zz}C >{\zz}C@{}}%Вертикальные полосы не используются принципиально, как и лишние горизонтальные (допускается по ГОСТ 2.105 пункт 4.4.5) % @{} позволяет прижиматься к краям
			\toprule     %%% верхняя линейка
			Веб-краулеры & Время загрузки 2000 страниц (мин.) & Время загрузки 30000 страниц (мин.) & Время загрузки 50000 страниц (мин.) & Время настройки Веб-краулера (мин.) \\
			\midrule %%% тонкий разделитель. Отделяет названия столбцов. Обязателен по ГОСТ 2.105 пункт 4.4.5
			Heritrix & 5.0 & 12.9 & 22.3 & 20.0 \\				
			OpenWebSpider & 5.5 & 13.0 & 25.0 & 34.0 \\
			Methanol & 4.2 & 13.5 & 26.5 & 30.0 \\			
			Обобщенное ядро с \textit{Module\_BDD} модулем & -- & -- & \textbf{10.5} & 4.0\\
			Обобщенное ядро с \textit{Module\_MDD} модулем & -- & \textbf{6.4} & -- & 4.0\\			
			Обобщенное ядро с \textit{Module\_SDD} модулем & \textbf{2.0} & -- & -- & 4.0\\		
			Обобщенное ядро с модулем анализа текста и ссылок & \textbf{3.5} & \textbf{9.0} & \textbf{13.0} & 4.0\\		
			\bottomrule %%% нижняя линейка
		\end{tabulary}%
	\end{SingleSpace}
\end{table}

\begin{itemize}
	\item получено небольшое время добавления внешних модулей (алгоритмы обхода Веб-пространства, синтаксические анализаторы, методы обновления индекса и др.) к обобщенному ядру и их конфигурирования (табл.~\cref{tab:webCrawlerComparison},~\cref{tab:crawlersEnglish}) по сравнению с зарубежными аналогами, архитектура которых не предусматривает возможность масштабирования и добавления в нее других модулей и время настройки в среднем занимает \textit{28 минут}.
	\item Веб-краулер с обобщенным ядром имеет возможность добавления в него специализированных модулей и их адаптации для тематического поиска (рис.~\cref{fig:crawlerModes}): классический (рекурсивно обходит все встречающиеся гиперссылки); интересо-ориентированный; ориентированный на популярность; ориентированный на интересы и популярность. В исследованиях \cite{BlekanovBondarenko1,BlekanovBondarenko2} проверялась качество поиска и полезность этих модулей, а также их интеграция с обобщенным ядром. В то время как Heritrix, OpenWebSpider и Methanol Web Crawler используют только классический режим обхода Веб-пространства и не обрабатывают текст источников информации в Вебе, но поддерживают возможность поиска музыкальных, видео файлов, файлов PDF.
	
	\begin{figure}[ht]
		\centerfloat{
			\includegraphics[scale=0.6]{crawlerModes}
		}
		\caption{Режимы поиска Веб-краулера с обобщенным ядром.}\label{fig:crawlerModes}
	\end{figure}
	
	\item Веб-краулеры на основе обобщенного ядра с разными модулями загрузки Веб-страниц (\textit{Module\_BDD}, \textit{Module\_MDD} и \textit{Module\_SDD}) в классическом режиме (обрабатывая информацию только о ссылках) показали высокую производительность среди зарубежных аналогов (табл.~\cref{tab:crawlersEnglish},~\cref{tab:crawlersRussian}), как в российском, так и зарубежном сегменте Веба. Вместе с этим поисковый робот на основе обобщенного ядра с модулем анализа текста не уступает по производительности зарубежным аналогам, которые не обрабатывают текстовую информацию Веб-страниц, в англоязычном сегменте Веба, а в российском Вебе зарубежные реализации Веб-краулеров сильно уступают в производительности. 
	\item зарубежные аналоги показали низкую производительность обработки информационных Веб-ресурсов в российском сегменте Веба (табл.~\cref{tab:crawlersEnglish}), а \textit{70\%} ресурсов из этого сегмента Веб-краулеры Heritrix, OpenWebSpider и Methanol Web Crawler не смогли загрузить и обработать. Данный факт объясняется низкой приспособленностью зарубежных аналогов к российскому Вебу.
	\item Среди зарубежных аналогов неплохую эффективность загрузки и обработки \textit{2000}, \textit{300000} и \textit{500000} Веб-страниц показал Веб-краулер Heritrix. Его можно использовать как неплохой инструмент для поиска (с классическим алгоритмом обхода Веба) разнородной информации и добавления ее в индекс.
\end{itemize}

Таким образом, из поставленного эксперимента следует, что Веб-краулер с обобщенным ядром эффективней собирает и обрабатывает информацию с Веб-ресурсов, чем зарубежные аналоги, которые сильно уступают по производительности, гибкости и масштабируемости архитектуры, а также приспособленностью к обработке информации в российском сегменте Веба.

В дальнейшем созданный Веб-краулер с обобщенным ядром планируется использовать в прикладных исследованиях по повышению вебометрических рейтингов университетов, научных институтов и других научных организаций.

\subsection{Веб-краулер как инструмент для вебометрических исследований на примере анализа Веб-пространства СПбГУ}\label{subsec:ch1/sec3/sub2}

\paragraph{Введение.} За последние несколько лет в области информационного Веб-поиска все чаще появляются задачи, связанные с развивающимся научным направлением вебометрика (webometrics) \cite{HolmbergThelwall,Pechnikov,PechnikovChirkovChuiko,BlekanovSergeevPechnikov}. К актуальным направлениям вебометрических исследований относятся задачи анализа и выявления гиперссылочных структур различных сегментов Веб-пространства (например, академический сегмент Веба, университетский, и др.), решение которых влияет на качество присутствия этих сегментов в Вебе, на результаты ранжирования поисковых машин (Google, Yandex и др.) или, в случае университетского Веба, на вебометрический рейтинг (http://www.webometrics.info) различных университетов мира \cite{BlekanovSergeevPechnikov}.

 Для получения и обработки больших объемов информации о веб-сайтах и их гиперссылках используются Веб-краулеры (поисковые роботы), общей задачей которых является специализированный обход Веба с целью сбора информации или определения гиперссылочной структуры и полезности каких-либо информационных ресурсов.
 
\paragraph{Эксперимент.} В эксперименте ставилась задача анализа и выявления гиперссылочной структуры Веб-пространства Санкт-Петербургского государственного университета (СПбГУ). 
 
 Для эксперимента использовался программный комплекс обобщенного ядра поискового робота, который обладает высокой гибкостью и масштабируемостью в сравнении с зарубежными аналогами, сильно уступающими в производительности собора и обработки веб-ресурсов и имеющими слабую приспособленность к анализу российского сегмента Веба \cite{BlekanovSergeevMartynenko}. 
 
 К Веб-краулеру с обобщенным ядром дополнительно был разработан и добавлен специализированный алгоритм обхода веб-страниц, который собирает и обрабатывает только страницы из Веб-пространства СПбГУ. В свою очередь пространство СПбГУ состоит из веб-сайта главного домена и сайтов всех его поддоменов (рис.~\cref{fig:spbuWebSpace}).
 
\begin{figure}[ht]
    \centerfloat{
	        \includegraphics[scale=0.5]{spbuWebSpace}
	    }
    \caption{Веб-пространство СПбГУ.}\label{fig:spbuWebSpace}
\end{figure}

Используя программный комплекс на основе обобщенного ядра поискового робота со специализированным алгоритмом, запущенного с начального множества веб-страниц, требовалось в автоматизированном режиме получить значения следующих показателей, характеризующих гиперссылочную структуру Веб-пространства СПбГУ:
\begin{itemize}
	\item объем Веб-пространства СПбГУ (количество всех различных веб-страниц из Веб-пространства СПбГУ);
	\item количество всех поддоменов из Веб-пространства СПбГУ; 
	\item количество тупиковых (не имеющих ссылок) веб-страниц; 
	\item количество неработающих гиперссылок; 
	\item количество гиперссылок на внешние веб-ресурсы; 
	\item количество поддоменов, связанных с «Главной страницей» главного домена; 
	\item количество поддоменов, несвязанных с «Главной страницей» главного домена; 
	\item количество страниц, имеющие гиперссылки на «Главную страницу» главного домена; 
	\item количество страниц, не имеющие гиперссылки на «Главную страницу» главного домена; 
	\item гиперссылочная структура Веб-пространства СПбГУ в виде матрицы смежности. 
\end{itemize}

В качестве начального множества веб-страниц, с которого Веб-краулер запускал процесс сбора и обработки веб-ресурсов, брался URL-адрес главного веб-сайта СПбГУ -- «http://www.spbu.ru/».

\paragraph{Результаты эксперимента.} В ходе эксперимента всего Веб-краулером было обработано и проанализировано \textit{6\space429\space963} гиперссылки, которые содержались на страницах Веб-пространства СПбГУ. Из них: объем ссылок на внешние источники информации равен \textit{507\space168}, а объем внутренних ссылок (на страницы главного домена сайта СПбГУ и его поддоменов) -- \textit{5\space922\space795}. Кроме того, были получены следующие результаты (табл.~\cref{tab:indicators}):

\begin{table} [htbp]%
	\centering
	\caption{}%
	\label{tab:indicators}% label всегда желательно идти после caption
	\renewcommand{\arraystretch}{1.5}%% Увеличение расстояния между рядами, для улучшения восприятия.
	\begin{SingleSpace}
		\begin{tabulary}{\textwidth}{@{}>{\zz}L >{\zz}C@{}} %Вертикальные полосы не используются принципиально, как и лишние горизонтальные (допускается по ГОСТ 2.105 пункт 4.4.5) % @{} позволяет прижиматься к краям
			\toprule     %%% верхняя линейка
			Показатель & Значение показателя  \\
			\midrule %%% тонкий разделитель. Отделяет названия столбцов. Обязателен по ГОСТ 2.105 пункт 4.4.5
			объем Веб-пространства СПбГУ & 71688 \\				
			количество всех поддоменов из Веб-пространства СПбГУ & 315 \\
			количество поддоменов, сильно-связанных с «Главной страницей» главного домена & 280 \\
			количество поддоменов, несвязанных с «Главной страницей» главного домена & 35 \\
			количество тупиковых веб-страниц & 12516 \\
			количество неработающих гиперссылок & 680 \\
			количество страниц, имеющие гиперссылки на «Главную страницу» главного домена & 10632 \\
			количество страниц, не имеющие гиперссылки на «Главную страницу» главного домена & 61056 \\
			\bottomrule %%% нижняя линейка
		\end{tabulary}%
	\end{SingleSpace}
\end{table}

Также стоит отметить, что для простоты исследований и получения численных значений показателей было предложено хранить выявленную гиперссылочную структуру Веб-пространства СПбГУ в виде матрицы смежности A размерностью [71688\(\times\)71688]. С помощью нее и ее транспонированной формы удобно определять связанность между страницами и в дальнейшем оптимизировать структуру ссылок. 

\paragraph{Общие выводы.} По результатам эксперимента видно, что \textit{17,5\%} веб-страниц из всех страниц Веб-пространства СПбГУ занимают тупиковые страницы, которые, например, не учитывает поисковая машина Google в своем алгоритме ранжирования PageRank, что негативно влияет на коммуникабельность всего Веб-пространство университета в целом. 

Также из эксперимента наблюдается плохая связанность страниц Веб-пространства СПбГУ и его поддоменов со страницами главного сайта университета: только \textit{14,8\%} всех веб-страниц имеет ссылки на страницы главного сайта, а \textit{35} доменов со всеми своими страницами и вовсе никак не связаны с ним. 

Стоит отметить и большое количество неработающих гиперссылок (\textit{680} уникальных гиперссылок), которые неоднократно повторяются по всему Веб-пространству СПбГУ, тем самым снижают доверие пользователей к ресурсам сайтов университета. 

Таким образом, проведенный эксперимент демонстрирует слабую связанность и коммуникабельность внутренних ресурсов Веб-пространства СПбГУ, что влечет и его слабую позицию в вебометрическом рейтинге сайтов университетов мира.

\subsection{Вебометрические исследования сегмента университетского Веба с помощью поискового робота}\label{subsec:ch1/sec3/sub3}

\paragraph{Введение.} За последние годы все больше российских ВУЗов следят за позициями своих веб-сайтов в вебометрическом рейтинге испанской группы Cybermetrics Lab \cite{RankingWeb}, который входит в пятерку мировых рейтингов наряду с такими рейтингами, как Times, QS, шанхайский и т. п. Этот рейтинг оценивает лишь качество сайта университета, но качество веб-сайта оказывает большое влияние на «видимость» университета в целом. Ведь по публикациям ученых видны только они, но не виден университет.

Для повышения рейтинга университетского сайта необходимо работать в таких направлениях, как увеличение объема, улучшение качества контента, увеличение количества гиперссылок на сайт с авторитетных веб-сайтов, а также оптимизация внутренней гиперссылочной структуры сайта. Исследованиям в этой области посвящены работы \cite{HolmbergThelwall,Pechnikov,Pechnikov2,PechnikovChirkovChuiko,BlekanovSergeevPechnikov,BlekanovSergeevMaksimov,BlekanovSergeev}. Мы же решили сосредоточиться на вопросах внутренней гиперссылочной структуры как целого Веб-пространства университета, так и отдельных сегментов, из которых оно состоит. Правильно организованная веб-структура может значительно увеличить коммуникабельность сайта и повысить его показатель цитируемости в поисковых системах, увеличить количество внешних ссылок, улучшить позиции университетского веб-сайта в вебометрическом рейтинге.

\paragraph{Эксперимент.} В эксперименте ставилась задача анализа и выявления гиперссылочной структуры одного из сегментов, составляющих Веб-пространство СПбГУ, -- веб-сайт факультета ПМ–ПУ, а также получения рекомендаций по улучшению этой структуры.

В качестве инструмента для достижения поставленной в эксперименте цели использовалась аналитическая система на основе специ- ализированного поискового робота, которая была успешно апробирована на исследованиях зарубежного и российского сегментов Веба \cite{BlekanovSergeevPechnikov,BlekanovSergeevMaksimov,BlekanovSergeev,BlekanovSergeevMartynenko}.

Специализация используемого поискового робота была сконфигурирована таким образом, чтобы робот собирал и обрабатывал только страницы в рамках сегмента Веба, содержащего веб-сайт факультета ПМ–ПУ СПбГУ, и гиперссылки с этого сайта на внешние источники информации (рис.~\cref{fig:spbuApmathWebSpace}).

\begin{figure}[ht]
	\centerfloat{
		\includegraphics[scale=1]{spbuApmathWebSpace}
	}
	\caption{Веб-пространство СПбГУ и включенный в него исследуемый сегмент Веба.}\label{fig:spbuApmathWebSpace}
\end{figure}

С помощью аналитической системы на основе специализированного поискового робота, запущенной с начальным URL-адресом http://www.apmath.spbu.ru/, требовалось выявить гиперссылочную структуру сайта ПМ–ПУ в виде матрицы смежности и определить значения параметров, приведенных в таблице~\cref{tab:researchedParameters}.

\begin{table} [htbp]%
	\centering
	\caption{Исследуемые параметры.}%
	\label{tab:researchedParameters}% label всегда желательно идти после caption
	\renewcommand{\arraystretch}{1.5}%% Увеличение расстояния между рядами, для улучшения восприятия.
	\begin{SingleSpace}
		\begin{tabulary}{\textwidth}{@{}>{\zz}L >{\zz}C@{}} %Вертикальные полосы не используются принципиально, как и лишние горизонтальные (допускается по ГОСТ 2.105 пункт 4.4.5) % @{} позволяет прижиматься к краям
			\toprule     %%% верхняя линейка
			ID параметра & Параметр  \\
			\midrule %%% тонкий разделитель. Отделяет названия столбцов. Обязателен по ГОСТ 2.105 пункт 4.4.5
			1 & Общее количество уникальных внутренних страниц сайта ПМ–ПУ \\				
			2 & Количество страниц, связанных с «Главной страницей» сайта ПМ–ПУ \\
			3 & Количество страниц, не связанных с «Главной страницей» сайта ПМ–ПУ \\
			4 & Количество тупиковых веб-страниц \\
			5 & Количество уникальных гиперссылок сайта ПМ–ПУ \\
			6 & Количество неработающих гиперссылок0 \\
			7 & Количество гиперссылок на внешние веб-ресурсы \\
			8 & Количество внутренних гиперссылок на документы форматов *.doc (*.docx), *.pdf, *excel, *.ppt \\
			9 & Количество гиперссылок с переадресацией \\
			\bottomrule %%% нижняя линейка
		\end{tabulary}%
	\end{SingleSpace}
\end{table}

\paragraph{Результаты эксперимента.} В ходе эксперимента поисковым роботом на веб-сайте факультета ПМ–ПУ была обработана и проанализирована \textit{2 258 521} гиперссылка. Из них \textit{2 026 285} гиперссылок на внутренние веб-страницы сайта ПМ–ПУ, \textit{119 931} ссылка на страницы из Веб-пространства СПбГУ и \textit{112 305} ссылок на внешние источники информации. Кроме того, с помощью построенной матрицы смежности и ее транспонированной формы для параметров из таблицы~\cref{tab:researchedParameters} были получены числовые значения в виде гистограмм и круговых диаграмм (рис.~\cref{fig:histogramOfTab1},~\cref{fig:linkDiagram}).

\begin{figure}[ht]
	\centerfloat{
		\includegraphics[scale=0.5]{histogramOfTab1}
	}
	\caption{Гистограммы значений показателей таблицы~\cref{tab:researchedParameters}.}\label{fig:histogramOfTab1}
\end{figure}

\begin{figure}[ht]
	\centerfloat{
		\includegraphics[scale=0.5]{linkDiagram}
	}
	\caption{Круговая диаграмма внешних и внутренних ссылок.}\label{fig:linkDiagram}
\end{figure}

\paragraph{Общие выводы.} По результатам эксперимента видно, что тупиковые страницы (не имеющие никаких других ссылок), не являющиеся медиа-контентом, отсутствуют. В свою очередь, 16,2\% страниц не связаны с главной веб-страницей главного сайта. Однако если отбросить страницы с документами, то это отношение составит всего около 1\%.
Несмотря на относительно неплохую связность страниц сайта и отсутствие тупиковых страниц, ситуацию ухудшает огромное число гиперссылок, переадресующих пользователя в основном на главную страницу сайта или ее иностранные аналоги. Кроме того, не все гиперссылки приводятся к единой нормальной форме на стороне сервера. Большое количество таких гиперссылок влекут вероятность некорректного цитирования информационных ресурсов веб-сайта ПМ–ПУ в поисковых системах, а также оказывают негативное влияние на вебометрический рейтинг университета в целом.
Стоит отметить и большое количество веб-страниц, при переходе на которые возникает ошибка, вызванная отсутствием предполагаемых данных в указанном месте сервера или неверной обработкой адреса. Наличие таких страниц снижает доверие пользователей к ресурсам сайта факультета.
Таким образом, для улучшения присутствия веб-сайта факультета ПМ–ПУ в Вебе и улучшения его показателя цитируемости в поисковых системах необходимо разрешить выявленные в эксперименте проблемы на практике.

\section{Анализ гиперссылочной структуры ресурсов}\label{sec:ch1/sec4}

\subsection{Исследование закономерностей в гиперссылочной структуре сайтов большой размерности (должна быть на англ.)}\label{subsec:ch1/sec4/sub5}

\paragraph{Актуальность.} В настоящее время редкая организация не имеет собственного сайта. Сайт, в определенном смысле, лицо организации и его качеству придается большое значение. Качество сайта складывается из многих составляющих. Это и количество страниц, и их дизайн, и содержание страниц, и структура сайта.

Представление структуры сайта в виде ориентированного графа, узлы которого -- документы, а дуги -- ссылки, является общепринятым \cite{BroderKumarMaghoul}. Глобальными структурными характеристиками информационных сетей в Вебе занимаются многие исследователи. Так, в работе Broder A. и Kumar F. \cite{BroderKumarMaghoul} крупные сообщества сайтов представляются в ввиде графовой компоненты сильной связности, компонент In, Out и Tubes; в работах \cite{Thelwall,ThelwallZuccala,ThelwallWilkinsonMusgrove,PechnikovNwohiri} рассматривается распределение внешних ссылок и индекс цитирования различных университетских сайтов; в \cite{BlekanovSergeevMaksimovBOWTIE} вводятся характеристики связности сайта; в \cite{KenekayoroBuckleyThelwall} исследуется метод автоматической классификации ссылок и страниц по их характеристикам; и др.

Цель настоящей работы – исследование структуры сайта. Точнее – установления вида функциональной зависимости, связывающей число страниц сайта с числом его внутренних ссылок.

\subsubsection{Проверка гипотезы}

\paragraph{Постановка задачи.} Для решения этой задачи использовался специальный поисковый робот \cite{BlekanovSergeevMartynenko}. В процессе обхода сайта робот составляет два списка: список найденных страниц (узлов веб-графа) и список ссылок, связывающих найденные страницы (дуг веб-графа).

Шагом алгоритма просмотра сайта будем считать акт нахождения \(e\) ссылок. То есть, в результате \(i\) шагов находятся \(E_i = ei\) ссылок. Количество найденных на \(i\) шагах страниц обозначим через \(v_i = v(E_i)\). Будем обозначать число всех ссылок сайта через \(e_0\), число всех страниц -- \(v(e_0) = v_0\) . Поисковый робот позволяет не только найти число страниц и ссылок сайта, но и получить график функции \(v(e)\)).

Очевидно, что \(v \leq e \). Причем \(v_0 < e_0\) (случай \(v_0 = e_0\) для сайтов -- экзотический). Очевидно, также, что рост \(v(e)\) постепенно замедляется. Это объясняется тем, что с ростом числа найденных страниц растет вероятность того, что очередная ссылка указывает на уже найденную страницу. 

Попытаемся подобрать аналитическую функцию, приближенно описывающую экспериментальную \(v(E)\), со следующими свойствами:
\begin{enumerate}
	\item На плоскости \((e, v)\) проходит через начало координат.
	\item Возрастает с ростом \(e\).
	\item Производная функции убывает.
	\item Функция имеет простой вид и зависит от небольшого числа параметров.
\end{enumerate}

\paragraph{Эксперимент.} Будем сопоставлять функцию с экспериментальным набором \(V_i, E_i (i = 1,2, \dots, N)\), где \(N = \lceil \frac{e_0}{e} \rceil\). Нами исследованы сайты четырех университетов и получены соответствующие наборы (табл.~\cref{tab:uniSites}).

\begin{table} [htbp]%
	\centering
	\caption{Сайты исследуемых университетов.}%
	\label{tab:uniSites}% label всегда желательно идти после caption
	\renewcommand{\arraystretch}{1.5}%% Увеличение расстояния между рядами, для улучшения восприятия.
	\def\tabularxcolumn#1{m{#1}}
	\begin{tabularx}{\textwidth}{@{}>{\raggedright}X >{\centering}m{3.5cm} >{\centering}m{2.5cm} >{\centering\arraybackslash}m{2.5cm}@{}}% Вертикальные полосы не используются принципиально, как и лишние горизонтальные (допускается по ГОСТ 2.105 пункт 4.4.5) % @{} позволяет прижиматься к краям
		\toprule     %%% верхняя линейка
		Университет & URL-адрес & Всего страниц & Всего ссылок \\
		\midrule %%% тонкий разделитель. Отделяет названия столбцов. Обязателен по ГОСТ 2.105 пункт 4.4.5
		%Санкт-Петербургский государственный университет & www.spbu.ru & 41 183 & 2 664 000 \\
		Московский государственный университет & www.msu.ru & 47 832 & 1 891 000 \\				
		Университет Айдзу & www.u-aizu.ac.jp & 4 161 & 49 900 \\
		Токийский университет & www.u-tokyo.ac.jp & > 17 000 & 240 000 \\			
		\bottomrule %%% нижняя линейка
	\end{tabularx}%
\end{table}

Для оценки качества приближения будем использовать среднюю относительную погрешность:
\[
\Delta = \frac{1}{N} \sum \frac{\lvert V_i - v_i \rvert}{v_i}
\]

Далее рассматриваются три варианта аппроксимирующей функции:
\[v^{(1)} = \alpha E^\beta ,\]
\[v^{(2)} = \frac{\alpha E}{E + \beta},\]
\[v^{(3)} = \alpha (\ln(1 + E))^\beta .\]

Все три функции удовлетворяют требованиям 1, 2 и 3 при \(\alpha > 0, \beta > 0\). \(v^{(1)}\) удовлетворяет требованию 3, если к тому же  \(\beta < 0\).

Рассмотрим их последовательно:

\begin{enumerate}
	\item Прологарифмируем уравнение: \[\ln v^{(1)} = \ln \alpha + \beta \ln E .\] Получаем систему линейных уравнений \[x + A_i y = B_i, i = \overline{1,N} ,\] где \(x = \ln \alpha, y = \beta, A_i = \ln E_i, B_i = \ln V_i\). Ее решение методом наименьших квадратов: \[\alpha = \exp \frac{C_1C_4 - C_2C_3}{N(C_1^2 - C_3)}, \beta = \frac{C_1C_2 - C_4}{C_1^2 - C_3},\] где \(C_1 = \sum_{1}^{N}A_i, C_2 = \sum_{1}^{N}B_i, C_3 = \sum_{1}^{N}A_i^2, C_4 = \sum_{1}^{N}A_i B_i\).
	\item Система линейных уравнений: \[\alpha E_i - \beta V_i = V_i E_i, i = \overline{1,N}.\] Ее решение методом наименьших квадратов: \[\alpha = \frac{a_3 a_4 - a_2 a_5}{a_1 a_4 - a_2^2}, \beta = \frac{a_2 a_3 = a_1 a_5}{a_1 a_4 - a_2^2}, \] где \(a_1 = \sum_{1}^{N} E_i^2, a_2 = \sum_{1}^{N} E_i V_i, a_3 = \sum_{1}^{N} E_i^2 V_i, a_4 = \sum_{1}^{N} V_i^2, a_5 = \sum_{1}^{N} V_i^2 E_i\).
	\item Логарифмируем \(v^{(3)}\): \[\ln v^{(3)} = \ln \alpha + \beta \ln \ln (1 + E).\] Получаем систему, почти совпадающую со случаем 1). Ее решение будет таким же, за исключением того, что в первом случае \(A_i = \ln E_i\), а в этом \[A_i = \ln \ln (1+E).\]
\end{enumerate}

\paragraph{Результаты эксперимента.} На рисунках~\cref{fig:msuFunc},~\cref{fig:aizuFunc},~\cref{fig:tuniFunc} представлены графики функций \(v^{(1)}, v^{(2)}, v^{(3)}\) и \(V\) для каждого из перечисленных университетов.

\begin{figure}[ht]
	\centerfloat{
		\includegraphics[scale=0.35]{msuFunc}
	}
	\caption{Функции \(v^{(1)}, v^{(2)}, v^{(3)}\) и \(V\) Московского государственного университета.}\label{fig:msuFunc}
\end{figure}

\begin{figure}[ht]
	\centerfloat{
		\includegraphics[scale=0.35]{aizuFunc}
	}
	\caption{Функции \(v^{(1)}, v^{(2)}, v^{(3)}\) и \(V\) для университета Айдзу.}\label{fig:aizuFunc}
\end{figure}

\begin{figure}[ht]
	\centerfloat{
		\includegraphics[scale=0.35]{tuniFunc}
	}
	\caption{Функции \(v^{(1)}, v^{(2)}, v^{(3)}\) и \(V\) для Токийского университета.}\label{fig:tuniFunc}
\end{figure}

В табл.~\cref{tab:uniErrors} приведены относительные погрешности каждой из формул для каждого из университетов.

\begin{table} [htbp]%
	\centering
	\caption{Средняя относительная погрешность функций \(v^{(1)}, v^{(2)}, v^{(3)}\) для каждого университета.}%
	\label{tab:uniErrors}% label всегда желательно идти после caption
	\renewcommand{\arraystretch}{1.5}%% Увеличение расстояния между рядами, для улучшения восприятия.
	\def\tabularxcolumn#1{m{#1}}
	\begin{tabularx}{\textwidth}{@{}>{\raggedright}X >{\centering}m{3.5cm} >{\centering}m{2.5cm} >{\centering\arraybackslash}m{2.5cm}@{}}% Вертикальные полосы не используются принципиально, как и лишние горизонтальные (допускается по ГОСТ 2.105 пункт 4.4.5) % @{} позволяет прижиматься к краям
		\toprule     %%% верхняя линейка
		Университет & \(v^{(1)}\) & \(v^{(2)}\) & \(v^{(3)}\) \\
		\midrule %%% тонкий разделитель. Отделяет названия столбцов. Обязателен по ГОСТ 2.105 пункт 4.4.5
		%Санкт-Петербургский государственный университет & www.spbu.ru & 41 183 & 2 664 000 \\
		Московский государственный университет & 0,092 & 0,037 & 0,014 \\				
		Университет Айдзу & 0,167&  0,209 &  0,229 \\
		Токийский университет & 0,068 & 0,076 & 0,062 \\			
		\bottomrule %%% нижняя линейка
	\end{tabularx}%
\end{table}

\paragraph{Выводы.} Рассмотрена проблема нахождения функции, аппроксимирующей экспериментальный график зависимости числа найденных страниц сайта от числа найденных ссылок. Предложены три варианта аппроксимации. Выяснено, что наилучшие приближения получаются для степенной и дробно-линейной функции. Предложенные функции могут быть использованы для изучения параметров сайтов и их кластеризации.

\subsection{Применение модифицированного алгоритма LSH для кластеризации внешнего окружения веб-пространства университетов}\label{subsec:ch1/sec4/sub7}

Для большинства крупных организаций немаловажное значение имеет их рейтинг, который рассчитывается в зависимости от параметров, связанных с их родом деятельности. В частности, на общий рейтинг ВУЗа большое влияние оказывает его вебометрический рейтинг. Известно, одним из главных показателей, который влияет на вебометрический рейтинг любой организации, в том числе и университета, является количество внешних ресурсов \cite{RankingWeb}, ссылающихся на сайт университета. Но кроме количества ссылающихся ресурсов также важно понимать качество этих ресурсов, их природу, определить к какой области относится тот или иной внешний ресурс. Данные веб-ресурсы образуют группы сайтов с одинаковым родом деятельности. Таким образом, возникает задача выявления этих групп -- кластеризации. Требуется определить степень влияния количества и размеров найденных групп на вебометрический рейтинг сайтов университетов. По найденным кластерам можно определить, с какими группами внешних ресурсов следует сайтам университетов выстраивать гиперссылочные взаимосвязи для повышения цитируемости.

Объектом данного исследования являются университетские сайты, которые имеют чрезвычайно большие размеры (например, сайт СПбГУ содержит более \textit{50} тыс. внутренних веб-страниц и около 5 млн. гиперссылок) \cite{BlekanovMoskalets,BlekanovSergeevMaksimovBOWTIE}. Окружение таких сайтов может составлять десятки – сотни тысяч страниц. Стандартными методами кластеризации такого объема внешних веб-ресурсов не обойтись, так как при работе с большими коллекциями документов многие из данных методов показывают крайне неудовлетворительные результаты в плане производительности \cite{EneImMoseley} (например, метод Single linkage позволяет создавать кластеры произвольной формы, однако имеет высокую трудоемкость -- \(O(n^2)\), где \(n\) -- число документов). Для того чтобы избавиться от «проклятия» размерности применяется вероятностный метод понижения размерности многомерных данных Locality-Sensitive Hashing (в дальнейшем LSH), основная идея которого состоит в подборе хэш-функций для некоторых измерений для того, чтобы похожие объекты попадали в одну корзину \cite{Buhler}.

\subsubsection{Уменьшение размерности для анализа больших коллекций текстовых документов}

Для преобразования текстовых документов в числовые множества в данной работе использовался метод Shingling \cite{Broder}, который разбивает каждый документ на небольшие множества по \(k\) слов в каждом. Далее, для сравнения документов, применялась технология Min Hashing, которая позволяет быстро сравнивать множества, содержащие большое число элементов.

\textit{Мера Жаккара.} В работе для определения похожести двух текстов применяется мера Жаккара, которая похожесть множеств определяет отношением числа элементов, входящих в пересечение двух множеств к числу элементов, входящих в объединение этих множеств \cite{SingthongchaiNiwattanakul}:
\[
\textit{Sim}(C_1, C_2) = \frac{\lvert C_1 \cap C_2\rvert}{\lvert C_1 \cup C_2\rvert}
\]

Так как каждый документ преобразуется во множества, состоящие из \(k\) слов, то набор документов можно представить в виде сильно разреженной булевой матрицы, где столбцы представляют собой документы, а строки -- элементы универсального множества (например, множество элементов, где каждый элемент представляет собой множество \(k\) слов). Элемент такой матрицы равен единице, если документ (столбец) содержит данное множество k-слов. В противном случае элемент матрицы равен нулю.

\textit{Minhashing.} Для большой коллекции документов булева матрица будет сильно разрежена. Соответственно, матрица, описывающая такую коллекцию, займет много места в памяти. Кроме того, дальнейшая ее обработка так же займет большое количество времени. Для того чтобы решить возникшие проблемы, данную матрицу преобразуем в матрицу, хранящую определенное количество хэш-функций и информацию о сходстве между похожими документами.

\textit{Minhash}-функция \(h(C)\) -- это номер первой строки для столбца \(C\) в булевой матрице, где строки перемешаны случайным образом \cite{ChumPerdochMatas}.

Как видим, число случайных перестановок задает число Minhash-функций. Например, можно использовать сто случайных перестановок для создания ста сигнатур для каждого столбца матрицы.

Сигнатуры могут быть записаны в другой матрице сигнатур (Signature matrix), чьи колонки представляют собой документы, а строки -- Minhash-значения.

Чем больше хэш-функций, тем выше вероятность того, что \(\textit{Sim}(C_1, C_2) = \textit{Sim}(M[C_1], M[C_2])\), где \(C_i\) -- столбец булевой матрицы, \(M[C_i]\) -- столбец матрицы сигнатур. Это важное свойство позволяет преобразовывать булевы матрицы с большим количеством строк в небольшие матрицы сигнатур, сохраняя сходство похожих множеств. Соответственно, похожесть двух столбцов определяется долей строк, в которых они равны.

Таким образом, каждый документ может быть представлен в виде вектора, число элементов которого равно количеству Minhash-функций.

\textit{Locality-Sensitive Hashing (LSH).} Основная идея: сгенерировать из большого множества документов маленькие списки пар документов, чья схожесть должна быть посчитана. Для сравнения двух документов (столбцов) устанавливается порог \(t\) \((t < 1)\). Пара документов считается похожей только в том случае, если доля одинаковых значений в матрице сигнатур больше \(t\). Для матрицы сигнатур необходимо несколько раз вычислить хэш-значения, и поместить документы с одинаковым значением в одну корзину (\textit{bucket}). Документы, которые хоть раз попали в одну корзину, будут рассмотрены как кандидаты на сравнение \cite{GionisIndykMotwani}.

Для множественного подсчета хэш-функций столбцов, необходимо разбить матрицу сигнатур \textit{M} на \(b\) частей по \(r\) строк в каждой. Для каждого \(b\) подсчитать хэш-значения столбцов и поместить столбцы с равным значением в одну корзину. Кандидатами на сравнение будут те столбцы, которые хоть раз попали в одну корзину. Для правильной работы алгоритма необходимо настроить \(r\) и \(b\) таким образом, чтобы похожие документы попадали в одну корзину, а непохожие -- в разные.

\subsubsection{Иерархическая кластеризация с использованием LSH}

Данный алгоритм кластеризации использует хэш-таблицы, которые были сформированы в результате LSH \cite{KogaIshibashiWatanabe}. Данный алгоритм кластеризации с высокой долей вероятности создает такие же кластеры, как и в методе Single linkage. Ниже представлено детальное описание алгоритма.

\textit{Предварительные условия:}
\begin{itemize}
	\item \(t < 1\) -- порог, задающий минимальную схожесть документов;
	\item \(r = 1\) -- начальное значение строк матрицы сигнатур в каждой группе;
	\item \(r_\textit{min}\) -- минимальное значение строк матрицы сигнатур в каждой группе;
	\item \(\Delta\) -- коэффициент уменьшения параметра \(r\).
\end{itemize}

Каждый документ представляет собой отдельный кластер.

\textit{Шаг 1.} Для каждой группы \(b\) в каждом столбце вычислить хэш-функции, и сохранить столбцы, у которых хэш-значения хотя бы раз попали в одну корзину, при условии, что в корзине должны находиться столбцы, принадлежащие разным кластерам. Если, например, какие то два столбца принадлежат одному кластеру, то один из случайно выбранных столбцов удаляется из корзины.

\textit{Шаг 2.} Для каждого из столбцов, входящих в одну корзину, отобрать пары кластеров, расстояние между которыми больше \(t\).

\textit{Шаг 3.} Пары кластеров, соответствующие парам, полученным на шаге 2, объединяются в один.

\textit{Шаг 4.} Если \(r \leq r_\textit{min}\), то алгоритм прекращает работу. Иначе, переход на шаг 5.

\textit{Шаг 5.} \(r = r - \Delta\). Переход на шаг 1.

\textit{Оценка качества кластеризации на основе LSH.} Для оценки качества модифицированного метода агломеративной кластеризации с использованием LSH в данной работе использовалась тестовая коллекция текстовых документов Reuters-21578, содержащая \textit{21 578} документов. На данной коллекции оценивались два метода иерархической кластеризации: стандартный алгоритм иерархической кластеризации single-link и модифицированный алгоритм single-link, основанный на применении LSH. В качества основных метрик качества в работе использовались точность (\(R\)) полнота (\(P\)), аккуратность (\textit{Acc}) и \(F\)-мера.

В табл.~\cref{tab:clasterizationReuters} приведены результаты оценки качества иерархической кластеризации для каждого из алгоритмов, а так же время работы каждого алгоритма при кластеризации \textit{1000}, \textit{10 000} и \textit{20 000} документов.

\begin{table}[ht]%
	\caption{Результаты оценки качества кластеризации на коллекции Reuters.}%
	\label{tab:clasterizationReuters}% label всегда желательно идти после caption
	\renewcommand{\arraystretch}{1.6}%% Увеличение расстояния между рядами, для улучшения восприятия.
	\def\tabularxcolumn#1{m{#1}}
	\begin{tabularx}{\textwidth}{@{}>{\raggedright}X >{\centering}m{3.2cm} >{\centering}m{1.5cm} >{\centering}m{1.5cm} >{\centering}m{1.5cm} >{\centering}m{2.3cm} >{\centering\arraybackslash}m{2cm}@{}}% Вертикальные полосы не используются принципиально, как и лишние горизонтальные (допускается по ГОСТ 2.105 пункт 4.4.5) % @{} позволяет прижиматься к краям
		\toprule     %%% верхняя линейка
		Алгоритм & Количество документов & \textit{Acc}, \% & \(R\), \% & \(P\), \% & \(F\)-мера, \% & Время работы\\
		\midrule %%% тонкий разделитель. Отделяет названия столбцов. Обязателен по ГОСТ 2.105 пункт 4.4.5
		 & 1000 & 75 & 79 & 60 & 68 & 4 с \\
		 Single-Link & 10 000 & 79 & 82 & 62 & 71 & 410 с \\
		 & 20 000 & -- & -- & -- & -- & > 1 ч \\
		 \midrule
		 & 1000 & 72 & 81 & 66 & 73 & 3 с \\
		 Single-Link + LSH & 10 000 & 72 & 80 & 69 & 74 & 41 с \\
		 & 20 000 & 78 & 83 & 64 & 72 & 90 с \\
		\bottomrule %%% нижняя линейка
	\end{tabularx}%
\end{table}

Как видно из таблицы, аккуратность и полнота метода Single-Link + LSH близки по значениям к аккуратности и полноте метода Single-Link. Точность же разработанного метода иногда превосходит точность метода Single-Link. \(F\)-меры у обоих методов примерно одинаковые.

На рис.~\cref{fig:clasterizationTime} приведена зависимость продолжительности работы алгоритмов кластеризации от числа входных документов.

\begin{figure}[ht]
	\centerfloat{
		\includegraphics[scale=0.7]{clasterizationTime}
	}
	\caption{Время работы алгоритмов кластеризации.}\label{fig:clasterizationTime}
\end{figure}

Графики показываю, что алгоритм кластеризации с использованием LSH работает гораздо быстрее, чем алгоритм без использования LSH. С увеличением числа документов время работы алгоритма с использованием LSH линейно растет.

\subsubsection{Эксперимент. Исследование сайтов университетов методом кластеризации на основе LSH}
\paragraph{Постановка эксперимента.} В эксперименте требовалось для заданного списка сайтов университетов России, США и Великобритании, занимающих по своим регионам ведущие позиции в вебометрическом рейтинге \cite{RankingWeb}, с помощью специализированного поискового робота \cite{BlekanovSergeevMartynenko} и базы данных Majestic \cite{Majestic} получить списки и содержимое всех внешних веб-страниц, которые их цитируют. А также с помощью разработанного авторами метода агломеративной кластеризации на основе алгоритма LSH выявить целевые группы найденных внешних веб-ресурсов с одинаковым родом деятельности и степень их влияния на вебометрический рейтинг сайтов исследуемых ВУЗов.

Для исследования выбраны следующие сайты ВУЗов:
\begin{itemize}
	\item Московский государственный университет им. М.В. Ломоносова (msu.ru), Россия;
	\item Санкт-Петербургский государственный университет (spbu.ru), Россия;
	\item Новосибирский государственный университет (nsu.ru), Россия;
	\item Массачусетский технологический институт (mit.edu), США;
	\item Гарвардский университет (harvard.edu), США;
	\item Стэнфордский университет (stanford.edu), США;
	\item Кембриджский университет (cam.ac.uk), Великобритания;
	\item Оксфордский университет (ox.ac.uk), Великобритания;
	\item Университетский колледж Лондона (ucl.ac.uk), Великобритания.
\end{itemize}

\paragraph{Результаты сбора.} В итоге сбора данных для анализа внешних веб-ресурсов сайтов исследуемых университетов были получены следующие результаты, представленные в табл.~\cref{tab:externalResources}.

\begin{table}[ht]%
	\caption{Количество внешних ресурсов, окружающих университеты.}%
	\label{tab:externalResources}% label всегда желательно идти после caption
	\renewcommand{\arraystretch}{1.6}%% Увеличение расстояния между рядами, для улучшения восприятия.
	\def\tabularxcolumn#1{m{#1}}
	\begin{tabularx}{\textwidth}{@{}>{\centering}X  >{\centering}m{2.6cm} >{\centering}m{2.6cm} >{\centering}m{2.6cm} >{\centering}m{2.6cm} >{\centering\arraybackslash}m{2.6cm}@{}}% Вертикальные полосы не используются принципиально, как и лишние горизонтальные (допускается по ГОСТ 2.105 пункт 4.4.5) % @{} позволяет прижиматься к краям
		\toprule     %%% верхняя линейка
		URL-адрес ВУЗа & Позиция в рейтинге Webometrics &  Внешние ссылки на сайте ВУЗа & Внешние домены на сайте ВУЗа & Количество цитирующих сайт ВУЗа веб-страниц & Количество цитирующих сайт ВУЗа доменов \\
		\midrule %%% тонкий разделитель. Отделяет названия столбцов. Обязателен по ГОСТ 2.105 пункт 4.4.5
		spbu.ru & 539 & 3 280 & 469 & 1 599 059 & 16 028 \\
		msu.ru & 129 & 817 & 262 & 7 039 127 & 39 416 \\
		nsu.ru & 616 & 5 747 & 921 & 926 004 & 15 186 \\
		harvard.edu & 1 & 918 & 163 & 75 994 723 & 319 445 \\
		stanford.edu & 2 & 175 & 79 & 29 551 130 & 311 148 \\
		mit.edu & 3 & 9 & 6 & 41 271 678 & 324 989 \\
		ox.ac.uk & 16 & 8 075 & 1 631 & 8 920 524 & 117 959 \\
		cam.ac.uk & 15 & 3 385 & 1 174 & 12 084 107 & 120 796 \\
		ucl.ac.uk & 24 & 25 218 & 5 683 & 4 733 566 & 66 035 \\
		\bottomrule %%% нижняя линейка
	\end{tabularx}%
\end{table}

Таблица показывает, что университеты США, занимающие первые три позиции вебометрического рейтинга, имеют значительно больше внешних ресурсов, ссылающихся на них, чем ведущие университеты России и Великобритании. Ведущие университеты Великобритании также опережают российские по общему количеству ссылок и доменов, ссылающихся на них. Количество внешних ресурсов напрямую влияет на такой вебометрический индикатор, как Impact ресурса в Вебе.

На показатель видимости ресурса в Вебе, помимо числа внешних ссылающихся ресурсов, также влияет и их качество. Для кластеризации внешних Веб-ресурсов были выбраны только англоязычные и русскоязычные страницы Веба. В табл.~\cref{tab:externalResourcesClasterization} приведено количество доменов для каждого университета, к которым применялся метод агломеративной кластеризации с использованием LSH.

\begin{table}[ht]%
	\centering
	\caption{Количество внешних ресурсов для кластеризации.}%
	\label{tab:externalResourcesClasterization}% label всегда желательно идти после caption
		\begin{tabular}{| c | c | c | c | c |}% Вертикальные полосы не используются принципиально, как и лишние горизонтальные (допускается по ГОСТ 2.105 пункт 4.4.5) % @{} позволяет прижиматься к краям
			\hline
			\makecell{\\Университет} & \multicolumn{2}{c|}{\makecell{Внешние домены на \\ сайте ВУЗа}} &  \multicolumn{2}{c|}{\makecell{Цитирующие сайт ВУЗа \\ домены}}\\
			\cline{2-5}
			& на русском & на английском & на русском & на английском\\
			\hline
			spbu.ru & 342 & 81 & 7 373 & 3 929  \\
			\hline
			msu.ru & 177 & 62 & 3 343 & 1 008 \\
			\hline
			nsu.ru & 535 & 242 & 2 508 & 1 428  \\
			\hline
			harvard.edu & 4 & 151 & 1 545 & 16 632  \\
			\hline
			stanford.edu & 5 & 65 & 409 & 6 088  \\
			\hline
			mit.edu & 3 & 3 & 1 625 & 24 930  \\
			\hline
			ox.ac.uk & 29 & 1 399 & 126 & 3 187  \\
			\hline
			cam.ac.uk & 91 & 4 668 & 671 & 9 196 \\
			\hline
			ucl.ac.uk & 14 & 1 054 & 162 & 3 112  \\
			\hline
		\end{tabular}%
\end{table}

Ниже приведены результаты кластеризации внешних ресурсов. Для определения их тематики извлекались наиболее часто встречающиеся множества слов (shingles), которые присутствовали в кластере. В результате анализа были выделены основные часто встречающиеся группы: университеты, научные сообщества, поисковые системы, социальные сети (включая различные блоги, форумы и т.д.), медиасфера, новостные порталы и сайты других организаций. Таким образом, определялась численность каждой группы для университета.

\paragraph{Результаты кластеризации внешних доменов на сайте ВУЗа.} Внешние домены -- это домены, на которые ссылаются исследуемые сайты университетов. Ввиду ограниченного числа внешних ссылок на сайтах университетов США, были рассмотрены только внешние ресурсы университетов России и Великобритании. Были получены кластеры для каждого из исследуемых университетов. Ниже представлена сравнительная гистограмма внешнего окружения университетов России и Великобритании (рис.~\cref{fig:histogramUKRU}).

\begin{figure}[ht]
	\centerfloat{
		\includegraphics[scale=0.7]{histogramUKRU}
	}
	\caption{Сравнительная гистограмма университетов России и Великобритании.}\label{fig:histogramUKRU}
\end{figure}

Рисунок показывает, что внешнее ресурсы, на которые ссылаются сайты университетов, практически совпадают. Университеты и научные сообщества составляют значительную долю внешних ресурсов, на которые ссылаются сайты университетов России и Великобритании.

\paragraph{Результаты кластеризации цитирующих сайт ВУЗа доменов.} Цитирующие веб-ресурсы -- это сайты, сгруппированные по доменам, которые ссылаются на исследуемые сайты университетов. Были получены кластеры для каждого из исследуемых университетов. Ниже приведена сравнительная гистограмма внешнего окружения университетов, представляющих страны России, США и Великобритании (рис.~\cref{fig:histogramUKRUUS}).

\begin{figure}[ht]
	\centerfloat{
		\includegraphics[scale=0.7]{histogramUKRUUS}
	}
	\caption{Сравнительная гистограмма ресурсов, ссылающихся на университеты России, США и Великобритании.}\label{fig:histogramUKRUUS}
\end{figure}

Анализ сайтов, ссылающихся на сайты ВУЗов, показал, что университеты из одной страны имеют схожее внешнее окружение во многих областях. У всех университетов преобладают такие группы, как Университеты и Научные сообщества. У университетов США и Великобритании эта доля выше, чем у России. Соответственно, в других областях доля российских университетов больше.

Можно сделать вывод о том, что чем выше доля университетов и научных сообществ среди внешних ресурсов университета, тем выше его вебометрический показатель Impact. Этот показатель зависит от количества и качества внешних ресурсов, ссылающихся на сайт. Внешние группы, содержащие университеты и научные сообщества, имеют гораздо больший вес для вебометрического рейтинга, чем другие группы внешних ресурсов.

\paragraph{Общие выводы.} Разработан алгоритм иерархической кластеризации с использованием вероятностного метода понижения размерности многомерных данных Locality-Sensitive Hashing, оптимально подходящий для кластеризации больших коллекций документов (massive datasets). Данный алгоритм был апробирован на тестовой коллекции текстовых документов Reuters. На тестовой коллекции алгоритм показал приемлемую точность (accuracy) и \(F\)-measure в сравнении с классическими методами иерархической кластеризации. Метод кластеризации с использованием LSH значительно превосходит по скорости работы классические методы иерархической кластеризации.

Разработанный алгоритм был использован для анализа веб-пространства нескольких университетов с их окружением.

Получены данные о внешнем окружении веб-пространства университетов России и Великобритании и о ресурсах, ссылающихся на сайты университетов России, США и Великобритании.

\subsection{Теоретико-графовые характеристики в вебометрических исследованиях внутренней топологии крупных сегментов Веба}\label{subsec:ch1/sec4/sub8}

\paragraph{Введение.} В настоящее время стремительно развивается молодое научное направление вебометрика \cite{AlmindIngwersen,Thelwall,Pechnikov}, которая занимается исследованиями количественных и качественных характеристик топологии (гиперссылочной структуры) различных веб-сегментов. Основоположниками таких исследований являются испанская группа Cybermetrics Lab, которая разработала вебометрический рейтинг сайтов различных крупных организаций \cite{RankingWeb} (таких как научно-образовательные учреждения, больницы, научно-исследовательские центры и т. п.).  

Одной из актуальных задач вебометрики является задача анализа внутренней топологии различных сайтов \cite{Thelwall,BlekanovSergeevMaksimov,MaksimovBlekanov,BlekanovSergeevMaksimovBOWTIE}, а также выявление критериев оценки качества веб-сегментов и сравнение этих сегментов по полученным оценкам.  

Для решения вышеуказанной задачи в статье авторами рассмотрены известные теоретико-графовые характеристики веб-графов и разработан комплекс программ на их основе. Данный программный комплекс используется для вычисления указанных выше характеристик, их сравнения на разных веб-сегментах большой размерности, а также для визуализации результатов сравнительного анализа.  

\paragraph{Эксперимент.} В работе ставился эксперимент, в котором разработанный авторами комплекс программ для выявления значимых страниц внутренней топологии крупных веб-сегментов (с помощью используемых теоретико-графовых характеристик) и их визуализации апробировался на крупных сайтах университетского Веба. 

Получение внутренней топологии веб-ресурсов в эксперименте выполнялось с помощью программно-аналитического комплекса для вебометрических исследований, основанного на обобщенном ядре поискового робота \cite{BlekanovSergeevMartynenko} и успешно апробированного в исследованиях \cite{BlekanovSergeevMaksimov,MaksimovBlekanov,BlekanovSergeevMaksimovBOWTIE}.

В качестве теоретико-графовых характеристик сайтов были выбраны несколько мер центральности и связности веб-графов, а именно:
\begin{enumerate}
	\item Степень вершины (Degree Centrality) -- показатель, указывающий для каждой страницы количество страниц, связанных с ней. В веб-графе вычисляются полустепень захода (indegree, количество входящих ссылок) и полустепень исхода (outdegree, количество исходящих ссылок) ссылок \cite{WassermanFaust,OrtegaAguillo}.
		\item Мера центральности Betweenneess -- показатель, указывающий насколько часто данная веб-страница лежит на кратчайшем пути между всеми парами страниц сайта \cite{WassermanFaust,OrtegaAguillo}.
		\item Мера центральности Closeness -- среднее расстояние от данной веб-страницы до всех остальных страниц сайта \cite{WassermanFaust}.
		\item Мера центральности PageRank -- показатель важности страницы. Чем выше ее показатель, тем она важнее \cite{PageBrinMotwani}.
		\item Мера связности p-Cliques -- подграф, который представляет собой полный граф. p -- количество вершин в данном подграфе \cite{OrtegaAguillo}.
		\item Мера связности k-Cores -- подграф, в котором каждая страница связана по крайне мере с k другими страницами в этом подграфе \cite{OrtegaAguillo}.
\end{enumerate}

А также использовались такие общие показатели топологии, как:
\begin{enumerate}
	\item Расстояние (Distance) -- средняя длина всех кратчайших путей в графе \cite{MaksimovBlekanov}. 
	\item Диаметр (Diameter) -- длина самого большого кратчайшего пути в графе \cite{MaksimovBlekanov}.
\end{enumerate}

В качестве сайтов университетского Веба были взяты следующие сайты из Мирового вебометрического рейтинга университетов \cite{RankingWeb}:
\begin{enumerate}
	\item Сайт, занимающий первое место в общем рейтинге:
		\begin{enumerate}
				\item Сайт Гарвардского университета -- ГУ (www.harvard.edu). 
			\end{enumerate}
	\item Сайты, занимающие первые места в рейтинге по Российской Федерации: 
		\begin{enumerate}
				\item Сайт Московского государственного университета -- МГУ (www.msu.ru). 
				\item Сайт Санкт-Петербургского государственного университета -- СПбГУ (www.spbu.ru).
			\end{enumerate}
\end{enumerate}

При получении внутренней топологии выше указанных сайтов учитывались только главные их домены (поддомены не рассматривались). 

Используя вышеуказанный программный комплекс, требовалось получить визуальное представление топологии сайтов университетов в виде композиции наилучших значимых веб-страниц по всем мерам центральности (по каждой мере выбирался топ-10 страниц с наилучшими весами), а также построить и оценить расстояние от главной страницы до этих страниц.

\paragraph{Результаты эксперимента.} В ходе эксперимента для выбранных сайтов университетского Веба были получены следующие значения мер связности и общих показателей топологии (табл.~\cref{tab:uniPagesInnerTopology}):

\begin{table} [htbp]%
	\centering
	\caption{Теоретико-графовые характеристики внутренней топологии университетских сайтов.}%
	\label{tab:uniPagesInnerTopology}% label всегда желательно идти после caption
	\renewcommand{\arraystretch}{1.5}%% Увеличение расстояния между рядами, для улучшения восприятия.
	\begin{SingleSpace}
		\begin{tabulary}{\textwidth}{@{}>{\zz}L >{\zz}C >{\zz}C >{\zz}C@{}} %Вертикальные полосы не используются принципиально, как и лишние горизонтальные (допускается по ГОСТ 2.105 пункт 4.4.5) % @{} позволяет прижиматься к краям
			\toprule     %%% верхняя линейка
			Теоретико-графовые характеристики & ГУ & МГУ & СПбГУ\\
			\midrule %%% тонкий разделитель. Отделяет названия столбцов. Обязателен по ГОСТ 2.105 пункт 4.4.5
			Количество страниц & 451 & 45 557 & 51 945 \\				
			Количество ссылок, принадлежащих главному домену & 27 615 & 2 043 221 & 4 862 750 \\
			Общее количество ссылок & 72 355 & 2 143 474 & 5 017 815 \\
			Расстояние & 3.04 & 7.19 & 6.02 \\
			Расстояние & 7 & 30 & 18 \\
			Среднее значение indegree, outdegree & 13.59 & 28.1 & 69.79 \\
			Количество вершин в p-Cliques & 26 & 10 & 144 \\
			\bottomrule %%% нижняя линейка
		\end{tabulary}%
	\end{SingleSpace}
\end{table}

Композиция всех наилучших (по заданным мерам центральности) значимых веб-страниц исследуемых сайтов представлена на рис.~\cref{fig:harvardUComposition}-\cref{fig:spbUComposition}. На рис.~\cref{fig:harvardUComposition}-\cref{fig:spbUComposition} введены следующие обозначения цветов: 
\begin{itemize}
	\item Красный -- главная страница сайта; 
	\item  Зеленый -- полустепень исхода; 
	\item Бирюзовый -- полустепень захода; 
	\item Синий -- мера Betweenness; 
	\item Розовый -- мера Closeness; 
	\item Оранжевый -- мера PageRank.
\end{itemize}

\begin{figure}[ht]
	\centerfloat{
		\includegraphics[scale=0.48]{harvardUComposition}
	}
	\caption{Композиция всех наилучших значимых веб-страниц сайта ГУ.}\label{fig:harvardUComposition}
\end{figure}

\begin{figure}[ht]
	\centerfloat{
		\includegraphics[scale=0.5]{moscowSUComposition}
	}
	\caption{Композиция всех наилучших значимых веб-страниц сайта МГУ.}\label{fig:moscowSUComposition}
\end{figure}

\begin{figure}[ht]
	\centerfloat{
		\includegraphics[scale=0.46]{spbUComposition}
	}
	\caption{Композиция всех наилучших значимых веб-страниц сайта СПбГУ.}\label{fig:spbUComposition}
\end{figure}

\paragraph{Общие выводы.} По результатам эксперимента можно сделать следующие выводы:
\begin{enumerate}
	\item Сайт ГУ представляет собой навигационный сайт, перенеся большую часть информации на свои поддомены (ссылки главного домена составляют \textit{38.2\%} от общего числа ссылок). Данное явление подтверждается тем, что веб-страницы по полустепени исхода у ГУ лежат близко к главной странице (глубина 2-4), в то время как у МГУ и СПбГУ они лежат на определенном удалении (8-11 у МГУ и 6-9 у СПбГУ).
	\item Показатели СПбГУ (такие как расстояние, диаметр, количество вершин в p-Cliques, положение значимых вершин относительно начальной страницы) лучше, чем таковые у МГУ. Это значит, что сайт СПбГУ показывает более высокую связность, нежели сайт МГУ.
	\item Авторитетные страницы расположены на следующем уровне глубины после главной страницы, о чем говорит их высокий вес PageRank.
	\item Значимые веб-страницы по мере Betweenness у всех сайтов лежат в радиусе пяти ссылок от главной страницы, однако у сайта МГУ они расположены последовательно. Это объясняется тем, что множество кратчайших путей проходит через эти страницы. При отказе одной из таких страниц увеличиваются размеры множества кратчайших путей, а также появляется риск полной потери связей с множеством страниц.
	\item Значимые веб-страницы по мере Closeness у ГУ лежат на уровнях 1 и 3, у МГУ -- на уровне 4 и 5, у СПбГУ -- на уровне 3. Это показывает, что рассмотренные сайты централизованы недалеко от главной страницы.
\end{enumerate}

Также эксперимент показал некоторую особенность главных страниц рассмотренных сайтов, а именно:
\begin{itemize}
	\item главная страница ГУ не попала ни в один из топ-10;
	\item главная страница МГУ попала лишь в топ-10 по степени захода;
	\item главная страница СПбГУ попала сразу в три топ-10: по степени захода, по мере Betweenness и по мере PageRank. 
\end{itemize}

В дальнейшем планируется расширить эксперимент, проверив эргономические параметры значимых веб-страниц, выделенных программным комплексом, у заданных сайтов. 

\FloatBarrier
           % Глава 1
\chapter{Web 2.0}\label{ch:ch2}

\subsubsection{2.1.1}

According to a number of scholars, the rise of Twitter publishing had to lead a new level of informational openness in social networks and to increase of press freedom. At the same time, some recent studies have shown significant criticism towards political efficacy of Twitter [7] (нужен источник). We analyze Twitter as a part of the hybridized media system [6] (нужен источник) suggesting after Adam and Pfetsch [1] (нужен источник) that political dimension of hybridization and spill-over effects between different media will be specific to the socio-political context. In terms of the media system modernization [15] (нужен источник) and hybridization context, Russian media system is the one bearing both unique and post-Communist-specific features. In our research, we aim at analyzing the role of Twitter in the process of agenda setting in the hybridized media system in Russia and finding out if the use of Twitter potentially leads to the formation of the ‘crossroads of opinions' or, in contrast, to encapsulation of political discussion and further fragmentation of public discourse.Our research focuses on structural and content aspects of discussion on anti-migrant bashings in Biryulyovo (Moscow) in October 2013. Methods of research include automated web crawling, framing analysis via semantic coding, and interpretation of descriptive statistics. Preliminary results suggest unexpectedly high level of mediatization of the discussion. The ‘crossroads' nature of the discussion showed up but in such a way that makes Twitter differ from elite communicative milieus like the Russian Facebook.

\subsubsection{2.4.1}

The community-based structure of communication on social networking sites has long been a focus of scholarly attention. However, the problem of discovery and description of hidden communities, including defining the proper level of user aggregation, remains an important problem not yet resolved. Studies of online communities have clear social implications, as they allow for assessment of preference-based user grouping and the detection of socially hazardous groups. The aim of this study is to comparatively assess the algorithms that effectively analyze large user networks and extract hidden user communities from them. The results we have obtained show the most suitable algorithms for Twitter datasets of different volumes (dozen thousands, hundred thousands, and millions of tweets). We show that the Infomap and Leiden algorithms provide for the best results overall, and we advise testing a combination of these algorithms for detecting discursive communities based on user traits or views. We also show that the generalized K-means algorithm does not apply to big datasets, while a range of other algorithms tend to prioritize the detection of just one big community instead of many that would mirror the reality better. For isolating overlapping communities, the GANXiS algorithm should be used, while OSLOM is not advised. 

\subsubsection{2.4.2}

Despite disputable possibility of extension of analysis of social relations on Twitter to real life, Twitter discussions are still being under attention of scholars studying structures and meanings of news- and issue-based ad-hoc public discourse. One of the socially relevant aspects of Twitter studies is that of influencers -- accounts that produce impact, either inside or outside Twitter. But there is still no agreement in the research community on how to defme and measure who is an influencer: either by 'absolute figures' or by network analysis metrics; this issue is even rarely discussed. Politically, today's mediatized public sphere where traditional media play the role of information hubs is highly uneven in terms of access to opinion expression; it privileges institutional players, including political elites, corporations, and media themselves. Hopes that Twitter would provide a more equal space for public deliberation are still not proven enough. Using web crawling and manual assessment of Twitter ad-hoc discussion on the Biryulyovo bashings of 2013, we show that users who post or even get commented most do not make it to the positions of most 'central' users by network metrics. We also demonstrate that usen that rank high by betweenness and pagerank centntrality form circles of reciprocal commenting that show the social cleavage wider than the discussion itself.

\subsubsection{2.4.3}

Today, a range of research approaches is used to define the so-called influencers in discussions in social media, and one can trace both conceptual and methodological differences in how influencers are defined and tracked. We distinguish between ‘marketing’ and ‘deliberative’ conceptualization of influencers and between metrics based on absolute figures and those from social network analytics; combining them leads to better understanding of user activity and connectivity measures in defining influential users. We add to the existing research by asking whether user activity necessarily leads to better connectivity and by what metrics in online ad hoc discussions, and try to compare the structure of influencers. To do this, we use comparable outbursts of discussions on inter-ethnic conflicts related to immigration. We collect Twitter data on violent conflicts between host and re-settled groups in Russia and Germany and look at top20 user lists by eight parameters of activity and connectivity to assess the structure of influencers in terms of pro/contra-migrant cleavages and institutional belonging. Our results show that, in both discussions, the number of users involved matters most for becoming an influencer by betweenness and pagerank centralities. Also, contrary to expectations, Russian top users all in all are, in general, more neutral, while Germans are more divided, but in both countries pro-migrant media oppose anti-migrant informal leaders.

\subsubsection{2.4.4}

Abstractive summarization is a technique that allows for extracting condensed meanings from long texts, with a variety of potential practical applications. Nonetheless, today’s abstractive summarization research is limited to testing the models on various types of data, which brings only marginal improvements and does not lead to massive practical employment of the method. In particular, abstractive summarization is not used for social media research, where it would be very useful for opinion and topic mining due to the complications that social media data create for other methods of textual analysis. Of all social media, Reddit is most frequently used for testing new neural models of text summarization on large-scale datasets in English, without further testing on real-world smaller-size data in various languages or from various other platforms. Moreover, for social media, summarizing pools of texts (one-author posts, comment threads, discussion cascades, etc.) may bring crucial results relevant for social studies, which have not yet been tested. However, the existing methods of abstractive summarization are not fine-tuned for social media data and have next-to-never been applied to data from platforms beyond Reddit, nor for comments or non-English user texts. We address these research gaps by fine-tuning the newest Transformer-based neural network models LongFormer and T5 and testing them against BART, and on real-world data from Reddit, with improvements of up to 2\%. Then, we apply the best model (fine-tuned T5) to pools of comments from Reddit and assess the similarity of post and comment summarizations. Further, to overcome the 500-token limitation of T5 for analyzing social media pools that are usually bigger, we apply LongFormer Large and T5 Large to pools of tweets from a large-scale discussion on the Charlie Hebdo massacre in three languages and prove that pool summarizations may be used for detecting micro-shifts in agendas of networked discussions. Our results show, however, that additional learning is definitely needed for German and French, as the results for these languages are non-satisfactory, and more fine-tuning is needed even in English for Twitter data. Thus, we show that a ‘one-for-all’ neural-network summarization model is still impossible to reach, while fine-tuning for platform affordances works well. We also show that fine-tuned T5 works best for small-scale social media data, but LongFormer is helpful for larger-scale pool summarizations.

\subsubsection{2.4.5}

\textit{Objectives.} Social media have become a place where the bulk of grassroots political discussion takes place. Today, the growing body of research is dedicated to cumulative patterns on online deliberation, the predecessors of which were the concept of the spiral of silence, the silent majority hypothesis, and influencer studies. However, when applied to the dissonant, disruptive, and discontinued online discussions of today where gatewatching is much less predictable and cumulation of support is often accompanied by communicative aggression, these concepts need to be reconsidered and re-tested. Also, the current public communication online is much more multi-level than before; even within one platform, several communication layers may be defined, and their inter-relations in terms of public opinion aggregation remain under-studied. \textit{Research goal.} This paper aims at discovering patterns of cumulative deliberation in online communication. We first discuss the umbrella concept of cumulative deliberation. Then, we test the dynamics of the discussion on Belarusian oppositional YouTube in terms of impact of cross-account commenting on growth of commenting within the cross-account community and the overall discussion. \textit{Method and sampling.} We have collected the data by YouTube crawling. The data include all user comments of 2018 for six salient Belarusian oppositional accounts. To define the cross-account commenters, we used Gephi-based web graph reconstruction. We manually coded user posts for interactivity, aggression, and criticism. Dependencies in the dynamics of the discussion were tested by correlational and cluster analysis. \textit{Results.} We have discovered that users diverged into two mutually exclusive modes of expression, namely aggressive-dialogical and (self-)critical. Of the cross-account commenters, several users demonstrated ‘cumulative’ behavior and personal opinion bubbles. While there were nearly no dependencies discovered in the dynamics of user posting, criticism and self-criticism show capacity of spurring/diminishing the dynamics of public communication online, even if commenting on the whole is cumulative, not dialogical.

\subsubsection{2.4.6}

Agendas in online media have become a scholarly focus nearly two decades ago, leading to shifting conceptualizations of what we see as agenda. Thus, agendas and agenda shifts \textit{inside} online discussions have shown its potential to influence offline deliberation, aggregate support, fuel protest, passing through and/or bypassing traditional media’s gatekeeping. Real-time (or nearly-real-time) learning about quick agenda movement inside globalized public debate might be particularly important for international organizations like UN or EU. However, we today lack both knowledge on how agendas move in such discussions and instruments on such analysis. In particular, we are next-to-unaware of to what extent globally relevant themes get contextualized within language-based discussion segments, as well as to what extent the latter depend on each other and lag behind each other in developing agendas and public opinion on quickly evolving issues or conflicts. In this paper, we propose a method of agenda detection based on neural-network text summarization and compare summaries of tweet packages across three languages within the Twitter hashtag \#jesuischarlie. We show that sentiment detection may allow for quality assessment of the text summaries, as compared to aggregated sentiment to the original tweets. We show that, outside France, agendas were more interpretational, abstract, and non-contextualized. The pattern of news changing to ‘issue outburt’ was simultaneous in dense discussion segments and lagged behind in a sparser one. We also show that, globally, main issues of the discussion may be spotted within the first hour.

\subsubsection{2.4.7}

With the emergence of discussion platforms like Twitter, the hopes rose that computermediated public sphere would become more even in access to discussion than massmediatized public sphere of the late 20th century. Scholars have argued that it will eventually form an ‘opinion crossroads’ where conflicts would be discussed by all the parties involved. But today, existing research provides mixed evidence on whether ordinary users, rather than mainstream media and institutional actors, can become influencers in discussions on current issues, e.g. relations between host and migrant communities. We focus on the Twitter discussion about an inter-ethnic conflict in Moscow’s Biryuliovo district in 2013, as well as the comparative ‘calm’ period in March 2014, and look at who were real influencers by reconstructing the discussion’s web graph, as well as analyzing and juxtaposing its metrics to figures indicating user activity. Our results show that ad hoc discussion differs dramatically from an issuebased one in terms of the influencer nature and composition; the role of active tweeting is questioned. We also show that nationalist accounts play a much bigger role than expected in both periods.

\subsubsection{2.5.1}

Конфликты с участием переселенцев «с глобального Юга» и их обсуждение в медиа и социальных сетях стали чертой современной публичной сферы в Европе и России. Сетевые дискуссии способны играть роль групп давления в формировании повестки обсуждения миграции и принятия политических решений, актуализируя как конфликтные, так и консенсусные общественные настроения. Важно знать, существуют ли закономерности в том, кого пользователи сети (в том числе иститутциональные) видят как виновников конфликтов и от кого ожидают его разрешения. Мы рассматриваем четыре конфликтные дискуссии в Твиттере в России и Германии 2013--2016 гг. и оцениваем связь статуса пользователя с паттернами присвоения вины и ответственности, а также сопоставляем дискурсивные стратегии обвинения в двух странах. Методология исследования включает веб-краулинг по заданному словарю, экспертную оценку аккаунтов пользователей, ручное кодирование твитов и описательную статистику. Итоги исследований говорят о том, что в целом обвинение направлено в первую очередь против трех групп: региональных властей, нацилнального лидера/правительства и иммигрантов; разрешения конфликтов пользователи ждут от федеральных властей, ответственных за изменение иммиграционной политики. Но если в Германии очевиден политический разлом между правыми (пронационалистическими) и левыми (феминистскими) стратегиями поиска виноватых, то в российском Твиттере доминирует обвинение и власти, и иммигрантов, а также распространены сниженные ожидания мирной политической разрядки межэтнической напряженности.

\subsubsection{2.5.2}

Recently, the growing role of social network users in content dissemination has brought to life the concept of secondary gatekeeping -- selection and republication of content already selected and published by traditional gatekeepers. Secondary gatekeeping is believed to be raising the media in-platform visibility, but it may also have negative effects such as adding to creation of echo chambers and deepening the gaps between conflicting views. Such studies are particularly relevant for emergencies or social conflicts where sharing relevant content may be crucial for lowering social unease. But till today the nature of secondary gatekeeping remains highly understudied. We have conducted a comparative study of three \textit{ad-hoc} Twitter discussions on heated ethnic/racial conflicts in the USA (Ferguson riots), Germany (Köln mass abuse), and Russia (Biryulyovo anti-migrant bashings) to assess the patterns of content sharing by active discussants. We used vocabulary-based web crawling and human coding of over 1,000 tweets in randomized samples. Our results show that, in all cases, there’s weak but significant correlation between the type of user and his/her attitude to minority with the attitudes expressed in content, while it is not always true that users prefer the same gatekeeper type, e.g. online or social media. As difference between individual users remains statistically significant, this may mean that the nature of heated \textit{ad-hoc} discussions facilitates formation of ‘individual-level filter bubbles’ in addition to bigger echo chambers.

\subsubsection{2.5.3}

With growth of Internet communication, hopes rose that online discussions will equalize ordinary users and institutional discussants. But what roles traditional media play in online discussions remains under-researched. We argue that mediatization of Twitter discourse is worth studying, as activity of registered media in online discussions may play a role in preserving the “offline” deliberative inequalities. To assess the roles of media in Twitter discussions, we look at two structural aspects of their presence: posting activity and users’ sharing of media content -- within heated \textit{ad hoc} Twitter discussions on inter-ethnic conflicts in Russia, the USA, and Germany. Our findings show huge national differences in mediatization patterns, including the roles of political media. But also we see that, in all the three cases, the discourse is shaped by “white majority media”, and outlets that would represent the oppressed minority are virtually absent in both posting and link sharing patterns.

\subsubsection{2.5.4}

Online communication platforms have rapidly become a substantial element of e-governance processes in Europe and beyond. Today, research has shown that, in cases of social unrest and/or emergency, political actors responsible for their resolution are able to efficiently use microblogging platforms (including Twitter) to promote the discourse of harmonization. But in today's Russia, where the growth of inter-ethnic conflicts between the re-settlers from the post-Soviet South (Central Asia and South Caucasus) and the host communities in cities and towns has coincided with the growth of online communication milieus and their radicalization, political actors as well as NGOs seem to play minor roles in online communication management, including the cases of social unrest. We explore two Twitter discussions on inter-ethnic conflicts in Moscow to describe the presence of political actors, their roles in conflict resolution, and the patterns of expectations of other users towards the politicians. We discover extremely low political participation, as well as the phenomenon of 'radical replacement' of the roles of political emergency managers by nationalist users.

\subsubsection{2.5.5}

Embeddedness of politicians and political organizations in a discussion defines its level of institutionalization and creates a public arena for collaboration between publics and institutional actors. Thus, testing whether traditional hierarchies (in terms of presence of politically institutionalized actors) show up in online discussions deserves scholarly research. Moreover, it is also important to see whether more democratic societies show patterns of public involvement of politically institutionalized users that would differ from those in more authoritarian contexts.

To assess the ‘influencer’ status of politically institutionalized actors on Twitter cross-culturally, we have selected conflictual Twitter discussions in Germany, the USA, and Russia, all based on violent inter-ethnic clashes. Using vocabulary-based web crawling, we collected data on them and formed samples of top users selected by four activity metrics and five network metrics, to assess the positions of political users in the top lists and correlations of user status with their top list ranks. To this, we added qualitative assessment of presence of political users in comparative perspective.

Our results show that, in all the cases, presence of political actors in online discussions is scarce; also, political actors tend to fail to link user groups or stay in the center of discussion. There is also meaningful divergence of Russia from the pattern that Germany and the USA show: while in these countries politicians gain user attention based on content, in Russia it is the status itself that matters, and political users tend to gain weight in the discussion structure despite low attention levels.

\subsubsection{2.5.6}

YouTube-based discussions are a growing area of academic attention. However, we still lack knowledge on whether YouTube provides for forming critical publics in countries with no established democratic tradition. To address this question, we study commenting to Belarusian oppositional YouTube blogs in advance of the major wave of Belarusian post-election protests of 2020. Based on the crawled data of the whole year of 2018 for six Belarusian political videoblogs, we define the structure of the commenters’ community, detect the core commenters, and assess their discourse for aggression, orientation of dialogue, direction of criticism, and antagonism/agonism. We show that, on Belarusian YouTube, the commenters represented a genuine adversarial self-critical public with cumulative patterns of solidarity formation and find markers of readiness for the protest spillover.

\section{Концептуальное назначение}\label{sec:ch2/sec1}

%\begin{figure}[ht]
%    \centerfloat{
%        \includegraphics[scale=0.27]{latex}
%    }
%    \caption{TeX.}\label{fig:latex}
%\end{figure}
%
%Для выравнивания изображения по-центру используется команда \verb+\centerfloat+, которая является во
%многом улучшенной версией встроенной команды \verb+\centering+.

\section{Принципы построения ресурсов}\label{sec:ch2/sect2}

%А это две картинки под общим номером и названием:
%\begin{figure}[ht]
%    \begin{minipage}[b][][b]{0.49\linewidth}\centering
%        \includegraphics[width=0.5\linewidth]{knuth1} \\ а)
%    \end{minipage}
%    \hfill
%    \begin{minipage}[b][][b]{0.49\linewidth}\centering
%        \includegraphics[width=0.5\linewidth]{knuth2} \\ б)
%    \end{minipage}
%    \caption{Очень длинная подпись к изображению,
%        на котором представлены две фотографии Дональда Кнута}
%    \label{fig:knuth}
%\end{figure}
%
%Те~же~две картинки под~общим номером и~названием,
%но с автоматизированной нумерацией подрисунков:
%\begin{figure}[ht]
%    \centerfloat{
%        \hfill
%        \subcaptionbox[List-of-Figures entry]{Первый подрисунок\label{fig:knuth_2-1}}{%
%            \includegraphics[width=0.25\linewidth]{knuth1}}
%        \hfill
%        \subcaptionbox{\label{fig:knuth_2-2}}{%
%            \includegraphics[width=0.25\linewidth]{knuth2}}
%        \hfill
%        \subcaptionbox{Третий подрисунок, подпись к которому
%            не~помещается на~одной строке}{%
%            \includegraphics[width=0.3\linewidth]{example-image-c}}
%        \hfill
%    }
%    \legend{Подрисуночный текст, описывающий обозначения, например. Согласно
%        ГОСТ 2.105, пункт 4.3.1, располагается перед наименованием рисунка.}
%    \caption[Этот текст попадает в названия рисунков в списке рисунков]{Очень
%        длинная подпись к второму изображению, на~котором представлены две
%        фотографии Дональда Кнута}\label{fig:knuth_2}
%\end{figure}
%
%На рисунке~\cref{fig:knuth_2-1} показан Дональд Кнут без головного убора.
%На рисунке~\cref{fig:knuth_2}\subcaptionref*{fig:knuth_2-2}
%показан Дональд Кнут в головном уборе.
%
%Возможно вставлять векторные картинки, рассчитываемые \LaTeX\ <<на~лету>>
%с~их~предварительной компиляцией. Надписи в таких рисунках будут выполнены
%тем же~шрифтом, который указан для документа в целом.
%На~рисунке~\cref{fig:tikz_example} на~странице~\pageref{fig:tikz_example}
%представлен пример схемы, рассчитываемой пакетом \verb|tikz| <<на~лету>>.
%Для ускорения компиляции, подобные рисунки могут быть <<кешированы>>, что
%определяется настройками в~\verb|common/setup.tex|.
%Причём имя предкомпилированного
%файла и~папка расположения таких файлов могут быть отдельно заданы,
%что удобно, если не~для подготовки диссертации,
%то~для подготовки научных публикаций.
%\begin{figure}[ht]
%    \centerfloat{
%        \ifdefmacro{\tikzsetnextfilename}{\tikzsetnextfilename{tikz_example_compiled}}{}% присваиваемое предкомпилированному pdf имя файла (не обязательно)
%        %\input{Dissertation/images/tikz_scheme.tikz}
%
%    }
%    \legend{}
%    \caption[Пример \texttt{tikz} схемы]{Пример рисунка, рассчитываемого
%        \texttt{tikz}, который может быть предкомпилирован}\label{fig:tikz_example}
%\end{figure}
%
%Множество программ имеют либо встроенную возможность экспортировать векторную
%графику кодом \verb|tikz|, либо соответствующий пакет расширения.
%Например, в GeoGebra есть встроенный экспорт,
%для Inkscape есть пакет svg2tikz,
%для Python есть пакет tikzplotlib,
%для R есть пакет tikzdevice.

\section{Сбор данных}\label{sec:ch2/sec3}

%\noindent Нумерованный список:
%\begin{enumerate}
%    \item Первый пункт.
%    \item Второй пункт.
%    \item Третий пункт.
%\end{enumerate}
%
%\noindent Маркированный список:
%\begin{itemize}
%    \item Первый пункт.
%    \item Второй пункт.
%    \item Третий пункт.
%\end{itemize}
%
%\noindent Вложенные списки:
%\begin{itemize}
%    \item Имеется маркированный список.
%          \begin{enumerate}
%              \item В нём лежит нумерованный список,
%              \item в котором
%                    \begin{itemize}
%                        \item лежит ещё один маркированный список.
%                    \end{itemize}
%          \end{enumerate}
%\end{itemize}
%
%\noindent Нумерованные вложенные списки:
%\begin{enumerate}
%    \item Первый пункт.
%    \item Второй пункт.
%    \item Вообще, по ГОСТ 2.105 первый уровень нумерации
%          (при необходимости ссылки в тексте документа на одно из перечислений)
%          идёт буквами русского или латинского алфавитов,
%          а второй "--- цифрами со~скобками.
%          Здесь отходим от ГОСТ.
%          \begin{enumerate}
%              \item в нём лежит нумерованный список,
%              \item в котором
%                    \begin{enumerate}
%                        \item ещё один нумерованный список,
%                        \item третий уровень нумерации не нормирован ГОСТ 2.105;
%                        \item обращаем внимание на строчность букв,
%                        \item в этом списке
%                              \begin{itemize}
%                                  \item лежит ещё один маркированный список.
%                              \end{itemize}
%                    \end{enumerate}
%
%          \end{enumerate}
%
%    \item Четвёртый пункт.
%\end{enumerate}

\section{Анализ пользовательских Web-графов}\label{sec:ch2/sec4}

\subsection{Detection of Hidden Communities in Twitter Discussions of Varying Volumes}\label{subsec:ch2/sec4/sub1}

\subsubsection{1. Introduction}

With the growth of the Internet, its complexity has also grown substantially. In particular, social networking has provided for formation of user clusters corresponding to an increasingly big number of communities at various levels, up to the cross-platform one. As previous research states, the formation of communities is based on user traits (e.g., group belonging, identity, spoken language, or location), discussion patterns (e.g., liking, commenting, mentioning, or reposting), and content features (e.g., political views, values, or cultural differences). Communities may also be situational, or ad hoc \cite{BrunsBurgess}, such as hashtag-based ones, or calculated, created by special curation \cite{BrunsBurgess2015}; in such cases, their borders are seen well enough. However, there are cases when hidden communities gradually accumulate online, e.g., within hashtagged discussions, and remain latent for non-involved observers; or, they are hidden on purpose, such as those created by terrorist groups. In any case, the detection of latent communities may serve not only as a metric for assessing the complexity of online discussions but also as a measure that can add to resolving burning social issues of security, political polarization, or cultural intolerance.

Despite the growing number of works dedicated to detection of intentionally hidden and/or latent self-cumulative communities on social media \cite{PerligerPedahzur,WangGongLiu}, a cumulation of communities in real-world discussions has so far mostly been an object of ‘one case/one model’ studies where tests on real data bear an illustrative function only, if at all. There is almost no comparative assessment of community detection models. Furthermore, the  and traits of discovered communities found by one model are not judged against those detected by other models, which, though, seems natural to do. Rare exceptions include works on several spectral algorithms for directed graphs \cite{VanLierdeDelvenneVanDooren} which consider clustering algorithms based on cyclic and acyclic patterns, as well as comparative evaluation of community detection algorithms \cite{GeorgeShujaeeKerwat} in which only two algorithms, Infomap and Louvain, are considered (also assessed in our paper -- see below). However, in these works, the authors do not assess the possibility of using methods for identifying communities on graphs of different sizes, especially large (around 100,000 nodes) and extra-large (more than 500,000 nodes), and also do not investigate user interactions in discussions. Another shortcoming of the current research on hidden communities is that the models are not compared on the datasets of varying size, despite the fact that not all the models work with large-scale data well enough.

Our research aims at partly covering the three aforementioned research gaps. The first gap is absence of comparative assessment of community detection models. The second one is selecting (or creating) the best model for real-world discussion data from social media. The third one is testing the models for the datasets belonging to one platform but significantly differing in size. Thus, the research design of this paper is the following: we test nine existing models of community detection in social media data on four datasets collected in 2013 to 2016 on Twitter by seven quality metrics that are considered precise in the previous literature. The old enough datasets are used, as they are virtually clean, in terms of botization and other types of computational propaganda, and are free from algorithmic impact upon the discussions cast by the platform itself.

The remaining part of this paper is organized as follows. Section 2 describes the difference between various clustering algorithms for networks seen as directed/undirected graphs. Section 3 provides the selection of community detection methods and metrics for their quality evaluation, while Section 4 tells of our research design. We demonstrate the results for four various datasets in Section 5 and discuss them in the Conclusion.

\subsubsection{2. Related Work}

\paragraph{2.1. Selecting Community Detection Algorithms} Community allocation methods can be divided into two categories.

The first category includes destiny-based clustering (DBSCAN, OPTICS, K-Means type) \cite{EsterKriegelSander,KriegelKrogerSander,Lloyd}. One of the advantages of these methods is the efficiency of the allocation of compacted clusters, but when applying these methods to graph analysis, the following problems arise:

\begin{enumerate}
	\item The need to use additional metrics to estimate the distance between vertices, and, in the case of oriented graphs, this choice is very narrow;
	\item The limitations of computing on large-dimensional graphs (with more than 50,000 nodes) (in particular, the speed of searching for new clusters is significantly reduced, and, on super-large graphs (with more than 500,000 nodes), this class of algorithms cannot get a result without additional optimization);
	\item Limited application to sparse data \cite{BodrunovaOrekhovBlekanov}.
\end{enumerate}

The second type includes graph methods for searching for hidden communities based on graph properties \cite{CauteruccioCorradiniTerracina} and the structure of connections of graph nodes \cite{RosvallAxelssonBergstrom,BlondelGuillaumeLambiotte,RaghavanAlbertKumara}. The advantages of these methods include the speed and the ability to calculate, including on large-scale graphs. However, when using such approaches, the number of hidden communities found in the graph can vary from several hundred to several thousand (although in many SNA tasks it is required to find up to several dozen). In addition, this class of algorithms is very sensitive to graph data. This is why a separate task within community detection in graphs is used to adapt existing methods to specific real data from the relevant subject area. In \cite{AgresteDeMeoFiumara}, the authors received high ratings of the efficiency of the Infomap \cite{RosvallAxelssonBergstrom} and LPA \cite{RaghavanAlbertKumara} algorithms on synthetic data and real web data of large size. In \cite{DengZhaiLv}, the authors compared their algorithm with the Infomap and OSLOM \cite{LancichinettiRadicchiRamasco} algorithms, where they obtained the best results in speed and floating quality results on billing network data. In \cite{AmatiAngeliniCruciani}, the authors used the Louvain \cite{BlondelGuillaumeLambiotte} algorithm as one of two algorithms for clustering large volumes of related data, and, in \cite{YuLiangJieJie,MotheMkhitaryanHaroutunian}, the authors also used graph algorithms for searching communities on real data, which showed high application efficiency. In studies \cite{DeitrickHu,ChenYinLi}, Infomap and SLPA (GANXiS) \cite{XieSzymanskiLiu} were also used in their work, but Infomap proved to be better.

As we can see from the above studies, there are few works dedicated to the analysis of user discussions on social networks. Individual tasks, such as topic modeling and sentiment analysis, are being investigated; however, user interaction within the discussion has not been investigated using these methods. Therefore, in this work, the main focus is on the application and testing of the most efficient methods for identifying hidden communities on real-data graphs of user discussions.

\paragraph{2.2. Communities in Networked Discussions: A Social Scientific Perspective} Testing algorithms and assessing their quality implies that they are tested either against known baseline results (by other methods or on marked-up data) or some pre-set ground truth. Establishing the latter requires a clear sociological task, e.g., for networked discussions, detection of politically polarized communities, sentiment-based grouping, and discursive or cultural divisions (as such a task would shape expectations on the number, size, and interaction patterns of the latent communities).

In this research, however, we cannot fix strict enough ground truth, as we compare the algorithms without setting a task. Yet, despite intentionally not setting the ground truth to compare to before testing, it is quite clear and commonsensical that social researchers who would like to use the algorithms for community detection usually aim at goals that imply a low number of communities. Thus, in political polarization studies or research on social group conflicts, the number of polarized views-based communities cannot be high \cite{BodrunovaBlekanovSmoliarova}. In general, studies of hidden communities are often necessary to detect the structure of conflict, such as (non-)binary political polarization, mutual hostility of social groups, opinion preferences towards burning issues on public agendas, recommender communities in marketing, or hidden chains of influencers \cite{BodrunovaBlekanovMaksimov}. In all these cases, it is expected that the discussion graphs are divided into a small number (2 to 5) of comparably large communities that represent social groups, political views, communities of practice, or center-periphery relations.

This is why, comparing the algorithms, we will not only look at the metrics calculated automatically (see below) but also judge qualitatively, by comparing the number and size of the detected communities, implying that algorithmic group detection should aim at finding a small number of robust communities.

\subsubsection{3. Models and Methods in SNA}

\paragraph{3.1. User Discussion Model} In this work, as a model of user discussion, we present a directed user web graph \(G\):

\begin{equation}
	\label{eqn:51}
	G_D = G(V,E),
\end{equation} where \(V = \{V_G,V_R, V_{GR}\}\) means all the users who took part in the discussion, \(V_G\) are the users who only generated content in the discussion \(D\), \(V_R\) are the users who only reacted to the content in discussion \(D\), and \(V_{GR}\) are the users who both generated content and reacted to other content. Furthermore, \(E\) refers to the connections that represent likes, comments, or reposts between the discussion participants. Visual representation of a networked discussion split by participation type according to \cref{eqn:51} is presented on Figure~\cref{fig:networkedDiscussionRepresentation}.

\begin{figure}[ht]
	\centerfloat{
		\includegraphics[scale=0.7]{networkedDiscussionRepresentation}
	}
	\caption{Visual representation of the structure of a networked discussion, by user type.}\label{fig:networkedDiscussionRepresentation}
\end{figure}

Let us assume that user \(A\) from \(V_G\) published a message on the social network on some event, and user \(B\) from \(V_R\) reacted to this post via a like, comment, or repost. Then there will be a connection \(B -- A\) (‘\(A\) caused by this publication a reaction from user \(B\)’) in an unweighted directed graph. 

Thus, in this paper, the search for communities in social graphs is the task of analyzing \cref{eqn:51} -- that is, finding sub-structures in a discursive community with three different participation strategies. The methods used for this analysis are discussed below.

\paragraph{3.2. Description of Community Detection Methods} This section will describe clustering techniques that have not yet been applied to \cref{eqn:51}.

\textit{3.2.1. The Directed Louvain Algorithm} Directed Louvain \cite{DuguePerez} is an extension for directed graphs of the greedy Louvain \cite{BlondelGuillaumeLambiotte} algorithm: starting from any set of vertices, the algorithm calculates the increase in modularity from moving vertices between communities. This increase is calculated using the following formula:
\[
	\Delta_{Q_d} = \frac{d^C_i}{m} - \left[\frac{d^{out}_i \cdot \sum_{tot}^{in} + d^{in}_i \sum_{tot}^{out}}{m^2}\right],
\] where \(d^C_i\) -- the degree of node \(i\) in community \(C\); \(d^{in}_i\) \((d^{out}_i)\) -- the indegree (outdegree) of node \(i\); \(sum^{in}_{tot}\) \((\sum_{tot}^{out})\) -- the incoming (outcoming) edges of community \(C\).

\textit{3.2.2. The Leiden Algorithm} The Leiden \cite{TraagWaltmanVanEck} algorithm is partially based on the smart local move algorithm, which can be seen as an improvement on the Louvain \cite{BlondelGuillaumeLambiotte} algorithm, and also uses the idea of speeding up local movement of nodes and the idea of moving nodes to random neighbors. The algorithm can use CPM or modularity as quality functions, and includes three stages:

\begin{enumerate}
	\item Local node movement;
	\item Improving partition;
	\item Enhanced partition-based network aggregation using a non-enhanced partition to
	create an initial partition for the aggregate network.
\end{enumerate}

\textit{3.2.3. The Directed Label Propagation Algorithm} The directed label propagation algorithm (DLPA) \cite{Li} is a development for directed graphs of the label propagation algorithm \cite{RaghavanAlbertKumara} method. LPA is one of the fastest community finding algorithms for undirected graphs. It can be used with large graphs, relies on topology, and is easy to implement:

\begin{enumerate}
	\item The algorithm assigns a unique label to each node;
	\item Each node selects a label among its neighbors based on the frequency of occurrence; 
	\item If the distribution of labels reaches a steady state, the algorithm stops, otherwise it returns to step 2;
\end{enumerate}

The principle of label selection looks like this:
\[
	l^{new}_v = \lvert N^l(v) \rvert,
\]
where \(v\) -- the node, \(l\) -- the node label, and \(N(v)\) -- neighbors of node \(v\).

DLPA differs from the original algorithm by the weighting rule for each edge:
\[
1 - \frac{E^{out}_S E^{in}_T}{k_S k_T}
\]
where \(E^{out}_S\) -- the outdegree of the source node, \(E^{in}_S\) -- the indegree of the target node, \(k_S\) -- the degree of the source node, and \(k_T\) -- the degree of the target node.

\textit{3.2.4. The Infomap Algorithm} The Infomap \cite{RosvallAxelssonBergstrom} algorithm works as follows: each node is assigned to a separate community. Then, in random sequential order, each node is moved to a neighboring com- munity, if this movement decreases the map equation value. The repetition is performed each time in a new random sequential order until no movement leads to a decrease in map equation. Then, the graph is rebuilt, and the communities of the last level are replaced by nodes at this level. Then, the procedure is repeated for this level. The network is rebuilt until the map equation result cannot be reduced any more.

\textit{3.2.5. The Generalized K-Means Algorithm} Generalized K-means using PageRank \cite{HajijSaidTodd} improves the k-means \cite{Lloyd} method, generalizing it to both directed and undirected graphs. Generalized k-means uses the PageRank algorithm as a measure of centrality and the Dijkstra’s algorithm as a metric for the distance between vertices for directed graphs.

Like the usual k-means, this algorithm consists of three stages:

\begin{enumerate}
	\item The initialization stage: k centroid vertices are randomly selected;
	\item The assignment stage: using Voronoi diagrams with centroids to divide the set of vertices into subsets;
	\item The update stage: building subgraphs, calculating PageRank for each subgraph, and updating the centroids.
\end{enumerate}

\textit{3.2.6. The Order Statistics Local Optimization Method} The order statistics local optimization method (OSLOM) \cite{LancichinettiRadicchiRamasco} is based on local optimization of the fitness function, which expresses the statistical significance of clusters in relation to random fluctuations, which is estimated using the extreme and order statistics tools. OSLOM can be used on its own or as a rework procedure for breaks/coverage provided by other methods. The method includes three phases:

\begin{enumerate}
	\item The search for significant clusters before convergence;
	\item Analysis of the resulting set of clusters to detect their internal structure or possible
	associations;
	\item Discovery of the hierarchical structure of clusters.
\end{enumerate}

\textit{3.2.7. The Speaker–Listener Propagation Algorithm} GANXiS, aka SLPA \cite{XieSzymanskiLiu}, is another extension of the LPA \cite{RaghavanAlbertKumara} algorithm that consists of three stages, as described below.

\begin{itemize}
	\item The memory of each node is initialized with the identifier of this node (unique label).
	\item The steps are repeated until the stopping criterion is met:
	\begin{enumerate}
		\item one node is selected as a listener;
		\item each neighbor of the selected node sends one tag following a certain conversation rule, e.g., choosing a random tag from its memory with a probability proportional to the frequency of occurrence of this tag in memory;
		\item the listener accepts a label from a collection of labels received from neighbors following a specific listening rule, e.g., choosing the most popular label from what it has observed at the current stage.
	\end{enumerate}
	\item Finally, post-processing based on in-memory labels of nodes is used to display communities.
\end{itemize}

SLPA utilizes an asynchronous update scheme, i.e., when updating a listener’s memory at time \(t\), some already-updated neighbors have memories of size \(t\) and some other neighbors still have memories of size \(t - 1\). SLPA reduces to LPA when the size of memory is limited to 1, and the stopping criterion is the convergence of all labels. Thus, SLPA is a joint version of the LPA algorithm that is suitable for detecting overlapping (fuzzy) clusters, while the disjoint version of the algorithm may be employed for finding non-fuzzy clusters and is used here for comparative purposes. We will name the joint algorithm GANXiSo, and the disjoint one GANXiSd.

\paragraph{3.3. Evaluation Metrics for Community Detection} In order to evaluate partitions, we decided to use various quality indicators that characterize how similar the structure of connections of a given network is to a community. The metrics are based on the idea that communities are collections of nodes with more connections inside and fewer connections outside. They were selected from a standard set of known metrics, which are commonly used to assess the graph properties of detected communities in a directed graph \cite{Rahman,LeichtNewman,Newman,Fagiolo,ChenNguyenSzymanski,KaurSinghKaushal}. Such properties, for example, include the ability to distinguish large communities in a graph, the density of internal connections within formed communities, the number of edges in a graph outside the community, and the ratio of incoming edges to the total number of community edges. In selecting the metrics, we also used the following logic: (1) the metrics must vary in terms of which graph properties they measure; (2) the metrics have to suit both the density-based and pattern-based clustering; and (3) taken together, the metrics must describe the main graph properties in terms of community detection. The metrics selected have no underlying ‘ground truth’ capacity; that is, none of them can tell whether a hidden community is found or not. However, taken together, the metrics provide for much better clarity on how well the algorithms cluster the graph nodes. The metrics, though, may provide for sociological meaning (e.g., stronger connection of nodes may reveal modularity based on opinion, group belonging, or expression sentiment), but the meaning of the result is task-dependent (e.g., we consider it better when a discussion is less modular, which would mean more equal spread of political opinion). It is not the peculiar sociological meaning that we demand from each metric; taken together, they are to provide for a bigger picture of how close a given algorithm is to finding a closer-to-life number of communities.

\begin{enumerate}
	\item The \textit{NEDindex} \cite{Rahman} value ranges from 0 to 1. Cluster nodes are more strongly con- nected to the entire graph if \textit{NEDindex} tends towards 1, and vice versa, cluster nodes are weakly connected if \textit{NEDindex} is close to 0. \[D(G) = 2 \sum_{1 \le i \le n}e_i,\] \[\textit{NED}(C) = \frac{\lvert V_c \rvert + \lvert E_c \rvert + D(C)}{\lvert V_c \rvert + \begin{pmatrix}
			\lvert V_c \rvert \\
			2 \\
	\end{pmatrix} + D(G,V_c)},\] \begin{equation} 
		\label{eqn:52} 
		\textit{NEDindex} = \sum_{1 \le i \le n} \frac{\textit{NED}(C_i) \cdot D(C_i)}{D(G)}. 
	\end{equation} This metric indicates the strength of the connection between the cluster vertices relative to the entire graph.
	
	\item Directed modularity \cite{LeichtNewman} is an extended version of the modularity metric \cite{Newman} for directed graphs. The higher the value, the better the result. \begin{equation}
		\label{eqn:53} 
		Q = \frac{1}{m} \sum_{ij}\left[A_{ij} - \frac{k^{in}_i k^{out}_j}{m}\right]\delta_{c_ic_j}
	\end{equation} Modularity reflects the concentration of edges within communities compared with random distribution of links between all nodes regardless of communities. Modularity also shows the effectiveness of the method to detect large communities.

	\item Clustering Coefficient is used in this paper as a version of the clustering coefficient \cite{Fagiolo} extended for directed and weighted graphs. \begin{equation}
		\label{eqn:54} 
		C^D = N^{-1} \sum_{i=1}^{N} C_i^D
	\end{equation} In this metric, the average value of the ratio of existing triangles based on the vertex \(i\) to all kinds of triangles based on this vertex is considered, that is, the completeness of the relationship between the vertices is considered. The closer the metric value is to 1, the better the community detection.

	\item Conductance \cite{ChenNguyenSzymanski} is the proportion of the total number of edges outside the commu- nity for unweighted networks or the proportion of the total weight of such edges for weighted networks. This metric allows you to know the “conductivity” of the resulting community. The closer the conductance value to 0, the better the quality of the community. \begin{equation}
		\label{eqn:55}
		\frac{\lvert E_c^{out} \rvert}{\lvert E_c^{in} \rvert + \lvert E_c^{out} \rvert}
	\end{equation}

	\item Contraction \cite{ChenNguyenSzymanski} measures the average number of edges per node within community C or the average weight per node of such edges. This metric shows how important this community is to the rest of the graph. The closer the contraction value to 1, the better the quality of the community. \begin{equation}
		\label{eqn:56}
		\frac{\lvert E_c^{in} \rvert}{c}
	\end{equation}

	\item Expansion \cite{ChenNguyenSzymanski} measures the average number of edges (per node) outside the community C, or the average weight per node of such edges. This metric shows how strongly this community is connected to the rest of the graph. The lower the expansion value, the better the quality of the community.\begin{equation}
		\label{eqn:57}
		\frac{\lvert E_c^{out} \rvert}{c}
	\end{equation}

	\item Community Fitness \cite{KaurSinghKaushal} calculates the ratio of the total indegree number of the community \(C\) to the total degree of \(\alpha\), where is \(\alpha\) a positive number that controls the size of communities. This allows you to find out the density of detected communities. The higher the community fitness value, the better result we get. \begin{equation}
		\label{eqn:58}
		\textit{comfit} = \frac{k_{in}^C}{(k_{in}^C + k_{out}^C)^\alpha}
	\end{equation}
\end{enumerate}

\subsubsection{4. Experiment}

\paragraph{4.1. Experiment Description} The methods suggested above will be applied to the four datasets, after which the resulting values will be evaluated by the metrics~\cref{eqn:52,eqn:53,eqn:54,eqn:55,eqn:56,eqn:57,eqn:58}. For completeness of the experiment, as test data, we take a graph of small size (about 10 thousand nodes), medium size (up to 50 thousand nodes), large size (up to 200 thousand nodes), and extra-large (more than 500 thousand nodes). Moving from a small dataset to large ones during testing, algorithms and metrics that do not perform efficiently on the task will be cut off. Next, two tables will be created for each dataset. The first table will contain the results of applying the metrics, highlighting the best results in each metric. The second table will contain the count of resulting clusters and the size of the seven largest of them. Based on the data obtained for each dataset, it will be possible to draw a conclusion on efficiency of the algorithms on the datasets of varying volume.

Our strategy differs from the ‘ground truth’ one: instead of having just one pre- analyzed discussion with known modularity as a ‘ground truth’ (which might lead to the situation of random closeness of this or that model to the content-based graph modularity), we take four discussions on various volume and set the sociological perspective for expectations, to see which algorithm(s) come closer to them on datasets of varying volume. This is why our design implies use of multiple metrics and multiple datasets.

\paragraph{4.2. The Datasets} To test and evaluate the extant tools, real-world datasets crawled from Twitter, a microblogging service with elements of a social network, were used. Each of these datasets was collected using the Twitter API during or immediately after high-profile, socially resonant events that triggered social unrest \cite{BodrunovaBlekanovSmoliarova,BodrunovaLitvinenkoBlekanov,BodrunovaBlekanov}. For accurate results, the datasets vary in size (the number of tweets), and testing will be carried out from the smallest to the largest one.

The mall size dataset (hereinafter referred to as “Biryulevo”) consists of the tweets posted on 17 to 31 October 2013, during the sharp phase of the riots in the Moscow district Biryulevo-Zapadnoe. Initially, it contains 20,106 links and 11,429 users, but, after cleaning from isolate users, 10,275 vertices and 20,093 edges remain.

The medium size dataset (hereinafter referred to as “Cologne”) contains the data collected for 1 to 31 January 2016, within the discussion on mass attacks on women in Cologne on the eve of 2016. Initially, it contained 40,117 users and 98,508 links, after cleaning from isolate users, the number of users decreased to 36,850, and the number of links to 96,244.

The large size dataset (hereinafter referred to as “Ferguson”) contains data from 22 to 31 August 2014, on the riots in Ferguson, USA, provoked by a murder of an African Ameri- can teenager by a white policeman. The dataset contains 169,676 users with 334,050 links, and, after cleaning from isolate users, 143,024 users with 325,369 links remain.

The extra-large dataset (hereinafter referred to as “Charlie Hebdo”) was obtained for the period of 7 to 10 January 2015 by collecting a response to a terrorist attack in the editorial office of the Charlie Hebdo magazine. The total number of participants is 719,503, connections are 981,131, and, after clearing, there are 617,041 users and 980,351 connections.

\subsubsection{5. Results}

As stated above, in this section we present the results of the experiment for four various datasets of different volumes.

The results of the task solution for the Biryulevo case are shown in Tables~\cref{tab:biryulevoMetricsEvaluation} and~\cref{tab:biryulevoCommunityNumber}. Each method has its own peculiarities:

\begin{itemize}
	\item The Infomap algorithm tends to highlight one large community, which can be used as a starting point for deeper research;
	\item Similar to the Infomap algorithm, GANXiS (both pure and overlapping versions) distinguishes one large society, but this set is smaller than for Infomap;
	\item The generalized K-means algorithm does not take into account the lack of connections between isolates, which is the reason for the unification of almost all small discussions (2-3-4 participants) into one large society;
	\item Directed Louvain shows good results. However, it was not able to get ahead in any metric, steadily holding on to the second or third place;
	\item Despite it only partially belonging to the directed graph clustering algorithms, Leiden (both the one that uses modularity and the one that uses the clique percolation method) shows good results, being ahead of the directed Louvain algorithm in almost every point. Its peculiarity is that it highlights the communities of average size and/or close to each other in size;
	\item DLPA allocates almost twice as many communities as other algorithms, all but the largest do not differ much in size;
	\item OSLOM identifies a large number of communities, which, moreover, can strongly overlap. 
\end{itemize}

\begin{table}[ht]%
	\centering
	\caption{Results of evaluations by metrics on the Biryulevo dataset.}%
	\label{tab:biryulevoMetricsEvaluation}% label всегда желательно идти после caption
	\begin{adjustbox}{width=1\textwidth}
		\small
		\begin{tabular}{ c  c  c  c  c  c  c  c }% Вертикальные полосы не используются принципиально, как и лишние горизонтальные (допускается по ГОСТ 2.105 пункт 4.4.5) % @{} позволяет прижиматься к краям
			\toprule
			& DirMod & ClusCoe & ComFit & NEDind & Conduc & Contrac & Expans\\
			\hline
			DLPA & 0.486 & \textbf{0.056} & 1.373 & 0.605 & 0.625 & \textbf{0.849} & 2.361 \\
			DirLouv & 0.574 & 0.039 & 6.548 & 0.793 & 0.286 & 0.766 & 0.715\\
			Infomap & 0.088 & 0.034 & 7.877 & \textbf{0.940} & \textbf{0.209} & 0.685 & \textbf{0.375}\\
			OSLOM & -- & 0.035 & 2.479 & 0.501 & 0.967 & 0.087 & 1.403 \\
			GKM & 0.390 & 0.0099 & 7.956 & 0.423 & 0.950 & 0.190 & 2.161 \\
			LeidMod & \textbf{0.584} & 0.039 & 6.686 & 0.806 & 0.278 & 0.761 & 0.686 \\
			LeidCPM & 0.571 & 0.040 & 5.052 & 0.676 & 0.365 & 0.828 & 1.001 \\
			GANXiSd & 0.406 & 0.040 & 2.538 & 0.661 & 0.526 & 0.841 & 1.681\\
			GANXiSo & -- & 0.049 & 3.295 & 0.643 & 0.550 & 0.848 & 2.282 \\
			\bottomrule
			\multicolumn{8}{@{}p{\textwidth}}{%
				\hspace*{2.5em}% абзацный отступ - требование ГОСТ 2.105
				Note. The best results are highlighted in bold.
			}\\
		\end{tabular}%
	\end{adjustbox}
\end{table}

\begin{table}[ht]%
	\centering
	\caption{The number of communities obtained on the Biryulevo dataset and the value of the seven largest of them for each algorithm.}%
	\label{tab:biryulevoCommunityNumber}% label всегда желательно идти после caption
%	\begin{adjustbox}{width=1\textwidth}
%		\small
		\begin{tabular}{ c  c  c  c  c  c  c  c  c }% Вертикальные полосы не используются принципиально, как и лишние горизонтальные (допускается по ГОСТ 2.105 пункт 4.4.5) % @{} позволяет прижиматься к краям
			\toprule
			& Count & 1 & 2 & 3 & 4 & 5 & 6 & 7\\
			\hline
			DLPA & 1159 & 741 & 205 & 171 & 115 & 111 & 108 & 105 \\
			DirLouv & 243 & 691 & 583 & 494 & 486 & 463 & 459 & 381 \\
			Infomap & 202 & 9276 & 232 & 86 & 41 & 33 & 28 & 26 \\
			OSLOM & 642 & 2637 & 1743 & 1425 & 847 & 763 & 737 & 619\\
			GKM & 200 & 3300 & 2815 & 1812 & 815 & 496 & 313 & 90 \\
			LeidMod & 238 & 716 & 464 & 458 & 451 & 416 & 388 & 356\\
			LeidCPM & 315 & 627 & 613 & 226 & 226 & 214 & 208 & 208 \\
			GANXiSd & 627 & 4881 & 190 & 146 & 129 & 113 & 103 & 101 \\
			GANXiSo & 647 & 4956 & 638 & 255 & 166 & 146 & 139 & 83 \\
			\bottomrule
		\end{tabular}%
%	\end{adjustbox}
\end{table}

The results for the Cologne case are shown in Tables~\cref{tab:cologneMetricsEvaluation} and~\cref{tab:cologneCommunityNumber}. At this dataset volume,~\cref{eqn:53} metric showed high computational complexity, so it was no longer used.

\begin{table}[ht]%
	\centering
	\caption{Results of evaluations by metrics on the Cologne dataset.}%
	\label{tab:cologneMetricsEvaluation}% label всегда желательно идти после caption
	\begin{adjustbox}{width=1\textwidth}
		\small
		\begin{tabular}{ c  c  c  c  c  c  c  c }% Вертикальные полосы не используются принципиально, как и лишние горизонтальные (допускается по ГОСТ 2.105 пункт 4.4.5) % @{} позволяет прижиматься к краям
			\toprule
			& DirMod & ClusCoe & ComFit & NEDind & Conduc & Contrac & Expans\\
			\hline
			DLPA & -- & \textbf{0.070} & 2.299 & 0.723 & 0.521 & 0.815 & 1.378 \\
			DirLouv & -- & 0.028 & 7.753 & 0.917 & 0.234 & 0.691 & 0.459 \\
			Infomap & -- & 0.024 & 8.468 & \textbf{0.949} & \textbf{0.201} & 0.653 & \textbf{0.344}\\
			OSLOM & -- & 0.034 & 3.380 & 0.495 & 0.963 & 0.098 & 1.249 \\
			GKM & -- & 0.013 & \textbf{28.498} & 0.422 & 0.934 & 0.209 & 2.600 \\
			LeidMod & -- & 0.028 & 7.764 & 0.916 & 0.233 & 0.693 & 0.459 \\
			LeidCPM & -- & 0.033 & 4.379 & 0.689 & 0.397 & 0.793 & 1.030 \\
			GANXiSd & -- & 0.031 & 3.057 & 0.693 & 0.501 & 0.807 & 1.798 \\
			GANXiSo & -- & 0.042 & 3.551 & 0.670 & 0.526 & \textbf{0.829} & 2.290 \\
			\bottomrule
			\multicolumn{8}{@{}p{\textwidth}}{%
				\hspace*{2.5em}% абзацный отступ - требование ГОСТ 2.105
				Note. The best results are highlighted in bold.
			}\\
		\end{tabular}%
	\end{adjustbox}
\end{table}

\begin{table}[ht]%
	\centering
	\caption{The number of communities obtained on the Cologne dataset and the value of the seven largest of them for each algorithm.}%
	\label{tab:cologneCommunityNumber}% label всегда желательно идти после caption
%	\begin{adjustbox}{width=1\textwidth}
%		\small
		\begin{tabular}{ c  c  c  c  c  c  c  c  c }% Вертикальные полосы не используются принципиально, как и лишние горизонтальные (допускается по ГОСТ 2.105 пункт 4.4.5) % @{} позволяет прижиматься к краям
			\toprule
			& Count & 1 & 2 & 3 & 4 & 5 & 6 & 7\\
			\hline
			DLPA & 2479 & 20,669 & 425 & 277 & 249 & 198 & 173 & 168 \\
			DirLouv & 735 & 4356 & 4105 & 4091 & 3604 & 3447 & 1933 & 1886  \\
			Infomap & 673 & 27,016 & 3726 & 1145 & 590 & 465 & 445 & 338 \\
			OSLOM & 1690 & 10,823 & 7336 & 3058 & 1881 & 1519 & 1466 & 921 \\
			GKM & 200 & 11,474 & 7737 & 4353 & 2971 & 2947 & 1859 & 1119  \\
			LeidMod & 734 & 4453 & 3882 & 3540 & 3212 & 3181 & 2686 & 1739 \\
			LeidCPM & 1221 & 2395 & 672 & 625 & 414 & 405 & 362 & 355 \\
			GANXiSd & 1864 & 24,911 & 128 & 110 & 108 & 94 & 89 & 80 \\
			GANXiSo & 1964 & 25,365 & 458 & 210 & 162 & 128 & 110 & 83 \\
			\bottomrule
		\end{tabular}%
%	\end{adjustbox}
\end{table}

For the mid-range dataset, the algorithms have shown the following results:

\begin{itemize}
	\item Infomap showed a more even selection of communities than on the Biryulevo dataset, nevertheless retaining the emphasis on the largest of them;
	\item GANXiS continued to highlight one large community, and the emphasis on this community grew;
	\item As in the case of the Biryulevo dataset, generalized K-means has combined minor discussions into one large group;
	Directed Louvain performed well again, continuing to hold on to second or third place;
	\item Unlike for the first dataset, the sizes of communities after the Leiden algorithm differ depending on the metric (CPM or modularity). The first option showed a large number of fairly small communities, while the second showed a smaller number of moderately large communities;
	\item DLPA on the Cologne dataset showed the risk of going into “overflow” of one community, not giving any other one chances to grow while, again, having twice as many	communities as other algorithms;
	\item OSLOM distinguished communities of not the best quality even in comparison with	GANXiS.
\end{itemize}

The results for the large dataset (Tables~\cref{tab:fergusonMetricsEvaluation} and~\cref{tab:fergusonCommunityNumber}) may be described the following way:

\begin{itemize}
	\item Infomap once again got four best metrics results, receiving, though, one large community;
	\item GANXiS kept the trend shown on small and medium datasets;
	\item The generalized K-Means algorithm has shown its inapplicability to sufficiently perform on large graphs, since both the computational cost and the memory cost made it impossible to use this algorithm;
	\item Directed Louvain retains its position relative to other algorithms;
	\item The Leiden algorithm kept the trend shown on the medium-sized dataset;
	\item DLPA identified more equal-sized communities than on the medium dataset, which demonstrates variability of the algorithm performance depending on the size of initial data;
	\item OSLOM got results worse than on medium dataset on every, except metric (8).
\end{itemize}

\begin{table}[ht]%
	\centering
	\caption{Results of evaluations by metrics on the Ferguson dataset.}%
	\label{tab:fergusonMetricsEvaluation}% label всегда желательно идти после caption
	\begin{adjustbox}{width=1\textwidth}
		\small
		\begin{tabular}{ c  c  c  c  c  c  c  c }% Вертикальные полосы не используются принципиально, как и лишние горизонтальные (допускается по ГОСТ 2.105 пункт 4.4.5) % @{} позволяет прижиматься к краям
			\toprule
			& DirMod & ClusCoe & ComFit & NEDind & Conduc & Contrac & Expans\\
			\hline
			DLPA & -- & \textbf{0.085} & 1.522 & 0.640 & 0.599 & \textbf{0.847} & 2.735 \\
			DirLouv & -- & 0.027 & 6.318 & 0.949 & 0.208 & 0.666 & 0.371 \\
			Infomap & -- & 0.026 & \textbf{6.476} & \textbf{0.959} & \textbf{0.199} & 0.655 & \textbf{0.340}\\
			OSLOM & -- & 0.034 & 2.426 & 0.493 & 0.964 & 0.096 & 1.305 \\
			GKM & -- & -- & -- & -- & -- & -- & -- \\
			LeidMod & -- & 0.027 & 6.244 & 0.945 & 0.212 & 0.670 & 0.380 \\
			LeidCPM & -- & 0.035 & 4.494 & 0.740 & 0.345 & 0.782 & 1.062 \\
			GANXiSd & -- & 0.031 & 3.303 & 0.772 & 0.422 & 0.773 & 1.260 \\
			GANXiSo & -- & 0.034 & 4.064 & 0.779 & 0.397 & 0.756 & 1.294 \\
			\bottomrule
			\multicolumn{8}{@{}p{\textwidth}}{%
				\hspace*{2.5em}% абзацный отступ - требование ГОСТ 2.105
				Note. The best results are highlighted in bold.
			}\\
		\end{tabular}%
	\end{adjustbox}
\end{table}

\begin{table}[ht]%
	\centering
	\caption{The number of communities obtained on the Ferguson dataset and the value of the seven largest of them for each algorithm.}%
	\label{tab:fergusonCommunityNumber}% label всегда желательно идти после caption
%	\begin{adjustbox}{width=1\textwidth}
%		\small
		\begin{tabular}{ c  c  c  c  c  c  c  c  c }% Вертикальные полосы не используются принципиально, как и лишние горизонтальные (допускается по ГОСТ 2.105 пункт 4.4.5) % @{} позволяет прижиматься к краям
			\toprule
			& Count & 1 & 2 & 3 & 4 & 5 & 6 & 7\\
			\hline
			DLPA & 18,389 & 3059 & 1224 & 1044 & 913 & 877 & 708 & 662 \\
			DirLouv & 4429 & 16,029 & 13,916 & 13,391 & 12,936 & 8827 & 4334 & 4053 \\
			Infomap & 4321 & 106,521 & 15,617 & 1651 & 1012 & 506 & 409 & 353\\
			OSLOM & 11,546 & 16,417 & 15,372 & 10,438 & 8874 & 6587 & 4081 & 3074 \\
			GKM & -- & -- & -- & -- & -- & -- & -- & -- \\
			LeidMod & 4481 & 15,804 & 12,715 & 12,212 & 11,385 & 9162 & 4355 & 4260 \\
			LeidCPM & 6226 & 1001 & 729 & 704 & 648 & 572 & 554 & 531 \\
			GANXiSd & 8472 & 98,344 & 211 & 184 & 159 & 104 & 100 & 97 \\
			GANXiSo & 7489 & 111,798 & 184 & 159 & 155 & 102 & 95 & 84\\
			\bottomrule
		\end{tabular}%
%	\end{adjustbox}
\end{table}

The results for the extra-large dataset (Tables~\cref{tab:charlieHebdoMetricsEvaluation} and~\cref{tab:charlieHebdoMetricsEvaluation}) are the following:

\begin{itemize}
	\item Infomap showed a better division into same-size communities, retaining the emphasis on the largest of them;
	\item It revealed that the GANXiS algorithm is not applicable to extra-large networks;
	\item As in the case of the large dataset, generalized K-means is not applicable to the networks of this size;
	\item The directed Louvain continued to evenly allocate communities;
	\item Leiden shown dramatically grown difference between the size of CPM and modularity metric-based communities;
	\item DLPA on an extra-large dataset has allocated a lot of small, smaller than before, communities;
	\item OSLOM has got even worse results than on the large dataset on every metric.
	
\end{itemize}

\begin{table}[ht]%
	\centering
	\caption{Results of evaluations by metrics on the Charlie Hebdo dataset.}%
	\label{tab:charlieHebdoMetricsEvaluation}% label всегда желательно идти после caption
	\begin{adjustbox}{width=1\textwidth}
		\small
		\begin{tabular}{ c  c  c  c  c  c  c  c }% Вертикальные полосы не используются принципиально, как и лишние горизонтальные (допускается по ГОСТ 2.105 пункт 4.4.5) % @{} позволяет прижиматься к краям
			\toprule
			& DirMod & ClusCoe & ComFit & NEDind & Conduc & Contrac & Expans\\
			\hline
			DLPA & -- & \textbf{0.027} & 1.375 & 0.672 & 0.514 & \textbf{0.764} & 1.550 \\
			DirLouv & -- & 0.012 & 4.395 & 0.953 & 0.163 & 0.604 & 0.256 \\
			Infomap & -- & 0.012 & 4.381 & 0.942 & 0.166 & 0.608 & 0.268 \\
			OSLOM & -- & 0.015 & 1.730 & 0.479 & 0.979 & 0.054 & 1.239 \\
			GKM & -- & -- & -- & -- & -- & -- & -- \\
			LeidMod & -- & 0.011 & \textbf{4.408} & \textbf{0.955} & \textbf{0.162} & 0.054 & 1.239 \\
			LeidCPM & -- & 0.016 & 3.351 & 0.758 & 0.276 & 0.713 & 0.663 \\
			GANXiSd & -- & -- & -- & -- & -- & -- & --  \\
			GANXiSo & -- & -- & -- & -- & -- & -- & --  \\
			\bottomrule
			\multicolumn{8}{@{}p{\textwidth}}{%
				\hspace*{2.5em}% абзацный отступ - требование ГОСТ 2.105
				Note. The best results are highlighted in bold.
			}\\
		\end{tabular}%
	\end{adjustbox}
\end{table}

\begin{table}[ht]%
	\centering
	\caption{The number of communities obtained on the Charlie Hebdo dataset and the value of the seven largest of them for each algorithm.}%
	\label{tab:charlieHebdoCommunityNumber}% label всегда желательно идти после caption
%	\begin{adjustbox}{width=1\textwidth}
%		\small
		\begin{tabular}{ c  c  c  c  c  c  c  c  c }% Вертикальные полосы не используются принципиально, как и лишние горизонтальные (допускается по ГОСТ 2.105 пункт 4.4.5) % @{} позволяет прижиматься к краям
			\toprule
			& Count & 1 & 2 & 3 & 4 & 5 & 6 & 7\\
			\hline
			DLPA & 84,409 & 1244 & 882 & 726 & 662 & 576 & 529 & 488  \\
			DirLouv & 26,415 & 68,217 & 45,678 & 44,734 & 40,369 & 26,807 & 22,583 & 18,774 \\
			Infomap & 26,497 & 164,839 & 33,635 & 32,111 & 25,038 & 20,098 & 17,226 & 14,033 \\
			OSLOM & 67,296 & 34,960 & 27,735 & 26,061 & 17,319 & 10,912 & 7614 & 6962 \\
			GKM & -- & -- & -- & -- & -- & -- & -- & -- \\
			LeidMod & 26,335 & 56,201 & 39,371 & 33,976 & 27,120 & 25,745 & 20,270 & 17,705 \\
			LeidCPM & 34,647 & 949 & 719 & 626 & 607 & 565 & 531 & 521 \\
			GANXiSd & -- & -- & -- & -- & -- & -- & -- & -- \\
			GANXiSo & -- & -- & -- & -- & -- & -- & -- & -- \\
			\bottomrule
		\end{tabular}%
%	\end{adjustbox}
\end{table}

\subsubsection{6. Discussion and Conclusions}

The experiments conducted have shown the relative mathematical efficiency of the selected methods for detecting user communities in the discussions of Twitter (the mi- croblogging service with elements of a social network) when the discussions are presented in the form of directed graphs of different sizes. Table~\cref{tab:modelSummaryEvaluation} clearly demonstrates that the Infomap algorithm can be called the best; however, due to the property of highlighting one large society, it can rather be used as a preparatory stage for the subsequent clustering by another algorithm, or the repeated application of Infomap. After Infomap comes the Leiden algorithm, which has shown the best results after Infomap, and does not have the ability to single out one large community. Therefore, these two algorithms can be used in conjunction. If it is necessary to isolate overlapping communities based on test results, the GANXiS algorithm should be used.

\begin{longtblr}[
	caption = {Results of evaluations by metrics on the Charlie Hebdo dataset.},
	label = {tab:modelSummaryEvaluation},
	remark{\hspace*{2.5em}Note} = {The best results for each case/algorithm are highlighted in bold.},
	]{
		colspec = {XXXXX}, 
		width = 1.0\linewidth,
		rowhead = 1,
	} 
			\toprule
			& Bir & Col & Fer & ChE \\
			\hline
			\SetCell[c=5]{c}DirMod \\
			\hline
			DLPA & 0.486 & -- & -- & -- \\
			DirLouv & 0.574 & -- & -- & --\\
			Infomap & 0.088 & -- & -- & --\\
			OSLOM & -- & -- & --  & --\\
			GKM & -- & -- & -- & --\\
			LeidMod & 0.39 & -- & -- & -- \\
			LeidCPM & \textbf{0.584} & -- & -- & --\\
			GANXiSd & 0.406 & -- & -- & -- \\
			GANXiSo & -- & -- & -- & -- \\
			\hline
			\SetCell[c=5]{c}ClusCoe \\
			\hline
			DLPA & \textbf{0.056} & \textbf{0.07} & \textbf{0.085} & \textbf{0.027} \\
			DirLouv & 0.039 & 0.028 & 0.027 & 0.012 \\
			Infomap & 0.034 & 0.024 & 0.026 & 0.012 \\
			OSLOM & 0.035 & 0.034 & 0.034 & 0.015\\
			GKM & 0.0099 & -- & -- & --\\
			LeidMod & 0.039 & 0.028 & 0.027 & 0.011 \\
			LeidCPM & 0.04 & 0.033 & 0.035 & 0.016 \\
			GANXiSd & 0.04 & 0.031 & 0.031 & -- \\
			GANXiSo & 0.049 & 0.042 & 0.034 & -- \\
			\hline
			\SetCell[c=5]{c}ComFit \\
			\hline
			DLPA & 1.373 & 2.299 & 1.522 & 1.375 \\
			DirLouv & 6.548 & 7.753 & 6.318 & 4.395 \\
			Infomap & 7.877 & 8.468 & \textbf{6.476} & 4.381\\
			OSLO & 2.479 & 3.38 & 2.426 & 1.73 \\
			GKM & \textbf{7.956} & \textbf{28.498} & -- & --\\
			LeidMod & 6.686 & 7.764 & 6.244 & \textbf{4.408} \\
			LeidCPM & 5.052 & 4.379 & 4.494 & 3.351 \\
			GANXiSd & 2.538 & 3.057 & 3.303 & -- \\
			GANXiSo & 3.295 & 3.551 & 4.064 & -- \\
			\hline
			\SetCell[c=5]{c}NEDind \\
			\hline
			DLPA & 0.605 & 0.723 & 0.64 & 0.672 \\
			DirLouv & 0.793 & 0.917 & 0.949 & 0.953 \\
			Infomap & \textbf{0.94} & \textbf{0.949} & \textbf{0.959} & 0.942 \\
			OSLO & 0.501 & 0.495 & 0.493 & 0.479 \\
			GKM & 0.423 & 0.422 & -- & --\\
			LeidMod & 0.806 & 0.916 & 0.945 & \textbf{0.955}  \\
			LeidCPM & 0.67 &6 0.689 & 0.74 & 0.758  \\
			GANXiSd & 0.661 & 0.693 & 0.772 & -- \\
			GANXiSo & 0.643 & 0.67 & 0.799 & -- \\
			\hline
			\SetCell[c=5]{c}Conduc \\
			\hline
			DLPA & 0.625 & 0.521 & 0.599 & 0.514 \\
			DirLouv & 0.286 & 0.234 & 0.208 & 0.163 \\
			Infomap & \textbf{0.209} & \textbf{0.201} & \textbf{0.199} & 0.166 \\
			OSLO & 0.967 & 0.963 & 0.964 & 0.979 \\
			GKM & 0.95 & 0.934 & -- & --\\
			LeidMod & 0.278 & 0.233 & 0.212 & \textbf{0.162} \\
			LeidCPM & 0.365 & 0.397 & 0.345 & 0.276  \\
			GANXiSd & 0.526 & 0.501 & 0.422 & -- \\
			GANXiSo & 0.55 & 0.526 & 0.397 & -- \\
			\hline
			\SetCell[c=5]{c}Contrac \\
			\hline
			DLPA & \textbf{0.849} & 0.815 & \textbf{0.847} & \textbf{0.764} \\
			DirLouv & 0.766 & 0.691 & 0.666 & 0.604  \\
			Infomap & 0.685 & 0.653 & 0.655 & 0.608  \\
			OSLO & 0.08 &7 0.098 & 0.096 & 0.054  \\
			GKM & 0.19 & 0.209 & -- & --\\
			LeidMod & 0.761 & 0.693 & 0.67 & 0.603 \\
			LeidCPM & 0.828 & 0.793 & 0.782 & 0.713   \\
			GANXiSd & 0.841 & 0.807 & 0.773 & -- \\
			GANXiSo & 0.848 & \textbf{0.829} & 0.756 & -- \\
			\hline
			\SetCell[c=5]{c}Expans \\
			\hline
			DLPA & 2.361 & 1.378 & 2.735 & 1.55\\
			DirLouv & 0.715 & 0.459 & 0.371 & 0.256\\
			Infomap & \textbf{0.375} & \textbf{0.344} & \textbf{0.34} & 0.268   \\
			OSLO & 1.403 & 1.249 & 1.305 & 1.239   \\
			GKM & 2.161 & 2.6 & -- & --\\
			LeidMod & 0.686 & 0.459 & 0.38 & \textbf{0.253}  \\
			LeidCPM & 1.001 & 1.030 & 1.062 & 0.663  \\
			GANXiSd & 1.681 & 1.798 & 1.26 & -- \\
			GANXiSo & 2.282 & 2.290 & 1.294 & -- \\
			\bottomrule
%		\end{tabular}%
%	\end{adjustbox}
\end{longtblr}

Table~\cref{tab:communitiesObtained} shows that the generalized K-means and GANXiS algorithms are not suitable for large-scale Twitter data. Infomap and Leiden (again), just as Direct Louvain, show very similar results as to the number of the detected communities, and the numbers are lower. Here, we need to repeat that, without a clear sociological task, we could not set the ground truth to which to compare the test results; the number of the detected communities (higher or lower) is better depending on the research goal. In general, the low number of robust communities is optimal. This is why we see these three algorithms as coming closer to possible ground truth in social research. Yet, it is also necessary to state that algorithmic results do not correspond to the logic of social studies designated in Section 2.2, as the methods come to either delineation of one large community or a much larger number of smaller communities detected. Of the three aforementioned algorithms with the best (lowest) number of communities, Infomap finds the largest groups, which might be considered the best result.

\begin{table}[ht]%
	\centering
	\caption{The number of communities obtained as compared by algorithm and the seven largest communities, summary table.}%
	\label{tab:communitiesObtained}% label всегда желательно идти после caption
	\begin{adjustbox}{width=1\textwidth}
		\small
		\begin{tabular}{ c  c  c  c  c  c  c  c  c  c  }% Вертикальные полосы не используются принципиально, как и лишние горизонтальные (допускается по ГОСТ 2.105 пункт 4.4.5) % @{} позволяет прижиматься к краям
			\toprule
			Algorithm & Case & Count & 1 & 2 & 3 & 4 & 5 & 6 & 7 \\
			\hline
			\multirow{4}{*}{DLPA} & Bir & 1159 & 741 & 205 & 171 & 115 & 111 & 108 & 105 \\
			& Col & 2479 & 20,669 & 425 & 277 & 249 & 198 & 173 & 168 \\
			& Fer & 18,389 & 3059 & 1224 & 1044 & 913 & 877 & 708 & 662\\
			& ChE & 84,409 & 1244 & 882 & 726 & 662 & 576 & 529 & 488 \\
			\hline
			\multirow{4}{*}{DirLouv} & Bir & 243 & 691 & 583 & 494 & 486 & 463 & 459 & 381 \\
			& Col & 735 & 4356 & 4105 & 4091 & 3604 & 3447 & 1933 & 1886 \\
			& Fer & 4429 & 16,029 & 13,916 & 13,391 & 12,936 & 8827 & 4334 & 4053 \\
			& ChE & 26,415 & 68,217 & 45,678 & 44,734 & 40,369 & 26,807 & 22,583 & 18,774 \\
			\hline
			\multirow{4}{*}{Infomap} & Bir & 202 & 9276 & 232 & 86 & 41 & 33 & 28 & 26 \\
			& Col & 673 & 27,016 & 3726 & 1145 & 590 & 465 & 445 & 338 \\
			& Fer & 4321 & 106,521 & 15,617 & 1651 & 1012 & 506 & 409 & 353 \\
			& ChE & 26,497 & 164,839 & 33,635 & 32,111 & 25,038 & 20,098 & 17,226 & 14,033 \\
			\hline
			\multirow{4}{*}{OSLOM} & Bir & 642 & 2637 & 1743 & 1425 & 847 & 763 & 737 & 619 \\
			& Col & 1690 & 10,823 & 7336 & 3058 & 1881 & 1519 & 1466 & 921 \\
			& Fer & 11,546 & 16,417 & 15,372 & 10,438 & 8874 & 6587 & 4081 & 3074 \\
			& ChE & 67,296 & 34,960 & 27,735 & 26,061 & 17,319 & 10,912 & 7614 & 6962 \\
			\hline
			\multirow{4}{*}{GKM} & Bir & 200 & 3300 & 2815 & 1812 & 815 & 496 & 313 & 90 \\
			& Col & 200 & 11,474 & 7737 & 4353 & 2971 & 2947 & 1859 & 1119 \\
			& Fer & --& --& --& --& --& --& --& --\\
			& ChE & --& --& --& --& --& --& --& --\\
			\hline
			\multirow{4}{*}{LeidMod} & Bir & 238 & 716 & 464 & 458 & 451 & 416 & 388 & 356\\
			& Col & 734 & 4453 & 3882 & 3540 & 3212 & 3181 & 2686 & 1739\\
			& Fer & 4481 & 15,804 & 12,715 & 12,212 & 11,385 & 9162 & 4355 & 4260  \\
			& ChE & 26,335 & 56,201 & 39,371 & 33,976 & 27,120 & 25,745 & 20,270 & 17,705 \\
			\hline
			\multirow{4}{*}{LeidCPM} & Bir & 315 & 627 & 613 & 226 & 226 & 214 & 208 & 208 \\
			& Col & 1221 & 2395 & 672 & 625 & 414 & 405 & 362 & 355 \\
			& Fer & 6226 & 1001 & 729 & 704 & 648 & 572 & 554 & 531 \\
			& ChE & 34,647 & 949 & 719 & 626 & 607 & 565 & 531 & 521 \\
			\hline
			\multirow{4}{*}{GANXiSd} & Bir & 627 & 4881 & 190 & 146 & 129 & 113 & 103 & 101 \\
			& Col & 1864 & 24,911 & 128 & 110 & 108 & 94 & 89 & 80 \\
			& Fer & 8472 & 98,344 & 211 & 184 & 159 & 104 & 100 & 97 \\
			& ChE & --& --& --& --& --& --& --& --\\
			\hline
			\multirow{4}{*}{GANXiSo} & Bir & 647 & 4956 & 638 & 146 & 139 & 83 & 255 & 166 \\
			& Col & 1964 & 25,365 & 458 & 128 & 110 & 83 & 210 & 162 \\
			& Fer & 7489 & 111,798 & 184 & 102 & 95 & 84 & 159 & 155 \\
			& ChE & --& --& --& --& --& --& --& --\\
			\bottomrule
		\end{tabular}%
	\end{adjustbox}
\end{table}

Yet, it would be too bold to call such a result satisfactory in sociological terms. Socio- logically, the use of the tested algorithms poses a general question on their applicability for current social science tasks, as well as the following questions:

\begin{itemize}
	\item How does the structure of hidden communities relate to the expected social, cultural, and/or political cleavages in the discussion?
	\item How can algorithmic detection of hidden communities come closer to detecting communities of views, as linked to communities of formal connections?	
\end{itemize}

Earlier, we had partly answered this question \cite{BodrunovaBlekanov} by using a complicated multi-step methodology of group views detection. However, the use of a combination of Infomap (to find the discussion core) and Leiden (to find modules within the core) might be a shorter way to find the communities based on social traits or views. This creates implications for further research that will show whether the formal community structure corresponds to the substantial divisions in the discourses.

Comparing our results to other works also shows the following. Despite the fact that, in our work, we did not test the algorithms on synthetic data, our performance on real data was better than that in \cite{VanLierdeDelvenneVanDooren} on synthetic data. However, since ground-truth testing also has its drawbacks on real data, further analysis of the clustering results is planned in future tests. Additionally, in the nearest future, we plan the following expansion of our work:

\begin{enumerate}
	\item Conduct semantic analysis of the quality of the results obtained using experts or NLP methods;
	\item Expand the work results by adding agglomerative clustering and Markov stopping moment for optimal clustering \cite{BodrunovaOrekhovBlekanov}.
\end{enumerate}

\subsection{Comparing influencers: activity vs. connectivity measures in defining key actors in twitter \textit{ad hoc} discussions on migrants in Germany and Russia}\label{subsec:ch2/sec4/sub3}

\subsubsection{1. Introduction}

Uneven representation of group interest in mediatized public discussions has been established in the research literature \cite{Nieminen} as one of the fundamental problems of public communication and public decision-making. Among the reasons for that, there is representation of newsmakers privileging institutional actors vs. ordinary citizens as \textit{vox populi} \cite{ScheufeleTewksbury}. Since Internet had emerged as a public communicative space less dependent on media, scholars expressed hopes that networked communication would provide for equalizing citizens with institutional actors within public discussions \cite{White1997} bypassing media who used to serve as gatekeepers of public agendas \cite{White1950}. But, till today, horizontalization of discursive relations online remains highly disputable \cite{Fuchs}; moreover, new societal cleavages emerge in hybrid media systems \cite{Chadwick} due to digital divide, interest- and value-based variance in media diets, and growing platform- oriented fragmentation of public arenas.

In online communicative milieus, the figures of \textit{newsmaker}, \textit{informer}, and \textit{opinion leader} are re-conceptualized as that of \textit{influencer} \cite{PattersonGrennyMaxfield}, partly based on an older idea of ‘influential’ \cite{Rogers}. Influencers combine beyond-the-average capacities of information dissemination with those of casting impact upon users’ opinions and formation of discussion circles often described as echo chambers \cite{Wallsten}, and thus are key structural elements of networked discussions \cite{Castells2007,BakshyRosennMarlow}.

Despite influencers’ expected crucial role in reshaping power relations between institutional and non-institutional participants of online discussions, they are, till today, under-studied in such aspects as dependence of influencer position upon user activity, institutional status, or taking sides in conflict. Social network analysis (SNA) tries to predict influencers technically, based on their activity and metadata, as well as on the discussion graph structure; other important works explore the interplay between the nature of the publics and constellations of influencers \cite{Habermas,Dahlgren,BrunsBurgess,Papacharissi,BrunsHighfeld2016}. Within this research cluster, Twitter as a microblogging platform has gained particular attention, but it is yet unclear whether this platform tends to democratize influencers in the so-called \textit{ad hoc} public discussions that rapidly rise and disseminate on events of high social relevance or on issues with high potential of social polarization.

To a large extent, this is due to the fact that very different approaches to defining and detecting influencers co-exist in computer science and communication disciplines. Earlier, we traced at least two concepts of influencer (based on user activity and user connectivity, respectively), as well as a methodological divide in detecting influencers via absolute-figure metrics and SNA metrics \cite{BodrunovaLitvinenkoBlekanov2016}; we also stated that few attempts had been made to juxtapose these ways of detecting influencers. Also, comparative studies beyond the Western and Arab Spring countries remain rare \cite{HladikStetka}.

Thus, the aim of this paper is twofold. First, we assess whether user activity necessarily leads to better connectivity, and by what metrics. Then, we try to compare the structure of influencers across countries in terms of their institutional belonging and pro-/anti-migrant stance. We do this by collecting and analyzing data on the Twitter discussion around anti-migrant bashings in Biryuliovo (Moscow) in 2013 and the one around the mass harassment in Cologne in 2016. To accomplish this, we have collected the discussion content, selected the metrics, applied them and formed user lists by activity and connectivity metrics (betweenness and pagerank centralities). We manually assessed the listed accounts to position them institutionally and politically.

Section 2 presents our conceptualization of ‘marketing’ vs. ‘deliberative’ influen- cers, while Sect. 3 reviews today’s approaches to defining influencers, including those based on user activity and connectivity metrics. Section 4 describes the cases, research hypotheses, and our methodology. Section 5 discusses our results.

\subsubsection{2. Actor Disparities in \textit{Ad Hoc} Twitter Discussions}

By 1990s, the public sphere theory had already stated that public discussions were arenas of high disparities in terms of who formed the opinions and influenced the discussion agendas. As mentioned above, institutional and elite representatives were naturally preferred by media; moreover, media themselves became the key nodes in information networks and performed agenda setting \cite{McCombsShaw,McCombs}. Another reason for criticism of media-based public spheres was their oppressive majority-oriented discourse \cite{Fraser,LaclauMouffe,FentonDowney,Dahlberg}; a lot of efforts have been put by countries in Europe and beyond to establish public media that would encompass at least some minority views.

With the rise of online platforms, hopes for better access of citizens to public discussion first rose \cite{Fuchs} and then faded, as both social \cite{Nakamura} and communicative \cite{Daniels} offline divides were accompanied by new disparities emerging due to divergent media consumption \cite{PfetschAdam,BodrunovaLitvinenko} and digital divide \cite{Norris,VanDeursenVanDijk}, among other reasons. Not even asking whether Twitter discussions have any impact upon real-world policymaking, scholars doubt even whether ‘Habermas is on Twitter’ \cite{BrunsHighfeld2016} \cite[p.~31]{Murthy}. Out of this, a range of research agendas have emerged on who become discussion leaders (influencers) and whether the disparities in influence persist. Also, we need to know how we define and detect the influencers, as their detection appears to be measure-dependent.

\textit{Defining an Influencer.} SNA is widely used to show deviant users in Twitter discussions. As we stated before \cite{BodrunovaLitvinenkoBlekanov2016}, there are at least three major divisions in SNA-based influencer studies that define influencers in differing ways.

Here, we will only shortly reconstruct our logic and show applicability of this logic to comparative studies. Thus, the three divisions may be conceptualized as follows. The first one is between ‘marketing’ and ‘deliberative’ influencers. The former generates a self-oriented ‘long tail’ of attention and support \cite[p.~1261]{DuboisGaffney} \cite{Aquino}; here, key characteristics of an influencer are \(N_{followers}\), the quantity and regularity of posting, and the vastness of ‘support waves’ of liking and retweeting. The latter, ‘deliberative’ influ- encer, helps in formation of a politically relevant and effective discussion by linking user groups with varying or even opposing views, as well as of intertwining topic-based echo chambers; also, such a user is linked to the maximum number of other users within the discussion by interacting with them. As inclusiveness and horizontality \cite{Papacharissi2010}, along with rationality and orientation to consensus, are key features of an effective ‘field of discursive connections’ \cite[p.~37]{Calhoun}, deliberative influencers are key for formation of ‘opinion crossroads’ \cite{BodrunovaLitvinenko,VanDeursenVanDijk} as a metaphor of an all-involving public discussion. Structurally, inter-linkage between clusters in a discussion and \(N_{users}\) involved in commenting and retweeting becomes the feature that defines an influencer. We consider these approaches mutually amplifying, as they both, in a way, are extensions of theory of two-step communication flow via opinion leaders \cite{Katz}.

To add, two more divisions may be traced: first, the one between user activity metrics (\(N_{posts}\), \(N_{likes}\), \(N_{retweets}\), \(N_{comments}\) left, \(N_{users}\) followed, \(N_{users}\) involved by a given user into any type of interaction) and user connectivity metrics (\(N_{likes}\), \(N_{retweets}\), \(N_{comments}\) received, \(N_{followers}\), \(N_{users}\) interacting with a given user, and centrality metrics that describe a user’s position in the web graph). And second, the same metrics are divided into absolute-figure ones measured for every user independently and graph-based metrics that, for every user, depend on the overall graph configuration \cite{BodrunovaLitvinenkoBlekanov2016}.

\textit{Conceptual Limitations in Twitter Studies of Influencers.} But before discussing particular ways of detecting influencers on Twitter, we need to mention that there are limitations for that; they are linked to the nature of the discussion, its level of rationality, and inherent Twitter mechanisms that technically privilege certain actors \cite{BodrunovaLitvinenkoBlekanov2016}.

In short, the first limitation is linked to the fact that ‘issue publics’ \cite[p.~422]{Habermas} \cite[p.~108]{VanDeursenVanDijk}, or \textit{ad hoc} publics \cite{BrunsBurgess}, become affective \cite{Papacharissi} and quickly rise and dissolve \cite[p.~74]{Dahlgren}. This, in its turn, raises two issues: (1) that of representability of \textit{ad hoc} discussion for stable discursive patterns outside the time of the event; (2) comparability of \textit{ad hoc} discussions in terms of their structure and the conclusions they allow for. In response to this, we may state that our experiments (work in progress) show that the structure of \textit{ad hoc} discussions changes the same way in six different discussions if isolated users are eliminated; that is, the patterns of \textit{ad hoc} discussions are, at least partly, comparable. In future, we will also test their comparability with stable discussions.

The issue of rationality has been debated among scholars since the appearance of Twitter itself. Twitter pessimists claim that the platform is home for depoliticized trivial content full of ‘white noise’ \cite{HartleyGreen} and subjected to slacktivist practices \cite{Morozov}. Other studies, though, show that migroblogging changes news agendas \cite{BroersmaGraham}, generates ‘sub-political’ discussion topics \cite{LindgrenLundstrom}, and may result into ‘self-generated public opinion’, as in long-text blogs \cite{KoltsovaKoltcov}; we share the latter opinion. Also, scholars have called Twitter the quickest platform for expression of public sentiment \cite{BrunsBurgessCrawford}; this is why we cannot dismiss the Twitter influencers’ potential of shaping the discussions.

The third limitation poses the question of structural limitations for all-involving discussion. Twitter networks resemble information-sharing ones and not offline social networks \cite[p.~264]{BastosRaimundoTravitzki} and, thus, privilege ‘gatewatchers’ \cite{Bruns} or ‘gateways’ \cite{BastosRaimundoTravitzki} who multiply and disseminate information from both outside the network and from influ- encers, as Twitter networks demonstrate ‘highly skewed distribution of followers and a low rate of reciprocated ties’ \cite[p.~263]{BastosRaimundoTravitzki}. But in our paper we try to see whether active users with a particular position towards migrants get to the influencer lists; later, we may check whether the structure of the network played a role in their promotion.

\subsubsection{3. Absolute Figures vs. Centrality Metrics in Detecting Influencers}

In our earlier case study, we have showed that existing research actually rarely links absolute-figure metrics to SNA-based graph-dependent metrics (centralities) \cite{BodrunovaBlekanovMaksimov}. Extremely wide SNA literature is dedicated to predicting the key nodes in discussion networks; a smaller bunch of works applies the network-based metrics to Twitter discussions (as examples, see \cite{DuboisGaffney,AlmindIngwersen}), using not only single metrics but also their combinations \cite{KwakLeePark,GonzalezBailonBorgeHolthoeferMoreno} and case-specific derivatives \cite{MairederWeeksDeZuniga}.

Several of these works have focused on the institutional nature of the key network nodes; mainly, researchers are looking at whether media continue to be information flow hubs -- and express significant doubts. Thus, authors \cite{BastosRaimundoTravitzki} have shown that it was content that mattered for generating ‘highly replicated messages... without relying on the activity of user hubs’ \cite[p.~260]{BastosRaimundoTravitzki}, and that the role of media outlets in forming retweet waves was much exaggerated \cite[p.~269]{AlmindIngwersen}. Other authors \cite{DuboisGaffney} have shown that media remained influencers only by indegree and eigenvalue metrics; another research group \cite{HilbertVasquezHalpern} has demonstrated that new groups of influential users join experts and media. Thus, we expect that media would still be among network-detected influencers but they will not be the leading ones.

But at the same time, research that uses absolute-figure metrics provides a more nuanced picture on who is labeled as influencer, which metrics to use for detecting them, and whether institutional (political, media, economic etc.) users remain among them. Earlier, we have shown that the majority of researchers have named \(N_{retweets}\) the most efficient metric to detect an influencer \cite{BodrunovaLitvinenkoBlekanov2016}. But other authors warn that \(N_{retweets}\) cannot actually help differentiate between ‘having a following’ due, e.g., a big number of tweets by a given user or a celebrity -- and ‘being seen as an expert’ whose tweets are genuinely shared more than those of other users \cite[p.~1263]{DuboisGaffney}; \(N_{tweets}\) has been shown to be a mediating factor for other metrics \cite{Jungherr}. This understanding corresponds to our ‘marketing’ vs. ‘deliberative influencers’ division.

Also, most of these works insist that institutionalized users remain highly influential in how discussions develop. Of course this partly due to another view on influencer as on ‘prestigious actor whose position is approved by the audience and who initiates more support than criticism’ \cite{Adam}, which does not take into account the user’s position in the network. In this line of research, several case studies have proved that Twitter strengthens the pre-existing hierarchies with media and political leaders \cite{WuHofmanMason,VaccariValerianiBarbera,JungherrJuergens}, as well as experts and long-established institutions \cite{FoxZickuhrSmith,Page}, still playing the key role in information dissemination. Using a composite measure named ‘mentions’ (that comprises several absolute-figure metrics), author \cite{Vis} shows that journalists and mainstream media were dominating the top100 accounts in the Twitter coverage of the UK 2011 riots. Similar results were received for New Zealand \cite{Bruns2014} where, of top16 Twitter accounts by retweet \& comment, 11 were institutional and included media. This may happen because journalists often retweet other journalists \cite{LotanGraeffAnanny}, but this can hardly influence the top lists selected out of several hundred thousand users
.
So far, only rare works tried to combine or juxtapose the absolute and network-based metrics \cite{BodrunovaLitvinenkoBlekanov2016,Adam,GruzdRoy,XuSangBlasiola}. To see the correlations between the two types of metrics, we will use the scheme we had elaborated earlier \cite{BodrunovaLitvinenkoBlekanov2016}. We will use both activity/connectivity and absolute/network-based divisions to describe the metrics we will juxtapose. Thus, the metrics we will use for top list formation are the following:

\begin{itemize}
	\item activity, absolute: \(N_{tweets}\);
	\item activity, network-based: outdegree centrality;
	\item connectivity, absolute: \(N_{retweets}\), \(N_{comments}\), \(N_{recom}\) -- retweets and comments combined (as it was conceptualized in \cite{Bruns2014});
	\item connectivity, network-based: indegree, betweenness, and pagerank centralities.
\end{itemize}

\subsubsection{4. The Cases, Research Hypotheses, and Methodology of the Study}

To formulate our research questions more precisely, we also need to take into con- sideration the context of the cases under scrutiny. The relevant aspects include the expectations from the Russian and German Twittersphere formulated in the existing research; the description of the cases; the societal cleavages inside it. This is done to help form our expectations of who would be the influencers within the discussions.

\paragraph{4.1 The Inter-ethnic Conflicts in Germany and Russia and Their Social and Communicative Context} 

To explore the issues described above, we have focused on comparable conflictual \textit{ad hoc} discussions. The topic of migrant crime and the following anti-migrant uprising provides cases that possess the following features: they have a rapid violent trigger, cause social polarization and street action, involve authorities, and get to national Twitter trending topics.

\textit{The German case.} According to statista.com, the number of regular Twitter users in Germany in 2015 was only 1.73 million (2\% of the population), with about twice that number using it occasionally \cite{Kissane}, and it seems not to grow since 2010 \cite{TumasjanSprengerSadner}.

The German media system belongs to the democratic corporatist model \cite{HallinMancini} with a strong tradition of freedom of expression combined with the tradition of corporatism, including the leading role of public TV. Also, the press market, despite the adherence to the notion of objectivity, is characterized by a degree of political polarization and media-political parallelism, as well as by powerful tabloids. The German Twitter, though being an undeniable news alert arena and one of the political facilitation tools in mass actions, has generated virtually no research on its structure and discussion features. Thus, our expectations are based on the overall structure of the media market, traditions of balanced reporting and public deliberation, and specific features of the German media market and civil society stated above. Thus, we expect German state actors and NGOs to be present in the discussion and perhaps even to become the discussion centers; supra-national mainstream media (like foreign newspapers or Euronews TV channel) will also be present.

The event under our scrutiny is the Köln mass harassment. During the New Year’s Eve 2015/2016 in Köln (Cologne), numerous sexual assaults were committed on women by groups of young men, allegedly mainly from the North African and Arab countries. The attacks triggered a new wave of far-right protests of the ultraconservative party ‘Alternative für Deutschland’ (AfD) and anti-migrants movement ‘Pegida’. Public support for refugee-welcoming politics of Angela Merkel has significantly dropped within several weeks \cite{Dearden}. The national media reported about the attacks with delay of several days, which led to new accusations of the mainstream media in pro-migrant bias and to escalation of the debate about the ‘lying press’ (‘Lügenpresse’) by AfD and Pegida. The Parliamentary Assembly of the Council of Europe reacted with a debate under urgent procedure on January 25, 2016, stating in Resolution 13961 that ‘media hold an important responsibility to report on objective facts without stigmatization. Partial, late or dishonest media reporting on crimes can feed in con- spiracy theories and fuel hatred against a part of the population. It can also contribute to mistrust in the authorities and the media’.

The Russian Case. After 1991, the Russian media system has seen fundamental transformation but, in political respect, it remained mostly post-Soviet \cite{Vartanova}. Today, the country’s media sphere is highly fractured along the lines of value-based cleavages between small cosmopolitan hyper-urban and huge mid-urban post-Soviet population clusters \cite{BodrunovaLitvinenko2013,BodrunovaLitvinenkoGavraYakunin}. Online, this division shows up in formation of platform-wide political echo chambers \cite{BodrunovaLitvinenkoGavraYakunin}, with the Russian Facebook serving as the best example.

Research on Russian Twitter is as well extremely scarce; it is hard even to estimate the overall use of Twitter in Russia. As for August 2015, figures varied from 8 to 11 mln subscribers, of which around 50\% seemed to be active users (used Twitter once a month or more, as estimated by TASS). The existing research on Russian Twitter provides mixed evidence on whether Twitter in Russia can play a role of an ‘opinion crossroads’. Several works have proved that political representation of pro- and anti-‘systemic’ actors on the Russian Twitter is virtually equal \cite{Greene,NikiporetsTakigawa}, but at the same time others stated that topic-based clusters with clear political bias were evident in earlier years \cite{BarashKelly}. Importantly, the latter work also stated the absence of any distinct nationalist clusters in the Russian blogs and on Twitter in particular. Except for our earlier works \cite{BodrunovaLitvinenkoBlekanov2016,BodrunovaBlekanovMaksimov}, there was no substantial attempt to study the nature and structural roles of influencers on the Russian Twitter. The newest work \cite{SanovichStukalPenfoldBrown} also proved extremely high ‘botization’ of political topics on the Russian Twitter; this is why we take as case the discussion of almost 4 years ago when it was not yet the case.

The events we analyze -- anti-migrant bashings in Biryuliovo district of Moscow -- happened in October 2013 and were in Twitter Trending Topics for over two days. The timeline included akilling of a Muscovite Egor Scherbakov by an Uzbek named Orkhan Zeinalov, the bashings at Biryuza trade center, its warehouse and the surroundings in Biryuliovo where the alleged killer should have resided along with many of his fellows, and the subsequent police street actions, several ‘gatherings’ of the locals, and arrest and trial of the suspect; the events were also accompanied by statements of federal and Moscow authorities. Thus, the actors that we may trace were: authorities (federal, Moscow, local); police; eyewitnesses; migrants. As the case was reported in federal and local media, we also expect high level of media involvement.

Expectations in Terms of Influencer Structure and Positioning. In both countries, we expect institutional actors to dominate top user lists, despite high levels of eyewitness posting. We expect national and local authorities, media, and police to be the main influencers; to a smaller extent, we expect NGOs and other pro-migrant speakers to form the lists. According to earlier research, we expect neither nationalists nor migrants to be highly influential in the Russian case, while we may expect anti-migrant citizens to show up in both cases; but taking into consideration the traditions of public discussion in Germany, we expect users and media to be mostly neutral.

\paragraph{4.2 Research Questions} 

Based on everything aforementioned, we have formulated four research questions.

\paragraph{RQ1. Do the users that post most become discussion centers in both absolute and network-based metrics? That is, does \(N_{tweets}\) significantly correlates to \(N_{retweets}\), \(N_{comments}\), \(N_{recom}\), outdegree, betweenness, and pagerank centralities in both cases?}

\paragraph{RQ2. Do institutionalized users dominate over ordinary users by both activity and connectivity metrics? Do the patterns of institutionalization differ a lot? We expect that, for Russia, pro- and anti-migrant users (like NGOs and nationalists) will be absent from both the lists of active (\(N_{tweets}\), indegree) and ‘central’ (betweenness and pagerank) top user lists; but in Germany we expect more political actors, social organizations, and NGOs to form the lists.}

\paragraph{RQ3. Do media occupy significant place in top user lists? Within the lists, are media of all views are represented?}

\paragraph{RQ4. Are institutional and most non-institutional top users neutral in terms of taking sides in the conflict?}

\paragraph{4.3 Methodology and Research Process} 

To collect the discussion bulk, we conducted vocabulary-based web crawling; then, we reconstructed the discussion web graphs. For this, we developed a specialized web crawler \cite{BlekanovSergeevMartynenko}. We used our own software to overcome limitations common for openly available API-based analogs; our algorithm is human-like, which allowed for unfolding of the discussion in the past and trespassing the time and quantity upload limits. To form the vocabularies, we first collected relevant keywords and hashtags at trendinalia.com and double-checked the lists on two other Twitter trending topics trackers.

Then, we added more hashtags based on manual snowballing of tweets in over 1,000 tweets for both cases. The vocabularies for Russia included 6 main hashtags/keywords, and for Germany -- 15 hashtags/keywords.

For Russia, the research period chosen was October 1 to 31, 2013, to capture the outburst of the discussion and its long tail. 3,574 users with 10,715 posts were identified as a result of crawling and formed the core dataset. One step further in crawling was made to identify those who commented or retweeted the collected tweets, to calculate properly the number of comments and retweets; this returned 12,040 users. For Germany, a similar strategy of uploading (January 1 to 31, 2016) discovered a significantly bigger discussion of 12,382 users involved with 64,874 posts posted; one step further returned 40,117 users.

For comparison, we used the user lists from the core datasets, but the data on commenting and retweeting for individual users are taken from the bigger datasets.

Then, we have conducted the following procedures:
\begin{enumerate}
	\item To calculate the SNA metrics, we reconstructed the discussion graphs. The graphs themselves were non-directed (as we were not interested in directions of interactions, only in numbers), but our data allowed for calculating in-/outdegrees independently.
	
	\item From the graphs, we received the values for the chosen variables: \(N_{tweets}\), \(N_{retweets}\), \(N_{comments}\), \(N_{recom}\), indegree, outdegree, betweenness, and pagerank for the core datasets.
	
	\item After that, we formed additional dataset of users with \(N_{tweets} \geq 10\) to include only those who actively participated in the discussion. This was done in order to exclude the discussion ‘long tails’ with large number of users who, though, posted only a few tweets each and, thus, would distort the results true for active users. For Germany, the list included 1,211 users; for Russia, only 178 users.
	
	\item Then, we conducted descriptive statistics (Spearman rho) to see to what extent the chosen metrics correlate in the core datasets and the datasets with \(N_{tweets} \ge 10\) (see Tables~\cref{tab:spearmanCorrelationRussiaCore} and~\cref{tab:spearmanCorrelationRussiaActive} for Russia and Tables~\cref{tab:spearmanCorrelationGermanyCore} and~\cref{tab:spearmanCorrelationGermanyActive} for Germany). We considered the use of Spearman’s rho appropriate despite we realized that absolute figures, including \(N_{tweets}\) and in-/outdegree values may play a role in formation of other centrality metrics, and we expect them to correlate, but it is the strength of correlation that we will be looking at. Also, as stated above, betweenness and pagerank are network-dependent, while in-/outdegree are calculated as absolute numbers of user interactions.
	
	\item We manually checked the user top lists, to assess the patterns user transposition from the lists by activity metrics to the lists by connectivity metrics, and those by absolute figures -- to those by network metrics; we also marked their institutional belonging and pro-/anti-migrant position (see Figs.~\cref{fig:topUsersRussia} and~\cref{fig:topUsersGermany}). To do so, we checked a user’s self-description, the collected tweets, and the user’s tweets closer to nowadays.
\end{enumerate}

The results assessed in comparative perspective are presented below.

\begin{table}[ht]%
	\centering
	\caption{Spearman’s correlation between activity and connectivity measures in Russia for the core dataset \((N_{users} = 3,574)\).}%
	\label{tab:spearmanCorrelationRussiaCore}% label всегда желательно идти после caption
	\begin{adjustbox}{width=1\textwidth}
		\small
		\begin{tabular}{ c  c  c  c  c  c  c  c  c }% Вертикальные полосы не используются принципиально, как и лишние горизонтальные (допускается по ГОСТ 2.105 пункт 4.4.5) % @{} позволяет прижиматься к краям
			\toprule
			& Tweets & Retweets & Comments & Recom & Indegree & Outdegree & BC & PRC \\
			\hline
			Tweets & 1,000 &  &  &  &  &  &  & \\
			Retweets & ,472** & 1,000 &  &  &  &  &  & \\
			Comments & ,408** & ,482** & 1,000 &  &  &  &  & \\
			Recom & ,489** & ,893** & ,753**  & 1,000 &  &  &  & \\
			Indegree & ,486** & ,226** & ,219** & ,238** & 1,000 &  &  & \\
			Outdegree & ,345** & ,168** & ,154** & ,179** & ,430** & 1,000 &  & \\
			Betweenness & ,410** & ,215** & ,200** & ,227** & ,493** & ,532** & 1,000 & \\
			Pagerank & ,403** & ,186** & ,185** & ,194** & ,808** & ,453** & ,513** & 1,000\\
			\bottomrule
		\end{tabular}%
	\end{adjustbox}
\end{table}

\begin{table}[ht]%
	\centering
	\caption{Spearman’s correlation between activity and connectivity measures in Russia for the
		dataset of active users, \(N_{tweets} \geq 10 (N_{users} = 178)\).}%
	\label{tab:spearmanCorrelationRussiaActive}% label всегда желательно идти после caption
	\begin{adjustbox}{width=1\textwidth}
		\small
		\begin{tabular}{ c  c  c  c  c  c  c  c  c }% Вертикальные полосы не используются принципиально, как и лишние горизонтальные (допускается по ГОСТ 2.105 пункт 4.4.5) % @{} позволяет прижиматься к краям
			\toprule
			& Tweets & Retweets & Comments & Recom & Indegree & Outdegree & BC & PRC \\
			\hline
			Tweets & 1,000 &  &  &  &  &  &  & \\
			Retweets & ,461** & 1,000 &  &  &  &  &  & \\
			Comments & ,417** & ,753** & 1,000 &  &  &  &  & \\
			Recom & ,453** & ,954** & ,893**  & 1,000 &  &  &  & \\
			Indegree & ,443** & ,444** & ,354** & ,429** & 1,000 &  &  & \\
			Outdegree & ,182* & ,190* & ,233** & ,208** & ,505** & 1,000 &  & \\
			Betweenness & ,335** & ,320** & ,344** & ,340** & ,646** & ,754**  & 1,000 & \\
			Pagerank & ,644** & ,414** & ,437** & ,357** & ,420** & ,873** & ,850** & 1,000\\
			\bottomrule
		\end{tabular}%
	\end{adjustbox}
\end{table}

\begin{table}[ht]%
	\centering
	\caption{Spearman’s correlation between activity and connectivity measures in Germany for the core dataset \((N_{users} = 12,382)\).}%
	\label{tab:spearmanCorrelationGermanyCore}% label всегда желательно идти после caption
	\begin{adjustbox}{width=1\textwidth}
		\small
		\begin{tabular}{ c  c  c  c  c  c  c  c  c }% Вертикальные полосы не используются принципиально, как и лишние горизонтальные (допускается по ГОСТ 2.105 пункт 4.4.5) % @{} позволяет прижиматься к краям
			\toprule
			& Tweets & Retweets & Comments & Recom & Indegree & Outdegree & BC & PRC \\
			\hline
			Tweets & 1,000 &  &  &  &  &  &  & \\
			Retweets & ,470** & 1,000 &  &  &  &  &  & \\
			Comments & ,481** & ,681**  & 1,000 &  &  &  &  & \\
			Recom & ,503** & ,941** & ,864**  & 1,000 &  &  &  & \\
			Indegree & ,485** & ,960** & ,777** & ,958** & 1,000 &  &  & \\
			Outdegree & ,368** & ,309** & ,431** & ,369** & ,378** & 1,000 &  & \\
			Betweenness & ,446** & ,608** & ,654** & ,662** & ,687** & ,748** & 1,000 & \\
			Pagerank & ,395** & ,786** &,767** & ,843** & ,861** & ,285** & ,629** & 1,000\\
			\bottomrule
		\end{tabular}%
	\end{adjustbox}
\end{table}

\begin{table}[ht]%
	\centering
	\caption{Spearman’s correlation between activity and connectivity measures in Germany for the
		dataset of active users, \(N_{tweets} \geq 10 (N_{users} = 1,211)\).}%
	\label{tab:spearmanCorrelationGermanyActive}% label всегда желательно идти после caption
	\begin{adjustbox}{width=1\textwidth}
		\small
		\begin{tabular}{ c  c  c  c  c  c  c  c  c }% Вертикальные полосы не используются принципиально, как и лишние горизонтальные (допускается по ГОСТ 2.105 пункт 4.4.5) % @{} позволяет прижиматься к краям
			\toprule
			& Tweets & Retweets & Comments & Recom & Indegree & Outdegree & BC & PRC \\
			\hline
			Tweets & 1,000 &  &  &  &  &  &  & \\
			Retweets & ,461** & 1,000 &  &  &  &  &  & \\
			Comments & ,417** & ,753** & 1,000 &  &  &  &  & \\
			Recom & ,453** & ,954** & ,893**  & 1,000 &  &  &  & \\
			Indegree & ,443** & ,444** & ,354** & ,429** & 1,000 &  &  & \\
			Outdegree & ,182* & ,190* & ,233** & ,208** & ,505** & 1,000 &  & \\
			Betweenness & ,335** & ,320** & ,344** & ,340** & ,646** & ,754**  & 1,000 & \\
			Pagerank & ,644** & ,414** & ,437** & ,357** & ,420** & ,873** & ,850** & 1,000\\
			\bottomrule
		\end{tabular}%
	\end{adjustbox}
\end{table}

\begin{figure}[ht]
	\centerfloat{
		\includegraphics[scale=1.0]{topUsersRussia}
	}
	\caption{Institutional belonging and pro/contra-migrant positioning of top users in Russia.}\label{fig:topUsersRussia}
\end{figure} 

\begin{figure}[ht]
	\centerfloat{
		\includegraphics[scale=1.0]{topUsersGermany}
	}
	\legend{
		Regular -- ‘ordinary user’ \\
		Bold -- institutional user/representative \\
		Bold italic -- media account/journalist \\
		Italic -- ‘Twitter media’ \\
		Green -- strong/institutional support of migrants (absent from picture)\\
		Light green -- weak support of migrants \\
		White -- neutral user \\
		Light orange -- weak anti-migrant attitude \\
		Orange -- strong anti-migrant attitude/nationalist account.
	}
	\caption{Institutional belonging and pro/contra-migrant positioning of top users in Germany.}\label{fig:topUsersGermany}
\end{figure} 

\subsubsection{5. Results}

\paragraph{RQ1. Number of Tweets and Discussion Centers.} For both countries, the more users post, the more likely they become both ‘marketing’ (by \(N_{retweets}\) and \(N_{comments}\)) and ‘deliberative’ (by betweenness and pagerank) influencers, despite the difference in the size of datasets. The correlations remain in place for all the metrics for full and active-user datasets, despite the fact that elimination of ‘the crowd’ significantly drops outdegree for the Russian discussion and high values for \(N_{comments}\) and betweenness centrality for Germany. All in all, strength of the correlations remains comparable for all four of our datasets (though higher for Germany for the two aforementioned metrics), which might be telling of the nature of \textit{ad hoc} discussions on Twitter, but can also support the idea of the mediating role of \(N_{tweets}\); this needs further exploration.

But if we look closer at correlation values, we will see that outdegree correlates more weakly with \(N_{tweets}\) than other metrics throughout our data; this might mean that tweeting a lot does not provide for necessarily becoming commented or retweeted; the value is never higher than 0,418, and thus, the correlation is weak enough. Also, only in the German case, \(N_{retweets}\) matters for getting higher betweenness and pagerank, while in Russia their correlation is much weaker. But what, instead, seems to be important for becoming an influencer is a user’s indegree, that is -- how many users have interacted with you. This parameter becomes more important for becoming an authoritative node than the number of tweets, retweets, or comments -- the metrics that many works stated as markers of influencers. For all four of our datasets, the strength of ties between indegree and pagerank is 0,804 or higher. That is, a successful strategy within an \textit{ad hoc} discussion might be not to comment many times or get into a long meaningful discussion but to make a bigger number of users comment on you, perhaps by commenting them as well. On the other hand, outdegree does not seem to matter much for both betweenness and pagerank; this brings us to the conclusion that attractive content (which makes users interact with it) may be more important within such discussions than user activity.

\paragraph{RQ2. Institutionalization of the Discussion.} In Figs.~\cref{fig:topUsersRussia} and~\cref{fig:topUsersGermany}, we have marked institutional users, including media, bold (media -- bold italic). In general, the picture is similar in both countries and may be described as ‘liberal media against individual nationalists’. What is similar (and striking) in both countries is the absence of the much-awaited national/regional authorities, as well as NGOs and human rights watchers. We cannot prove dominance of institutional accounts over ‘ordinary users’, unlike in previous research \cite{XuSangBlasiola}, as we find only two politicians and one account of Public Advisory Chamber among the Russian top users, and two NGO-like organizations in the German top user lists.

Also, we cannot prove absence of nationalists as top users: in both countries, they are not only present but demonstrate blooming activity, and if in Russia they are most active in commenting, in Germany they lead both tweeting and commenting activities, all of them being non-institutionalized.

We call the picture similar, despite that, from Figs.~\cref{fig:topUsersRussia} and~\cref{fig:topUsersGermany}, the German discussion seems highly radicalized and the Russian one shows a lot of neutral users. But this may be explained by two factors. First, Biryuliovo happened over three years ago, and assessment of the accounts of top users does not bring over a lot of anti-migrant posts; this has cast its impact upon our allocation of users as neutral or biased. But we need to state that we have discovered a dominant mood among Russian ‘ordinary people’ which may be described as ‘angry patriotism’: today, such users (over a dozen in the Russian top lists) express ‘patriotic’ views like supporting Donbass population or tweeting on national pride (e.g. on leading industries like aviation, military equipment etc.) but at the same demonize the current country’s leaders as corrupt and ineffective. Thus, these users are highly politicized, no less than the German users; it is just not always possible to deduce their attitude from the tweets of the discussion (especially if they retweeted other accounts) and today’s tweets. Second, some media accounts in the Russian list (like lifenews\_ru, izvestia\_ru, or RT\_russian) were marked neutral, as we could not find direct proof of their anti-migrant bias, but their overall tone is pro-establishment, and thus their position fluctuates from supporting the state views on open visa regime for Central Asian post-Soviet ethnicities to populism of attaching social and cultural threats to their communities. Having said this, we can consider the situation similar indeed in both cases, as it is highly polarized, full of political criticism, and intolerant.

\paragraph{RQ3. The Place of Media in the Discussions.} Media, indeed, occupy a significant place in the discussion and represent a variety of political views and positions. Unlike on Russian Facebook, in this discussion both pro-establishment and highly oppositional media (ru\_rbc, GraniTweet), as well as foreign liberal media and journalists (MaloverjanBBC, fulelo, SvobodaRadio) are present, and the liberal-oppositional media show their efficacy, as they become retweeted and commented by many people without tweeting a lot. In Germany, it is mainstream media, and mostly newspapers, that also become influencers without posting a lot; they get retweeted and commented in general and by a lot of users in particular, and gain high pageranks. But even if so, media do not outperform nationalist users, and they do not get high betweenness centrality, which means that they do not play the role of ‘information mini-hubs’ as the basic nodes of the online public sphere. They remain authoritative (especially in Germany), but the niche of ‘deliberative connectors’ remains free and is occupied by the most polarized users. Thus, the ‘opinion crossroads’ may be there in terms of representation of views within the whole discussion but it is still a question whether the opposing views actually have a chance to meet.

\paragraph{RQ4. Neutrality of Top Users.} As already stated above, neutrality of users cannot be proven, especially in case of Germany. In Russia, general negative politicization of the audience goes along with nationalist and pro-nationalist views, and in Germany the discussion after a major public harassment is shaped not by the forces countering intolerance but by openly anti-migrant discussants; in both cases, it is individuals that polarize the discourse against re-settlers and media that counter this -- even if due to different reasons. Thus, for most of the German media, supporting immigrants is a non-valent issue, and expressing an alternative position would amount to a scandal. In Russia, the division between pro-establishment and oppositional media is also true for the migration issues, and thus liberal-oppositional media support their political standing by expressing pro-migrant views.

\subsubsection{6. Conclusion}

We have looked at two \textit{ad hoc} discussions on violent inter-ethnic conflicts, namely the Biryuliovo bashings in 2013 and mass harassment in Cologne in 2016, to see whether in such discussions user activity leads to higher positions within the discussion network and higher connectivity. Along with this, we assessed the substantial features of top users, such as their institutional status and opinion positioning.

Despite the differences in samples, we have managed to show that comparing influencers is possible, and there are patterns in the structure of influencing that are similar. The main methodological finding is that, in both discussions, the number of users involved mattered more for becoming an influencer (by BC and PRC) than the number of actual tweets, retweets, or comments received by a user.

Though direct comparisons were not always possible by our methodology, we have found more similarities in the two cases than we had expected. Thus, in both countries, the situation may be described as opposition ‘liberal media vs. nationalist users’, and the absence of both authorities and NGOs is striking. Media do become influencers, but in terms of authority (or ‘marketing’ approach) rather than in deliberative terms, as they do not get high betweenness centrality and thus may have difficulties in performing the roles of shapers of information flows. In both countries, the discussion was highly opinionated and emotionally heated even within several weeks after the events, and seemingly higher neutrality of the Russian users was compensated by their overall politicization and rebuttal modus.

Thus, as to the question of whether and how ‘opinion crossroads’ are forming, there is evidence that, in general, left-right or in-system/oppositional views are well represented by media within the discussions. But virtual absence of pro-migrant institutions and opposition of liberal media to pro-nationalist ‘ordinary users’ shows that, in both countries, the discussion is far from being balanced, rational, and inclusive.

\subsection{Public Opinion Dynamics in Online Discussions: Cumulative Commenting and Micro-level Spirals of Silence}\label{subsec:ch2/sec4/sub5}

\subsubsection{1. Introduction}

Social media have become a place where the bulk of grassroots political discussion takes place. Social networking sites, online platforms, and messengers have become a focus of scholarly attention in their relation to public opinion formation and spread, which, in its turn, directly affects the quality of offline democratic deliberation.

The models that described opinion formation in the era of traditional media have been applied to online discussions, with varying degree of proof and support. Thus, the idea of opinion leaders \cite{KatzLazarsfeld} has been re-interpreted to describe the practices of influencers \cite{Gillin,BakshyHofmanMason,BodrunovaLitvinenkoBlekanov2017}, and the older idea of echo chambers \cite{Key} has transformed into filter bubbles \cite{Bruns2019}.

Along with this, today’s public spheres are, in general, described as dissonant \cite{Pfetsch} and discontinued \cite{SmoliarovaBodrunovaBlekanov2020}. The dynamics of formation of dominant opinions have been conceptualized in several theoretical models, among them, most famously, in the idea of spiral of silence \cite{NoelleNeuman1974}. In this conceptualization, under pressure of fear to be excluded, supporters of one candidate become gradually silenced by the supporters of the other one, and the latter are capable of doing this just by expressing their opinion more massively.

However, when applied to the dissonant and disruptive online discussions of today where cumulation of support is often accompanied by communicative aggression \cite{Sidorov,BodrunovaLitvinenkoBlekanov2021}, Noelle-Neumann’s concept and similar conceptualizations needs to be reconsidered and re-tested. This is especially because there is growing evidence that online discussions are the place where silent majority flourishes \cite{MustafarajFinnWhitlock,ChenLiYao,McKeeverMcKeeverHolton,MaiShanBai}.

Today, instead of dialogue and polylogue in the forms of opinion exchange, we witness cumulative modes of discussion prevail \cite{BodrunovaLitvinenkoBlekanov2021}. To re-conceptualize the ways public communication exists and develops in online milieus, we suggest the idea of \textit{cumulative deliberation} on which we elaborate below (see Sect. 2). Our previous works, as well as many works before ours, have shown that cumulation may be regarded as a universal mechanism of public opinion growth and polarization \cite{BodrunovaBlekanovSmoliarova}. E.g., in the form of echo chambering, it works on many levels in globalized discussions, such as language, hashtagging, and sentiment \cite{BodrunovaBlekanovKukarkinCH}. We have also demonstrated that, on Russian oppositional YouTube, users hardly form what might be called ‘a discussion’; they, rather, express themselves \textit{urbi et orbi} without entering any dialogue with particular users, despite their high polarization in political terms and an extremely high level of aggressive and hate speech discovered \cite{BodrunovaLitvinenkoBlekanov2021}.

The current paper further fosters the idea of discovering patterns of cumulative deliberation in oppositional discourse and, this time, uses a year-long dataset from Belarusian oppositional YouTube of 2018. We have collected comments left by users under the videos of six salient Bealrusian oppositional accounts (one of which, Nexta, has become world-renowned during the Belarusian protest of 2020). We have identified users who linked these accounts by posing to at least five of them, and have analysed this core of commenting for the features important for the deliberative process, namely orientation to dialogue \cite{Ksiazek}, aggression \cite{BodrunovaLitvinenkoBlekanov2021,KsiazekPeerZivic}, and political criticism. In the latter we have taken in mind the difference between uncritical, policy-critical, and leadership critical publics, as it was suggested and explained for authoritarian societies \cite{Toepfl}.

In particular, we focus on two research questions. First, we ask whether (and how exactly) the deliberative features mentioned above accumulate in the discussion and affect its dynamics in the core of the discussion (which is, in the Belarusian case, the cross-commenter conglomerate) on various time spans. Second, we look at whether the deliberative feature of the discussion core affect the bulk of the discussion. This paper is only a trial that shows the necessity of further exploration of cumulative patterns of online discussions and points out to possible bigger cumulative effects that deserve exploration and explanation.

The remainder of the paper is organized as follows. Section 2 elaborates on cumulative deliberation as a concept and suggests how dynamics of the discussions might be affected by cumulative patterns of opinion formation. Section 3 poses the research questions. Section 4 describes sampling, data collection, dataset limitations, and methods of data assessment. Section 5 provides the results; Sect. 6 discusses them and concludes the paper.

\subsubsection{2. Cumulative Foundations and Deliberative Features of Online Opinion Formation}

\paragraph{2.1 Cumulative Patterns in Online Opinion Formation}
As already stated above, the idea of cumulation as a basic principle of how online public opinion forms has come from two streams of literature.

\paragraph{Cumulative Effects in Online Opinion Formation.} The first one focuses upon patterns	of winning the public debate -- that is, majority/minority formation in public opinion and electoral research. In political sociology, the works that have demonstrated that people may change opinion just knowing or supposing what others might think have been numerous \cite{NowakSzamrejLatane}, and they have posed a question of whether opinion change is a result of persuasive external arguments, including mediated ones.

However, despite their evident high impact upon the rational choice debate, relatively few important works linked individual/public opinion shifts to communication and deliberation patterns on aggregated levels. Well-known are, e.g., the works on ‘spiral of silence’ by Elizabeth Noelle-Neumann \cite{NoelleNeuman1974} or Latané’s social impact theory that emerged, i.a., from studying interaction patterns \cite{NowakSzamrejLatane,Latane}. But it was not before the advent of social media that the scholars have received a chance to study the multi-level communicative interactions.

There is one significant factor that distinguishes the pre-Internet era from today’s communication. In online communication, the iceberg of what could previously remain an oral utterance exchange and disappear forever after being pronounced, remains written and accessible by the newcomers who join the communicative milieu later in time. Potentially, this casts impact upon the newcomers and their personal, seconds-long opinion formation, in accordance with the social impact theory; this remains highly understudied. At the same time, there is a substantially bigger corpus of works upon various discussion-level effects of opinion cumulation, such as echo chambering and filter bubbles studies or research on influencers who become such by cumulation of user support. (Of these studies, we will focus on user comments studies -- see below).

Thus, cumulative effects in communication have long been demanding generalization
and systematization -- or at least a proper acknowledgement of the role of cumulation in
opinion formation online.

\paragraph{Dissonant Public Spheres: Getting Rid of Excessive Normativity.} The second stream of literature is that on public sphere as a normative space of (arguably) consensusoriented discussion of conflictual issues in public domain, as well as on deliberation as a special form of public communication, an imagined round-table process of collection and adjustment of individual, group, and institutionalized opinions that leads to decision on a given issue. For several decades, public sphere studies have almost exclusively, either adhering or critically, oriented to the works by Juergen Habermas.

The Habermasian view on deliberation is known for its high normativity, as communicative actions need to be rational, consensus-oriented, and civil in order to count for public decision-making. Despite the overwhelming amount of academic criticism from both left and right parts of the political spectrum, as well as the shifts in Habermas’s own conceptualization of public sphere and deliberation, the normative notion of orientation to consensus (or to struggle, in the leftist view \cite{Mouffe}) and high \textit{deliberative quality} of public communication was not challenged.

Again, it was not before the advent of blogs and social networking sites that the scholars questioned the belief that people come into a discussion either ‘to agree or to argue’ \cite{YardiBoyd} -- that is, with a certain deliberative goal. Today, this notion has been shaken by description of online public spheres as dissonant -- that is, not aiming to consonance and consensus; public spheres where myriads of actors express themselves simultaneously \cite{Pfetsch}. There are also works that describe public spheres as discontinued, whose actors constantly change and re-aggregate \cite{SmoliarovaBodrunovaBlekanov2020}, and, in general, disrupted \cite{BennettPfetsch}.

Moreover, recent research suggests that users might participate in dialogical forms of communication much more rarely than deliberation theory suggests \cite{BodrunovaLitvinenkoBlekanov2021}, and that their motivation is neither to agree nor to argue in order to come to a conclusion, but more to express themselves, solidarize with an earlier opinion or negate it without having in mind any longer-term effects. Yet, even such communicative activity, ‘irrelevant, aggressive, and stupid’ \cite[p.~659]{SchultesDornerLehner}, seemingly aimless in deliberative terms and dismissed as slacktivism \cite{Morozov} and ‘sofa warrior’ behavior, acquires importance if seen from the cumulative-effects viewpoint.

The plea for finding new grounds for assessment of public spheres influenced by massive micro-level communication, united with the view upon cumulation as an overarching pattern of how opinion aggregates and dissipates online, has led us to the concept of \textit{cumulative deliberation}. This concept is, in essence, an alternative vision on public deliberation that covers by one umbrella many cumulative effects scattered around academic works on older and newer media and opinion formation via them. Cumulative deliberation allows for a closer-to-reality look at how people talk online, as it lets avoid the excessive normativity and merit the smallest forms of online user activity: each individual comment, like, share, or mention matters for cumulative deliberation. It also aims at discovering patterns of opinion aggregation on various levels and explaining/predicting critical thresholds in cumulative communication. And the patterns of opinion cumulation, not only user intentions, should be subjected to critical assessment within the democratic perspective of public communication.

\paragraph{2.2 User Comments on YouTube and Other Social Media: The Issue of Deliberative Quality of Communicative Micro-action}

Thus, user commenting seen as deliberative activity may be put into the cumulative deliberation perspective -- that is, cumulative patterns and their effects in user commenting need to be detected and assessed, without putting the blame to users for not being rational or purposeful. Previous research shows that users’ comment ‘exchange is socially and not deliberatively motivated’, and, in reading comments, ‘entertainment dimension -- a dimension that is not usually considered to be linked to deliberation processes -- is the more stable one’ \cite[p.~798]{SpringerEngelmannPfaffinger}. This points out to the necessity for the academe to move towards accepting these motivations as normal and assessing the cumulative effects of such online posting.

Close enough to this notion, computational social scientists have tried to detect statistical effects in comment threads and discussion graph structure, such as various types of distributions like lognormal \cite{GomezKappenKaltenbrunner}, negative binomial \cite{TsagkiasWeerkampDeRijke}, Zipf’s law \cite{KumarMahdianMcGlohon}, or power law \cite{BodrunovaBlekanov}. However, cumulative effects may be of a more nuanced nature, as they combine statistical features of communication flows and contents of user utterances (for a short review on early works that include content features into statistical modelling of commenting dynamics, see \cite{WangYeHuberman}). Even successful modelling of discussion growth \cite{WangYeHuberman} often sees comments as discreet items, not cumulative conglomerates, and fails to encompass deliberative features of comments (sometimes focusing on one selected). On the other hand, communication scholars who focus on deliberative quality of commenting \cite{Rowe} are less attentive to various formal levels on which cumulation might shape the discussion, as well as to statistical features of discussions linked to platform affordances and human capacities that define posting and reading.

In our current research, we do not aim at creating overarching models of user commenting; we just want to demonstrate easily discoverable impact of deliberative features of user posts upon various structural levels and temporal fragments of online discussions. In previous studies, several user-dependent features of comments and their deliberative potential have been discussed. In a more systematic view, assessment of democratic qual- ity of deliberation online moves from evaluating the constellation of major institutional discussants to single user and features of their statements.

\paragraph{Orientation to Dialogue.} Addressing the other is an inevitable pre-requisite for any sort of dialogue, public or private. Interactivity has been conceptualized as capable of ‘contribut[ing] to shaping a democratically valuable and vivid interpersonal discourse on topics of public interest’ \cite[p.~1112]{ZiegeleBreinerQuiring}; see also \cite{BoczkowskiMitchelstein,Freelon,RuizDomingoMico}. However, user comments do not necessarily have this feature. Thus, an important division has been introduced between user-to-content (e.g., news video) and user-to-user orientation in commenting \cite{KsiazekPeerZivic}, and user-to-content comments, as a rule, do not involve other users to interaction.

While we praise this important distinction, we prefer to see formal interactivity (often marked by user mentions) as just a part of wider orientation to dialogue. Thus, our earlier study \cite{BodrunovaLitvinenkoBlekanov2021} shows that, on Russian-speaking YouTube, ‘discussion’ in the form of dia- logue hardly takes place at all, despite the political topicality of the assessed videos, and formal and substantial markers of dialogue may not correspond. E.g., even when users formally reply to other users’ comments, they construct the phrases the way that does not invite their correspondents to answer; they just solidarize with the expressed opinion or share their own experience that, in most cases, supports the view in the primary post. Or, vice versa, a comment may be openly dialogical in substance (interrogative, sarcastic, or denialist) without any formal interactivity features. In normative view, dialogue features should oppose.

\paragraph{Aggression.} When a dialogue does take place, it may easily take an anti-normative direction, as users attach each other aggressively, rather than interact in civil manner. Studies of user sentiment, emotions, affect, and aggression in social media content are overwhelming in number. Of emotional substance of comments, various types of aggressive speech are the sharpest; this, i.a., includes hate speech, obscene and tabooed lexicons, politically motivated offence, calls for violence, cyberbullying, humiliation, and harsh shaming. However, we have demonstrated \cite{BodrunovaLitvinenkoBlekanov2021,BodrunovaNigmatullinaBlekanov} that communicative aggression may play constructive roles in online discussions, such as influencing its dynamics and performing individual- and group-level psychological functions.

Today, surprisingly few works focus on how aggressive and emotional speech relates to opinion aggregation and fueling discussions. This is especially true for languages like Russian and Belarusian where obscene lexicon has grown into a highly developed sub-language.

\paragraph{Criticism.} Aggression needs to be distinguished from criticism, another pragmatic dimension of user speech. Critical assessment of reality has for centuries been acknowledged as a pre-requisite for substantial debate on nearly any issue. Political criticism may be directed to individual/group ideological opponents, as well as to the state, public bodies, persons, and policies. Recently, three types of publics have been defined for authoritarian public spheres, based on what is subjected to criticism \cite{Toepfl}, namely uncritical, policy-critical, and leadership-critical publics. Criticism to policies is expected to be milder, more discussion-oriented, and perhaps a bit more toothless, as it follows from the paper. Criticism towards the authorities brings more dangers to the speaker and, thus, demands passing a certain psychological threshold. This might imply that leadership-critical users could use aggressive speech and be less willing to take part in meaningful discussion. Both critical stances, however, may foster cumulative effects in discussions, as earlier critical comments may help diminish participation thresholds for users who comment later, and thus create avalanche effects.

Taken together, these three deliberative features may work on various levels of online discussion architectonics, from single comments to platform-based publics. We will look at their potential to foster cumulative effects in online communication on Belarusian YouTube.

\subsubsection{3. Research Questions}
Having said that, we have formulated the research questions for our study.

\paragraph{RQ1. Are the three deliberative features relate to each other in the discussion, and how?}

\paragraph{RQ2. Do the three deliberative features influence the discussion dynamics, and on what levels and time spans?}

\paragraph{RQ3. Are there any distinctive patterns of cumulative deliberation that we can point to, and on what levels and time spans?}

We do not pose any exact hypotheses, as our study is exploratory; below, we will summarize our findings according to the RQs stated above.

\subsubsection{4. Data Collection and Methods}

\paragraph{4.1 Case Description: Belarusian Oppositional YouTube}

To address the research questions, we have used the dataset of comments published within the year of 2018 to six YouTube accounts from Belarus, all of them politically oppositional, most salient on Belarusian political YouTube in 2018--2019.

This choice is explainable in many respects. First, Belarus remains one of most understudied media landscapes of Europe; as a recognized autocracy, it represents a case for analysis of authoritarian public sphere. Second, the dataset is not hashtag-based, which allows for avoiding the dissipative and affective effects \cite{Papacharissi} in how it develops; the dataset represents a ‘chronic’ discourse on the oppositional YouTube. Its oppositional orientation allows for comparing the results with our earlier work on Russia \cite{BodrunovaLitvinenkoBlekanov2021}; the pre-set political standing of channels sets the conditions for potentially conflictual and polarized opinion exchange. Al the same time, the volume of discussion is by an order of magnitude smaller than in Russia and the most studied Western countries, due to lower levels of YouTube consumption and smaller populace in general; thus, a year-long dataset remains feasible for both automated and manual research. In this respect, Belarusian YouTube is sort of an ‘ideal gas’ of studying online political communication in autocracies. To this, we add that, to avoid geographical echo chambering, we have focused not only on the accounts based in Minsk but also on those produced in other Belarusian regions. The accounts are Belsat (Belarusian oppositional satellite TV channel sponsored from abroad), Nexta (by 2021, a world-famous Belarusian blogger and protest facilitator), Garantiy NET (‘No guarantee’, Gomel; by 2021, the videos of 2018 are deleted), Narodny Reporter (‘Popular reporter’, Brest region), The Belarusian Experimental Field (positioned as all-Belarusian), and Rudabelskaya pakazuha (‘Window dressing in Rudabelka’, Gomel region).

It is worth mentioning that, in 2020, Belarus has lived through a months-long political protest unprecedented in the newest Belarusian history. The data of a year and a half before the protests erupted demonstrate the mood of Belarusians and are relatively highly politicized.

\paragraph{4.2 Data Collection and the Dataset}

YouTube crawling was performed by the web crawler created and patented by our research group \cite{BlekanovSergeevMartynenko}. The overall number of user comments in the dataset is 120,412 posted by 3000+ users.

\paragraph{4.3 Research Procedures and Methods Used}

The logic of our research demanded several steps of working upon the dataset.

\paragraph{Defining the Discussion Core.} After preliminary studying of comment structure and their appearance on actual YouTube pages, we have divided the dataset into the core and periphery. The core users were those who performed an important deliberative function of cross-account commenting, thus transforming the accounts and the discussion into one discursive field. Reconstructing the web graph of the discussion by the YifanHu Gephi algorithm, we have identified users who posted in at least five of six accounts (76 users with 8,610 comments). Answering the RQs was linked to both the discussion core itself and to relations between the core and the periphery of the discussion.

\paragraph{Forming Sub-Datasets.} We have identified the levels on which we search for cumulative patterns. We have used seven levels/time spans formed from the bulk data. The logic of it is to ‘zoom in’ from the whole discussion to the level of one particular day (step-by-step via other time spans), one video, and one user. The sub-datasets used for this paper include the following (see Table~\cref{tab:projectSubdatasets}).

\begin{table}[ht]%
	\centering
	\caption{Sub-datasets of the project.}%
	\label{tab:projectSubdatasets}% label всегда желательно идти после caption
	\begin{adjustbox}{width=1\textwidth}
		\small
		\begin{tabular}{ l  l  l  l  l }% Вертикальные полосы не используются принципиально, как и лишние горизонтальные (допускается по ГОСТ 2.105 пункт 4.4.5) % @{} позволяет прижиматься к краям
			\toprule
			\# & Time span & \makecell[l]{Unit of analysis\\/ Granger test increment} & Type of data & N of comments \\
			\hline
			(1) & Year & Any & \makecell[l]{The whole discussion\\bulk (used to detect\\the core and form\\other datasets)} & 120,412 \\
			(2) & Year & Week & \makecell[l]{The whole discussio\\core} & 8,610 \\
			(3) &  \makecell[l]{5 weeks, February 26\\to April 1} & 24 h & The most rapidly growing discussion segment around a key event &  \makecell[l]{(3a) 1,164 -- core vs\\(3b) 17,222 -- bulk} \\
			(4) &  \makecell[l]{Week, March 26\\to April 1} & 6 h &  \makecell[l]{The most commented\\week around a key\\event} & 287 \\
			(5) & Day, March 26 & 1 h &  \makecell[l]{The day next to the\\key event} &  \makecell[l]{(5a) 113 -- core vs. (5b) 1,269 -- bulk} \\
			(6) & Year & Year &  \makecell[l]{Comment\\conglomerates by each\\cross-commenter} & 8,610 \\
			(7) & Year & Year &  \makecell[l]{Comments under a\\particular video of\\March 25} & 127 \\
			\bottomrule
			\multicolumn{5}{@{}p{\textwidth}}{%
				\hspace*{2.5em}% абзацный отступ - требование ГОСТ 2.105
				\textit{Note}. The sub-datasets (3a), (4), (5a) are formed out of the sub-dataset (2) and, thus, were coded within the coding process for the sub-dataset (2).
			}\\
		\end{tabular}%
	\end{adjustbox}
\end{table}

Zooming in was performed after detecting the most actively growing segment of the year-long discussion. This segment was 5 weeks from February 26 to April 1, 2018, structured around an event of high importance to the opposition, namely the celebration of 100 years of the so-called 1918 Belarusian Republic (March 25, most active commenting taking part on March 26).

\paragraph{Coding Deliberative Features of Comments.} The contents of the sub-datasets (2) including (3a), (4), and (5a), as well as (6) and (7), were coded for the three deliberative features described above, that is, orientation to dialogue (‘dialogue’ from now on), aggression, and criticism.

For the latter, we have introduced additional variables after preliminary reading of the data and following Toepfl’s argument \cite{Toepfl}. Thus, three more coding variables were added: criticism towards authorities/security services (corresponding to ‘leadership-critical public’), and self-criticism. We have discovered this phenomenon of Belarusians criticizing themselves and self-blaming for the political situation in the country and have coded this as a separate variable. Thus, ‘criticism’ was coded as comprising the three types of critique found in the comments.

The coding was performed by two coders (Cohen’s kappa = 0,81), both speaking native Russian and Belarusian, as the discussion was bilingual.

\paragraph{Testing the Dependencies.} For testing the inter-relation between the variables, descriptive statistics (Cramer’s \(V\)) were used. For demonstration of discursive divergence of the cross-commenters, k-means clustering with silhouette quality metric was applied. For assessment of the possible impact of the deliberative features to discussion dynamics, Granger testing was used.

\subsubsection{5. Results}

\paragraph{RQ1.} \textit{Inter-relation between the deliberative features and their impact to discussion structure.} After we have coded the data, it became evident that dialogue, aggression, and criticism were inter-related in our data in specific ways that might relate to cumulation of support or criticism during the discussion.

From the first correlation tests we have seen that, counter-intuitively, aggression is linked not to criticism but to dialogue, and, moreover, this inter-relation intensifies on the most intense commented days (see Tables~\cref{tab:commentInterRelationYearly} and~\cref{tab:commentInterRelationWeekly} for year-long and day-long spans, respectively). Dialogue and aggression correlate strong enough, while none of criticism types correlates with dialogue or aggression, and criticism to power and self.

\begin{table}[ht]%
	\centering
	\caption{Inter-relation between deliberative features in the comments of cross-commenters on the year-long span (Cramer’s V, N = 8,610).}%
	\label{tab:commentInterRelationYearly}% label всегда желательно идти после caption
	\begin{adjustbox}{width=1\textwidth}
		\small
		\begin{tabular}{ l  l  l  l  l  l  l }% Вертикальные полосы не используются принципиально, как и лишние горизонтальные (допускается по ГОСТ 2.105 пункт 4.4.5) % @{} позволяет прижиматься к краям
			\toprule
			Deliberative feature & Dialogue & Aggresion & Criticism & Criticism power & Criticism policy & Criticism self \\
			\hline
			Dialogue & 1.000 & & & & & \\
			Aggresion & 0.289*** & 1.000 & & & & \\
			Criticism & \(-0.177\)*** & \(-0.074\)*** & 1.000 & & & \\
			CriticismPower & \(-0.151\) & \(-0.014\) & n/a & 1.000 & & \\
			CriticismPolicy & \(-0.048\)*** & \(-0.052\)*** & n/a & \(-0.107\)* & 1.000 & \\
			CriticismSelf & \(-0.043\)*** & \(-0.076\)*** & n/a & \(-0.156\)*** & \(-0.054\)*** & 1.000 \\
			\bottomrule
		\end{tabular}%
	\end{adjustbox}
\end{table}

\begin{table}[ht]%
	\centering
	\caption{Inter-relationbetweendeliberativefeaturesinthecommentsofcross-commentersona weekly span, March 26 to April 1 (Cramer’s \(V\), \(N = 287\))*.}%
	\label{tab:commentInterRelationWeekly}% label всегда желательно идти после caption
	\begin{adjustbox}{width=1\textwidth}
		\small
		\begin{tabular}{ l  l  l  l  l  l }% Вертикальные полосы не используются принципиально, как и лишние горизонтальные (допускается по ГОСТ 2.105 пункт 4.4.5) % @{} позволяет прижиматься к краям
			\toprule
			Deliberative feature & Dialogue & Aggresion & Criticism & Criticism power & Criticism self \\
			\hline
			Dialogue & 1.000 & & & & \\
			Aggresion & 0.470*** & 1.000 & & & \\
			Criticism & \(-0.347\)*** & \(-0.243\)*** & 1.000 & & \\
			CriticismPower & \(-0.299\)*** & \(-0.205\) & n/a & 1.000 & \\
			CriticismSelf & \(-0.153\)* & \(-0.095\) & n/a & \(-0.143\)* & 1.000 \\
			\bottomrule
			\multicolumn{6}{@{}p{\textwidth}}{%
				\hspace*{2.5em}% абзацный отступ - требование ГОСТ 2.105
				*\textit{Note}. The correlations between criticism to policy and other variables was not tested due to scarcity of data.
			}\\
		\end{tabular}%
	\end{adjustbox}
\end{table}

Based on this finding and on the insights from the coding process, we have tested whether the discovered correlations matter on the user level -- that is, whether there are distinct discursive clusters that would diverge, e.g., the ‘dialogical and aggressive’ users and critical users. For this, we have used \(k\)-means clustering and silhouette as a quality metric (see earlier accounts on this method used for social media data in \cite{BodrunovaBlekanovKukarkin2018}).

Indeed, we have discovered that two clusters form well enough (\(S > 0,6\); see Fig.~\cref{fig:twoClusterDivergence}); however, the two clusters diverged only on aggression \((p = 0.000)\) and dialogue \((p = 0.001)\), not on criticism \((p = 0.3)\). The best clustering option (of over two dozen probed, from 2 to 5 clusters, from 3 to 5 variables) is presented on Fig.~\cref{fig:threeClusterDivergence}, it has three clusters and is based on the three initial variables (for all variables, \(p < 0.0002\)). On Fig.~\cref{fig:threeClusterDivergence}, the aggressive-dialogical Cluster 3 is formed by 4 users only, with the rest of the users distributing nearly equally between two ‘calm’ clusters with varying level of criticism. However, if we use another sorting option (with observations taken as fixed intervals), their number grows to 15.

\begin{figure}[ht]
	\centerfloat{
		\includegraphics[scale=1.0]{twoClusterDivergence}
	}
	\caption{Plot of means for two-cluster divergence of cross-commenters.}\label{fig:twoClusterDivergence}
\end{figure} 

\begin{figure}[ht]
	\centerfloat{
		\includegraphics[scale=1.0]{threeClusterDivergence}
	}
	\caption{Plot of means for three-cluster divergence of cross-commenters.}\label{fig:threeClusterDivergence}
\end{figure} 

This supports what we have seen during coding. There are two modes used by cross- commenters of our dataset. Mostly, users self-express ‘urbi et orbi’ or, more rarely, address their comments to the authors of the discussed videos. There were nine users among 76 who never addressed anyone and just commented by expressing themselves; for 42 of 76 users (55\%), dialogue-oriented comments were fewer than 10\%. Such users seem to form a personal ‘opinion bubble’ which they stretch across the YouTube channels they comment in. The dialogue appears to ‘belong’ to other users that enter aggressive interrogations with other users (whom they often dismiss as trolls), and these micro-skirmishes evolve into microscopic ‘spirals of silencing’ one of the opponents.

Another interesting, even if intuitively understandable, finding is that self-criticism is inversely proportional to criticism to authorities, as well as to policy criticism (despite the fact that all options could be coded for one comment). However, unlike criticism on the whole, self-criticism does not help users diverge. We have tested self-criticism in various combinations of variables for 2, 3, and 4 clusters, and self-criticism was never significant. Thus, this modus is spread all over the discussion core, but in small proportions which are not enough to influence individual user strategies; only 12 of 76 users (circa 16\%) have not expressed self-criticism in their comments.

\paragraph{RQ2.} \textit{Impact of the deliberative features to the discussion dynamics.} To assess whether the deliberative features cast impact upon the discussion dynamics, we have tested the relations between the overall number of posts and the number of posts with deliberative features for various time spans (see Table~\cref{tab:projectSubdatasets}). After our findings for the Russian YouTube of 2019 [11] where we have seen that aggression highly influences discussion dynamics, as well as vice versa, we expected this to be also true for the Belarusian YouTube.

However, of all the Granger tests we have conducted for weekly, daily, 6-h, and hourly spans, only two cases detected dependencies. Thus, by Granger tests, dynamics of self-criticism shaped the dynamics of the discussion core in 6-h increment, both for the overall number of posts (6.3*, \textit{p = 0.021}) and for the posts with self-critical ones excluded (5.64*, \textit{p = 0.03}). In its turn, criticism in the discussion core affects the number of comments in the whole dataset within one day (1-h increment). This may point out to a cumulative pattern of core-to-periphery influencing: when criticism in the core discussion grows or diminishes, intensity of the whole discussion changes accordingly. This, of course, needs more testing to detect stricter patterns of influence; however, for the discussion under scrutiny, we see that the macro-levels (yearly and monthly spans) are too loose in terms of number of posts to create meaningful patterns of cumulation, while, within hours, the discussion fluctuates depending on how critical it is.

Another interesting finding, perhaps, is that neither the overall number of comments casts impact upon the number of comments with distinct deliberative features, which is counter-intuitive. We think that this effect is related to the fact that the discussion is too loose in terms of number of posts in time. In non-dense discussions, each week and day start anew: what was said before, even immediately before, neither spurs nor slows down the pace of dispersed user commenting. This is also supported by how the number of comments in the core and the whole discussion inter-relate: their correlation is very high (R Pearson 0.880***), but Granger test does not show any mutual impact (which was expected, at least in terms of impact of the whole discussion upon its core).

This, combined with the knowledge of low orientation to dialogue, creates a picture of dispersed, self-expression-oriented, and non-deliberative \textit{cumulative commenting}. Despite formally present user-to-user comments, they do not involve users into opinion argumentation and (dis)agreement. What they do is perhaps lowering entry thresholds for those who would like to express their dissent towards authorities and the country.

\paragraph{RQ3.} \textit{Qualitative assessment of the results.} Thus, what we have seen as cumulative patterns are the following:
\begin{enumerate}
	\item The personal bubbles of opinion expression. Such users do not enter any dialogue with other commenters; when put into the position of influencing, they may spread their influence to others who reply them in various YouTube accounts.
	
	\item Micro-spirals of silencing. The aggressive-dialogical users enter discussion episodes and silence their opponents in skirmish-like communication. Due to this, a bigger ‘spiral of support’ forms and grows under particular videos.
	
	\item Micro-impact of criticism expressed by users. While the level of aggression was low enough not to cast any significant impact upon the dynamic of user talk, criticism took its place and cast impact within 1-h and 6-h time spans.
	
	\item Cumulative commenting. This might as well describe what happens on the macro-level of the discussion bulk. Cumulative commenting happens when comments are relatively rare and dispersed, previous talk does not shape the current one even within several days, and dialogue does not form. However, opinion (in this case, clearly oppositional, critical towards the authorities, and self-critical to a large extent) forms anyway, which points out to the cumulative effects of how we perceive comments online.
\end{enumerate}

\subsubsection{6. Discussion and Conclusion}

In the discussion on Belarusian oppositional YouTube of 2018, we have discovered two modes of expression that were mutually exclusive -- namely, the aggressive-dialogical and (self-)critical ones. This was true on both the level of a single post and user deliberative strategy. We have seen that these discursive strategies play a role in formation of cumulative patterns within the discussion, forming personal opinion bubbles and micro-spirals of silencing. Both strategies deserve further exploration, as they, e.g., contradict conventional view on filter bubbles as formed by many users. Opinion bubbles of the users that stand in the middle of discussions may involve other users (who respond to their comments) into their orbits, thus forming hardly detectable but clearly important influence spreading patterns. Proving this goes beyond the immediate goal of this paper, but definitely deserves exploring.

We have also shown that Toepfl’s theory \cite{Toepfl} needs to be amplified by the notion of self-critical public. Criticism and self-criticism have become the only deliberative features of the discussion under our scrutiny that, at least to some extent, shaped the overall / core discussion dynamics.

We have spotted patterns of cumulative nature that work each on its own level of the discussion. Thus, micro-spirals of silencing happen within small interactions, and thanks to this a bigger ‘spiral of support’ forms easier under individual videos; personal opinion bubbles exist on the user level; multiple deliberative posts affect the immediate development of the discussion; and cumulative commenting is how it looks in general.

Unlike in Russia \cite{BodrunovaLitvinenkoBlekanov2021}, aggression, due to its non-saliency, did not play a meaningful role in discussion dynamics. This points out to cultural differences in how authoritarian public spheres work, even in countries as culturally close as Russia and Belarus.

\section{Анализ пользовательских дискуссий}\label{sec:ch2/sec5}

\subsection{Кто Виноват? Паттерны присвоения вины и ответственности в сетевых дискуссиях об иммигрантах в России и Германии}\label{subsec:ch2/sec5/sub1}

Конфликты с участием переселенцев по оси «север "--- юг» и их обсуждение в медиа и социальных сетях стали чертой современной публичной сферы в Европе и России. По данным ООН \cite{UN2013}, Россия и Германия остаются одними из основных мест привлечения глобальных миграционных потоков. Это порождает в двух странах напряженность, заметную по росту числа крупных и малых конфликтов на межэтнической почве, а также по поляризации и борьбе в сфере регулирования миграции.

Сетевые дискуссии могут играть роль групп давления в формировании повестки дня и миграционной политики \cite{TrottierFuchs}, но непрямым образом \cite{Tufekci}. Дискуссии актуализируют конфликтные настроения, демонстрируя социальную поляризацию, которая может не совпадать с партийным распределением \cite{BodrunovaBlekanovSmoliarova}, и, несмотря на ограничения исследований виртуального пространства с точки зрения социальной репрезентации \cite{Daniels}, изучение сетевого контента способно выявить аттитюды, в т. ч. описать участников конфликта и отношение к ним.

Важную роль и в разрешении конфликтов, и в обсуждении миграции играет то, какие акторы признаются виновниками конфликта и на кого возлагается ответственность за выход из кризисной ситуации. Остается вопросом, есть ли в присвоении вины и ответственности кросс-культурные закономерности или же они являются национально-обусловленными. Для выявления паттернов вины и ответственности мы изучили дискуссии в Твиттере о четырех сходных конфликтах в России и Германии 2013–2016 гг. Это конфликты с насильственным триггером, в которых межличностный конфликт «вскрывает» групповую напряженность между «принимающим сообществом» и иммигрантами. В разных политических контекстах и медиасистемах паттерны присвоения вины и ответственности различаются, несмотря на схожесть иных аспектов конфликта. При этом исследований, где они бы сравнивались, нет, а русскоязычные сравнительные исследования дискурса о мигрантах в России и Германии публикуются реже, чем раз в десять лет \cite{Malakhov,Svinkina,Didenko}. Сравнительные работы политических лингвистов не задействуют либо российский \cite{NesterovaBurova}, либо немецкий \cite{Kalygina} контексты. Игнорирование российского контекста также характерно для немецких исследований, за исключением редких случаев \cite{Pottker}. Наше исследование призвано частично восполнить этот пробел. Ниже мы оцениваем связь статуса пользователя с паттернами присвоения вины и ответственности, а также описываем направленность вины и ответственности в Твиттере двух стран.

\subsubsection{Социальные медиа и конфликт: присвоение вины и ответственности}

Сегодня ученые исследуют роль новых механизмов общественной дискуссии в политике, в т. ч. роль социальных медиа как в разрешении кризисов \cite{Patrona}, так и в нагнетании конфликтности \cite{HeverinZach}. При этом современные работы менее оптимистичны, чем ранее: они указывают на «нормализацию» пользования социальными медиа со стороны властных акторов \cite{LasorsaLewisHolton} и перенос оффлайновых иерархий в сетевое общение \cite{Daniels}, что снижает потенциал общественного давления.

Несмотря на актуальность этой проблемы, изучение стратегий обвинения в текстах социальных медиа практически не проводится. В когнитивистике, впрочем, нет согласия по поводу того, насколько обвинение в принципе является нормой в ситуации конфликта \cite{PisarevaGritsenko}. Так, в 1990-е гг. считалось, что обвинение и самообвинение связаны с психическим состоянием респондентов \cite{TennenAffleck}, но позже было показано, что перекладывание вины на «систему» помогает физиологической адаптации \cite{LaVeistSellersNeighbors}, т. е. является ожидаемым ответом на конфликт. При этом в научно-правовом дискурсе доминирует иное содержание термина «вина» "--- как доказанной персональной причинности при нарушении закона \cite{Dmitrieva}.

Англоязычная литература во многом сфокусирована на феномене обвинения жертв конфликта (victim blaming), который особенно заметен в дискуссиях, ценностно поляризующих общество \cite{StubbsRichardsonRaderCosby}. Обвинения в соцсетях колеблются между тремя разными по природе полюсами: обсуждением вины акторов конфликтов как политического явления, упоминанием персональной вины преступников и обвинением жертв. Поэтому мы будем использовать анализ заявлений, \textit{claim-making analysis} \cite{KoopmansStatham}, как метод оценки текста, позволяющий включить в рассмотрение все три вида обвинений.

\paragraph{Иммигранты в России и Германии и их обсуждение в традиционных и социальных медиа}

В фокус нашего внимания попали конфликты с насильственным триггером, в которых межличностный конфликт «вскрывает» межгрупповой конфликт, имеющий этнорелигиозную и мигрантскую составляющие. В силу этнитизации и исламизации темы миграции в публичной дискуссии \cite{Malakhov,BodrunovaKoltsovaKoltcov} в таких конфликтах межэтническая напряженность раскрывается в многообразии ее контекстов.

Российские конфликты, выбранные нами, включают убийство Егора Щербакова Орханом Зейналовым и последующие народные волнения в Бирюлево (Москва) в 2013 г. и убийство ребенка няней-узбечкой Гюльчехрой Бобокуловой в 2016 г. Конфликты в Германии включают массовое насилие над женщинами в новогоднюю ночь 2015--2016 г. в Кельне и наезд автобуса на рождественскую ярмарку в Берлине в 2016 г. Во всех этих случаях триггер был создан представителями иммигрантских сообществ, понимаемых как отличных по фенотипу, языку и религиозным предпочтениям. За этим последовала реакция принимающего сообщества, в т. ч. с выходом на улицы и/или обсуждением в соцсетях (тематика попадала в мировые и / или национальные тренды Твиттера). Дискуссии в соцсетях включали вопросы о миграционной политике (закрытие рынков, введение виз для иммигрантов, политика приема беженцев, политика безопасности для городской среды) и вовлекала представителей власти.

\paragraph{Иммигранты в медиа России и Германии: недопредставленность и демонизация}

На иммигрантские сообщества в силу их стигматизации потенциально направлено повышенное внимание обвинителей. А неполная включенность иммигрантов в публичную сферу \cite{Zvereva2014} не оставляет шансов на равный ответ \cite{BleichBloemraadDeGraauw}. В постсоветской России публичный дискурс о мигрантах складывался из властного дискурса (которому часто противостоял дискурс чиновный), масс-медийного и редко актуализируемого академического. К началу 2010-х гг. они переплелись, и на фоне роста неприятия иммигрантов \cite{Bessudnov} доминантами стали: указание на этничность иммигрантов; миф о большом вкладе «приезжих» из Центральной Азии в рост преступности \cite{Mukomel}; дуализм в речах федеральных и региональных чиновников "--- нужда в рабочих руках vs. необходимость введения виз \cite{BodrunovaKoltsovaKoltcov}; замалчи- вание конфликтов и общая неуверенность в том, как освещать в СМИ межэтнические кризисы \cite{HutchinsonTolz}.

Еще одной чертой миграционного дискурса стало нарастание межэтнической проблематики в блогах и социальных медиа \cite{BodrunovaLitvinenkoBlekanov2017}, где формируется собственная повестка обсуждения \cite{Mukomel,ApishevKoltcovKoltsova16}. При этом в публичном поле сложились две коалиции в сфере политики о миграции: во-первых, заинтересованные в субсидиях коммерческие структуры и активисты-«квазиправозащитники», а, во-вторых, немногочисленные активисты, работающие с реальными проблемами иммигрантов \cite{Kondakov}. Исследователи подчеркивают жесткость дискурса о мигрантах, в котором очевидны обвинительная модальность высказываемой позиции \cite{VikulovaSerebrennikova}, стабильность негативных аттитюдов даже в ситуации контакта с иммигрантами \cite{Sokolov}, рост поляризации по оси «свой "--- чужой» \cite{Mukomel2014}, роль политических акторов и СМИ в поддержании настороженного отношения к иммигрантам \cite{KugaiKovaleva,HutchinsonTolz}.

В Германии к 2015--2016 гг. повестка дня в публичных дебатах о миграции изменилась в связи с «миграционным кризисом». В 2000-е гг. внимание общества привлекали выходцы из Турции и стран бывшего СССР, а также репрезентации ислама в публичном дискурсе в целом \cite{BonfadelliMoser}; «арабская весна» и сирийский кризис сместили повестку в сторону обсуждения ближневосточных групп переселенцев \cite{HolmesCastaneda}. Еще недавно дискурс об иммигрантах формировался как дискурс толерантности и о толерантности \cite{Didenko}, но «кризис беженцев» выявил поляризацию населения и складывание анти-мигрантской «темной Германии» \cite{BodrunovaLitvinenkoBlekanov2017,SchelterKunegis}. При этом СМИ хранили подчеркнутый нейтралитет: многие редакционные гайдлайны прямо запрещали называть национальность правонарушителей в криминальных сводках.

\paragraph{Контекст обсуждения конфликтов с участием иммигрантов}

Материалы нашего исследования анализируют онлайн-тексты о конфликтах и/или миграции с применением автоматизированных методов \cite{BikkulovBershadskaya,ApishevKoltcovKoltsova16,Gabrielova,BodrunovaLitvinenkoBlekanov2016,BodrunovaLitvinenkoBlekanov2017}. Показано, что следует принять во внимание дискурсивные и медийные факторы. И в Германии, и в России обсуждение вопросов миграции связано с дискуссией о национализме \cite{Elwert,Drobizheva}, так как связь иммиграции с исламом является предметом религиозно-обусловленной нетолерантности, угроз, призывов к насилию от сторонников «чистоты нации» \cite{Awan,BodrunovaLitvinenkoBlekanov2021}. Также в обеих странах наблюдается напряженность между федеральным и региональным уровнем формирования миграционной политики.

Но в Германии иммигрантские сообщества на несколько лет раньше начали пользоваться гаджетами и выходить в социальные медиа \cite{HeppBozdagSuna,Hinkelbein,KuzhelevaSagan}. В Германии интенсивность пользования Твиттером политиками, госорганами и НКО выше, чем в России, и мы ожидаем большее присутствие иммигрантского сообщества и институтов в немецком Твиттере. К тому же разная степень гражданственности общества в России и Германии позволяет ожидать, что в России промигрантские голоса не будут связаны с НКО, тогда как в Германии иммигрантов, говоря простым языком, будет кому защищать.

\subsubsection{Гипотезы и методология исследования}

Исходя из вышеизложенного, мы сформулировали пять гипотез.
\begin{itemize}
	\item H1.Объемтвитовсуказаниемнавозложениевиныи/илиответственности
	за конфликт выше в российском Твиттере, чем в немецком.
	\item H2. Статус пользователя связан с избираемыми им паттернами вины
	и ответственности.
	\item Н3. В России основным объектом обвинения выступают иммигранты
	(индивидуально и как сообщество), а в Германии "--- государственные
	органы.
	\item Н4. В обеих странах обвинения будут звучать в адрес федеральных
	властей, а разрешения конфликта люди будут ждать от региональных
	властей.
	\item Н5. В обеих странах объем обвинения жертв будет пренебрежимо
	мал.
\end{itemize}

Общая методология проекта описана нами ранее \cite{BodrunovaLitvinenkoBlekanov2016,BodrunovaLitvinenkoBlekanov2017}, здесь мы опишем лишь ход данного исследования. Оно включало следующие этапы: (1) отбор кейсов (с помощью сервиса \textit{trendinalia.com}); (2) формирование наборов ключевых слов (словарей) для веб-краулинга; (3) веб-краулинг на основе словарей. На последней стадии загружались метаданные и твиты пользователей (шаг 1), а также метаданные пользователей, комментировавших, ретвитивших или лайкавших собранные твиты, для дальнейшего определения влиятельных пользователей в выборке шага 1 (шаг 2).

Пользователи шага 1 ранжировались по числу твитов, и для каждого набора данных вычислялся порог активности пользователей K, с тем чтобы сформировать выборки твитов сравнимого объема и репрезентативности для обработки в ручном режиме. После препроцессинга и очистки от нерелевантных твитов получены выборки по конфликтным событиям: Бирюлево "--- 1099, Бобокулова "--- 1513, Берлин "--- 1437, Кельн "--- 2525 твитов. В случае Кельна выборка построена путем обогащения сокращенной выборки твитов «обычных» людей твитами медиа из числа влиятельных пользователей, т. к. первичная выборка слишком велика для ручного анализа.

Выборки закодированы вручную (\textit{Cohen’s Kappa for two coders}). Каждый твит кодировался по пяти переменным: (1) статус пользователя: системный политический актор, внеститемный актор (НКО / представитель радикальной группы / представитель диаспоры), медиа, обычный пользователь; (2) наличие фрейма обвинения или (3) фрейма ответственности; (4) адресат обвинения или (5) возложения ответственности: федеральная власть и национальный лидер, региональная власть и полиция, имми- грантское сообщество, иммигрант(ы) "--- триггер(ы) конфликта, жертва или все общество, медиа, политические группы, иное. Спорные случаи обсуждались рабочей группой.

Для оценки связи переменных применялась корреляция Спирмена. Выборки твитов подвергались экспертному чтению. Ограничения метода включали невозможность гарантировать принадлежность твитов к географическому региону и невозможность собрать твиты, размещенные «под замком»; но мы заинтересованы в изучении публичного дискурса, а подавляющее большинство твитов на немецком и русском принадлежит жителям Германии и России, поэтому данные ограничения несущественны.

В определении стилистики высказывания, указывающего на вину/ ответственность, мы пользовались наработками \textit{claim-making analysis} и теории фрейминга \cite{Fuersich}. Так, принималось в расчет только активное возложение вины и/или ответственности, отличное от критики акторов и призыва к действию. Брались в расчет фреймы вины и ответственности «за пределами» триггера, разъясняющие причины триггерного события, предъявляющие претензии к персонам или социальным группам, сравнивающие триггерное событие с иными подобными. Не кодировались попытки снятия вины и защиты триггеров конфликта и иных акторов.

\subsubsection{Результаты: паттерны вины и ответственности в Твиттере России и Германии}

Обвинение и высказывание ожиданий разрешения конфликтов присутствует, как и ожидалось, в Твиттере в обеих странах. Гипотеза 1 касалась объема таких высказываний; (Рисунок~\cref{fig:blameAndResponsibilityTweets}): как видно Гипотеза 1 должна быть отвергнута, т. к. объем таких твитов в Германии почти в полтора раза выше, чем в России. Интересно сопоставить кейсы с групповым (Бирюлево и Кельн) и индивидуальным триггером (Бобокулова и Берлин). Групповые конфликты демонстрируют сопоставимый паттерн, а в берлинском кейсе обсуждение ответственных занимает намного больше места относительно и российских, и кельнского кейса. Подробное чтение твитов показывает, что разница обусловлена доминированием нейтрального новостного дискурса в кейсе Бобокуловой и тем, что пользователи в России часто констатируют бессилие властей.

\begin{figure}[ht]
	\centerfloat{
		\includegraphics[scale=0.5]{blameAndResponsibilityTweets}
	}
	\caption{Объем твитов с фреймами возложения вины и ответственности, в \% от числа твитов по кейсу/стране}\label{fig:blameAndResponsibilityTweets}
\end{figure}

Гипотеза 2 касалась связи статуса пользователя и адресата обвинения. Как видно (Таблица~\cref{tab:blameAndResponsibilityLink}), в двух из четырех кейсов статус пользователя слабо коррелирует с его/ее мнением о том, кто несет вину за конфликт. Иначе говоря, в половине случаев НКО, государственные акторы и обычные пользователи действительно направляют гнев и обвинения на разных субъектов. Малый объем найденных твитов не позволяет подтвердить Гипотезу 2, но искомая связь очевидна в половине кейсов минимум, при том что в анализ вошли только случаи прямого обвинения и возложения ответственности. Обвинительный дискурс «разлит» в Твиттере в обеих странах, но все же по-разному. Так, в кейсе Бобокуловой, твиты пропитаны дискурсом о вине почти без ее возложения на конкретных акторов. В марте большинство пользователей, включая медиа, называли Бобокулову «няня-убийца» еще до решения суда; в октябре ситуация повторилась, хотя суд отказался назначить Бобокуловой уголовное наказание. При этом в берлинском кейсе водители автобуса назывались подозреваемыми, а не убийцами, даже когда вина тунисца Амира Амри казалась очевидной, что позволило немецким СМИ сохранить лицо при смене подозреваемого.


\begin{table}[ht]%
	\caption{Связь статуса пользователя и адресата возложения вины и ответственности, корреляция Спирмена}%
	\label{tab:blameAndResponsibilityLink}% label всегда желательно идти после caption
	\renewcommand{\arraystretch}{1.6}%% Увеличение расстояния между рядами, для улучшения восприятия.
	\def\tabularxcolumn#1{m{#1}}
	\begin{tabularx}{\textwidth}{@{}>{\centering}m{5cm} >{\centering}m{5cm} >{\centering\arraybackslash}m{6.5cm}@{}}% Вертикальные полосы не используются принципиально, как и лишние горизонтальные (допускается по ГОСТ 2.105 пункт 4.4.5) % @{} позволяет прижиматься к краям
		\toprule     %%% верхняя линейка
		Кейс / страна & Паттерн вины & Паттерн ответственности\\
		\midrule %%% тонкий разделитель. Отделяет названия столбцов. Обязателен по ГОСТ 2.105 пункт 4.4.5
		Бирюлево & 0,010 & "---\\
		Бобокулова & 0,237** & "---\\
		Кельн & 0,136* & 0,182\\
		Берлин & 0,067 & 0,111\\
		Россия в целом & 0,141 (sig. 0,069) & 0,048\\
		Германия в целом & 0,095 (sig. 0,065) & 0,171*\\
		\bottomrule %%% нижняя линейка
		\multicolumn{3}{@{}p{\textwidth}}{%
			\hspace*{2.5em}% абзацный отступ - требование ГОСТ 2.105
			Примечание. Прочерк означает недостаточное количество данных для анализа.
		}\\
	\end{tabularx}%
\end{table}

Гипотезы 3, 4 и 5 касались адресатов обвинения и возложения ответствен- ности в разных странах. Результаты анализа адресатов обвинения и возложения вины представлены для России (Рисунок~\cref{fig:blameAndResponsibilityRussia}) и для Германии (Рисунок~\cref{fig:blameAndResponsibilityGermany}). При этом Гипотеза 3 о том, что в России основным объектом обвинения выступают мигранты, а в Германии "--- госорганы, должна быть отвергнута.

\begin{figure}[ht]
	\centerfloat{
		\includegraphics[scale=0.8]{blameAndResponsibilityRussia}
	}
	\caption{Распределение паттернов вины и ответственности, российские кейсы}\label{fig:blameAndResponsibilityRussia}
\end{figure}

\begin{figure}[ht]
	\centerfloat{
		\includegraphics[scale=0.8]{blameAndResponsibilityGermany}
	}
	\caption{Распределение паттернов вины и ответственности, немецкие кейсы}\label{fig:blameAndResponsibilityGermany}
\end{figure}

В конфликтах с групповым триггером наблюдаются сходные паттерны присвоения вины: в обеих странах лидируют региональные власти и по- лиция, затем идет федеральная власть, затем мигранты в целом. Обвинения ложатся либо на лидера страны (особенно на Ангелу Меркель, персонифицирующую политику толерантности), либо на правительство в целом («Правительство готово трусливо спрятаться за овощей-оккупантов, подло предав свой народ "--- лишь бы сохранить власть»). Но наблюдается разница в объеме агрессии по отношению к иммигрантам: в российском Твиттере присутствуют и сарказм, и ноты ненависти, а в немецком "--- сарказм соседствует с признанием своей вины политиками и призывами выяснить, кто все-таки виноват. Также имеется разница в присвоении вины иным акторам. Если в России почти нет обвинений ни СМИ, ни политических групп, то в Германии их объем сравним с объемом обвинений в сторону триггера. В случае Кельна это связано с молчанием СМИ по поводу конфликта и ожиданиями граждан, что СМИ будут более жестко и правдиво освещать роль иммигрантов в нападениях на женщин. Также политическая поляризация и ее отражение в паттернах обвинения очевидно отличает немецкий Твиттер от российского: так, участок дискуссии под феминистским хэштегом \#ausnahmslos («без исключения») снимает вину с мигрантов и перекладывает ее на сексистское общество, в то время как радикальный хэштег \#einearmlaenge, появившийся после рекомендации одного из политиков женщинам держаться «на расстоянии вытянутой руки» от мужчин, использовался для критики не только правительства и политики толерантности, но также феминистских групп. По ходу дискуссии оба хэштега обросли саркастическим дискурсом, противоположным по смыслу. Обвинялись также политические партии "--- как ХДС/ХСС, так и СПД.

В случае Кельна мы могли сравнить медиадискурс с дискурсом «обычных пользователей» (Рисунок~\cref{fig:blameAndResponsibilityKoln}). Видно, что последний является гораздо более жестким с точки зрения наличия обвинений и отличается от первого по структуре. Наибольшую вину пользователи возлагают на политических игроков: федеральное правительство, региональные власти, политические партии, тогда как медиа больше обвиняют полицию (хотя есть и твиты о ее эффективной работе), а также пишут о том, что среди обвиняемых есть иммигранты. Именно пользователи налагают вину на медиа и политические группы, а СМИ обвиняют скорее местные власти и не налагают ни вину, ни ответственность на федеральное правительство.

\begin{figure}[ht]
	\centerfloat{
		\includegraphics[scale=0.8]{blameAndResponsibilityKoln}
	}
	\caption{Паттерны вины и ответственности в кейсе Кельна: сопоставление дискурса СМИ и «обычных пользователей»}\label{fig:blameAndResponsibilityKoln}
\end{figure}

В кейсе Кельна медийные твиты абсолютно не отражают картину поляризованных настроений пользователей, ограничиваясь констатациями и репортажами о подробностях конфликта. Но и между газетами есть различия: так, только таблоид \textit{Bild} четко указывает, что в насилии в Кельне виноваты иностранцы; СМИ левого спектра, такие как \textit{Sueddeutsche Zeitung} и \textit{Der Freitag}, осторожны в формулировках и нейтральны к иммигрантам.

В конфликтах с индивидуальным триггером паттерн вины иной: основной объем обвинений переложен на виновника/виновницу конфликта, и большую роль в этом в обеих странах играет медийный дискурс, освещающий детали расследования и ареста. Но следует обратить внимание на намного больший объем обвинений в сторону властей в России и в сторону иммигрантских групп "--- в Германии.

Следует признать Гипотезы 3 и 4 опровергнутыми. Реальность дискуссии оказалась противоположна ожиданиям в обеих странах: вина ложится на региональные власти и полицию, а разрешения конфликта люди ожидают от федеральных властей, связывая конфликты с ошибками в миграционной политике. Подтверждена Гипотеза 5: твиты, обвиняющие жертв, во всех конфликтах оказались единичными, а вина и ответственность ложатся на общество, а не на пострадавших. При этом ответственность общества формулируется по-разному. В обеих странах присутствуют радикальные призывы к объединению и «наведению порядка»: «Отвратительные лицемеры! Закончите уже эту \#Asylpolitik! [\#политикуубежища!»] или «Ору папе БЕЙ ХАЧЕЙ! папа в ответ орет из другой комнаты СПАСАЙ РОССИЮ». Однако в Германии также присутствует контрдискурс о необходимости толерантности и снижения уровня сексизма в обществе.

\subsubsection{Заключение}

Вопреки ожиданиям, паттерны наложения вины и ответственности в российском и немецком Твиттере в целом сходны, причем объем обвинений в немецком Твиттере больше. В групповых конфликтах виноватыми признаются федеральная и местная власти и в меньшей степени иммигранты, в конфликтах с персональным триггером "--- зачинщики конфликта. В том, на кого возлагаются вина и ответственность отражается состояние каждого из обществ. Так, в Германии, где медиа и партии активно влияют на миграционную политику, мы видим политическую поляризацию обвинений и выпады в сторону СМИ; но дискурс защиты беженцев и медийный нейтралитет отчасти компенсируют разлом между радикальными феминистками и сторонниками неоконсервативного движения. В российском Твиттере обвинители более солидарны и более агрессивны, а призывы к объединению против неугодных иммигрантских групп ничем не компенсируются, в то время как ожидания, что ситуация разрешится политически ниже в сравнении с Германией.

\subsection{Content Sharing in Conflictual \textit{Ad-Hoc} Twitter Discussions: National Patterns or Universal Trends?}\label{subsec:ch2/sec5/sub2}

\subsubsection{1. Introduction}

In the last decade, online content selection and news spreading proved to be ‘a phenomenon of growing social, economic and political importance’ \cite[p.~331]{LeeMa} capable, arguably, of subversion of traditional media gatekeeping \cite{Carr}. In 2010s, research on factors influencing content sharing has dealt with user motivations and content exposure concepts, among which selective exposure theory shows significant explanatory power; but, till today, these studies provide mixed evidence on which particular factors influence more the selection of shared content, and they remain case-oriented and lack comparative perspectives.

In neighboring research zones, content selection has been studied for decades; thus, of media effects studies, gatekeeping theory is the closest to the studies of content sharing, as it has for long been based on similar sets of research questions. Today, this theory expands to a much greater variety of gatekeeper types, and multi-layer (primary vs. secondary, as well as other divisions) gatekeeping is discussed. But these studies suffer from the same shortcomings; also, the studies do not involve the context of media systems as a variable, while it inevitably casts impact upon the spectra of available content.

This paper aims to partially cover these gaps by analyzing secondary gatekeeping patterns within three Twitter discussions on recent ethnic/racial conflicts in the USA, Russia, and Germany. Here, comparability comes from their \textit{ad hoc} nature \cite{BrunsBurgess} as well as from similarity of the cases in terms of political polarization and policy implications.

Using a specially developed web crawler, we collected the bulk of the discussions, ranged the users by the number of tweets, and assessed the links to outer content in 1,000+ tweets by active users for each case by six variables. In Sect. 1, we present the overview of the research approaches, including primary vs. secondary gatekeeping. In Sect. 2, the cases and the hypotheses are presented. In Sect. 3, we describe data collection, sampling, and coding. In Sect. 4, we provide the results and test the hypotheses. In Sect. 5, we discuss our contribution to the existing research and the limitations of the study.

\subsubsection{2. Explaining Patterns of Content Sharing: Research Approaches}

Today, content sharing encompasses sharing the same content within a platform (e.g. retweeting) and selection and re-publication of media content from outside the platform -- or both, with no distinction in many of published works. Since early 2000s, the researchers have tried to conceptualize and empirically assess the motivations and patterns of content sharing on social media platforms. One of the approaches was uses and gratifications theory which focused on user motives for sharing, finding out that these were, primarily, information seeking, socializing, entertainment, and status seeking  \cite{LeeMa}, especially on online networks as content sharing communities \cite{LeeAntoniadisSalamatian}. Later, network-based approaches added to this evidence that weakly tied users were more likely to share content if they are tied uni-directedly, rather than bi-directedly, thus forming an evidence of content sharing hierarchies in social media \cite{ShiRuiWhinston}

Another research stream tries to identify the factors that influence content selection beyond user motivation. Recently, a range of works tried to check whether the selective exposure theory \cite{SearsPreedman} works for social media -- that is, whether information shared by the users is in line with their values, beliefs, and biases. So far, there is mixed evidence whether this theory works well for content sharing, e.g. on Twitter. For example, authors \cite{ThorsonWells} have shown that, during London riots of 2011, users shared tweets that were in line with their beliefs; in experiment-based research, users do tend to demonstrate bias in news selection \cite{MunsonResnick} -- to the extent of formation of echo chambers \cite{Garett} with the same views reproduced. A group of Italian and American researchers have demonstrated that the selective exposure theory, especially when it comes to personal traits and beliefs, can explain the polarizing effects in echo chambering \cite{QuattrociocchiCaldarelliScala}, consumption of (fake) news \cite{BessiCaldarelliDelVicario}, and in formation of scientific networks \cite{QuattrociocchiAmblardGaleota}. Most echo chambers research is, though, focused on studying in-platform content exchange and its effects (for a review, see \cite{Bessi}), including spreading of online behavior \cite{Centola} and raising issues via collective intelligence \cite{Levy,MaloneKlein}. On the other hand, authors \cite{MorganLampeShafiq} have proved that in sharing news from outside Twitter, users tend to include into their shared links media of various political stands; similar findings were reported for Youtube \cite{HalveyKeane}.

But in both cases, the role of content attracted from outside the platform in formation of user echo chambers remains understudied; we still do not know how this or that piece of content is selected and what makes it attractive enough to break the cross-platform boundaries. Also, being well-grounded in theory, this research remains case-oriented and often lacks comparative component; we try to partly cover this gap.

\paragraph{Content Selection and Gatekeeping.} Mechanisms of content selection for (re)publication were traditionally studied in media effects research, and in particular -- in studies of gatekeeping \cite{White1950}. In general, gatekeeping helps reduce the complexity of the world turning the incoming unstructured information flow into limited and structured sets of comprehensible messages \cite{Singer}. Also, as early as in 1975 \cite{Janowitz}, gatekeeping was associated with objectivity model of journalism based on impartiality and balance, as opposed to advocacy model. But from the very early days of gatekeeping research, it was also considered an oppressive mechanism allowing non-professional intentions in news selection \cite{Mills}.

Media have for decades been studied as major societal information gatekeepers. In 1980s to 1990s, the gatekeeping theory focused on structural constraints of information selection \cite{DonahueOlienTichenor}, of which time and space pressures, professional values, and organizational routines seemed to be more important than the type of platform, thus forming cross-media understanding of news selection priorities. Also, for international gatekeeping, news values theory was used to explain the choice of local/national vs. international news \cite{RobertsBantimaroudis}.

The growth of Internet has drastically changed the gatekeeping landscape. Media institutions felt they were losing monopoly in gatekeeping -- that is, the control over the relationship with the users they had previously prospered with. Industry predictions ranged from alarmist claims of the ‘collapse of gatekeeping’ \cite{WilliamsDeliCarpini} to transformationalism. Among the latter, several ideas deserve mentioning. Thus, the ‘gatewatching’ concept was suggested \cite{Bruns} assigning the users a capacity to transmit information from media without actual participation in content production. Other concepts were ‘network gatekeeping’ \cite{BarzilaiNahon} and ‘audience gatekeeping’ \cite{ShoemakerVos} that described the growing role of information consumers in dissemination of news; soon, these ideas expanded beyond the industry-based structural understanding of ‘gatekeeping layers’, and \textit{secondary gatekeeping} started to be discussed in media studies \cite[p.~6]{ShoemakerVos} \cite{RobertsBantimaroudis,White1950}. Also, today, the gatekeeping theory meets social network analysis where key nodes (users) who either disseminate information themselves or link discussion zones are considered new gatekeepers \cite{JurgensJungherrSchoen} or, rather, gatekeepers with smaller gateways \cite{BastosRaimundoTravitzki}.

Today, three gatekeeping models are most elaborated (for an overview, see [19]). For the first one, authors define five levels of analysis: individual, that of communication routines, organizational, social-institutional, and social-systemic \cite{ShoemakerVos}. For the second \cite{Nielsen}, primary gatekeepers are those ‘who by filtering information and publishing decide what is specifically what we commonly understand as news’; secondary gatekeepers ‘are those that filter already available news content’. Thus, it directly links primary gatekeeping to editorial work, and secondary gatekeeping to content sharing (link-, affinity-, or community-based). The third one \cite{ThorsonWells,WellsThroson} shifts from selection to ‘curation’ of content and suggest five types of curators: journalists, strategic communicators, ‘social others’, computer algorithms, and a reader’s self \cite[p.~35]{WellsThroson}
.
We, though, need to state that we see the distribution of primary/secondary gatekeeping role to be much more complicated than it comes from the current academic debate, and a ‘two-step gatekeeping process’ \cite{Singer} reflecting the idea of two-step communication flow \cite{Katz}, is just one of options. To reduce the complexity, we define them as content producer (‘primary gatekeeper’), content re-publisher (‘secondary gatekeeper’), and intermediary (‘tertiary gatekeeper’). Primary gatekeepers are those who produce pieces of content originally; the content may as well be user-generated. In contrast to this, secondary gatekeepers, when re-publishing, are considering what might be interesting for their core information cocoon. This division seems to be a key one, as their final goals may be similar, as well as factors that shape the individual level of decision-making (from ideology to personality \cite{ShoemakerReese}), ‘resulting in choices that are organizationally efficient and culturally acceptable inside and outside the newsroom’ \cite[p.~56]{Singer}. Tertiary, or ‘algorithmic’ \cite{Napoli} gatekeepers, in their turn, use storage, aggregation, and affinity-based algorithms to shape the overall picture.

What differentiates ‘new’ and ‘old’ gatekeepers is selection principles as part of the ‘content curation’ process. For all of them, the issue of subjectivity in selection has never been resolved: in the 50s and 60s, while it was argued that editors were stuck into the machinery of editorial choice based on vigorous criteria \cite{Gieber}, other authors \cite{Snider,White1950} argued that subjectivity was the basic principle for content selection. With mediatization of society \cite{EsserStromback,Hjarvard}, this gap in understanding has only deepened, as media logic itself has become more subjected to various social pressures, some seen as external for editorial offices and some explicating in subjective editorial decisions. But in online user-level content sharing, user decisions become even more individualized and hectic, also due to combining random/purposeful decisions based on personal traits \cite{Bessi} with the impact of tertiary gatekeepers.

\subsubsection{3. Twitter Discussions as Research Object: The Three Cases}

Of all social media, Twitter has been under researchers’ particular scrutiny due to its bigger feasibility, availability of research software, and platform features. Several works \cite{BakshyRosennMarlow,ShiRuiWhinston} have stated that Twitter is for content dissemination more than other social networks are. Here, we need to discuss how peculiarities of Twitter discussions \cite{ChaBenevenutoHaddadi} and relations between Twitter and the outer media system affect content sharing.

First, we need to acknowledge that replicability of our findings to other discussions may be low due to the \textit{ad hoc} nature of the event-based discussions \cite{BrunsBurgess}. This may lower the comparative potential and be viewed as a natural limitation of this research, even if we have no empirical proof of non-replicability.

Second, the network structure of the discussion, including echo chambers and influencers \cite{BodrunovaBlekanovMaksimov}, may also play a role in content selection -- or, rather, content selection helps shape echo chambers, as it is natural to suggest that the logic of consonant reality \cite{NoelleNeumann} is shows up in all user decisions, including content sharing. In our previous \cite{BodrunovaLitvinenkoGavraYakunin,BodrunovaLitvinenkoBlekanov2016} and ongoing research on Germany and Russia, we have noticed that echo chambers are not only based on attitudes towards immigrants but also on user type. Thus, in Germany in 2014 and 2016, pro-migrant camp is formed around left-liberal media and NGOs, while the pro-nationalist one -- around nationalist leaders in Germany and Switzerland who look like ‘ordinary users’. In the Russian case, under the cover of intense tweeting by media, we discovered a chain of similar pro-nationalist ‘ordinary users’ who also acted as gateways for each other’s content. We may, thus, hypothesize that cross-platform content sharing is linked to the user type, as well as to his/her attitude towards the minority under question.

Third, what may also matter is the difference of media systems and the role of Twitter in them. Russia differs from both Germany and the US in terms of low trust to print media, competitive social media market, and low hybridization of TV. The US and Russian media markets have a distinct online-only segment, while in Germany dominant positions on offline market survive online. All the media systems have relatively high media-political parallelism, and one can trace political alienations in media content and endorsement: progressive/conservative in the US, left/right in Germany, and pro-/anti-systemic in Russia \cite{BodrunovaLitvinenko2013,BodrunovaLitvinenkoGavraYakunin}. As to the Twitter’s position on the media markets, its market niche is a bit bigger in the US than in Germany and Russia; there is lack of substantial research on the comparable user profiling in the three countries.

Thus, we formulate the following RQs: (1) is the pattern of content sharing national-bound or universal? Do established media remain the leading primary gatekeeper for Twitter users in \textit{ad hoc} discussions? Do we see any differentiating impact of the outer media systems upon content sharing? (2) Does the type of secondary gatekeeper correlate with the type of primary gatekeeper? (3) Do personal biases of the secondary gatekeeping correlate with the type of the primary gatekeeper and the views in the attracted content?

\paragraph{The Cases Under Scrutiny.} To better formulate our hypotheses, we will also shortly describe the cases we look at.

In Russia, in October 2013, a Muscovite Egor Scherbakov was, allegedly, killed by an Uzbek immigrant Orkhan Zeynalov; later, Zeynalov was arrested and found guilty. Before the trial, local residents committed repeated bashings of a major warehouse in the Moscow district of Biryulyovo where the migrant community dwelled and worked; they raised against police and formed ‘people’s brigades’ to monitor the local area. The case became the most resonant conflict between Russians and Southern re-settlers in the recent Russian history.

In the USA, in August 2014, a white police officer Darren Wilson shot dead an unarmed African American youngster Michael Brown. Peaceful protest and violent riots rose repeatedly in Ferguson, Missouri, as well as all over the US. The case has polarized American society not only along white/non-white lines but also across political affiliations. 

In Germany, in the New Year night of 2016, mass sexual harassment of women by alleged migrants from Middle East and North Africa took place in Cologne. This made people in many cities protest, mostly under the PEGIDA movement slogans. National media were practically silent on the case, as editorial guidelines prevent journalists from covering ethnics origin of the criminals, while local media reported the case.

\paragraph{Research Hypotheses.} Having stated all the mentioned above, we have formulated the following hypotheses:
\begin{itemize}
	\item H1. Across the cases, an ‘onion’ pattern of content sharing will emerge: social media will be shared most; online-only media will follow, traditional media will come third. In general, structures of media systems will have no visible impact upon content sharing patterns.
	\item H2. The type of secondary gatekeeper will correspond with that of primary gatekeeper in all the three cases.
	\item H3. Personal preferences (pro-/anti-minority) of users will correlate with the type of primary gatekeeper, as well as with pro-/anti-minority bias in shared content.
\end{itemize}

\subsubsection{4. Methodology}

To test the hypotheses, we have collected the discussion content, formed tweets datasets for manual coding, coded them, and made calculations for H1 and descriptive statistics (Cramer’s V) for H2 and H3.

\paragraph{Data Collection.} We have used vocabulary-based web crawling for data collection. To form the vocabularies, we have collected and manually assessed relevant keywords and hashtags at trendinalia.com; all the cases were in Twitter trending topics for at least two days. We have also added keywords found via manual snowballing in 1,000+ popular tweets per case. We used the crawler developed earlier with modules modified for our purposes. The uploading periods included October 1 to 31, 2013 for Russia, August 22 to 31, 2014 for the USA (as the peaking dates of the conflict, August 14 to 21, contained an unfeasible number of tweets for our software), and January 1 to 30, 2016 for Germany. The uploads returned 10715 tweets by 3574 users for Russia, 193812 tweets by 70018 users for the US, and 64874 tweets by 12382 users for Germany.

\paragraph{Sampling of Datasets for Data Annotation.} As we were interested in the active users’ content sharing patterns, we have formed feasible datasets of tweets for assessment based on user lists ranged by quantity of the published tweets. The collections had to be feasible for manual coding, and thus we were seeking the datasets containing around 1000 tweets; for further research, the datasets are to be enlarged. We had to cut the long tail of the discussion, which may have affected the results in terms of the variety of content shared; but the reason for looking at most active users was that they tend to have more stable patterns of posting and content sharing. Thus, our sample is not representative, but randomized sampling would have even distorted the results in a way, as active users are also those who are retweeted more than others \cite{BodrunovaBlekanovMaksimov}, and, thus, the content they share creates impact that goes beyond average.

Cutting the user ‘long tail’ implied establishing a threshold of user activity -- that is, the minimum number of tweets posted in the discussion \cite{Chadwick,MunsonResnick}. Due to very differing user activity, these thresholds (\textit{K} = \textit{Ntweets} per user) differed much: for Russia, \(K = 20\), for the USA, \(K = 66\); for Germany, \(K = 10\); \(K\) was defined as the point where the graph that ranged users stopped to drop significantly from user to user. After that, randomized samples (certain percentage of tweets) from each user were taken depending on the overall user activity in the discussion, to preserve the comparable volume of the datasets: for Russia, 40\%, for the USA, 5\%, for Germany, 30\%. After cutting out irrelevant tweets, replicas, and foreign-language tweets, the final datasets included: for Russia, 1120 tweets; for the USA, 1095 tweets; for Germany, 911 tweets (since many tweets were in English and had to be cut out). The relatively small sample is explained by the fact that reading of media texts could not be automatized, as so far no sentiment analysis or machine-learning techniques could provide for clear understanding of user attitudes.

\paragraph{Coding.} 
Coding. Coding of the datasets was performed by experienced experts and native speakers. Four variables were coded: user type = secondary gatekeeper type (‘ordinary user’/political activist/political actor/hybrid media/online-only media/‘media’ on Twitter/‘culture’/other or non-defined); user’s attitude to minority (neutral/pro/anti/non-defined); type of primary gatekeeper (‘social’, online-only, hybrid media, politics, ‘culture’, other/non-defined); attitude to minority in media content (neutral/pro/anti/non-defined). Links to own content were excluded from the analysis. Also, we excluded the links to Twitter-based photo content (pic.twitter.com). In future, bigger samples may be coded.

\subsubsection{5. Results and Discussion}

On completing the aforementioned procedures, we have received the following results.

\textit{H1.} We have calculated the number of links to content outside Twitter in the three
datasets in absolute figures and in \% (see Table~\cref{tab:datasetLinkNumber}).

\begin{table}[ht]%
	\centering
	\caption{Number of links in the datasets, absolute \(N\).}%
	\label{tab:datasetLinkNumber}% label всегда желательно идти после caption
	\begin{adjustbox}{width=1\textwidth}
		\small
		\begin{tabular}{ c  c  c  c  c  c  c  c  c  c  c }% Вертикальные полосы не используются принципиально, как и лишние горизонтальные (допускается по ГОСТ 2.105 пункт 4.4.5) % @{} позволяет прижиматься к краям
			\toprule
			Country & \multicolumn{2}{l}{\makecell{All links, of all tweets}} & \multicolumn{2}{l}{\makecell{Social media}} & \multicolumn{2}{l}{\makecell{Online-only media}} & \multicolumn{2}{l}{\makecell{Hybrid media}} & \multicolumn{2}{l}{\makecell{Other}} \\
			\hline
			Russia & 159 & \textit{14.2\%} & 62 & \textit{39\%} & 29 & \textit{18.2\%} & 62 & \textit{39\%} & 6 & \textit{3.8\% } \\
			USA & 280 & \textit{25.5\%} & 70 & \textit{25\%} & 96 & \textit{34.3\%} & 67 & \textit{23.9\%} & 46 & \textit{16.4\%}       \\
			Germany & 317 & \textit{34.8\%} & 41 & \textit{12.9\%} & 41 & \textit{12.9\%} & 197 & \textit{62.1\%} & 38 & \textit{12\%} \\
			\bottomrule
		\end{tabular}%
	\end{adjustbox}
\end{table}

In Table~\cref{tab:datasetLinkNumber}, we show political, economic, and cultural links all together, as they are few; everywhere, media are, indeed, the dominant primary gatekeeper, but with very different roles of various types of media. This is why we describe the cases in more detail here.

In the Russian case, we see a picture of rivalry between social and hybrid media -- despite low structural online/offline parallelism and low trust to legacy media in Russia. Thus, the idea of the ‘onion’ pattern of content sharing is not supported; hybrid media are shared more than expected. Social media are represented by blog platforms (Livejournal, Blogspot), the local social network Vkontakte and global Facebook and Instagram, and video hostings (Youtube, Ustream). More interestingly, if we look at online-only and hybrid media in terms of their political bias (and to several political links), we will see that, on Russian Twitter, pro-elite and liberal-oppositional content is shared almost equally, making Twitter much closer an ‘opinion crossroads’ than, e.g., the Russian Facebook, a distinct platform-wide echo chamber, as stated earlier \cite{BodrunovaLitvinenkoGavraYakunin}.

For the American case, we see a picture opposite to the Russian one. Here, only 25\% of the links is to social media and other 24\% -- to hybrid media, while online-only comprise over one third. For America, the ‘onion’ pattern, again, is irrelevant. Also, there are several other features characteristic for the US case. First, all the variety of the mainstream media is shared. The leading print media included \textit{Los Angeles Times}, \textit{The Nation}, \textit{The Atlantic}, \textit{Chicago Tribune}, \textit{USA Today}, \textit{Time}, \textit{The Washington Post}, and \textit{The New Yorker}; of TV, MSNBC, CBS, CNN, and Fox got 2 to 3 links each. Regional media market is well-represented, especially by media of the St.Louis county. Second, there is a striking diversity in formats and political positions of online-only media, many of them being personal or group providers of opinionated writing, grassroots political reporting, and alternative news (like \textit{RawStory.com}, \textit{Madiaite.com}, \textit{Vox}, or \textit{Breitbart}). Some outlets balance between politics and journalism (\textit{TheConservativeFreeHouse}, \textit{ThinkProgress}, \textit{TalkingPointsMemo}, \textit{AmericanThinker}). This diversity is supported by that in political links (8\% of all links) belonging to activist, human-rights, and petition websites. Third, we have found virtually no African American media. Unlike in the tweets, the ‘attracted’ discourse, be it pro- or anti-Black, seems to be shaped by ‘white’ media -- except for Youtube that features eyewitnesses’ views, African American commentators, and rap music in support of Mike Brown.

The German case is the only that proves H1 right, but, again, the ‘onion’ logic does not fully work, as social and online media get equal quantity of links -- both almost five times less than the hybrid media. 26 links, or 8,2\%, belong to political organizations or activists. Surprisingly, 84 of 186 (45\%) of German hybrid media links are to advocacy outlets, but this may be explained by the fact that 70\% of them lead to a right-wing magazine \textit{Zuerst} and are posted by the same user. Of the rest, 46 are supra-regional general-interest and 50 are regional and local general-interest media. The regional distribution is representative enough; the links are not Cologne-centered. Nearly all the big players are also present: newspapers (\textit{FAZ}, \textit{Süddeutsche Zeitung}, and \textit{Die Welt} to \textit{Handelsblatt} and the tabloid \textit{Bild}), news magazines (\textit{Spiegel}, \textit{Focus}), public and commercial TV, and public radio. Due to this, left-right media market polarization reconstructs in the discussion, with right-wing media clearly winning. Online-only media referenced to are also highly politicized. More than half of them (24 out of 41) are the links to advocacy media fostering agendas and news frames alternative to the mainstream media, like \textit{Alternative Dresden News}, \textit{DortmundEcho}, or \textit{Pi-news}.

Thus, our H1 is fully rejected: the ‘onion’ pattern of content sharing is absent. Also, national media systems and the state of civil society seem to have a bigger impact than expected upon the sharing strategies of the users. While in the US and Germany 9\% and 8\% of the links, respectively, belong to political websites, Russian content sharing is depoliticized. What unites all the cases, though, is an ‘opinion crossroads’ in terms of political bias present via media links platform-wide, but it is yet unclear whether it is present within ‘echo chambers’.

\textit{H2 and H3.} To assess correlations between secondary gatekeepers’ types and attitudes to minorities vs. primary gatekeepers’ types and attitudes in their content, we have conducted Cramer’s V for the 4 variables described above. The results are shown in Tables~\cref{tab:gatekeeperTypeVS} and~\cref{tab:gatekeeperAttitideVS}.

\begin{table}[ht]%
	\centering
	\caption{Number of links in the datasets, absolute \(N\).}%
	\label{tab:gatekeeperTypeVS}% label всегда желательно идти после caption
	\begin{adjustbox}{width=1\textwidth}
		\small
		\begin{tabular}{ c  c  c }% Вертикальные полосы не используются принципиально, как и лишние горизонтальные (допускается по ГОСТ 2.105 пункт 4.4.5) % @{} позволяет прижиматься к краям
			\toprule
			& $\ldots$primary gatekeeper type & $\ldots$attitude to minority in primary gatekeeper’s text \\
			\hline
			Russia & 0.220 & \textbf{0.248*} \\
			USA & \textbf{0.183**} & \textbf{0.208**} \\
			Germany &\textbf{ 0.219*** } & \textbf{0.372***} \\
			\bottomrule
			\multicolumn{3}{@{}p{\textwidth}}{%
%				\vspace*{-4ex}% этим подтягиваем повыше
				\hspace*{2.5em}% абзацный отступ - требование ГОСТ 2.105
				Sig.: \(*p \le 0.05\); \(**0.05 < p \le 0.001\); \(***p < 0.001\).
			}\\
		\end{tabular}%
	\end{adjustbox}
\end{table}

\begin{table}[ht]%
	\centering
	\caption{Number of links in the datasets, absolute \(N\).}%
	\label{tab:gatekeeperAttitideVS}% label всегда желательно идти после caption
	\begin{adjustbox}{width=1\textwidth}
		\small
		\begin{tabular}{ c  c  c }% Вертикальные полосы не используются принципиально, как и лишние горизонтальные (допускается по ГОСТ 2.105 пункт 4.4.5) % @{} позволяет прижиматься к краям
			\toprule
			& $\ldots$primary gatekeeper type & $\ldots$attitude to minority in primary gatekeeper’s text \\
			\hline
			Russia & \textbf{0.288***} & \textbf{0.290***} \\
			USA & 0.150 & \textbf{0.350***}  \\
			Germany & \textbf{0.183***} & \textbf{0.401***} \\
			\bottomrule
			\multicolumn{3}{@{}p{\textwidth}}{%
%				\vspace*{-4ex}% этим подтягиваем повыше
				\hspace*{2.5em}% абзацный отступ - требование ГОСТ 2.105
				Sig.: \(*p \le 0.05\); \(**0.05 < p \le 0.001\); \(***p < 0.001\).
			}\\
		\end{tabular}%
	\end{adjustbox}
\end{table}

As we see from Tables~\cref{tab:gatekeeperTypeVS} and~\cref{tab:gatekeeperAttitideVS}, in all the cases, type and attitude of the secondary gatekeeper (that is, Twitter user) weakly but significantly correlates with bias in the shared text -- that is, the hypothesis of selective choice of content is supported in cross-country perspective, and in Germany these ties are the strongest. But at the same time the type of primary gatekeeper (that is, in most cases, media platform) has varying relevance. Thus, in Russia, users of various nature all rely on a variety of media, from hybrid to social, but still choose the content that supports their views; this means they look more at content than stick to their preferred media outlets. In the US, the position of a user towards the minority does not correlate with where a user goes for his/her information supply; it is the particular attitude that will be sought after, and this also may be a sign of occasion-based content consumption and sharing, as there is no stability in linking to the same source. That is, echo chambers in content sharing seem to exist in all the cases, but they are formed via different mechanisms. The type of links seems to be all in all less important for the users than the content, and this supports the idea of active seeking of information rather than ‘information cocooning’ and passivity in content consumption. But, as already stated, this may also support the idea of occasional sharing as soon as the users run into the content that supports their views. This needs further research.

What is also seen is that our idea of collision of the individual and organizational levels of gatekeeping may be supported in future. What we have definitely seen in case of Germany and for several users in the US and Russian datasets is what one may call ‘individual-level filter bubbles’ formed not only via posting but also via content sharing. This may amplify the echo chamber concept, as not only network-level inter-user information exchange but also individual-level filtering leads to the rise of closed-up communicative zones. In our future work, we will try to see in detail which gatekeepers were on which side of the conflict, with bigger samples allowing for better conduct of regression analysis.

\paragraph{Discussion and Limitations.} Several methodological issues rose in the process of dataset formation and coding. Thus, representativeness of the datasets was considered against feasibility, and the decision was to have more users, rather than more tweets from each user, in the datasets, especially because many most active users were posting or retweeting the same tweets multiple times and had to be almost fully eliminated from the datasets. Also, comparability of the cases and the codebooks was discussed; as stated above, the cases were considered comparable due to their political nature, time structure, and main actors involved. But still one needs to realize that the understanding of political cleavages differs much in the three countries, and in some cases bipolarism of the political spectra does not bring a proper solution in judging the content.

We could not always code the ethnic belonging of the users, and this variable was not used in the research. Manual assessment of the datasets shows much bigger presence of African American population on Twitter than migrant population in Germany or Russia (which is nearly 0), but, as we see from the results, this did not affect the results too much -- perhaps due to the structure of the media systems, since they all lack significant representation of the minorities among registered media, and this means that the users share what is available.

And, of course, it was hard not to allow a certain level of subjectivity in coding, but we hope that similarity of the codebooks, cross-checks, and advice from native speakers helped to overcome these limitations.

\subsubsection{6. Conclusion}
In this paper, we have shown that the complexity of gatekeeping in today’s commu- nication flows may be reduced to primary (content producers), secondary (content disseminators), and tertiary (‘algorithmic’) gatekeepers; thus, content sharing patterns shed light onto the relations between primary and secondary gatekeeping.

We have found that national media markets and the state of civil society still help break the cross-national patterns of content sharing on Twitter, and all the variety of views is presented. But at the same time, across cultures, users share the content they solidarize with in their views; but if in Russia and Germany, user bias is linked to the media they consume, in the USA users are more spontaneous in their choice. Thus, the idea of selective content sharing due to bias and user status is supported cross-culturally, but what is shared depends highly on what national arenas offer.

\subsection{Mediatization of twitter? Traditional and online media in ad hoc discussions on inter-ethnic conflicts}\label{subsec:ch2/sec5/sub3}

By 1990s, the scholarly community established consensual understanding that mediatized public sphere \cite{Schulz}. Mediatized public discussions were continuously highly uneven in terms of access to opinion expression \cite{Calhoun}. Among other features, it was privileging institutional players, including political elites and corporations, not only as main newsmakers but also as representatives of conflict sides. Also, media themselves became privileged as well, within the “fourth estate” role conceptualization and due to them being the main nodes in societal communication flows -- and also due to a range of media effects that bring significant distortions and biases to public discussions, gatekeeping \cite{White1950,ShoemakerVos} being one of them. Emergence of Internet seems to have equalized ordinary users to media and other institutional actors in their potential access to online deliberation, but not in their social and political capital \cite{Fuchs}, diminishing offline disparities \cite{Daniels}, political capacities \cite{Chadwick}, or users’ chances of becoming influencers \cite{BodrunovaBlekanovMaksimov,BodrunovaLitvinenkoBlekanov2016} or discussion gateways (Bastos, Raimundo, Travitzki, 2013)\cite{BastosRaimundoTravitzki}.

At the same time, several research streams have been focusing on the concept of mediatization of society, which provides another angle on the roles of established media in wider society. Mediatization research (among many others, see \cite{Hjarvard,EsserStromback2014,HeppKrotz}) shows that the concept of mediatization implies \textit{presence} of media as actors, gatekeepers/gatewatchers, or communication nodes in public discussions as well as \textit{impact} of media content and media logic upon a wide range of social practices, or arenas of mediatization \cite[p.~253]{ThimmDangAnhEinspanner}.

We argue that online user-generated discussion may also be viewed as such a social practice, or arena; not only politics on the whole or religious practices but also the communicative exchange among social media users may find itself under specific impact of established media. Thus, it can be mediatized -- that is, traditional media via their accounts may significantly alter the discussion itself (its topicality, boundaries, length, intensity, or user involvement), reshape agendas, and provide framing. But this, first of all, implies that media need to play significant structural roles within the discussion to be able to alter it. E.g., earlier, media-based public discussions may make us expect that media would be present and active within the discussion and would eventually become influentials/influencers \cite{PattersonGrennyMaxfield,ChaHaddadiBenevenuto} -- that is, highly authoritative users who either create waves of support or link micro-groups in discussion as key gateways \cite{BodrunovaBlekanovMaksimov}.

Another aspect of media presence in the online user discussions is presence via media content dragged into user tweets from outside Twitter. Here, both media hybridization theory \cite{Chadwick} and inter-media agenda-setting theory \cite[p.~113--117]{McCombs2004} allow for posing a question on which media (by type, segment, reach, or location) are popular content providers for Twitter users. In comparative perspective, this would provide some evidence on whether national or universal patterns of media presence dominate. Here, mediatization research neighbors content sharing studies \cite{BakshyRosennMarlow,Singer} and secondary gatekeeping studies \cite{Nielsen}, and we may fruitfully borrow from the latter to enrich the analysis of media presence in Twitter discussions.

So far, structural presence of media on Twitter has been studied in a much more thorough way than the types of content dragged into Twitter discussions. But we see it crucially important to study media’s presence that comes both “from within” and “from outside” Twitter, to be able to measure it in a more efficient way. As several research groups have shown that active tweeting does not always lead to being retweeted and commented, we need to look at other latent ways of casting influence, including sharing content from outside sources.

Such analysis becomes even more important for the discussions of high social relevance like conflicts, natural and anthropogenic disasters, or terrorist attacks, as in such cases people tend to follow news media at least to the same extent as -- or even more than -- authorities’ or rescue agencies’ accounts \cite{ChewEysenbach,Bruns2014}. These situations gather in social media, including Twitter, the so-called \textit{ad hoc} publics where influencer status is hardly predictable, but looking at the two aspects of presence may help define media’s real roles in the discussions.

We use a mixed research design that combines comparative assessment of media activity and media content that the users link to. We look at three similar \textit{ad hoc} discussions in English, German, and Russian (with an in-depth case study of the Russian discussion) to see whether activity of media leads to the influencer status in terms of retweeting and several types of centrality. Our initial assumption is that, first, media will dominate the discussions as active tweeters and, second, that the pattern of dominance will be the same in all the three cases: the biggest mainstream hybrid media will be dominant in both active tweeting and shared content, and the most active media will also become influential by connectivity metrics.

The topicality of the chosen cases is inter-ethnic conflicts, namely Moscow anti-migrant bashings of 2013, Ferguson riots in the US in 2014, and Köln mass harassment in 2016. These cases were selected due to the fact that inter-ethnic conflicts in different countries provide very similar frameworks for online discussions, and, thus, \textit{ad hoc} discussions become comparable, at least in terms of patterns of promoting users to influencer status and content sharing. Also, these conflicts are similar to what Cottle \cite{Cottle2006} has described as “mediatized public crises”; these conflicts all have violent triggers but then they move to public sphere and develop there. So far, the research on media and political conflicts \cite{Cottle,Cottle2006,Horsti,EskjaerHjarvardMortensen} touched ethnicity in several various ways, but mainly it explained how mainstream and diaspora media helped certain social groups in acculturation and integration into host societies, as well as how media frames on war and major conflicts worked in national and translational perspective. But it is still quite rare that smaller-scale and event-based inter-ethnic conflicts become the focus of such analysis, and it is even more rare for comparative studies of online communication.

Another line of research, namely online media and migration studies, provides insights in terms of how media diets of “mediatized migrants” \cite[p.~1]{HeppBozdagSuna} are structured, as well as how media appropriation strategies and media cultures \cite{BaileyGeorgiouHarindranath} may shape the presence of minority voices on social networking platforms. But, again, these studies mostly focus on everyday communicative disparities and mediatization effects \cite{Bozdag}, rather than on whether and how minority\&diaspora actors, including media, succeed (or fail) to shape online public discussions in emergency or conflict situations. Of these discussions, those on Twitter have the shortest reaction-to-event time and the biggest breaking news alert potential, thanks to the brevity of messages and newly-grown habits of institutional tweeting.

Today, Twitter as a communicative platform seems to have a disputable reputation in terms of its political relevance \cite{Fuchs} due to its high volumes of de-problematized discourse and “botization” \cite{SanovichStukalPenfoldBrown}. But still, at least for some researchers, Twitter “can be seen as highly relevant for political exchange and public discourse” \cite[p.~254]{ThimmDangAnhEinspanner}, not only because it works as a news alerts platform, as aforementioned, butalso due to its rapid massive response to crisis events, which suits our research goals. Also, due tofeasibility of Twitter research, there is already a body of case studies of influencers on Twitter (foran overview, see \cite{BodrunovaBlekanovMaksimov}) and of content sharing on Twitter, againstwhich we can assess our cross-national results

To assess the selected cases, we use vocabulary-based web crawling, web graph analytics based on the collected data, and qualitative assessment of selected Twitter accounts and tweet collections. The theoretical premises and limitations, methodology of the study, and our results are presented below.

\subsubsection{Review of Literature}

\paragraph{Twitter publics and media presence in them}
As stated above, it is hard to predict whether traditional disparities in access to opinion expression and deliberative influence are reproduced within \textit{ad hoc} discussions -- and whether media preserve their place as important users within them. In our research, we deal with “issue publics” \cite[p.~422]{Habermas}\cite[p.~108]{BrunsHighfeld2016} -- that is, \textit{ad hoc} discursive constellations \cite{BrunsBurgess} that may grow and dissolve within one heated discussion \cite[p.~74]{Dahlgren} becoming affective \cite{Papacharissi}. Discussions by such publics may not be representative enough for the “calm” periods of Twitter discourse on ethnic tensions.

We analyze the cases that, on one hand, foster ad hoc discussion outbursts by ordinary users based on emotions of both social solidarity and intolerance, thus de-privileging institutional players with their rational discourse. But, at the same time, these cases raise the necessity of verified on-the-ground information, which would empower media (and eyewitness users) within the discussions.

Cha et al. \cite{ChaHaddadiBenevenuto} have drawn attention to the change in understanding of mechanisms of communicative influence, namely to the shift from the classic two-step communication flow theory \cite{KatzLazarsfeld,Rogers} to a more modern approach that emphasizes interpersonal and group factors, including readiness for innovation \cite{DomingosRichardson,WattsDodds}. While both approaches still remain theories lacking solid empirical evidence, the same work by Cha et al. \cite{ChaHaddadiBenevenuto} states that the two-step communication flow idea finds proof on Twitter at large, thus emphasizing the role of conversational mechanisms against content-related mechanisms of communicative influence. This, in effect, may mean that media as established nodes of communication flows should retain their influencer positions despite the quality of audience and discussion itself.

But in some works we find evidence that media disappoint the researchers in terms of their overall presence on Twitter, not speaking of influence. Thus, Bastos and colleagues \cite{BastosRaimundoTravitzki} note that impact of media accounts on Twitter is exaggerated, as “only a small portion of tweets received by ordinary users comes from media outlets” \cite[p.~269]{BastosRaimundoTravitzki}. Cha et al. \cite{ChaHaddadiBenevenuto} note that overall presence of media in Twitter is really low (0.01\%). However, the same authors note that presence of media within the Twitter information flows on major news events is two to six times higher (still remaining low, though). Several studies have shown that, in times of crises, their importance grows further, as Twitter users retweet news sources rather than authority sources or expert opinions. E.g., Chew and Eysenbach \cite{ChewEysenbach} have shown that news websites were the most popular sources of retweets during the H1N1 pandemic. Bruns \cite{Bruns2014}, researching on the Twitter discussion on a major earthquake in New Zealand, shows that, out of top 16 user accounts most commented and retweeted, 11 belonged to institutions, that is, to media, authorities, and “utilities” (mobile operators). Despite the list of top influencers defined this way changes in the immediate aftermath of a tragedy, the list of top influencers remains packed with media and authorities’ accounts. Other works cast some light on why this may happen: thus, Lotan et al. \cite{LotanGraeffAnanny} show that journalists often retweeted their colleagues. At the same time, established media outlets, as An et al. \cite{AnChaGummadi} prove, retain the role of publishing news and stories without much interaction with readers; thus, it is users who actively drag media into the discussion by sharing their content -- either via retweeting tweets posted by media or via linking to media content published outside Twitter. Also, media show the biggest information production capacity, but it is celebrities who are followed most \cite{WuHofmanMason}.

\paragraph{Content sharing and media content shared}
Another stream of literature explores inter-media content distribution and its implications for agenda setting and qualities of public discussions -- from the viewpoint of inter-media impact upon agendas \cite{McCombs2004,HarringtonHighfieldBruns}, in-platform visibility of content \cite{Singer}, and new layers of societal gatekeeping of publicly relevant issues resulting from content sharing practices, or secondary gatekeeping \cite{Nielsen}. But we need to state that current works on Twitter do not pay enough attention to shared media content. Many works are dedicated to estimation of efficacy of news sharing on Twitter by media themselves but not to user strategies and configuration of media presence via shared content. Another bunch of works classifies the core content of tweets or users, without assessing shared content from the viewpoint of media presence.

The literature on content sharing focuses today more on linking user parameters to sharing outcomes and on explaining the driving factors of knowledge diffusion on social platforms \cite{Osatuyi}, including users’ motives for selection of particular content, including media content \cite{LeeMa}. Inside this field, a stream of research examines inter-relations between television or video viewing and Twitter. One work \cite{WohnNa} focuses on the types of content shared simultaneously with TV viewing. The work closest to ours is by Cha et al. \cite{ChaBenevenutoHaddadi} who try to classify content of tweets, types of shared media content serving as a classification parameter. But we have not found any work that would look at media hybridity (web 2.0 / online-only web 1.0 / hybrid) in shared content.

In absence of works that would provide preliminary knowledge on which types of media are more active and popular on Twitter, we have on it two assumptions countering each other. On one hand, mainstream media of hybrid nature (that is, traditional media also present online; see \cite{Chadwick}) usually possess bigger resources than online-only news outlets and, thus, provide exclusive and verified information faster, which gains them dominant positions against online-only media. Then, we know that patterns of power distribution on the Internet tend to more and more reproduce the offline ones. On the other hand, the online audience profile (based on younger, more cosmopolitan and tech-advanced user groups) may favor online media as more relevant to their interests and viewpoints. In the latter case, the structure of media presence would be the “onion” one: social media / online-only media / hybrid media, in descending order. This idea is partly supported, e.g., in the Russian case by the monthly statistics of media use online, as produced by Medialogia rating agency: hybrid media occupy two to three positions only within the top 10 titles most read online -- practically in all monthly lists during last three years. This situation may be described as low structural (online/offline) parallelism on the media market. But anyway we will hypothesize that hybrid media are more referenced to than the online-only ones, as this may be true for the US and German media systems with more established democratic media and higher structural parallelism.

\paragraph{Research hypotheses}
Thus, our hypotheses on mediatization of Twitter discussions of inter-ethnic conflicts are the following:
\begin{itemize}
	\item H1. Media in all the three ad hoc discussions will occupy significant (and comparable) places in the lists of top posting users within the discussions.
	\item H2. Content of mainstream hybrid media will be shared more than that of online-only and social media.
	\item H3. The patterns of media activity and content sharing will be the same in all the three cases.
	\item H4. In all the cases, media that were active on Twitter were also leaders of content sharing.
\end{itemize}

\subsubsection{Methodology}
To test our hypotheses, we needed to select comparable cases, upload the data, and assess the data in several ways. Our judgment on the cases was based on several aspects characteristic for all of them. These aspects included:
\begin{itemize}
	\item a clear inter-ethnic / inter-racial cleavage where a “majority” and a “minority” could be defined;
	\item presence of situational violence (a killing, harassment, street violence etc.);
	\item popular reaction to violence (mass demonstrations, people’s gatherings etc.);
	\item public involvement of federal/state-level/region-level authorities;
	\item public discussion wide enough to reach mainstream media and become national/world Twitter trending topics.
\end{itemize}

Thus, the three cases from Russia (Moscow anti-migrant bashings in 2013, October 1 to 31), the USA (Ferguson riots in 2014, August 15 to 30), and Germany (massive harassment in Cologne in 2016, January 1 to 30) were selected; in case of Ferguson, the exact dates were selected by pre-testing due to feasibility reasons, as the discussion was much bigger than the two others. To collect the Twitter data on the discussions on these cases, we used vocabulary-based web crawling and developed a focused web crawler with changeable modules. As its technical parameters had been described before \cite{BlekanovSergeevMartynenko}, we will not discuss them here. What we need to additionally state is, first, that our crawler is not language-dependent, and thus we could collect data in many languages; second, we consider our crawler to perform better than commonly used software for collection of data, as our software can “retreat” in time beyond the threshold of 900 available tweets set by Twitter. Also, our data collection algorithm ensures that tweets marked as “popular” by Twitter cast no impact upon the result and overcome other well-known limitations of Twitter data collection such as limited number of tweets and limited number of requests to the server per second.

The resulting collections of tweets were represented within non-directed web graphs where users were nodes and linkages between them (retweets or comments) were the edges. Our web crawler collects all public tweets marked by a keyword/hashtag, unlike Twitter Search API-based crawlers freely available on the market, as it uses an algorithm analogous to human unfolding of tweet rolls.

To address the hypotheses, we used the following methods:

\textit{H1. Media in all the three ad-hoc discussions will occupy significant (and comparable) places in the lists of top posting users within the discussions.}

We comparatively assessed the user lists ranged by quantity of tweets and described media found in top20 tweeters for each case.

\textit{H2. Content of hybrid media will be shared more than that of online-only and social media.}

Based on user lists ranged by quantity, we have formed feasible datasets of tweets for assessment. We were interested in the most active users, as they were those who tweeted most. The collections had to be feasible for manual coding, and thus we were seeking the datasets containing around 1000 tweets; later on, the datasets may be enlarged. The sampling procedure implied establishing a threshold of user activity – that is, the minimum of tweets posted within the discussion \cite{LevyWindahl,ChaHaddadiBenevenuto}, but due to very differing user activity these thresholds (\(K = \textit{Ntweets}\) per user) were different in each case: for Russia, \(K = 20\), for the USA, \(K = 67\); for Germany, \(K = 10\). After that, certain percentage of tweets from each user was taken depending on the overall volume of the collected data, to preserve the representativeness of the datasets: for Russia, 40\%, for the USA, 6\%, for Germany, 30\%. After leaving out irrelevant tweets, replicas, and foreign-language tweets, the final datasets included: for Russia -- 1120 tweets, for the USA -- 1060 tweets, for Germany -- 911 tweets.

The datasets were examined on the matter of what content (that from media and other outside sources) was referenced to in the tweets. Self-referential links (that is, links to the user’s own content on other platforms), either by media or by other users, were excluded from the analysis. The media whose content was detected were grouped into three groups: social media, online-only media and web 1.0 news portals, and hybrid media (offline media with online representation portals).

\textit{H3. The patterns of media activity and content sharing will be the same in all the three cases;}

\textit{H4. In all the cases, media that were active on Twitter were also leaders of content sharing.}

Results for H1 and H2 were qualitatively assessed for comparison.

\subsubsection{Results}

\paragraph{H1. Patterns of media activity}
Media in the lists of top users. The lists of top 20 most active users are presented in Table~\cref{tab:top20InConflictDiscussions} where media/journalist accounts are marked bold, while “Twitter media” (accounts created with recognizable “media” purposes) are italicized.

\begin{table}[ht]%
	\centering
	\caption{The top20 most active users in the discussions on inter-ethnic conflicts in Russia, the USA, and Germany, with \textit{Ntweets} indicated for each user}%
	\label{tab:top20InConflictDiscussions}% label всегда желательно идти после caption
	\begin{adjustbox}{width=1\textwidth}
		\small
		\begin{tabular}{ c  c  c  c  c  c  c }% Вертикальные полосы не используются принципиально, как и лишние горизонтальные (допускается по ГОСТ 2.105 пункт 4.4.5) % @{} позволяет прижиматься к краям
			\toprule
			\# & \multicolumn{2}{c}{RUSSIA} &  \multicolumn{2}{c}{GERMANY} & \multicolumn{2}{c}{USA}\\
			\hline
			1 & mynameisphilipp & 187 & Chat\_Atkins & 218 & \textit{\textbf{disetv}} & 4880 \\
			2 & \textbf{ifenews\_ru} & 165 & JoeyGerlach & 137 & gprince1110 & 2199 \\
			3 & BorisALV & 103 & SvenMFGN & 129 & iamemix & 2104 \\
			4 & White\_technolog & 97& DennisEntenmann & 126 & djayyy7 & 1080 \\
			5 & topoprf & 89 & LoveBeatsHB & 80 & deray & 1031\\
			6 & ruspoker & 80 & Juliet777777 & 51 & \textit{HandsUpTO} & 988\\
			7 & masfka & 76 & renben1 & 43 & Nowhydrogen & 568\\
			8 & volya\_naroda & 71 & Greg389 & 38 & Team\_LIBer8 & 549\\
			9 & \textbf{MetroRussia} & 69 & gavthebrexit & 37 & TRJustPassingBy & 493\\
			10 & \textbf{ruvr\_ru} & 67 & \textbf{MatthiasMeisner} & 34 & jlangdale & 419\\
			11 & ipotechniy & 57 & Laatzenietlopen & 34 & \textit{\textbf{News\_LNK}} & 383\\
			12 & \textbf{Mir24TV} & 55 & Rabenzauber & 32 & \textbf{CaccioppoliMike} & 369\\
			13 & \textbf{news\_kavkazcen} & 55 & pumaspucke & 30 & jeffwired & 321\\
			14 & MedvedRu & 51 & KoelleunTschues & 28 & M0HAMMEDWASIQ & 313\\
			15 & NoviniRosii & 49 & LupusLotarius & 27 & KennethaScott & 298\\
			16 & \textbf{krgzr} & 48 & vkenneth\_com & 26 & \textbf{\textit{rightnowio\_feed}} & 294\\
			17 & Estraniero & 45 & montagsdemoGIDA & 24 & Disciple4Lif & 281\\
			18 & ciperovich & 44 & \textbf{2n1f} & 23 & justinstoned & 260\\
			19 & \textbf{\textit{urannews}} & 44 & Antonvlimburg & 23 & ABLVCKGOD & 254\\
			20 & \textbf{\textit{istina}} & 41 & Aluhut\_fuer\_Ken & 22 & AnonFatCat & 252\\
			\bottomrule
		\end{tabular}%
	\end{adjustbox}
\end{table}

Table~\cref{tab:top20InConflictDiscussions} shows the following. The patterns of presence of media in top user lists are very different from country to country. In Russia, hybrid mainstream media (\@lifenews\_ru -- a major corporation of low-mid-market and tabloid media, \@MetroRussia -- the Russian version of the free-of-charge Metro newspaper, TV channel \@Mir24TV and state-funded radio \@ruvr\_ru) are very active, but online-only media (\@news\_kavkazcen, the Twitter representation of the famous pro-Caucasian Kavkaz Center, and \@krgzr -- a social-media-based project with formally “general-interest” but mostly politicized content) are neighboring them. In Germany and the US, to a big contrast, not media but individual journalists (\@MattiasMeisner and \@2n1f for Germany, \@CaccioppoliMike for the US) are most active, and mainstream media lose in activity to all other types of accounts (eyewitnesses, activists, religious and cultural commentators etc.).

\paragraph{“Twitter media” in the lists of top users.} What we have also discovered is that in Russia and the USA, unlike in Germany, “Twitter media” were quite active. As stated above, we call “Twitter media” the accounts that were created to perform media-like functions, e.g. dissemination of information, and have no other platform outside Twitter. As we see, the American Twittersphere has its own media that actively tried to shape the discussion on Ferguson; in Russia, similar accounts are found within the top 20 user list.

\paragraph{H2. Shared content}
As stated above, for qualitative assessment of the links discovered, we grouped them into four categories: social media, online-only media, hybrid media, others. Self-referential links were omitted from the analysis; by self-referential links, we mean clearly identifiable links by a user to his/her/its own content on other platforms. In ambiguous cases, the links were not considered self-referential. Links to Twitter were also omitted from our analysis, as well as links to photos posted via Twitter (pic.twitter.com).

Comparing the number of links discovered in the datasets, we can say that Russian Twitter was virtually twice less mediatized than Germany and the USA (159 vs. 280 and 317 links, or 14,2\% vs. 25,5\% and 34,8\%, respectively). In our cases, more established democracies showed bigger willingness of users to rely on external links. But at the same time the number of non-media links tells another story. Thus, non-media links in Russia were only 6 vs. 46 in Germany and 38 in the USA (3,8\% vs. 16,4\% and 12\%, respectively); that is, the Russian users, when they share content, seem to be a bit more media-dependent than Germans and Americans.

Comparing the structure of media content sharing in the three cases, we see that H2 should be completely rejected, since no country has shown the pattern we expected to see (hybrid media shared most, online-only media second, social media third). Thus, Russia, contrary to expectations, showed that social and hybrid media dominated the users’ attention (62, or 39\%, vs. 29, or 18,2\%, vs. 62, or 39\%). This happened despite the fact that Russia was characterized by relatively low online/offline parallelism and, in 2013, there were quite many online-only media which fostered political agendas. In Germany, the proportion was the opposite (70, or 25\%, vs. 96, or 34,3\%, vs. 67, or 23,9\%), despite firm establishment of hybrid media and their dominance within the online media market. In the US, the constellation, again, differs from the expected one, as social media and online-only media equally (almost 5 times each!) lag behind hybrid media (41, or 12,9\%, vs. 41, or 12,9\%, vs. 197, or 62,1\%).

To explain the differences we have found, we will go deeper into the cases with qualitative analysis. We will describe content sharing in the three cases very shortly, as our main focus is comparative analysis.

\textbf{The Russian case.} In terms of the content referenced to, social media dominate. Users reference to posts in the local social network Vkontakte vs. global Facebook and Instagram, blogging platforms (Livejournal, Blogspot), and video hostings (Ustream, Youtube). A more nuanced picture appears if we assess online-only and hybrid media in terms of their political stance and bias, as well as to political links. We clearly see that, on Russian Twitter, both pro-establishment and liberal-oppositional content coexist; media of both political camps are shared almost equally. This brings the Russian Twitter closer to the idea of “opinion crossroads” than, e.g., the Russian Facebook which works like an anti-establishment (or liberal-oppositional) echo chamber \cite{BodrunovaLitvinenko}. Among pro-establishment/state-owned hybrid media referenced to, one can find TASS information agency, Russia today, Vesti (the news bulletin of the state-owned nation-wide TV channel Rossiya 1), tabloid \textit{Moskovski Komsomolets} and others, while the oppositional/independent side is represented by an even bigger number of media, namely Gazeta.ru, BBC Russian and Slon (online-only) as well as hybrid Radio Liberty, business media holding RBC, \textit{Novaya gazeta}, Echo Moskvy radio station, and TV channel Dozhd (“Rain”), among others.

Politically neutral media are also present, being mostly regional information agencies and papers (\textit{The Moscow News}, Rosbalt, Fontanka, Kolomna Times). It is hard to say which exact media are the most popular, as the figures are too low to search for linkage leaders (those aforementioned reach up to 4 tweets each). Political links form a small but important cluster where the same cleavage between establishment and opposition clearly shows: Kremlin.ru and oprf.ru (the account of Public Advisory Chamber) neighbor Amnesty International, Kasparov.ru and Imrussia.org (sponsored by Mikhail Khodorkovsky). What we also see is that the same Kavkaz Center that was posting actively is linked to (but only once), and that there are no other pro-immigrant voices that users would reference to.

\textbf{The American case.} As shown above, in the Ferguson case, only one fourth of the links belongs to social media and another fourth to mainstream media, while online-only web initiatives encompass over one third of them. For the USA, H2 is also wrong, as, again, the hypothesized scheme of linking to hybrid media seems not to work.

As to the contents, there are several features that stand out in the Ferguson case. First, the most popular media that are linked to are Youtube (31 links), \textit{The Huffington Post} (12), \textit{The Gateway Pundit} (12), and \textit{The Daily Kos} (11); the most popular traditional media, \textit{The Washington Post} and \textit{The New York Times}, make up only to 10 and 6 references, respectively. Facebook and Instagram (9 and 7 links, respectively) add to dominance of online vs. hybrid media.

Second, there is an overall high diversity of media that are represented on the American Twitter by user shares. Thus, all the variety of the national media establishment was the basis for content sharing on the American Twitter. The leading print media were, \textit{i.a., Los Angeles Times, The Atlantic, The Nation, USA Today, Chicago Tribune, The Washington Post, Time, and The New Yorker}. On the other hand, TV channels and networks (CNN, CBS, MSNBC, and Fox) all gained unexpectedly low number of links (2 to 3 each). On the contrary, regional media market was represented quite well: local media from St.Louis county where Ferguson is located received 12 links. Also, 5 links belonged to non-US media (\textit{The Guardian}, BBC, and Russia Today).

Third, there is an even more striking diversity of online-only media, many of them being personal or group initiatives (\textit{JuanCole.com, DaleNapier.com, PatDollard.com, Steyn Online, Taki’s Magazine} or even \textit{3chicpolitico.com}), as well as providers of grassroots opinionated writing and alternative news. Given this, the abovementioned leaders in referencing look not like “someday-grassroots-now-established success media stories” but like “first among many equals”, perhaps just less opinionated than \textit{AmericanThinker, Alternet, ThinkProgress, TheConservativeFreeHouse, PoliticusUsa}, or \textit{TalkingPointsMemo} who try to balance between political advocacy and journalism. Alternative news providers are exemplified, \textit{i.a}., by \textit{Madiaite.com, Vox, RawStory.com}, or, of course, \textit{BreitBart.com}. All in all, online-only media are multiple and diversified indeed in terms of left-right divisions; we witness a highly developed market of online-only opinionated writing initiatives and close-to-journalism practices. This diversity is mirrored by that in political links (23, or 8\%) that belong overwhelmingly to activist, human-rights, and petition websites.

Fourth, despite the diversity of views described above, no African American websites that would provide the viewpoint of this community. There are pro-Middle East and pro-Indian websites linked to, but, unlike in the tweets themselves, representation of African American media is \textit{de-facto} 0 in all the segments. The shared content -- and, thus, the discourse around them -- seems to be shaped by “white” media (of course produced by multi-racial teams). The only exception is Youtube which channels the eyewitnesses’ views as well as opinions by African American commentators. We also cannot help mentioning high visibility of cultural products (like rap compositions to support Mike Brown) within the dataset.

\textbf{The German case.} For Cologne, the logic of H2 works only halfway. As shown above, as online-only outlets and social media gain absolutely equal number of links -- 41, and this figure is almost 5 times smaller than for the hybrid media. 26 links, or 8,2\%, belong to political institutions or activist portals. In 2,5\% of cases, users shared a link to news aggregators, quotes collections, Wikipedia, or Google Maps. Within 197 links to the hybrid media, 11 landing pages belong to non-German media (6 -- Switzerland, 4 -- Austria, 1 -- The New York Times, USA).

Contrary to expectations, 84 of 186 (45\%) of links to the German hybrid media belong to advocacy media outlets; they all position themselves as nationwide sources. But this needs to be judged as a case peculiarity, not the regular picture for Germany, as almost 70\% of these 84 links belong to a right-wing journal titled \textit{Zuerst}. Right-wing views are also presented by the outlets published by Institute of State policy (Institut für Staatspolitik) such as \textit{Sezession} magazine and weekly newspaper \textit{Junge Freiheit} (6 and 7 links out of 84, respectively). All the links to \textit{Zuerst} and \textit{Sezession} are published by one Twitter user who posted 88 tweets with external media links. Every link points to an advocacy media (65 links to two hybrid media, 7 links to two online-only media) or a website of political organization (12 links).

Among the rest, 46 may be recognized as leading to supra-regional media of general interest, and 50 as leading to regional and local general-interest outlets, thus forming a nearly perfect federal/local balance very different from the two other cases. Geography of the links covers a lot of regions, and it is worth noting that the links are not Cologne-centered.

On the supra-regional level, nearly all the big players are found among the shared links. Thus, among newspapers one finds \textit{FAZ, Die Welt, Süddeutsche Zeitung, Die Zeit. taz} and \textit{Handelsblatt}, as well as the tabloid \textit{Bild}. News magazines include \textit{Spiegel}, \textit{Focus}, and \textit{Stern}; public and commercial TV range from ARD and ZDF to N-TV, N24, and 3Sat; public radio links lead to Deutschlandsfunk and Deutschlandsradio Kultur. As it is widely known, the German market of supra-regional mainstream media that we recognize as those that follow the objectivity standards still shows considerable left-right polarization. This polarization reconstructs in the discussion on Cologne, but with right-wing media winning: 9 media are neutral, 4 left-wing and 4 right-wing, but the number of links is 12:13:21, and the media with explicit political preferences get more shares per edition than the neutral ones.

Left-wing advocacy media are, thus, outnumbered: the established left magazine Compact has only 5 links, the newspaper \textit{Neues Deutschland} -- only one. 6 more links point at two sources: Epoch Times, newspaper of Chinese political emigrants (4 out of 84), and a new liberal- conservative magazine Tichy’s Einblick run by a CEO the left-wing Ludwig Erhard Foundation. Online-only media links are also unexpectedly highly politicized: 24 out of 41 of them lead to advocacy media (like \textit{Alternative Dresden News, DortmundEcho}, or \textit{Pi-news}) fostering alternative news agendas.

In Germany, we see an interesting example of how inter-connected chains of opinionated sources are reproduced on Twitter via linkage. Thus, the media network of right-wing political movement “Neue Rechte” (New Right) includes websites \textit{Widerstand} and \textit{DortmundEcho} (7 out of 24 advocacy online media links) and hybrid magazines \textit{Zuerst} and \textit{Sezession}. The share of these New Right media is disproportionally big in our sample but all of them are introduced by only one user, one of the most active ones.

\paragraph{H3. Similarity of patterns of media activity and link sharing}
Both H1 and H2 provide for the conclusion that national context is still the primary definer of the mediatization patterns in social media, including Twitter; the patterns are fundamentally different in several aspects, including percentage of tweets with media links, presence of media in top posting users, and aspects of media market representation. E.g. in Russia media links are present in less than 15\% of tweets, while in Germany they are found in over one third of them.

In terms of media activity, Russia seems more mediatized than Germany and the US, and Germany lacks “Twitter media”, while in case of Ferguson only such “media” were among the most active users. But in link sharing, Russian users are least active, and this shows the picture of media being active but not attractive, and this corresponds to relatively low trust to textual media in the country. Germany is a bit more mediatized than the US, and its pattern definitely favors hybrid media (they comprise over 60\% of all links, and journalists are among the most posting users), while in the US the pattern favors online media. Then, despite all the three countries have relatively high levels of regional autonomy, in Germany we see the balance in representation of nation-wide/local media, while in the US major media dominate and only St.Louis outlets are relevant, and in Russia nation-wide media dominate.

In the US, political opinion is definitely present more in online-only outlets, but in Germany almost half of hybrid media are politicized, and in Russia politicization follows the national pro-/anti-establishment polarization in media that formally declare neutrality, while political opinion media are practically non-existent. As to purely political links, Russia showed remarkable absence of activist and discussion portals, while in the Western democracies they were present in approximately each 35 to 50 tweets. The overall structure of link sharing seems to reflect two main factors: the structure of media trust and the state of civil society in the given country, but this needs more thorough research.

Thus, H3 is also rejected -- but only partly. We can spot several important similarities that may point to the nature of Twitter discussions. First, commercial and cultural portals were insignificant everywhere (despite that in the US musicians were among the most active users), which made mediatization go along with politicization. Second, representation of left-right political divisions was quite significant in all the three countries, which makes Twitter an “opinion crossroads” -- maybe not in terms of the discussion but at least in terms of representation of views. Third, there were practically no links to the media that would represent “the oppressed minority” in each case. Thus, the dominance of “media of the white majority” in shaping the Twitter discourse comes not only through posting (which is contested in the American case) but also through mechanisms of link sharing. Fourth, TV was remarkably absent from both top posting users and from top shared media, while Youtube and other video sharing sites, especially in the US, were among the highly cited media.

\paragraph{H4. Media activity vs. content sharing}
As we clearly see from our results, in no case, active tweeting by media is related to active link sharing. This is even true for Russia where media are among the top tweeters. Thus, doubts must be expressed on whether active tweeting leads to raising the number of links shared by Twitter users. The two aspects of mediatization seem to have no relation to each other. Links are shared according to the traditional structures of media systems and activity of civil society agents.

\subsubsection{Discussion and Conclusions}

In this paper, we have argued that online discussions may also be viewed as mediatized (in terms of the place and impact of established media within the discussion), and thus, studies of content sharing and discursive influencers may be united with mediatization research. If so, mediatization needs to be interpreted two-fold: 1) as presence of media and media-like outlets as actors technically equal to other users, and, thus, their activity, connectivity and SNA-based efficacy parameters may be measured; 2) as presence of media via link sharing -- and, thus, media representation may be measured. In this paper, we have focused on only two metrics, namely posting activity and structure of link sharing. Our main aim was to juxtapose the results to see whether the comparative perspective was at all possible, and to assess the results in comparative perspective.

What we have seen is dominance of national modes of media activity and media link sharing. But there are also interesting similarities in discursive power distribution when we look at right-left and “majority”/“minority” representation. Thus, we have seen that active mediatization was much higher in Russia, while in the USA “Twitter media” were on the rise. In terms of content sharing, no similar pattern was discovered. But in all the cases we have seen “opinion crossroads” to form in political terms, with right media outperforming left in Germany due to activity of just one user. Also, we have seen that active tweeting has nothing to do with what links are shared; this poses a question of the Twitter strategy for media entities.

In conclusion, we could add two things on the limitations of this research. First of all, we recognize well the limitations in sampling and interpretation. Most data collection limitations were overcome by the design of our web crawler, but the qualitative interpretation of the discovered links could have had human-related limitations. But beyond this, we have realized that a conceptual limitation exists, and we will work on this in our further research. When assessing social media links, we were assessing channels, not contents, and sometimes ran into representations of online-only media via, e.g., Facebook. We counted this in as a Facebook link, as for our research the channel mattered more; but for possible content-oriented research it seems necessary to develop an understanding of distinguishing institutional(ized) users from “ordinary people” -- which might be complicated in case of popular bloggers or activists. This is especially true because we have realized that studying structural things does not tell it all, and mediatization of discussions on social platforms may (and should) be studied qualitatively in terms of power disparities and the role of particular channels in opinion formation, but not the way it is often studied (presence of one channel at another one), but within a bigger picture depicting all channels within a discussion. Studying link sharing is one of the ways to reach this goal. We hope that, in future, research on structure of media representation could be linked to studying content, to see what content-related factors shape the sharing patterns on particular channels in cross-country perspective.

\subsection{Politicians Driving Online Discussions: Are Institutionalized Influencers Top Twitter Users?}\label{subsec:ch2/sec5/sub5}

\subsubsection{1 Introduction}

Twitter has been often described in the literature as a user-driven platform where horizontality of user connections puts into question traditional hierarchies of actors who influence the debate; e.g., Twitter dramatically individualized participation in collective actions and social movements \cite{BennettSegerberg,BastosMercea}. According to the studies of digital protest and ‘hashtag activism’, Twitter has enabled longitudinal campaigning that changes structure and content of public discussions \cite{BonillaRosa}. Case studies of Twitter discussions on selected conflicts paid more attention to the content of debate and the character of public communication rather than users participating in the discussion \cite{BonillaRosa,GroshekTandoc}. Also, several important works focused on Twitter activities and strategies of single categories of users -- predominantly communicative elites, like politicians and journalists \cite{BaumannFabianLessmann,LasorsaLewisHolton}. But there is still lack of research on what place such users occupy in \textit{ad hoc} discussions on urgent public matters, whether they regain user attention and authority in the new communicative environment, whether they occupy positions that help organize meaningful discussions, and what lies behind their popularity.

Moreover, there is no consensus among researchers on what kind of users becomes influential and whom to consider an influencer at all, as well as which user metrics one needs to look at to collect a sample of influencers. Our earlier research deals with this issue and offers a framework for selecting top users \cite{BodrunovaLitvinenkoNigmatullina,BodrunovaLitvinenkoBlekanov2016}. In this paper, we apply this approach to assess the political institutionalization of online discussions in cross-national perspective -- that task is, again, under-addressed in current research on social media platforms.

Today’s proliferation of social media comes across cultures and political regimes, and it poses another research question -- whether national traditions of public discussions (or at least some of their features) show up on similar occasions online, or the platform features and limitations make them irrelevant. In this paper, we address this question by comparatively assessing the discussions in Germany (a European democracy), the USA (a non-European democracy), and Russia which shows authoritarian trends in media development, along with other areas \cite{Toepfl2010,Sherstobitov,SmythOates,Denisova,BodrunovaLitvinenko}.

In search of comparable discussions, we looked at three cases of inter-ethnic conflicts that gained national attention and triggered extensive networks of Twitter interactions: the Cologne mass harassment of 2016, Germany, the Ferguson unrest of 2014, USA, and the Biryulyovo bashings of 2013, Russia. All these cases were among national Twitter trending topics and all of them caused nation-wide discussion and social polarization.

For all the cases, we assessed the extent to which the offline user status may change the potential of becoming a Twitter influencer (top user by several parameters); we also looked at the nature of the users that were present in the discussions. In Sect. 2, we present the research framework and literature overview; Sect. 3 provides our research hypotheses; Sect. 4 is dedicated to the methodology and data collection; Sect. 5 pre- sents and discusses our research findings.

\subsubsection{2 Influencers in a Horizontal Network: User Status vs. User Rank in Twitter Discussions}

\paragraph{2.1 Twitter: A Platform for Reciprocal Communication?} The network theory applied in the Twitter studies draws attention to the structure of the Twitter discussion based on weak links and cluster dynamics, the formation of horizontal influencers and various approaches to their analysis, as well as the relationship of the structure and content of tweets in their relation to the type of user.

In comparison to other social networks, Twitter is often characterized by nonlinear dissemination of information \cite{BrunsBurgees2012} and low level of reciprocal relationships \cite{BakshyHofmanMason,Habermas,KwakLeePark} creating of Twitter a picture of quite a loose network of non-reciprocated communication. At the same time, a lot of research has focused on finding communicative disparities and hierarchies of various kinds in this seemingly equal and loose discursive milieu.

Despite the initial horizontal character of Twitter communication, scholars revealed several patterns of hierarchy in relation to user influencer status.

First, users differ in terms of activity: ‘the more active Twitter users frequently engage in gathering and sharing what they perceive to be relevant materials, for example tweeting links to further information (or retweeting relevant posts of other users) to their own followers or to Twitter communities formed around topical hashtags’ \cite{BrunsBurgees2012}. Other authors \cite{BastosRaimundoTravitzki} found that users with higher level of message communication have stronger and more selective influence on the information in Twitter.

Second, although Twitter is seen by scholars as ‘a platform that works on the idea of reciprocation’ \cite{GroshekTandoc}, they have found various types of user inequality in the patterns of reciprocation. Thus, it has been shown that users of one type (for e.g. Twitter accounts of news organizations) usually gain attention from users of the same type (other news organizations) \cite{WuHofmanMason} rather than expand their network to ordinary users. Follower-followee relationships on Twitter may form a kind of relatively stable, long-term network that tends to represent the long-term interests of users \cite{BrunsBurgees2012}.

Third, offline group representation patterns may influence differences in user impact upon how the discussion goes. For instance, as minorities are less connected to social movements, they get lower visibility among the influencers \cite{Juris,Clark,BodrunovaLitvinenko2013}; or, ‘young people from historically underrepresented communities neither experience themselves as members of a singular community on- or off-line, nor uniformly choose to follow or interact online with elites who have emerged as political leaders’ \cite{Clark}. Nevertheless, the optimistic view that is shared by a large group of scholars says that Twitter enables communities by ‘artefacts of engagement -- comments, photos, videos, tweets, and news stories’ \cite[p. 239]{Clark} which, as the same author claims after studying involvement of young people into Ferguson events and online discussions, “provided an avenue into what Dahlgren called the ‘proto-political’” \cite{Dahlgren}\cite[p. 244]{Clark}. Participation in online discussions made this social group also visible to other parts of the community; this was ‘important, as both strong collective identities and strongly felt shared grievances are known to play a significant role in spurring members of one’s extended social networks to consider participating for the first time’ \cite[p. 244]{Clark}. Clark also identifies ‘Black Twitter’ as ‘a space in which counter-publics may form as people find and follow one another, engage in discussions about the meanings of Blackness, and discuss strategies for engaging in political action that arise from those meanings’ \cite[p. 239]{Clark}. But at the same time, this implies that ‘Black Twitter’ in some respect opposes the wider discussion environment and lacks political institutionalization.

Fourth, a lot of research is focused on how the platform features of Twitter and structural peculiarities of the discussions bring certain users to the positions of influencers. The notions of gateway formation \cite{BastosRaimundoTravitzki} and gatewatching \cite{Bruns2011,Bruns2005} demonstrate that networking patterns strongly influence the potential discussion hierarchies. But this research is not that well linked to the offline nature of the users, which till today remains a research gap. Moreover, the nature of Twitter discussions has been also described as ad hoc, which implies that the findings for one discussion are not applicable to another one. But both research on Twitter gateway formation \cite{BastosRaimundoTravitzki} and our own research shows that technical nature of Twitter discussions has stable features that enable comparative research.

And, perhaps most important, several studies of Twitter have linked user rankings and influencer status to their offline metadata. Thus, studies on Twitter segments that belong to the USA and Sweden have shown that experts, professionals, and old-established organizations exert more considerable influence within discussions \cite{FoxZickuhrSmith,Ruth}. In case of emergency, e.g. natural disaster or a terror attack, the so-called utility accounts may start to play a bigger role \cite{Bruns2011}. The same author has also shown that media and politicians belong to influencers on Twitter \cite{Bruns2011}; other works have shown that, for example, Norwegian politicians and journalists share patterns of Twitter activity \cite{EnliSimonsen} but not in all features. E.g. politicians posted significantly more retweets with hashtags than the journalists, and the politicians’ promoted hashtags through retweets, while journalists used hashtags to promote original content \cite{EnliSimonsen}.

Thus, among the various factors that are, arguably, expected to shape growth of a user’s influencer status, beyond the platform and network nature, the most relevant for institutionalized users would be their activity, their belonging to certain types of institutions, their grassroots/systemic status, and their belonging to the majority/ minority. But this has not yet been tested in comparative perspective, including the democratic/non-democratic states.

As it is known from a lot of research on the aftermath of the Arab Spring \cite{Lynch,TufekciWilson}, as well as from other studies on Twitter and authoritarianism \cite{TufekciWilson,HowardAgarwalHussain}, post-communist, semi-authoritarian and authoritarian states may have differing patterns of political use of Twitter, as well as regime responses to it. This creates a perspective for comparing online discussion patterns in democratic, non-democratic, and hybrid or semi-authoritarian countries, to learn how much they actually differ.

\paragraph{2.2 Measuring Influence in Twitter Discussions} 
To dig a bit deeper into the nature of linkages between the offline user status and his/her online influencer status, we need to look at two major research lines that deal with: (1) how influencers are defined; and (2) how online behavior patterns of political actors is linked to their reaching the influencer positions.

As noted above, hierarchization of Twitter users has been documented in many studies, but not that many of them did it in relation to political actors. Moreover, as soon as it concerns political topics on Twitter, the overwhelming majority of the studies include into ‘politically relevant actors’ not only political institutions and parties, individual politicians, and grassroots movements and their leaders, but also media of various kind. We will follow this tradition and further discuss how offline dominance of these communicative actors translates into their online influencer status.

\textit{Defining Influencers.} As we have addressed this issue in our previous research papers \cite{BodrunovaLitvinenkoNigmatullina,BodrunovaLitvinenkoBlekanov2016}, we will now focus on the aspects relevant for the topic of this research only, as well as on the newest literature in the field.

As stated previously, today’s research on Twitter influencers may be divided in two streams. The first group of scholars includes an account or a group of users into influencers on the basis of non-network-dependent metrics such as the number of followers, comments, or retweets. Within this stream of literature, retweets have been marked as the preferred measure for defining who the influencers are -- these users are retweeted both most frequently and most widely \cite{FrebergGrahamMcGaughey,Aquino,SajuriaVanHeerdeHudsonHudson}. ‘An analysis of responses and retweets provides a useful indication of the overall visibility of each account: as discussed, retweets are a means of amplifying the reach of a tweet, and thus of increasing the visibility of a tweet and its sender’ \cite{BrunsBurgees2012,KwakLeePark,BoydGolderLotan,ChaHaddadiBenevenuto}. There even is a classification introduced by authors \cite{LarssonMoe} for the most active users that includes politicians, bloggers and journalists as well as anonymous users and is based on two types of retweeting -- that is, mentions via @ and RT networks.

Diffusion of messages on Twitter has different velocity depending on retweet activity \cite{BastosRaimundoTravitzki} and user status \cite{LarssonMoe,CollianderMarderFalkman}. What is also important is that a retweet is considered to be ‘a form of further disseminating someone else’s message to one’s own network’ \cite{GroshekTandoc,MascaroGoggins}. Measuring number of retweets is about ‘both influence in diffusion and value in market of a tweet’ \cite[p. 792]{LiuShiChen}. A simple retweet is, in most cases, regarded as the sign of solidarization with the position expressed in the original tweet, but we also know that the retweeting practices vary considerably \cite{BoydGolderLotan}. ‘Retweeting users may even see themselves as information brokers, bridging distinct communities of interest by passing on tweets from one network cluster to another’ \cite{BrunsBurgees2012}. In other works, comments and likes, as well as their combinations, are also considered important as influencer markers, but this line of research is definitely smaller.

The second approach measures connectivity (involvement into discussion or in a network) instead of activity. According to this approach, the most influential are those users who have potential in uniting echo chambers in Twitter because they form big nodes in the discussion networks. Most often, to measure this user potential, the metrics of SNA are employed. The clear difference between two approaches is rooted in fact that such SNA metrics as betweenness and pagerank centralities are network-dependent. “Degree centrality is conceptually the simplest one, which is defined as the number of links incident upon a node. Meanwhile, betweenness centrality and closeness centrality are global metrics which more complex but it can better to identify the influential nodes” \cite{Maharani}. A lot of combinations of network measures have been offered by researchers to define the top users; we rely most on (in/out)degree, betweenness, and pagerank metrics for ad hoc discussions \cite{BodrunovaLitvinenkoNigmatullina}.

In our previous research, we have pioneered juxtaposing absolute metrics and network-based metrics, to show that these correlate \cite{BodrunovaLitvinenkoNigmatullina,BodrunovaLitvinenkoBlekanov2016}. In the most recent studies, there have also been attempts to correlate individual indicators, but no more than three \cite{GroshekTandoc}. But this time we aim at linking this research method to the offline user status that we define via user metadata and user descriptions, as well as the content of the accounts.

\textit{User behavior and their influencer status.} Several studies \cite{EnliSimonsen,GrahamBroersmaHazelhoff,LarssonIhlen} have studied peculiar behavior that distinguishes them from average user profiles. Thus, it was shown that politicians and journalists form isolated networks and tend to interact within these circles – which also may be caused by the nature of the Twitter networks themselves, as other studies have found it for other types of users \cite{BodrunovaLitvinenkoNigmatullina,BodrunovaLitvinenkoBlekanov2016}. Thus, more than two thirds of German politicians that actively use Twitter referred to other politicians in their tweets \cite{BaumannFabianLessmann}. On the other hand, for wider audiences, ‘official’ accounts are able to establish themselves as authoritative sources of information, even in the open environments of social media’ \cite{BrunsBurgees2012}. We also know that behavior matters: witty and sharply formulated political tweets, for example, have a great potential of spreading around the network \cite{AelstVanErkelDheer,ParmeleeBichard}. Also, candidates who have a prominent position in the media are generally also the ones who are more popular on Twitter \cite{AelstVanErkelDheer}.

Similar patterns may be found in online behavior of media and journalists. In case of Ferguson, despite having vast following, on average more than a million followers, legacy media journalists and organizations reciprocated by following only small numbers of users and were the least engaging in terms of retweeting and commenting other users \cite{GroshekTandoc}. These signs of traditional communicative hierarchies let journalists and legacy media remain within top popular users as defined via following but in several days were no longer the most influential in terms of generating discussions and gatekeeping information about the Ferguson decision. This is partly explained by their passive activity on Twitter, according to the authors.

Thus, we will try to explore whether patterns of political actors (including media) are similar in cross-cultural perspective.

\paragraph{2.3 Establishing Comparative Perspective: Twitter and Politics in Germany, the USA, and Russia}

To form better our hypotheses, we need to establish a comparative perspective. For this, we will describe the political use of Twitter in the three countries and show that the cases under our scrutiny are comparable.

\textit{Twitter and politics in Germany, the USA, and Russia.} Although Twitter technologically is the same all over the world, societal structure of users, as well as political clusterization and political use of Twitter, differs from country to country.

24\% of US population used Twitter in 2016 (Pew Research Center). It is hard to identify the leaders of the US Twitter for the American audience, as the data that are provided are not country-based; e.g. the third poplar profile belongs to Barack Obama and is followed by 89.5 Mln of users who obviously live all over the world. According to the open data, most popular accounts of US citizens belong to the show business celebrities.

Throughout the recent years, online polarization has been increasing on the US Twitter \cite{GarimellaWeber,CohenRuths}. Support for the Democratic party is more visible on Twitter than that of Republicans (31\% vs. 19\%) \cite{GreenwoodPerrinDuggan}. Almost 60\% of social media users feel irritated while participating in online discussions with users of different political position \cite{DugganSmith}. A very small number of large organizations that have a very large online network of Twitter followers are concentrating the online political voice \cite{HongNadler}. Among members of the US House of Representatives, politicians with extreme ideological positions tend to have more Twitter followers, both left- and right-wing extremists have larger Twitter readership than their moderate peers \cite{HongKim}. Twitter is actively used in the USA in the election periods, as well as in times of nationwide events.

The divergence of users based on race is more significant in the US Twitter than elsewhere. Also, studies showed that education and ethnic heterogeneity are associated with uncivil communication on the US Twitter, while annual household income and unemployment rate were not related to Twitter incivility \cite{VargoHopp}.

Russian Twitter was launched in 2011 and is less popular than the national social networks Vkontakte (desktop monthly reach in June 2016 46 Mln) and Odnoklassniki (desktop monthly reach in June 2016 30 Mln) \cite{TNSWebIndex2016} as well as Facebook (monthly reach in June 2016 14.4 Mln) \cite{Frolova}. Monthly reach of Twitter was about 12 Mln \cite{Frolova,BrandAnalytics}. Kelly et al. have revealed several political clusters within the Russian Twitter in 2010 and 2011. The cleavages between clusters correlated with topics of discussion rather than with social stratification, which is a sign of presence of echo chambers in Russian Twitter \cite{KellyBarashAlexanyan}. At the same time, both pro-government and oppositional political voices are represented in Russian Twitter almost equally, with slight dominance of pro-establishment voices \cite{Greene,NikiporetsTakigawa}. Leaders of the Russian Twitter are TV anchors and politicians: prime minister Dmitry Medvedev and TV showman Ivan Urgant have about 5.5 Mln followers each. Old-established news organizations (\textit{Channel One} and \textit{Vesti}) are followed by 3.7 Mln of Twitter users each \cite{SocialBakers}. Parody microblogs are also highly visible: the spoof account @KermlinRussia gained over 1.5 Mln followers becoming more popular than the official presidential account the spoof was parodying \cite{Denisova2017}. While not a big share of Russian population uses Twitter the platform attracts more politically active users \cite{BodrunovaLitvinenko2013,BodrunovaLitvinenko2015}.

In general, the Russian segment of Twitter is under-researched. The evidence from previous research provides contradictory data on whether Russian Twitter tends to form echo chambers \cite{KellyBarashAlexanyan} or play rather a role of mediator between different views \cite{BodrunovaLitvinenko2013,BodrunovaLitvinenko2015,BarashKelly}. Also, very recently, the number of bots has risen substantially, especially in discussions on political matters, but this was still irrelevant in 2013.

Twitter is also not very popular in Germany. While more than 30\% of German population have created once a Twitter account, in 2015 only about 10\% of German population used this micro-blogging platform actively \cite{MoscaQuaranta}. Leaders of the German Twitter are celebrities -- football and show business stars with about 4 Mln of followers. In comparison to them, most popular media accounts gained about 2 Mln of followers and accounts of politicians attracted less than 400,000 followers.

Among German Twitter users who were politically active in the 2013 elections, almost every third prefers to participate in supportive networks where users always or often agree with others \cite{VaccariValerianiBarbera2016}; thus, German Twitter is tending to form echo chambers and user chains. Political journalists unite in a journalism-centered bubble where political accounts are included since journalists prefer to follow institutionalized users as sources and do not use Twitter for communication with the audience \cite{Nuernbergk}.

This description shows that we may expect higher presence of political actors in Twitter of democratic countries, and these users will be of different nature (institutional as well as grassroots, perhaps more institutional in Germany than in the US). We can also expect presence of all major political powers on Twitter, but their interaction with users will be scarce, as previous research shows.

\textit{The conflicts under scrutiny.} The three conflicts that we have chosen for our analysis include the Biryulyovo bashings of 2013, Russia, the Ferguson unrest of 2014, USA, and the Cologne mass harassment of 2016, Germany. These conflicts were selected due to the fact that they contained clear political elements, were all ethnic/racial-based, provoked massive discussions on Twitter (reaching the national Twitter trending topics), and had similar ad hoc nature in terms of technical organization. Political criteria for selection included the following: the cases had a violent trigger, evoked or revealed major political polarization in the society, caused peaceful protest, and involved institutional political actors and police forces.

In Russia, major anti-migrant bashings happened in September 2013 in the Moscow district of Biryulevo, triggered by an alleged killing of a Muscovite Egor Sviridov by an Uzbek Orkhan Zeinalov and directed to a warehouse where a lot of Central Asian immigrants (mostly illegal) lived and traded. Police and additional military forces had to gather in the area to stop violence; there were claims on the part of attackers that police over-reacted to the events. Later that autumn, peaceful gatherings of local dwellers continued; there were no direct links in change of Moscow policy towards immigration, but the markets and warehouses in the Russian capital have since then been under particular scrutiny of police and immigration officers.

In the USA, an African American teenager Mike Brown was killed unarmed by a white police officer named Darren Wilson. This event provoked major riots not only in the town of Ferguson where the killing happened but also in other cities and towns; peaceful protest spread up to Washington, DC. The case has polarized the audiences; several media reported that more money was raised for Wilson than for the Brown family. For several weeks, protests continued and had minor policy implications on the local level.

In Germany, in the New Year eve of 2016, mass harassment of female celebrators happened on the main square of Cologne; over a 1,000 women self-reported to be affected, and several risked their lives due to firework pieces being directed to them. As police proved, most of the harassers belonged to the re-settled communities from Middle East and North African countries. The conflict provoked nationwide debates but was also characterized by the silence of media for the first few days, as editorial guidelines did not let papers and TV report on the nationality of harassers.

\subsubsection{3 Method}

\paragraph{3.1 The Research Questions and the Hypotheses} Based on what was said above, we formed the following research questions:

\paragraph{RQ1. Is the structure of presence of politically relevant (institutionalized political, grassroots political, and media) users among the influencers similar in all the three cases?}
\paragraph{RQ2. Do politically institutionalized users play bigger role in the discussion than other politically aligned users?}
\paragraph{RQ3. Is offline status of the users linked to their online influencer status, and by which parameters?}

According to the research questions, we have formulated three hypotheses:

\begin{enumerate}
	\item H1. Due to the nature of the civil society in the countries, there is bigger presence of grassroots political users and NGOs in the USA and Germany in the top user lists than in Russia.
	
	\item H2. Due to importance of political actors during conflicts, as well as inertia of traditional hierarchies, users who belong to the current national-scale political institutions form big nodes in the discussion networks -- that is, they are ranked high by degree, betweenness, and pagerank centralities.
	
	\item H3. The position of users in the top lists correlates with the users’ political status. Thus, the more politically institutionalized a user is, the more he is ‘liked’, retweeted, and commented; also, (s)he occupies a higher place in the structural hierarchy of the discussion web graph (higher InDegree, Degree, Betweenness, and Pagerank). As to the OutDegree, it inversely correlates with political institutionalization, as traditional hierarchies are reproduced in the online discussions.
\end{enumerate}

\paragraph{3.2 The Research Methodology} To answer the research questions, we have done the following:
\begin{enumerate}
	\item We conducted vocabulary-based web crawling to collect the discussion content and to analyze their structure. The keywords were selected at www.trendinalia.com, in accordance with hashtags high-rated by this web service that monitor Twitter trending topics around the world in real time environment and diachronically. A special web crawler has been developed to bypass the limitations of Twitter API;
	
	\item We have reconstructed the discussion web graphs to get the data on network-related metrics;
	
	\item We sampled 50 top users based on nine parameters (number of tweets; number of interactions (likes, comments and retweets); centralities -- degree, indegree, out-degree, betweenness, and pagerank); then we merged them in general lists of top users. Thus, if a user was found in at least one single-metric list, he/she was to be found in the merged top user list. As many users were within top lists by many parameters, the final top lists included 205 users for Russia, 205 users for Germany, and 230 users for the USA. The total number of users in the datasets included 3574 users for Russia, 12382 users for Germany, and 70018 users for the USA;
	
	\item We manually coded the users in the lists for their offline status, aiming at showing their growing political institutionalization (from ‘ordinary users’ to ‘institutional political actors’);
	
	\item We used descriptive statistics (Spearman’s rho and Cramer’s V) to see whether the offline user status correlates with the online metrics;
	
	\item We qualitatively assessed the place of political users within the top lists.
\end{enumerate}

\subsubsection{4 Findings}


\textit{H1. Due to the nature of the civil society in the countries, there is bigger presence of grassroots political users and NGOs in the USA and Germany in the top user lists than in Russia} (Table~\cref{tab:topTierListUserStatus}).

\begin{table}[ht]%
	\centering
	\caption{User status in the top user lists}%
	\label{tab:topTierListUserStatus}% label всегда желательно идти после caption
	\begin{adjustbox}{width=1\textwidth}
		\small
		\begin{tabular}{ l  l  l  l }% Вертикальные полосы не используются принципиально, как и лишние горизонтальные (допускается по ГОСТ 2.105 пункт 4.4.5) % @{} позволяет прижиматься к краям
			\toprule
			& \makecell[l]{Germany\\\(N=230\)} & \makecell[l]{Russia\\\(N=205\)} & \makecell[l]{USA\\\(N=205\)}\\
			\hline
			Ordinary people & 48.70\% & 36.59\% & 27.32\% \\
			Activists and politically active bloggers & 10.87\% & 13.66\% & 23.41\% \\
			Media & 23.48\% & 22.93\% & 17.56\% \\
			Political actors outside state institutions & 1.74\% & 3.41\% & 3.41\% \\
			\makecell[l]{State authorities of any level and their individual\\representatives} & 3.48\% & 6.83\% & 0.49\%\\
			Arts people, celebrities, and experts & 6.96\% & 3.41\% & 13.66\% \\
			Other/non-defined/irrelevant & 4.78\% & 13.17\% & 14.15\% \\ 
			\bottomrule
		\end{tabular}%
	\end{adjustbox}
\end{table}

The biggest share of grassroots political users is present in the USA -- almost 25\% of our sample in comparison to 11\% in Germany and 14\% in Russia. In Germany, many ordinary users have a clear political position that can be understood from their tweets but they don’t define themselves as activists. Despite the share of authorities is bigger in Russia than in Germany, the level of authorities should be taken into consideration. While the sample of Russian political users includes low-level politicians -- deputies of local councils or executives, as well as institutionalized party members outside the State Duma, the list of influential German Twitter users includes a member of the European Parliament, former Federal Minister of Justice and Consumer Protection, and members of parliamentary parties. 50\% of politicians belonging to the German Twitter influencers represent the right-wing party Alternative for Germany (AfD). Thus, H1 is partly supported for the US and Russia.

\textit{H2. Due to importance of political actors during conflicts, as well as inertia of traditional hierarchies, users who belong to the current national-scale political institutions form big nodes in the discussion networks -- that is, they are ranked high by degree, betweenness, and pagerank centralities.}

H2 is not supported for the United States, since there are no state authorities of any level and their individual representatives among users ranked high by degree, betweenness, or pagerank centralities. For Germany, it is partly supported, since there are representatives of ‘Alternative for Germany’ (rating 11 out of 50 by betweeness and 44 by pagerank), ‘Die Linke’ (The Left Party -- rating 33 by degree) and the Social Democratic Party of Germany (rating 41 for degree).

Representatives of state actors are much more authoritative on the Russian Twitter. There are 7 of them that were ranked high by degree (50\% of this type of users in the whole sample), 4 ranked high by betweenness and 8 ranked high by pagerank

Among these users, one is the member of Committee on Drug Prevention by the Federation Council, upper house of the Federal Assembly of Russia (6 by degree, 10 by betweenness, 15 by pagerank). Beside him, the three users also were ranked high by all the three metrics: a member of the Council of the youth wing of the ‘United Russia’ party, an account of the Civic Chamber of the Russian Federation and a member of a populist Liberal Democratic Party of Russia.

\textit{H3. The position of users in the top lists correlates with the users’ political status. Thus, the more politically institutionalized a user is, the more he is ‘liked’, retweeted, and commented; also, (s)he occupies a higher place in the structural hierarchy of the discussion web graph (higher InDegree, Degree, Betweenness, and Pagerank). As to the OutDegree, it inversely correlates with political institutionalization, as traditional hierarchies are reproduced in the online discussions.}

What we see from Table~\cref{tab:userStatusRhoCorrelation} is the following. First, for Germany and the USA, we see a similar pattern. The growth of the level of political institutionalization correlates with absolute figures of user activity, namely likes and retweets -- that is, political figures tend to maintain authority and keep the audience’s wish to interact with their content. This is a sign of traditional hierarchies still being in place in online discussions; belonging to an institution linked to politics slightly raises the chances to be more times liked and retweeted -- but not commented, which also shows that the nature of the discussion is not dialogue and the politically institutionalized users do not tend to enter discussions and generate mutual commenting. This is also partially supported by the fact that there is no correlation between user status and betweenness centrality -- that is, politically institutionalized users do not tend to be linking groups of users more than others, which they could do if talked to people more.

\begin{table}[ht]%
	\centering
	\caption{Spearman’s rho correlating the user status to…}%
	\label{tab:userStatusRhoCorrelation}% label всегда желательно идти после caption
	\begin{adjustbox}{width=1\textwidth}
		\small
		\begin{tabular}{ c  c  c  c  c  c  c  c  c }% Вертикальные полосы не используются принципиально, как и лишние горизонтальные (допускается по ГОСТ 2.105 пункт 4.4.5) % @{} позволяет прижиматься к краям
			\toprule
			& …likes & …retweets & …comments & …Indegree & …OutDegree &…Degree & …Betweenness & …Pagerank \\
			\hline
			USA & 0.299** & 0.235** & -- & 0.214** & \(-0.179\)* & -- & -- & -- \\
			Germany & 0.145* & 0.231** & -- & 0.150* & \(-0.227\)* & -- & -- & 0.189** \\
			Russia & -- & -- & -- & 0.257** & -- & 0.231** & 0.153* & 0.288** \\
			\bottomrule
		\end{tabular}%
	\end{adjustbox}
\end{table}

Also, for the USA and Germany, we see clear support of the idea that political institutionalization raises the number of users who interact with the given user (by InDegree metric) and diminishes the number of users who are interacted with on the part of the given user. In other words, if you belong to a political institution, you post information that people tend to like and share, and you remain attractive for more users than average, but you tend to be passive in communicating with them. This is another sign of traditional hierarchies rising in the ad hoc discussions where all the users are seemingly equal. In Germany, political users also show higher authority among the most authoritative users (by pagerank centrality), but we explain this by the nature of the case and the top user list -- while in Ferguson there were almost no politicians at all, in Germany their presence was bigger and they could enter the ‘authority chain’ of users with high pagerank centrality.

In Russia, by sharp contrast, we see that political users, despite their relatively massive presence, do not create the picture of being different from ordinary users in terms of user attention (likes, retweets, or comments), but the structure of the discussion still makes them different from the average user. They are not only referenced to by more people (not because of content, it seems, but because of their status), but the people who reference to political actors also tend to be more authoritative. It even seems that institutionalized users form something like a self-referential cluster of the discussion, with external audience seeing also referencing to it. This looks like another form of keeping hierarchies – based not on political individual but on political status itself.

But anyway, in all the three countries, politically institutionalized users do not stand in the position of linking people in the discussions and forming circles of interest around them – only in Russia their betweenness centrality is slightly connected to the user status but the correlation is really weak. This may be the sign of the situation when, while traditional hierarchies still work in people’s minds, the real role of political actors in public discussions is not that significant.

These findings are partly supported also by Cramer’s V metric (see Table~\cref{tab:userStatusVCorrelation}) -- with the difference for Germany where there are practically no significant correlations that would show up. But for the USA and Russia correlations not only are there but are really strong, showing the differences between content-based and status-based nature of influencing the discussion by politically institutionalized actors in the US and Russia, respectively.

\begin{table}[ht]%
	\centering
	\caption{Spearman’s rho correlating the user status to…}%
	\label{tab:userStatusVCorrelation}% label всегда желательно идти после caption
	\begin{adjustbox}{width=1\textwidth}
		\small
		\begin{tabular}{ c  c  c  c  c  c  c  c  c }% Вертикальные полосы не используются принципиально, как и лишние горизонтальные (допускается по ГОСТ 2.105 пункт 4.4.5) % @{} позволяет прижиматься к краям
			\toprule
			& …likes & …retweets & …comments & …Indegree & …OutDegree &…Degree & …Betweenness & …Pagerank \\
			\hline
			USA & 0.925** & 0.944** & 0.915** & -- & -- & -- & -- & -- \\
			Germany & -- & -- & -- & 0.813** \textit{(sig. 0.055)} & -- & -- & -- & -- \\
			Russia & -- & -- & -- & 0.716*** & -- & 0.747*** & -- & 0.980*** \\
			\bottomrule
		\end{tabular}%
	\end{adjustbox}
\end{table}

Thus, we see our hypothesis is partly supported for Germany and the USA; there is the pattern of attention growth for more politically institutionalized users there, based on the number of users who pay attention to politicians, as well as on their content which users find sharable in the crisis situations. In Russia, on the contrary, the pattern is not there, and the discussion first seems more deprived of hierarchies, but the structural metrics reveal that the status itself matters -- arguably, more than content these user post. This leads us to the conclusion that political presence in Twitter in democratic countries differs from that in post-authoritarian ones, especially in terms of what creates the authority.

\subsubsection{5 Conclusion}

Thus, we see that even preliminary research like ours shows that there are significant gaps in political representation within conflictual discussions on Twitter; there is room for significant improvement in terms of discussion engagement of political actors, including media and grassroots politicians.

We also see that the USA and Germany tend to differ from Russia in terms of how exactly the hierarchies are reproduced -- in Russia, it is the status itself that matters, and the more politically institutionalized a user is, the bigger chance is that he/she will form an important network node, while in the US and Germany we see that content matters more, but political users fail anyway to be discussion centers.

This provides new grounds for further research on behavior and role of politically aligned users on Twitter.

Of course we see many limitations in this kind of research; also, we have not provided in-depth analysis of the strategies of political actors and the causes that lead them to the detected (and relative) failure in uniting the user groups and discussion echo chambers. But this lies a bit beyond the scope of this paper and demands further research. What we would like to underline is that the divergent patterns of political influencing may be grounded not only in the nature of the national publics but also the nature of the national political arenas.

\subsection{A Self-Critical Public: Cumulation of Opinion on Belarusian Oppositional YouTube before the 2020 Protests}\label{subsec:ch2/sec5/sub6}

\begin{displayquote}
	My friend! You’re a major, aren’t you! You get it all! Stop the guys! To our beautiful sky! Forests! Fields! <…> I love my country so much! You too! You understand everything! Stop the guys! We have such a chance! We all have! We are for the truth! Power is in the truth! Stop the guys! And we will owe you! We, like you, are all Belarusians! Let's try to live the truth!
\end{displayquote}

These words by an anonymous protester addressed to a special forces major were pronounced during the Belarusian street protest rally on the night of 10 August 2020, after the presidential elections, won by official vote count for the sixth time by Aliaksandr Lukashenka. The quote was publicized by \textit{Nexta}, a Belarusian oppositional media now existing as a complex of Telegram channels but, in 2018, mostly active on YouTube. Starting from the election eve, the popular protest rose in Belarus; from that moment on, \textit{Nexta}’s number of subscribers on Telegram grew in a rolling manner in real time for several days, soon reaching nearly 2 million users. This, for a 9.5 million Belarusians (even with the Belarusian diasporas and \textit{Nexta}’s international followers), was an absolute record.

Such growth did not start from scratch. In the 2010s, Belarusian opposition was perceived by many as relatively disunited and existing only “off-system,” without casting any impact upon political decision-making. The formally oppositional parties in the parliament were built into the system of power and mostly supported the status quo. The non-systemic opposition found its shelter on social media, including \textit{Odnoklassniki.ru} (“Classmates”), Twitter, and later YouTube and Telegram. In 2018 and 2019, Belarusian YouTube contained several Minsk-based and regional oppositional channels that contributed to preparing the soil for the post-election protest.

The quote above mirrored a strikingly peaceful mood of the Belarusian protesters who continuously emphasized that the protests were against a particular person and the regime associated with this person, not against fellow citizens. During the rallies, the protesters carried both the white-red-white symbols of the independent Belarus of 100years ago and the current Belarusian post-Soviet flags, underlined the right of the regime supporters to hold their own rallies, brought flowers to the police, and took off shoes if they climbed upon benches. This made both international and domestic observers wonder what the Belarusian oppositional public and the consensus behind it were.

The protest seemed to touch various social groups, from plant workers to students, housewives, and owners of small businesses. Despite this, the public behind it was not amorphous, mostly because it was networked. On social media including YouTube and Telegram, the protesters discussed electoral agendas and protest facilitation, as well as criticized organizational patterns of everyday life and “mundane” discontent related to them \cite{SmoliarovaBodrunova}, and such discussions were clearly making the discussants a critical public \cite{Toepfl}. However, their discourse differed substantially from both policy-critical and leadership- critical publics, as conceptualized by Toepfl for (semi-) authoritarian public spheres.

In this article, we aim at partly covering a gap that exists in studying autocratic public spheres and publics, including those of Belarus. To find out the features of the oppositional public that existed in Belarus before the protests, we collect and analyze YouTube comments to the popular Belarusian oppositional accounts and connections between their authors for the whole of 2018. The six YouTube accounts, \textit{Nexta} among them, were selected with the help of YouTube statistics, expert evaluation by Dr Aliaksandr Herasimenka (University of Oxford, United Kingdom, as for 2021), and discussions with local activists. The number of the collected comments was 120,411 and the number of commenters was 34,770.

First, we assess the commenters’ interconnectedness, in order to define their formal structure as a closed-up community, random-user conglomerate, or a group with a more complex structure. Second, we look at the core group of cross-commenters and discursive features of their comments, including orientation to dialogue, aggression, and antagonism/agonism. Third, we ask whether this discussant group, allegedly the most critical to the regime, formed a policy- critical, leadership-critical, or other type of critical public.

The remainder of the article is organized as follows. The second section sheds light on the theory of critical publics in (semi-)autocracies and describes the recent developments in the Belarusian public sphere and politics. The third section goes deeper into the research questions, data collection, and methods of data analysis. The fourth section presents our results, followed by the concluding remarks.

\subsubsection{The Belarusian Public Sphere and Social Media in the Late 2010s}

\paragraph{Critical Publics in Non-Democracies}
Today, studies of political publicness in non-democracies \cite{GurievTreisman,StockmannLuoShen} show that public discussion takes place there just as intensely as in democratic states, even if the balance between state-affiliated and “non-systemic” actors is significantly distorted.

The very presence of publics in autocracies, including those on social media, is not in question any longer; however, we know little about the nature of such publics and their potential for social change.

Toepfl \cite{Toepfl} has suggested that autocratic publics vary by the level and direction of political criticism they express; this, in effect, means that, to become functional, a public needs to be critical. Distinguishing between uncritical, policy-critical, and leadership-critical publics, Toepfl also linked publics to their residential milieus, like networking platforms or media audiences, thus making place/environment a significant attribute of an authoritarian public.

At a birds-eye view, these milieus tend to be less intersecting than in democracies. In post-Soviet autocracies, political cleavages differ in their nature from the Western left–right patterns \cite{BodrunovaLitvinenko2015,Manaev} and, spurred by platformization, tend to create “thinking ghettoes.” One such cleavage is “the system” versus “non-systemic” actors \cite{Ledeneva}. The former includes the state, its security services, and state-affiliated political, media, and non-governmental organization (NGO) entities, including, for example, formally oppositional parliamentary parties. The latter are political parties outside parliaments, political activists, oppositional and underground media, and human-rights-watch and social-aid NGOs often aided from abroad, which allows for denouncing them as foreign agents \cite{EgorovShutovKatsuk}. Today, projects and discussion milieus on social media have joined this cohort.

In Karol Jakubowicz’s terms, the media system of Belarus has, till today, been “atavistic” with respect to the communist media systems \cite[p.~21]{Jakubowicz}. Thus, public affairs newspapers are state-run for 93\%, and TV channels are for 68\% \cite[p.~210]{Manaev}. They enjoy a major market share and preferential treatment, thus clearly defining public agendas \cite{Shirokanova}. In response, they face “natural” and “adequate” restrictions in news selection \cite{SchimpfosslYablokov}. As Jaromilek \cite[p.~89, emphasis in original]{Jaromilek} notes, “the postulated \textit{plurality of openness and access to information} could not be reached at all. As a consequence, plurality of political meanings and \textit{communicative feedback} was badly affected.” Also, Belarusian TV content has been significantly complemented by Russian channels and TV series throughout the post-Soviet time.

Along with that, autocratic strategies of working with dissent vary from repression to cooptation to “gardening” \cite{LitvinenkoToepfl} when moderate oppositional milieus/publics are tolerated or weakly supported to demonstrate liberalism. In such conditions, independent media cannot provide for substantial change, as they, with very rare exceptions, do not form wide enough publics around them. Non-systemic media actors choose between radicalization and formation of “parallel” \cite{Kiriya2012} or “alternative-agenda” \cite{BodrunovaLitvinenko2015} discussion enclaves limited to their core interested audiences and lacking communication with the authorities \cite{Shirokanova}. Portals and whole platforms may become such closed-up milieus, but we have only scarce data on whether their users feel and behave as publics and can insist on self-assembly and self-expression \cite{Shirky}. Moreover, we lack knowledge on the nature of criticism and public dialogue on such milieus; in particular, whether they are antagonistic, agonistic/adversarial, or consensual \cite{Mouffe}, and to whom they address their criticism.

\paragraph{The “Egalitarian Nationalism” and Belarusian Identity in the 2010s}
If publics depend on identities, Belarusians had for long put this under question. By the 2000s, Belarus was described as a “denationalized nation,” for which nation-building and country development did not depend upon national(ist) identity \cite{Marples}, and a country in between great powers, whose statehood experienced several false starts \cite[p.~xi]{Wilson2021}. By then, a wave of national revival of the early 1990s was reversed by Lukashenka’s policies planted on fertile ground: in the 1980s, Belarus was “the most ‘Soviet’ of the USSR (Union of Soviet Socialist Republics) republics” \cite[p.~337]{Leshchenko}, with the highest standard of living among them. Successful exploitation of post-Soviet patterns in policing and economy was fostered by tolerance of Belarusians to the USSR and historic remoteness of the ideal of independent republic with its symbols, including the white-red-white flag and the Pahonia shield of the Dutchy of Lithuania of centuries ago. In as late as 2010, Belarus still “seem[ed] better suited to a more inclusive civic identity than an exclusive ethnic one” \cite{BuhrShadurskiHoffman}.

As Leshchenko \cite{Leshchenko2008} argued, in Belarus under Lukashenka, weakness of ethnic identity was used extensively by the regime. A special sort of national ideology -- “egalitarian nationalism” -- became a cornerstone of autocratic rule, failure of democratization, and authoritarian consolidation in Belarus. For domestic use, it “discarded Belarusian ethnic references, such as the language, and employed the ethical values of a collectivist repertoire instead” \cite[p.~1420]{Leshchenko2008} (see also \cite{Kazharski}). Restored Soviet policies, language (Russian), and symbols, as well as romanticization of Soviet-style Belarusianness \cite{Gapova} and Lukashenka’s image of personal caretaker of the country, helped reach a feeling of continuity of the Soviet lifestyle for many Belarusians. Emphasis on the “national mentality traits” like non-adherence to materialism, consumerism, and individualism, has been combined with showcases of jailing of “corrupt” professorship, while disappearance of journalists and activists critical toward the regime remained under-covered in national media. East–West differences similar to Ukrainian were less salient, though not non-extant. Orthodoxy, Catholicism, and Protestantism’s peaceful co-existence contributed to depoliticization of the general public, despite the state’s growing hostility to the Roman Catholic Church and simultaneous political involvement of the Belarusian Exarchate of the Russian Orthodox Church \cite{Vasilevich}. The national language (excluding dialects) was taught at schools but spoken by less than 5\% of the population; only after the Euromaidan events in Ukraine, it spread online, mostly by language enthusiasts.

As to the international arena, Lukashenka’s ability to maneuver politically and financially between the European Union and Russia became world-infamous and finally backfired in the late 2010s \cite{Nizhnikau}. This also became part of the forking national self-consciousness that fluctuated between East and West, just as between past and future; these dimensions were not always directly associated with each other. Bekus \cite{Bekus} described not one but two pro-European philosophies of Belarusians. Thus, “Belarus is Europe” linked Belarusians to the Western neighbors by classic cultural, political, and values-based heritage. Alternatively, moving “from Russia towards Europe [in order] to reach neutrality,” meant Belarus to become a “meeting point of civilizations” and “a neutral and self-sufficient country” (p.~2). This avoidance of alienation and polarity is also mirrored in how Belarusians intertwine the Russian and Belarusian language \cite{Bekus2014} not only in the widespread dialects like \textit{trasyanka} but also on language-patchworked newspaper pages.

After Euromaidan in Ukraine and Russian integrational rhetoric of the late 2010s, the second option started to win the public mind. Instead of polarizing toward East/West, Belarusians started to talk of their own way and “letting Belarus alone,” choosing neither the “color revolution” scenario nor integration with Russia. This, of course, would be an over-generalization to say that a public consensus emerged; but, by 2018, many small-scale discussions on social media raised the issue of political, linguistic, and historic identity of Belarusians.

\paragraph{Aliaksandr Lukashenka: Between Cynicism and Fears of Turmoil}
Toepfl \cite{Toepfl} tells of leadership-critical publics, and Belarus has had a leader that definitely evoked one. Lukashenka has retained his post since 1994, partly thanks to amendments of Constitution by referenda and pushing out oppositional parties \cite{KorostelevaLawsonMarsh}. However, the grassroots support to Bats’ka (“father” in Belarusian, a popular Lukashenka’s nickname) who embodied the post-Soviet values of paternalism and “sturdy householding” must not be underestimated \cite{OhanaGenerale}.

According to Jaromilek \cite[p.~87]{Jaromilek}, Lukashenka’s regime fully consolidated by 2009; after that, Belarus ceased to be the European Union’s primary post-Soviet target country \cite{Marples2009}. Wilson \cite{Wilson2021} draws attention to how successfully Lukashenka has combined political repression (never formally proven as coming from him directly), elimination of challengers from political scene, and “authoritarian public goods” like clean cities or absence of open corruption (reached, we would add, by Draconian measures like show-off anti-corruption campaigns in universities or extremely high punishments for small bribes). Within the authoritarian-populist welfare state \cite{HortZakharov}, the policies supporting byudzhetniki (people on state-funded salaries) have ensured him a wide frontier of support before the 2010s. However, the misbalances in the Belarusian econ- omy created by this approach could not allow for economic growth, and, after the Donbass conflict started, Lukashenka has changed the social contract from “power in exchange for growth of real income” to “power in exchange for peace and security” \cite{Nikolyuk}.

Yet, despite being labeled “the last dictator of Europe” \cite{Wilson2021}, Lukashenka has experienced a peculiar treatment by Belarusians, as, for many, he was tolerated rather than supported. Similar to popular jokes and late-Soviet anecdotes, he was cynically described as a part of “the system,” a distant but unavoidable curse, in clear accordance with the post-Communist mind-set \cite{Schopflin}. This signaled the rise of a public that did not link criticism toward leadership to constructive policy criticism.

\paragraph{Oppositional Political Parties Versus Being Oppositional}
Weakness of the Belarusian opposition, both party-based and unorganized, dates back to the Brezhnev era, as the dissident movement was weaker in Belarus than in Russia or Ukraine, and even the Chernobyl disaster has not created any signifi- cant oppositional consolidation \cite{KorostelevaLawsonMarsh}. Since the early 2000s, Belarusian parties, despite the wide political spectrum they represented, experienced extremely low support and were detached from ruling legal restraints for peaceful oppositional action, lack of resources, and inability to reach intra-party consensus.

Belarusian Popular Front (BNF) and several other non-systemic parties survived throughout the Lukashenka era; but they have been marginal, just as the systemic oppositional parties were toothless. Selective repression was accompanied by allowing some personal freedom to citizens subject to their abstaining from politics, to further discourage activism and label the opposition as “noisy” \cite[p.~381]{Bedford}, freak, and renegade. Mostly due to this, being oppo- sitional started to be seen as independent from parties, while institutionalized oppositional leaders were marginalized by various means \cite[p.~5]{Kazharski}. However, “ghettoization” of activist opposition \cite{BedfordVinatier} does not equal to hushing popular political criticism. In Belarus, grassroots opposition to the “system” beyond “systemic” and “non-systemic” parties needs to be taken into account. Thus, led by oppositional leaders from time to time only, political protest has accompanied Lukashenka’s presidency nearly throughout its five terms, with peaks at elections and also in 2010, 2011, and 2017 \cite[p.~6]{Kazharski}.

\paragraph{Internet in Belarus: A State-Controlled but Contested Terrain}
Internet, a commonly recognized soil for free assembly, reached, by optimistic measurement, 79\% \cite{Smorgunov} by 2018. In official polls of 2018, Internet rocketed to 60.4\% as the second main information source defeated only by TV which, though, was losing its positions (72\% against 85.7\% in 2016), just as print press and radio were \cite[p.~190]{Hradziushka}. Also, 2018 was the year when, on one hand, YouTube peaked in popularity and, on the other hand, the state enforced a new wave of Internet regulation.

Internet control in Belarus has been infamous for its strictness and absurdity, like online access by passport in Internet cafes. This was a result of a complex of laws that gradually tightened control over “illegal content” without providing a definition for it. Thus, since 2010, all Internet sites created in Belarus had to register and store their data in Belarus, and these were used to partially block, for example, Gmail, VK.com (ex-VKontakte.ru), or political news portals. In 2014, “a new edition of the Law on media <…> raised the status of blogs to that of media” \cite[p.~15]{Shirokanova}; in 2015, banning web pages after three warnings was introduced. However, there had been no multiple legal prosecutions for online posts (unlike in Russia), and Chinese-like firewalls were never implemented. American platforms, including YouTube, operated in a climate where restrictive laws were used arbitrarily, and pirated content flourished \cite[p.~625]{Iosifidis}. Commenting on YouTube was not restricted beyond the platform limitations and the necessity to register your identity when getting online -- which, though, did not prevent creating nicknames for YouTube accounts.

\paragraph{The Belarusian Public Sphere and Social Media}
The public sphere of Belarus has mostly been studied within the “post-communist media in transition” paradigm \cite{Sparks}. It has been critically distorted due to unhindered dominance of state-owned media and state actors, as well as intrusion of ideology into education: “[D]e jure and de facto dominance of the presidency <…> entrenched the state’s discursive hegemony in the public sphere” \cite[p.~463]{Burkhardt}. Thus, for a long time, the development of the public sphere in Belarus followed the polarizational “pro/contra system” pattern described above. Scholars mostly describe its actors as belonging to either pro-state or oppositional realms, not those that bridge political positions \cite{Lesnikova,OhanaGenerale}.

Since the fall of the USSR, the Belarusian public sphere has passed several stages of development, on which various actors played the leading roles. Thus, since late 1990s, oppositional media either resided abroad like \textit{Charter97.org} or were a new breed of \textit{samizdat}, or undercover publishing, and played both informational and facilitational roles creating networks of co-thinkers around themselves \cite[p.~68]{Lesnikova}. By 2008, the activist part of the public sphere consolidated as internally polarized \cite{Jaromilek}, and most pro-Western NGOs and oppositional media were relegated to nearly underground existence \cite{OhanaGenerale} becoming largely irrelevant for the general public.

Despite that, street protests, as stated above, have been there nearly throughout the Lukashenka era. They started to be facilitated by e-mail no later than 2006 \cite[p.~29]{Shirky2011}. With the advent of social networks, the latter became facilitators of oppositional street actions, either linked to nationalism or otherwise. Such events included Livejournal-based protest flash mobs of 2006 \cite{Shirokanova}, the applauding flash mob of 2011 discussed in Vkontakte.ru and the following protest rallies \cite[p.~69]{Lesnikova} \cite{Shirokanova}, or celebrations of the 100th anniversary of Belarusian Republic in 2018 discussed on Twitter and YouTube.

In between street rallies, platforms like VKontakte.ru and Facebook gave a shelter to activism and wider politicized discourse, mostly within the younger population strata over-represented in Bynet. Many older-age Belarusians remained non-involved in political discussion beyond chatting in kitchens, and political agendas alternative to TV, as Shirokanova \cite{Shirokanova} notes, remained nearly invisible for the majority of the population. Also, apolitical Odnoklassniki.ru quickly became very popular in Belarusian countryside, contributing to platformization of the society and depoliticization of online presence of Belarusians, even if cleverly used by the opposition for hidden network building \cite[p.~236]{Herasimenka}. Thus, social media have been home for oppositional discourse, but at the same time contributed to scatteredness of oppositional powers and platform-based divides between social strata. By 2018, no platform could be called a home of opposition.

\paragraph{By 2018: An Emergent Consensus}
We also need to add that the public mood in Lukashenka’s Belarus was closer to consensual than in Putin’s Russia or post-Maidan Ukraine. Before 2020, oppositional rallies did not gather thousands, and non-systemic parties could not pro- vide for an alternative vision of future shared by many; thus, no large-scale “gardening” was necessary. In the society with no huge rich–poor gap, widespread critique of the political regime at “kitchen chats” was, before a certain moment, counter-balanced by an equally widespread feeling of “anything but not war.” In this feeling, the unprecedentedly huge trauma of World War II (when, by Soviet-time calculations, circa 25\% of Belarusians were killed) reinstated itself in fears of Ukraine-like political turmoil dismissed as irrational, economically maleficent, and leading to a real war. To our viewpoint, social peace and feelings of shared (mis-)fortune, rather than attractive prospects for future, prevented radicalism, nationalism, and “non-systemic” politics to become the main force for the oppositional consensus before 2020.

However, it is not that, before 2020, Belarusians remained constantly tolerant to Lukashenka and its policies. Thus, the first “anti-systemic” turn in public opinion was noticed in the after-recession 2009–2010, along with a decline in the standard of living; at that moment, agreement with the state policies dropped, and preferences toward Europe rose. This movement was, presumably, partly reversed due to the Crimea and Donbass crisis. By 2017, a growing feeling of uncertainty and instability was properly addressed neither by the authorities nor by the opposition \cite{Manaev2017}; the “light-version” reforms that started in 2017 stopped in 2018 \cite{BelarusianYearbook}.

And, as stated above, by 2018, a consensus of another sort had started to form in Belarus. Wilson \cite{Wilson} points out that it emphasized the concept of sovereignty. Neither Russia “rising from its knees” nor Europe widely consid- ered as intermingling into the Ukrainian politics provided an attractive model of future. It was a gradual but rather fundamental turn since the mid-2000s when only 7\% rejected closer ties with both Russia and the European Union \cite[p.~225]{RontoyanniKorosteleva}. Below, looking at YouTube talk of 2018, we will try to show what the essence of this consensus was and whether there was a public that was its bearer.

\subsubsection{The Research Questions and Methodology}

\paragraph{The Research Questions}
As said above, and in accordance with the literature review, we have formulated the following research questions (RQs).

\paragraph{RQ1. Was the pool of commenters to oppositional chan- nels on the Belarusian YouTube a closed-up community, a group of random visitors, or a conglomerate with a more complex structure?}

\paragraph{RQ2. Was there a core group of commenters and, if yes, was it dialogue-oriented? And what was the nature of the dialogue -- was it antagonistic (confrontational), adver- sarial (recognizing differences), or consensual (express- ing agreement and/or movement to another political position)? What was the constellation of these modes in commenting?}

\paragraph{RQ3. Did this discussant group, allegedly openly critical to the regime, formed a critical public? Was it policy-critical, leadership-critical, or another critical public?}

We do not state any research hypotheses, as our research is exploratory.

\paragraph{Data Collection}
In selecting the YouTube channels for our dataset, we have used the following logic. First, not only Minsk-based but also regional channels had to be represented. Second, it had to be media-like channels, as those of oppositional parties were more oriented to party members; we also wanted to see how people comment on media content alternative to official TV, as media-like content continues to be publicly attractive and is perceived as pursuing public interest, not private polit- ical interests. After initial monitoring, we also followed the expert advice by Aliaksandr Herasimenka, a Belarus-born researcher on Belarusian opposition and media.

As a result, we have chosen the following six channels:
\begin{itemize}
	\item \textit{Belsat}, a Polish satellite channel in Belarusian, of clearly oppositional stance;
	\item \textit{Nexta} (Minsk);
	\item \textit{Garantiy net} (No guarantee, Homel);
	\item \textit{Narodny reportyor} (Popular reporter, Brest);
	\item \textit{Leave the Vagon!} / \textit{The Belarusian Experimental Field} (rural area);
	\item \textit{Rudabelskaya pakazukha} (\textit{Rudabelka} show-off, rural area).
\end{itemize}

The dataset was collected and pre-processed via a specialized web crawler with changeable modules \cite{BodrunovaLitvinenkoBlekanov2017} (for the use on YouTube, see the study by Bodrunova, Litvinenko, Blekanov, \& Nepiyushchikh \cite{BodrunovaLitvinenkoBlekanov2021}). The number of comments was 120,411, and the number of commenters was 34,770; beside this, data for cross-account commenting and data for graph metrics were uploaded. The only limitation was that the markers of the position of comments in threads (root comment/response) could not be uploaded; however, this was a minor limitation for us, as we were interested more in the discursive reciprocity than in formal interactivity.

\paragraph{Data Analysis}
To address RQ1, we have used web graph reconstruction (the Gephi library, ForceAtlas2, and OpenOrd algorithms). Within the graph, we have detected how many users have commented on several accounts, thus interlinking the accounts and creating an all-encompassing discussion core. The latter comprised the users who commented on at least five of six accounts. We have also qualitatively assessed the graph structure, the user metrics curves, and the key commenters to see whether the discussion core was separated from its periphery and corresponded to the structure depicted by Belarusian authorities.

To address RQ2, we have coded all the posts of the discussion core for dialogue markers. First, we defined presence of appeal to another user in the comments -- as stated above, it differed from formal belonging to comment threads. Second, we coded the antagonistic OR agonistic/adversary OR consensual nature of comments. After pre-reading, we coded:
\begin{enumerate}
	\item As antagonism (including calls for unification) both against the authorities/media/security services and against the fellow commenters;
	\item As willingness to discuss. questions to commenters and mentions of their arguments;
	\item As agreement with statements, expressions of approval and support to authors of videos were excluded if they lacked substantial claims.
\end{enumerate}

Third, we have coded the posts for presence of aggression and have checked the correlations (Spearman’s rho and Pearson’s) between aggression, on one hand, and overall orientation to dialogue, antagonism, agonism, and consensus, on the other hand. For coding posts, pairs of coders were trained and their agreement rate tested (Cohen’s kappa ranging from .72 to .93 for all the variables).

For RQ3, after preliminary self-informing reading and suggestion on the types of critical publics, we have coded the posts in the discussion core for the addressees of criticism. Correlations between the type of addressee and the RQ2 variables were also tested. Interpretive reading and elements of discourse analysis were employed for qualitative assessment of user statements and referrals \cite{Evolvi}. K-means clustering was used to find sub-publics and assess the relative contribution of the discursive features into types of sub-publics.

The results of our study are presented below.

\subsubsection{The Results: A Self-Critical Public}

\paragraph{RQ1: A Community or a Public?}
As mentioned above, we have uploaded the users’ comments together with their metadata, created a web graph of commenting (see Figure~\cref{fig:discussionWebGraph-1}), and assessed the resulting user metrics.

%\begin{figure}[ht]
%	\begin{minipage}[b][][b]{0.49\linewidth}\centering
%		\includegraphics[width=0.5\linewidth]{knuth1} \\ a)
%	\end{minipage}
%	\hfill
%	\begin{minipage}[b][][b]{0.49\linewidth}\centering
%		\includegraphics[width=0.5\linewidth]{knuth2} \\ b)
%	\end{minipage}
%	\caption{The discussion web graph: (a) by the ForceAtlas2 algorithm and (b) by the OpenOrd algorithm}
%	\label{fig:discussionWebGraph}
%\end{figure}

\begin{figure}[ht]
	\centerfloat{
		\hfill
		\subcaptionbox{\label{fig:discussionWebGraph-1}}{%
			\includegraphics[width=0.5\linewidth]{discussionWebGraph1}}
		\hfill
		\subcaptionbox{\label{fig:discussionWebGraph-2}}{%
			\includegraphics[width=0.5\linewidth]{discussionWebGraph2}}
		\hfill
	}
	\legend{lilac: commenters to \textit{Nexta};
		gray: users who commented in 2 to 5 accounts; 
		orange: commenters to \textit{Belsat};
		green: commenters to \textit{No guarantee};
		light blue: commenters to \textit{Popular reporter}; 
		yellow: commenters to \textit{Leave the Vagon!};
		blue: commenters to \textit{Rudabelka show-off};
		red: users who commented in all accounts.}
	\caption[Этот текст попадает в названия рисунков в списке рисунков]{The discussion web graph: (a) by the ForceAtlas2 algorithm and (б) by the OpenOrd algorithm}\label{fig:discussionWebGraph}
\end{figure}

As Figure~\cref{fig:discussionWebGraph} shows, attractiveness of the channels was unequal. Six of 10 commenters commented to \textit{Nexta}, and \textit{Nexta} only (59.34\%). \textit{Belsat} (11.54\%) and \textit{No guarantee} (5.35\%) came next, while \textit{Popular reporter} (3.62\%), \textit{Leave the Vagon!} (2.14\%), and \textit{Rudabelka show-off} (0.86\%) were least popular. This is explainable, given the latter channels’ local status.

Beside these users, circa 17\% of commenters united the channels and formed an inter-channel discussion core, as they commented to more than one channel. Thus, 7,076 users (16.96\%) commented on two to five channels; 75 users commented on all the six channels. The latter sub-sample will be used in our studies of the nature of Belarusian oppositional public.

To further clarify the core/periphery constellation, we have reconstructed the graph by another algorithm, OpenOrd (Figure~\cref{fig:discussionWebGraph-2}). Figure~\cref{fig:crossCommentersMetricsCurves} shows two modules, one comprising four channels and one with the two remaining ones, with gray and red users in between bridging densely the two nebulae. The periphery links mostly to Nexta and Belsat, the two channels closest to classic media by format. Activist chan- nels create denser commenting publics, but they are all interlinked.

\begin{figure}[ht]
	\centerfloat{
		\includegraphics[scale=0.75]{crossCommentersMetricsCurves}
	}
	\caption{Cross-commenters’ metrics curves: (a) betweenness centrality (\(K\)), user long tail cut at user \#553 (\(K = 0.0001\)); (b) pagerank centrality (\(M\)), user long tail cut at user \#410 (\(M = 0.00001\)); (c) indegree centrality (\(N\)), user long tail cut at user \#610 (\(N = 10\)). The six commented accounts are excluded.}\label{fig:crossCommentersMetricsCurves}
\end{figure}

To avoid algorithms-induced distortions in our judgment, we have also assessed the curves of user graph metrics, namely betweenness, pagerank, and degree centralities (Figure~\cref{fig:crossCommentersMetricsCurves}). Graphs of betweenness (“user-as-crossroads” importance) and indegree (number of incoming comments) centralities show smooth decline, which means that the borders of the core are not sharp. The pagerank graph shows five important users (all being ordi- nary citizens) and a group of \(\sim 70\) users of secondary importance. Overall, the graphs reveal that the commenters were not a closed-up community the way that Belarusian authorities often depicted them but a public with a cross-commenting core and a wide but active periphery (3+ comments per user on average).

Also, we have assessed the top 100 users by the three centralities, to see whether they represented foreign institutions or citizens whose presence would support the claim of “foreign agents” used by the authorities to discredit the opposition. To our best possible judgment, within the three top 100 lists, there were several media-like channels (\textit{Belarusian world, Real Belarus! Homel Society}, and \textit{Selvestor Vivat} of oppositional stance, \textit{West Polesian} focusing on folk culture, \textit{fanzone fanzone} on sports, and \textit{RadioDestroyer} on self-made radio stations) and three foreign users (one male from Poland, one mum with a little daughter from Germany, one unidenti- fied). One alleged “foreign-agent-like” account, \textit{Biełarus z-za miažy} (Polish “Belarus from abroad”), indeed, was present within top 100 pagerank pages, but we could not learn its stance, as it was completely unavailable in 2021. One Polish account was also found among 76 cross-commenters.

Thus, the active public we have discovered was genuinely Belarusian, with only minor presence of foreign users and media-like initiatives.

\paragraph{RQ2: The Conflictual Nature of the Public}
Our coding of 8,644 posts shows that antagonism was highly salient (3,131 posts or 36.2\%); thus, the discourse was, indeed, permeated by conflict. However, agonistic/adversarial mood was also salient enough (1,112 posts or \(\sim13\)\%); but it manifested only in readiness to discuss (posing questions and reacting to statements in a responsive mode or partial agreement), not in readiness to recognize differences. Consensual claims beyond emotional support (198 posts or 2.9\%) related exclusively to agreement with co-thinkers and authors of the videos, not to agreeing with ideological counterparts.

Two other dimensions of agonism/antagonism were pres- ence of dialogue markers (2,087 posts or 24.1\%) and aggression (883 posts or 10.2\%). The share of aggressive posts resembles that on the Russian political YouTube \cite{BodrunovaLitvinenkoBlekanov2021}. Interestingly, correlations (see Table~\cref{tab:discursiveFeaturesCorrelationsComment}) show that dialogue is linked to both antagonism and agonism. For antagonism, this is due to aggressive rebuttals toward alleged pro-Russian trolls; and this explains why dialogue is also linked to aggression. For agonism, it is genuine dialogue with questions to fellow commenters. Expectedly, for individual comments, aggression directly correlated with antag- onism and inversely with agonism, thus proving that posing substantial questions lowers aggression (even slightly more than agreement does!).

\begin{table}[ht]%
	\centering
	\caption{Correlations of discursive features: on the level of a comment}%
	\label{tab:discursiveFeaturesCorrelationsComment}% label всегда желательно идти после caption
	\begin{adjustbox}{width=1\textwidth}
		\small
		\begin{tabular}{ l  l  l  l  l  l  l  l  l }% Вертикальные полосы не используются принципиально, как и лишние горизонтальные (допускается по ГОСТ 2.105 пункт 4.4.5) % @{} позволяет прижиматься к краям
			\toprule
			& Aggression & Criticism & \makecell[l]{Criticism:\\Leadership} & \makecell[l]{Criticism:\\Policy} & \makecell[l]{Criticism:\\Self} & Antagonism &  Agonism & Agreement \\
			\hline
			Dialogue & .\textbf{242**} & \textbf{\(-.213\)**} & \textbf{\(-.166\)**} & \(-.073\)** & \(-.056\)** & .143** & \textbf{.190**} & \\
			Aggression & 1 & \(-.044\)** & \textbf{.049**} & \textbf{\(-.080\)**} & \(-.084\)** & \textbf{.394**} & \textbf{\(-.099\)**} & \(-.049\)**\\
			Criticism &  & 1 & \textbf{.707**} & \textbf{.329**} & \textbf{.372**} & \textbf{.247**} & \(-.055\)** & \(-.070\)**\\
			Criticism: Leadership & & & 1 & \(-.102\)** & \textbf{\(-.146\)**} & \textbf{.361**} & \textbf{\(-.086\)**} & \(-.065\)**\\
			Criticism: Policy & & & & 1 & \(-.087\)** & \(-.061\)** & & \(-.024\)**\\
			Criticism: Self & & & & & 1 & \(-.071\)** & .033** & \\
			Antagonism & & & & & & 1 & \textbf{\(-.289\)**} & \textbf{\(-.112\)**}\\
			Agonism & & & & & & & 1 & \(-.059\)**\\
			Agreement & & & & & & & & 1\\
			\bottomrule
			\multicolumn{9}{@{}p{\textwidth}}{%
				%				\vspace*{-4ex}% этим подтягиваем повыше
				\hspace*{2.5em}% абзацный отступ - требование ГОСТ 2.105
				Note. *\(p \le .05\); **\(p \le .01\). Correlations supported on the aggregate level are in bold.
			}\\
		\end{tabular}%
	\end{adjustbox}
\end{table}

To check whether discursive features depended on the number of comments by one user, we aggregated the comments by user, calculated percentages of comments for each feature, and checked the dependencies (see Table~\cref{tab:discursiveFeaturesCorrelationsAggregate}). The main correlations were supported and only involvement into dialogue grew substantially when the number of comments by a user grew.

\begin{table}[ht]%
	\centering
	\caption{Correlations of discursive features: on the aggregate level}%
	\label{tab:discursiveFeaturesCorrelationsAggregate}% label всегда желательно идти после caption
	\begin{adjustbox}{width=1\textwidth}
		\small
		\begin{tabular}{ l  l  l  l  l  l  l  l  l  l  l }% Вертикальные полосы не используются принципиально, как и лишние горизонтальные (допускается по ГОСТ 2.105 пункт 4.4.5) % @{} позволяет прижиматься к краям
			\toprule
			& Dialogue & Aggression & Criticism & \makecell[l]{Criticism:\\Leadership} & \makecell[l]{Criticism:\\Policy} & \makecell[l]{Criticism:\\Self} & Antagonism &  Agonism & Agreement & \makecell[l]{Agonism +\\Agreement} \\
			\hline
			Number of comments & .416** & & &  & & & & & & \\
			
			Dialogue & 1 & .\textbf{299**} & \textbf{\(-.238\)**} & \textbf{\(-.249\)**} & & & & \textbf{.479**} & & .442**\\
			
			Aggression & & 1 & & \textbf{.425**} & \textbf{\(-.304\)**} & & \textbf{.790**} & \textbf{\(-.298\)**} &  & \(-.323\)**\\
			
			Criticism & & & 1 & \textbf{.766**} & \textbf{.305**} & \textbf{.327**} & \textbf{.511**} & & & \(-.055\)** \\
			
			Criticism: Leadership & & & & 1 & \(-.102\)** & & \textbf{\(-.268\)**} & \textbf{.629**} & \textbf{\(-.311\)**} & \(-.339\)**\\
			
			Criticism: Policy & & & & & 1 & \(-.087\)** & & & \(.229\)** & \\
			
			Criticism: Self & & & & & & 1 & & & & \\
			
			Antagonism & & & & & & & 1 & \textbf{\(-.340\)**} & \textbf{\(-.225\)**} & \(-.378**\)\\
			Agonism & & & & & & & & 1 & & \(.966\)**\\
			Agreement & & & & & & & & & 1& .439**\\
			\makecell[l]{Agonism +\\Agreement} & & & & & & & & & & 1\\
			\bottomrule
			\multicolumn{11}{@{}p{\textwidth}}{%
				%				\vspace*{-4ex}% этим подтягиваем повыше
				\hspace*{2.5em}% абзацный отступ - требование ГОСТ 2.105
				Note. *\(p \le .05\); **\(p \le .01\). Correlations supported on the aggregate level are in bold.
			}\\
		\end{tabular}%
	\end{adjustbox}
\end{table}

Median percentage for dialogical comments per author was less than 10\% (9.6\%), even if 25 of 75 users had \(\ge 20\%\) of dialogue-oriented comments. Aggression was much lower (median \(= 4.5\%\)), with 10 users posting \(\ge 20\%\) aggressive comments. However, antagonism was high (median \(= 33.4\%\)), with 64 users allowing it in \(\ge 20\%\) comments and 13 reaching \(\ge 50\%\), only partly compensated by agonism and agreement (combined median \(= 11.9\%\)).

Altogether, the shape of conflict on Belarusian oppositional YouTube was antagonistic, with commenters’ firm conviction of particular views, not seeking agreement with opponents, and dialogue mostly aiming at rebuttals toward alleged trolls. However, it was relatively non-aggressive and, in many cases, showed readiness to pose questions to opponents and co-thinkers.

\paragraph{RQ3: Direction of Criticism and Discursive Clusters}
However, the true main feature of the discourse was criticism. It differed from aggression (“The Ministry of Education and the education departments [of local administrations] are destroying our Belarusianness”), and was not always antagonistic.

Following Toepfl’s \cite{Toepfl} division of authoritarian publics into uncritical, policy-critical, and leadership-critical, we have coded presence of criticism and its direction. After preliminary reading, however, we have added one more category, which was self-criticism -- critique addressed to “us ourselves,” Belarus as a country/society, or oneself personally. Strikingly salient in the dataset, self-criticism sharply distinguished the Belarusian discourse from, for example, the Russian one of 2019. For the reasons of objectivity, we have also coded support to authorities as a counterbalance.

The latter, though, was expectedly next-to-absent, found in 17 comments, with 3 among 494 comments by one user as a maximum. Criticism, on the contrary, was present in 3,801 comments (\(\sim 44\%\)), with only one user having \(\le 20\%\) of critical posts, and 36 users of 75 being overwhelmingly critical (\(\ge 50\%\) comments). However, the direction/addressee of this criticism varied.

The highest amount of criticism was evoked by the “system” \cite{Ledeneva}, with the median \(= 27.7\%\). This included Aliaksandr Lukashenka himself; authorities in general and civil servants in particular; militiamen; and pro-Russian trolls. Lukashenka is described as “tsar” and \textit{peresident} (“over-sitter”), \textit{ne Bat’ka} (as “parents are not elected”), \textit{lukashescu} (reminiscent to Ceausescu), and \textit{lukavy} (“cunning”) afraid of his own people. More troubling, though, are rare descriptions of lukanomics and demands for dismantling of presidency as an institution, along with support of fair elections and freedoms:

\begin{displayquote}
	…Turnover of power, or, better, absence of presidentship in Belarus is key to normal statehood.
\end{displayquote}

The authorities are described as arrogant, shame-evoking, thieves who “feed themselves at the trough” while lying and imitating work, oppositional to people, and powerless to decide. Militiamen are compared to jackals and Hitlerjugend. Often, together they are compared to occupants, and the country is called “occupied” and people “enslaved.” Antagonism in the dataset was mostly linked to descriptions of power as insane, incompetent, barbaric, showing-off, mistrusted, unduly rich, incapable of negotiation, and laughable at. The hopeless mood of “the further the worse” was dominant.

Many comments are addressed to alleged pro-Russian commenters, while, in the core, there was only one pro-Russian user. Pro-Russians were shamed (sometimes aggressively) as trolls and bots, while recognizable responses to (pro-)European or (pro-)American commenters were extremely rare (six of all). This might be a sign of pro-Russian “trollization” of the periphery of the discussion, oppo- site to the authorities’ claims of European impact upon online communication. In contrast to trolls, oppositional bloggers are a “punch in the gut” to the “system.”

In Toepfl’s theory, leadership-critical publics, presumably, appear where policy-critical publics already exist. However, we have discovered a leadership-critical public without policy criticism. For a Western observer, this is a paradox; for autocratic publics, though, cursing “the system” without assessing its individual shortcomings is characteris- tic, due to absence of traditions of publicly discussing policies. For policy criticism, the median \(= 5\%\); the users only mentioned some examples of bad policy decisions, without giving reasons or suggesting alternatives.

Self-criticism was more impressive. It stably spread through the dataset, with the median \(= 9\%\), and included sarcasm toward official slogans like “Belarus -- a country for living,” criticizing \textit{pamyarkounasc}’ - forbearance and long-suffering viewed as traditional for Belarusians, and, most often, stating citizen’s responsibility for letting the country into its current state:

\begin{displayquote}
	Quite livable here -- for an American pensioner.
\end{displayquote}

\begin{displayquote}
	[I feel] shame, because we are all responsible that our country is ruled by exactly this [person].
\end{displayquote}

\begin{displayquote}
	Everyone wants change but no one wants to do anything.
\end{displayquote}

\begin{displayquote}
	We need, in general, an alternative Belarusian national life.
\end{displayquote}

The narrative of self-victimization intertwined with two more narratives -- the necessity of unified efforts and deeds instead of words within legal boundaries -- 

\begin{displayquote}
	I agree this is scary, but law needs to be our weapon, if we do not follow the law, we will become them.
\end{displayquote}

\begin{displayquote}
	Alas, most probably, no way via the courts, as this is their arms against ordinary people. <…> our only armament against the system is solidarity.
\end{displayquote}

\begin{displayquote}
	We need to find other ways. It will anyway depend fully on people’s activity.
\end{displayquote}

\begin{displayquote}
	Even a critical mass of amoebae will never explode.
\end{displayquote}

-- and the narrative of the aforementioned “third way” for Belarusians:

\begin{displayquote}
	We will go another way. Not to NATO, not to the Russian Federation. How about this option? Neutralism.
\end{displayquote}

The users state that they want to “make friends, not enemies” with “all the countries of the civilized world,” including Europe, the United States, and Russia -- which, regretfully for the commenters, choses more and more a way to self-isolation. In such statements, neither policy criticism nor self-criticism is linked to aggression (see Table~\cref{tab:discursiveFeaturesCorrelationsAggregate}), and the more leadership criticism is present, the less self-criticism is expressed by the users.

Another addition to Toepfl’s view is that discursive features like dialogue or aggression/antagonism need to be used in combination with criticism to define the types of autocratic publics. By using k-means clustering, we have shown that these features allow for seeing difference between actively aggressive-critical, passive-critical, non-aggressive dialogical modes of discussion (see Figure~\cref{label}).

\begin{figure}[ht]
	\centerfloat{
		\includegraphics[scale=0.75]{crossCommentersSubClusters}
	}
	\caption{Sub-clusters of cross-commenters by discursive features. Three clusters represent the best decision by the silhouette metric; \(p \le .0000\) for all the clusters}\label{fig:crossCommentersSubClusters}
\end{figure}

\subsubsection{Discussion and Conclusion}

In general, what we have seen on Belarusian oppositional YouTube of 2018 was, to our viewpoint, the graduate cumulation of opinion that, in August 2020, reached a threshold and spilled over to Belarusian streets.

We have shown that, at least for 2018, the oppositional public on Belarusian YouTube was open but non-random; genuine but already highly alert against trolls and bots; detached from political parties but reaching the critical point of politicization. It bore the markers of readiness to massive protest, such as acknowledgment of people’s own guilt and calls for solidarity against the “system.” Without much dialogue on policing, criticism toward the systemic features of the state combined with unusually high self-blaming and non-acceptance of Belarus and Belarusians themselves in their long-suffering. The commenters’ positions reflected what later formed a countrywide consensus and partly laid a foundation for the 2020 protests.

Studies like ours allow for capturing the public mood better than polls or political party research; we have detected a combination of readiness for change, “adversarial antagonism,” and self-criticism. The Belarusian case demonstrates that antagonism on social media is a double-edged sword: it allows uniting in hatred but eliminates chances for deliberation of participants with opposing standpoints.

Our additions to the theory of autocratic publics include presence of leadership-critical publics without policy criticism in their discourse, lack of policy-and-leadership criticism, importance of self-critical narratives, and multi-dimensionality of discourses and (if publics are characterized by discourses) of publics. The linkages between antagonism/agonism and criticism of various directions need to be further explored, too. And, last but not the least, tracing today’s changes in the Belarusian oppositional discourse online would provide a unique chance to see how such discourses mutate after a relative failure (or victory?) of a nationwide resistance wave.

\FloatBarrier
           % Глава 2
\chapter{Сетевая коммуникация}\label{ch:ch3}

\subsubsection{3.1.1}

Recently, communication scholars have paid attention to the growing dissonant and dissipative character of the public spheres, especially in their connection to networked discursive spaces. While substantial dissonance of the discussions is well addressed, structural discontinuity of public discussion remains under-explored. Reproducibility of the discussions on similar issues or events in time, we argue, needs to be seen as a marker of stability of public spheres. In this paper, we compare the user and influencer structure of two similar discussions on German Twitter of 2016 (the Cologne mass harassment) and 2019 (the Chemnitz killing). We show that the overall reproducibility of the discussions is extremely low, and the only structural element that reproduces are influential media, mostly of national reach. But even the stability of media presence must be questioned, as both intensity of their presence within the discussion and user engagement with their tweets varies much from one discussion to the other. Thus, one may conclude that the structural stability of public discussion of similar events on Twitter is not reached.

\subsubsection{3.2.1}

Предметом исследования является использование пиктограмм для выражения эмоций и чувств в сетевой коммуникации. Объектом исследования выступили комментарии в нескольких популярных у молодежи сообществах социальной сети «ВКонтакте». Целью исследования было сравнение предпочитаемых пиктограмм-эмодзи, которыми выражают свои эмоции и чувства пользователи, состоящие в сообществах с разной частотой использования обсценной лексики. Гипотезой исследования выступило предположение, что группы с разной частотой использования обсценной лексики отличаются эмоциональной тональностью и используемыми пиктограммами-эмодзи для выражения эмоций и чувств. Методами исследования выступили опрос и контент-анализ. Для реализации первого метода была составлена анкета, в которой респондентам были заданы вопросы об особенностях их интернет-предпочтений, наиболее привлекательных для них интернет-сообществах в социальных сетях. Проведен анкетный опрос на выборке 854 человека, осуществлен сбор данных о частотности использования пиктограмм-эмодзи в 14 выделенных для детального анализа сообществах социальной сети ВКонтакте. Выявлено, что в сообществах, где пользователи реже используют ненормативную лексику, рейтинг эмодзи, обозначающих положительные (гедонические и сближающие) эмоции выше. В сообществах, где нецензурная брань используется чаще, пользователи наравне с позитивными, часто используют пиктограммы выражающие негативные эмоции -- удаляющие (гнев, злость) и меланхолические (печаль, тоска). Результаты исследования могут быть применены в разработке методов мониторинга настроений участников сетевой коммуникации.

\section{Принципы взаимодействия и распространения информации}\label{sec:ch3/sec1}

\subsection{Discontinued Public Spheres? Reproducibility of User Structure in Twitter Discussions on Inter-ethnic Conflicts}\label{subsec:ch3/sec1/sub1}

\subsubsection{1. Introduction}\label{subsubsec:ch3/sec1/sub1/subsub1}

Recently, communication scholars have stated that today’s public spheres \cite{Habermas} have become increasingly dissonant \cite{Pfetsch} and disintegrated. They rarely have consensus as a goal and, to a large extent, consist of ad hoc discussions \cite{BrunsBurgess} that have no continuity and quickly dissipate. Disconnection also emerges through ‘ever-more-fiercely negative campaigns, increasing political polarization, and public debates filled with prejudices and false assumptions’ \cite[p.~59]{Pfetsch}. One dimension of this disintegration has remained virtually unexplored, which is time. In traditionally mediated public spheres \cite{Calhoun}, the main structure of information flows organized by media and institutions remained stable in years or even decades; but we hardly know whether the user and influencer structure of networked discussions remains stable or changes with time, and to what extent. How can we expect public discussions to come to definitive conclusions stably accepted or at least discussed, if the very participant structure is unstable?

We argue that reproducibility of the discussion structure may be viewed as a sign of the long-term continuity of the public spheres and, thus, work as their quality metric. Long-term studies of networked discussions (\textit{e.g.} their polarization) remain rare despite the measurement’s tools accessible at the platform such as networking, tweeting and content-producing behavior of users \cite{GarimellaWeber}. ‘Most longitudinal network studies have confounded the processes of new tie formation and old tie maintenance, resulting in an incomplete understanding of the processes of network change’ \cite{Fu}. Studying Twitter data in timelines and investigating timeline narratives longitudinally are, till today, atypical approaches to data gathering and analysis \cite{BrookerVinesBarnett}; comparisons between structures of similar discussions of various times are next-to-absent.

Conflictual discussions online are vivid forms of expression of the public sphere \cite{BodrunovaBlekanovSmoliarova}. Among them, Twitter discussions are the most rapidly growing and, sometimes, most quickly dissipating. Our previous findings have shown that Twitter ‘is more complicated than the imaginary cocooned talk in echo chambers, especially for issues beyond elec- tions and direct policing’ \cite[p.~130]{BodrunovaBlekanovSmoliarova}. Moreover, according to the studies of digital protest and ‘hashtag activism’, Twitter may enable longitudinal campaigning that changes structure and content of public discussions for long enough time periods \cite{BonillaRosa}. Technological affordances equally contribute to the formation of social movements and the spread of false beliefs in the ‘unedited public sphere’ \cite{BimberDeZuniga} that emerges on social media. They also simultaneously allow for keeping the talk alive and forgetting it as soon as it goes beyond the Twitter scroll.

Case studies of Twitter discussions on selected conflicts paid more attention to the content of debate and the character of public communication than to who and how long participates in the discussion \cite{GroshekTandoc}. Also, comparative approaches have mostly been applied to simultaneous cases in different countries, not to similar discussions that hap- pened within one national context in varying times, with few exceptions when the scholars talk about a social/political movement with its evolving dynamics of activities (e.g. \#blacklivesmatter).

This study aims at comparing two national-level discussions in German-language Twitter about ethnicity-related conflicts: the one on the 2016 New Year’s Eve sexual assaults in Cologne and that on the 2018 protests in Chemnitz that took place after the death of a Cuban-German man. We ask to what extent the structure of influencers (ordinary users, media, and institutions) has remained similar, since the issue and public polarization behind it were the same. Comparison within the national context reduces bias present in cross-national studies, as the variety of influencers cannot be explained by varying national political cultures.

The structure of this paper is organized as follows. The next sections presents our research questions and methodology (Sect.~2) and the findings of our study (Sect.~3). We conclude with a short summary of our results (Sect.~4).

\subsubsection{2. Research Questions and Methodology}

As stated above, the study aims at comparing two national-level discussions in German- language Twitter about ethnicity-related conflicts: the one on the 2016 New Year’s Eve sexual assaults in Cologne and that on the 2018 protests in Chemnitz that took place after the death of a Cuban-German man. Our general assumption is that the structure of participation, as well as the structure of influencers, should remain similar, as the issue and public polarization behind the discussions were the same.

We conducted vocabulary-based web crawling to collect the discussion content and to analyze their structure. A special web crawler has been developed to bypass the limitations of Twitter API \cite{BlekanovSergeevMartynenko}. The total number of users in the two datasets included 12,382 users for Cologne and 22,973 users for Chemnitz.

As we moved from step to step in our research, we corrected the research questions and hypotheses based on the findings of the previous steps. Here, we will describe the RQs and methods used for each step, but the hypotheses are attached to the description of our results.

\paragraph{RQ1. Did Twitter users who discussed the Cologne case join the Twitter discussion about the Chemnitz case?}

To answer RQ1, we assessed the general reproducibility of the discussions by defining the number of all users and influencers who participated in both discussions.

\paragraph{RQ2. Does the structure of politically relevant (institutionalized political, grassroots political, and media) influencers repeat from Cologne to Chemnitz?}

To answer RQ2, we assessed the structure of influencers. For this, we sampled 50 top users for each of eight user metrics: number of tweets; number of interactions (likes, comments, and retweets); centralities "--- indegree, outdegree, betweenness, and pagerank. Then we merged them in aggregate lists of top users and eliminated the duplicates. As many users were within the top lists by many metrics, the final top lists included 230 users for the Cologne case and 207 users for the Chemnitz case. We manually coded the users in the final lists for their offline status in the following way. We coded an account as ‘ordinary user’ if its owner mentioned neither an institutional status nor political positioning in the user self-description at the top of the Twitter blog. If any political positioning or support of political values have been mentioned in the description, the account has been coded as ‘politically active blogger/activist.’ As ‘media’ we coded the accounts that either have a clear connection to a media project outside Twitter or position themselves as Twitter-only media projects. Institutional political actors included the accounts run by state authorities of any level and their individual representatives; political parties; individual politicians. Any other possible classification based on the description in the user account was marked as ‘other’.

\paragraph{RQ3. Does the salience of intersecting influencers differ in the two cases?}

To answer RQ3, we have defined the list of relevant intersecting influencers and have calculated their activity ratio and user engagement ratio. The former is the number of tweets published by the account in the first case divided by the number of tweets published in the second case; the latter is the ratio of the according numbers of user comments.

These ratios show whether the status of these influencers, both actively pursued and user-supported, has been stable in time.

\subsubsection{3. Findings}

\paragraph{3.1 Reproducibility of the General User Structure in the Discussions}

\paragraph{RQ1. Did Twitter users who discussed the Cologne case join the Twitter discussion about the Chemnitz case?}

\begin{itemize}
	\item H1a. \textit{Ad hoc} publics that emerged around the Cologne and Chemnitz cases repeat to a significant share (over 30\% of the participants who discussed Cologne also discussed Chemnitz).
	\item H1b. The list of top users changed less than the general sample of users taking part in the discussion, as influencers are the structural carcass of the issue-based debates.
\end{itemize}

To prove H1a we compared two general samples crawled for the Cologne case and for the Chemnitz case. From the whole dataset, we excluded those users who have not posted at least once and only interacted with those who tweeted by liking or sharing. (These users were crawled because they were necessary for the graph reconstruction of the discussions; we consider them irrelevant for RQ1.) The datasets for RQ1 included 12,382 users who tweeted about Cologne and 22,973 users who tweeted about Chemnitz. Only 1,735 users (circa 14\% of Cologne sample and 7,5\% of Chemnitz sample) participated in both discussions posting at least one tweet with at least one hashtag that defines the corpora of each discussion. 17\% of them published the same amount of posts (from 1 to 7 tweets in each discussion), 57\% tweeted more about Cologne and 26\% were more active in posting about Chemnitz. The total number of unique users amounts to 33,620 users; thus, the share of intersecting users is 5,2\%. If we exclude those who posted less than three tweets in both cases, the share of users participated in both discussions decreases even to 1,2\%. Therefore, H1a has to be rejected: the \textit{ad hoc} publics the emerged around the Cologne and Chemnitz events vary by 95\%.

To check H1b, we selected the top users (influencers) for both cases by the aforementioned procedure, which resulted into 222 users for the Cologne case and 207 users for the Chemnitz case, 413 unique users in total. Only 16 top users "--- 7\% from the Cologne top list and 7,7\% from the Chemnitz top list "--- have participated in both discussions (3,9\% of the total number of unique users; for the accounts, see Table~\cref{tab:intersectingInfluencers}). This also means that the list of top users changed slightly more than the general sample, but the difference seems not to be incredibly significant. Thus, H1ba is rejected, too.

\begin{table} [htbp]%
	\centering
	\caption{Intersecting influencers.}%
	\label{tab:intersectingInfluencers}% label всегда желательно идти после caption
	\renewcommand{\arraystretch}{1.5}%% Увеличение расстояния между рядами, для улучшения восприятия.
	\begin{SingleSpace}
		\begin{tabulary}{\textwidth}{@{}>{\zz}L >{\zz}C >{\zz}C@{}} %Вертикальные полосы не используются принципиально, как и лишние горизонтальные (допускается по ГОСТ 2.105 пункт 4.4.5) % @{} позволяет прижиматься к краям
			\toprule     %%% верхняя линейка
			Institutional type & Accounts & Number of users \\
			\midrule %%% тонкий разделитель. Отделяет названия столбцов. Обязателен по ГОСТ 2.105 пункт 4.4.5
			National media & BILD, DLFNachrichten, faznet, ndaktuell, SPIEGELONLINE, SZ, tagesschau, tazgezwitscher, welt, ZDFheute, zeitonline & 11 \\
			Regional media & WDR and ZDFnrw & 2 \\
			Journalists & MatthiasMeisner & 1 \\
			Bloggers & Korallenherz & 1 \\
			Politicians & HeikoMaas & 1 \\
			\bottomrule %%% нижняя линейка
		\end{tabulary}%
	\end{SingleSpace}
\end{table}

\paragraph{3.2 Reproducibility of the Influencer Structure: Politicization and Political Polarization}

\paragraph{RQ2. Is the structure of politically relevant (institutionalized political, grassroots political, and media) users among the influencers similar in both cases?}

\begin{itemize}
	\item H2a. The media segment of the influencer structure is the most stable (as media are the key actors in the mediated public sphere).
	\item H2b. The political segment (both institutional and grassroots) is more salient in the Cologne case (due to the political resonance of the case) (Table~\cref{tab:influencerCharacter}).
\end{itemize}

\begin{table} [htbp]%
	\centering
	\caption{Institutional character of influencers.}%
	\label{tab:influencerCharacter}% label всегда желательно идти после caption
	\renewcommand{\arraystretch}{1.5}%% Увеличение расстояния между рядами, для улучшения восприятия.
	\begin{SingleSpace}
		\begin{tabulary}{\textwidth}{@{}>{\zz}L >{\zz}C >{\zz}C@{}} %Вертикальные полосы не используются принципиально, как и лишние горизонтальные (допускается по ГОСТ 2.105 пункт 4.4.5) % @{} позволяет прижиматься к краям
			\toprule     %%% верхняя линейка
			Institutional type & Cologne & Chemnitz \\
			\midrule %%% тонкий разделитель. Отделяет названия столбцов. Обязателен по ГОСТ 2.105 пункт 4.4.5
			Institutional political actors & 5,22\% & 14,98\% \\
			Media & 23,48\% & 24,15\% \\
			Activists and politically active bloggers & 11,26\% & 27,05\% \\
			Ordinary & 50,00\% & 18,36\% \\
			Other/irrelevant & 11,74\% & 14,49\% \\
			Total N & 230 & 207 \\
			\bottomrule %%% нижняя линейка
		\end{tabulary}%
	\end{SingleSpace}
\end{table}

H2a is fully supported: the share of media accounts in both cases remained the same. Moreover, as shown in Table~\cref{tab:intersectingInfluencers}, 13 out 16 users that were high-ranked in both discussions are editorial mass media that operate mostly on the national level in Germany.


As for H2b, our findings demonstrate that political actors were much more visible in the discussion around events in Chemnitz that in Cologne: the share of institutionalized politicians has almost tripled. One might explain this tendency by the growth of Twitter activity of German politicians in general. SPD, die Linke, and FDP, as well as AfD, were fighting for users’ attention and gaining authority in an online discussion.

The salience of the AfD representatives has doubled. Four members of the right-wing party are present among the high-ranked users driving Twitter-discussion about events in Cologne in 2016, among them Björn Höcke, one of the founders of AfD Thuringia, the speaker of the parliamentary group of the AfD and the spokesman of the Thuringia Regional Association. In the discussion about Chemnitz in 2018, 8 accounts associated with AfD are found on the list of top users. All of them are ranked high by pagerank; 4 of 8 accounts are also ranked high by indegree and 1 by retweets.

Also, the rise of the grassroots politicization is clear. In our previous research on the Cologne case, we have shown that ‘many ordinary users have a clear political position that can be understood from their tweets, but they don’t define themselves as activists’ \cite[p.~141]{SmoliarovaBodrunovaBlekanov}. The Chemnitz case shows a clear difference: the share of ordinary people who explicitly declare their political position or another official status in their account descriptions has increased from 11\% to 27\%. This is another sign of politicization and political polarization during the case, as well as the sign of structural instability. Thus, H2b has not been supported.

\paragraph{3.3 Reproducibility of Intersecting Users and Their Roles: Media as the Carcass of Networked Discussions}

\paragraph{RQ3. Does salience of the intersecting accounts differ in the two cases?}

\begin{itemize}
	\item H3a. The activity ratios of the intersecting users remain stable (in between 0,8 and 1,2).
	\item H3b. The user engagement ratios of the intersecting users remain stable (in between 0,8 and 1,2).
\end{itemize}

We have shown above that media influencer accounts were, en masse, the only group of influencers that stably repeated in the two discussions. Hence, we have decided to assess the 13 media accounts that were discovered in both top user lists in terms of their activity and user engagement ratios.

Almost all media outlets tweeted more about Cologne than about Chemnitz (see Table~\cref{tab:mediaInfluencerRatios}). For public service TV, the difference is the most significant: \textit{@ZDFnrw} posted in January 2016 almost 12 times more tweets than in August-September 2018. Nation-wide news program \textit{@tagesschau} is the second account that was much more active in January 2016: they posted 7 times more tweets about Cologne. Interestingly, two media outlets that paid more attention to the Chemnitz case than to the Cologne events "--- \textit{@welt} and \textit{@tazgezwitscher} "--- represent the two opposite sides of the political spectrum (right-wing and left-wing, respectively), which, again, is a sign of political polarization. Thus, H3a is rejected.

\begin{table} [htbp]%
	\centering
	\caption{Activity and user engagement ratios for media influencers in the two discussions.}%
	\label{tab:mediaInfluencerRatios}% label всегда желательно идти после caption
	\renewcommand{\arraystretch}{1.6}%% Увеличение расстояния между рядами, для улучшения восприятия.
	\def\tabularxcolumn#1{m{#1}}
	\begin{tabularx}{\textwidth}{@{}>{\raggedright}X >{\centering}m{2.5cm} >{\centering}m{2.5cm} >{\centering}m{2.5cm} >{\centering\arraybackslash}m{2.5cm}@{}}% Вертикальные полосы не используются принципиально, как и лишние горизонтальные (допускается по ГОСТ 2.105 пункт 4.4.5) % @{} позволяет прижиматься к краям
			\toprule     %%% верхняя линейка
			& Cologne & Chemnitz & Activity ratio (Cologne to Chemnitz) & User engagement ratio (Cologne to Chemnitz) \\
			\midrule %%% тонкий разделитель. Отделяет названия столбцов. Обязателен по ГОСТ 2.105 пункт 4.4.5
			ZDFnrw & 83 & 7 & 11,86 & 0,07 \\
			Tagesschau & 123 & 18 & 6,83 & 0,2 \\ 
			Faznet  & 66 & 19 & 3,47 & 0,77  \\
			DLFNachrichten & 52 & 15 & 3,47 & 0,21 \\
			BILD & 16 & 5 & 3,2 & 3,79 \\
			SZ & 21 & 7 & 3  & 0,49 \\
			WDR & 31 &  15 & 2,066 & 0,08 \\
			SPIEGELONLINE & 43 & 24 & 1,79 & 0,42 \\
			Ndaktuell & 54 & 31 & 1,74 & 0,54 \\
			Zeitonline & 40 & 23 & 1,74 & 2\\
			ZDFheute & 32 & 28 & 1,14 & 0,36 \\
			Tazgezwitscher & 16 & 29 & 0,55 & 1,07 \\
			welt & 10 & 21 & 0,48 & 1,55 \\
			\bottomrule %%% нижняя линейка
	\end{tabularx}%
\end{table}

Contrary to the fact that media outlets tweeted significantly more about Cologne, the users’ involvement seems to follow the opposite trend. The biggest difference between users’ engagement rate between two cases is revealed for \textit{@ZDFnrw}, despite this media account has tweeted 12 times less about Chemnitz than about Cologne. Users of \textit{@tagesschau} were 5 times more involved into the discussion of about Chemnitz than by tweets about Cologne. The same gap is observed for another public broadcaster, the national radio station \textit{@DLFNachrichten}. Thus, national/regional PSB has lost its positions in terms of user engagement, despite the growing efforts in Twitter reporting.

Unlike the PSB TV, media with clear political positioning like the conservative \textit{@BILD} and \textit{@welt} and the left-wing \textit{@Tazgezwitscher}, received higher user attention in the Chemnitz case "--- again, which tells of user polarization. There is no media account except for \textit{@Zeitonline} that would fit into our expected ratio divergence values. Thus, both H3a and H3b need to be rejected. This, in its turn, is a sign of unstable positioning of the only segment of influencers that continued to the second Twitter discussion.

\subsubsection{4. Conclusion}

Our findings suggest that the level of general reproduction of the structure of \textit{ad hoc} conflictual discussions is extremely low. Comparing the two general samples crawled for Cologne and Chemnitz we revealed that the lists of users who have posted at least once match by 5\% only. The list of top users changed even more significantly, contrary to expectations. Among the top users, the media segment has been the most stable, which supports the idea of media remaining the key actors in the mediated public sphere \cite{Calhoun}. But even this media cluster did not preserve its positioning, neither in terms of activity nor in terms of user engagement. While media accounts tweeted significantly more about Cologne, the users’ involvement seems to follow the opposite trend. Political actors were much more visible in the Twitter discussion on Chemnitz than on Cologne: the share of institutionalized politicians has almost tripled, and the number of grassroots activists has more than doubled. This shows that politicization of the discussion became more formal, while the issue itself became relevant for non-politicized people. We can conclude that media remain the carcass of the public spheres which discontinue around them.

\section{Событийное формирование контента}\label{sec:ch3/sec2}

\subsection{Особенности использования пиктограмм-эмодзи в интернет-сообществах с разной частотой использования обсценной лексики}\label{subsec:ch3/sec2/sub1}

\paragraph{Постановка проблемы.}

Рост «киберпопуляции», наблюдаемый в последнее время из-за сложившихся эпидемиологических условий, затрудняет или делает невозможным непосредственное общение, приводит к возрастанию роли интернет-коммуникации. Для более полного выражения своих переживаний пользователи прибегают к различным формам вербальной и невербальной экспрессивности: использованию пиктограмм-эмодзи, стикеров, голосовых сообщений. Эмоционально окрашенные сообщения (как правило, негативного характера) часто сопровождаются использованием обсценной лексики. Подобные выражения допускаются в сообществах, имеющих менее строгий регламент и вольную модерацию, в них нередко унижают участников, провоцируют конфликты и разные формы киберагрессии. Сопоставление разных форм проявления эмоциональности (вербальной и невербальной) в социальных сетях помогает выявить специфику общения в группах с разным отношением к обсценной лексике. Анализ пиктограмм-эмодзи и отнесение их к определенным типам переживаний открывает возможности для более детального и многомерного изучения проявляемых в сети эмоций и чувств.

Повышенное значение эмодзи в сетевой коммуникации подтверждается тем, что пиктограммы постепенно вытесняют ряд сленговых слов в языке молодежи. Это, в частности, подтверждается результатами исследования социальной сети Instagram, приведенными в аналитическом обзоре Д. А. Войнова \cite{Voinov}. Еще в 2015 году Оксфордский словарь признал эмодзи «Face with Tears of Joy» (\facewithtearsofjoy) словом года из-за исключительно частого его употребления в сетевой коммуникации \cite{WordOfTheYear2015}. Т.Л. Копусь сообщает о полученных в Великобритании данных, свидетельствующих, что в современном интернет-общении 80\% сообщений содержат эмодзи и что 40\% всех текстов состоят только из пиктограмм \cite{Kopus}. С другой стороны, многие авторы, анализируя алфавиты эмодзи, приходят к выводу, что слишком многие символы могут толковаться по-разному и быть в значительной степени связанными с культурными кодами, принятыми жестами и символами. Смысл эмодзи сильно зависит от контекста, «негласных договоренностей, окказиональных значений часто используемых символов» \cite{Krylov}. Кроме того, технические особенности использования эмодзи заключаются в том, что пользователю в первую очередь предлагаются символы, которые он часто применяет в коммуникации. Они со временем составляют «активный словарь», к которому прибегает автор сообщений, в этом наборе отражаются индивидуальные характеристики языковой личности \cite{Krylov}. Это осложняет коммуникацию, ведь таким образом понимание затруднено и может стать источником разночтений, но стоит учесть, что использование пиктограмм инициирует творческие процессы и позволяет создавать уникальную сетевую идентичность, играть значениями слов и символов.

Кроме использования как некоторого средства самовыражения, следует отметить также и иные важные функции эмодзи:
\begin{itemize}
	\item экономия времени и сил на написание сообщений, передача фактуальной информации в сжатой, сокращенной форме \cite{Krylov,Pervukhina,FrolovFrolova};
	\item экономия места, замещение слова предметным изображением (заместительно-номинативная функция) \cite{Kosmarskaya}; 
	\item подчеркивание коммуникативной значимости передаваемой информации \cite{Tverdokhleb}; 
	\item универсальность передачи сообщения, улучшение взаимопонимания между участниками коммуникации \cite{Vinogradova,Inozemtseva2010,Inozemtseva2009}; 
	\item выражение настроения, эмоций, чувств, “оживление” коммуникации \cite{FrolovFrolova}; 
	\item эстетическая (украшение текста визуализацией) \cite{Kosmarskaya}; 
	\item сглаживание противоречий в общении, усиление эффекта речи, особенно в ситуации поздравления и извинения \cite{Pervukhina,Sampietro}; 
	\item управление настроениями читателей, мобилизация, активизация, манипуляция общественным мнением, в том числе и наиболее активной социальной группы "--- молодежи \cite{Voinov}; 
	\item привлечение внимания, усиление привлекательности текстовых сообщений \cite{ShapovalovaGusarovaDobrenko}; 
	\item внешняя и внутренняя идентификация, использование символов для обозначения ценностных ориентаций \cite{Voinov}; 
	\item компенсация “эмоциональной недостаточности” интернет-дискурса \cite{RyabkoFlyug}; 
	\item манифестация эмоций, их кодирование и репрезентация \cite{Mozgovaya}.
\end{itemize}

Таким образом, анализ современного сетевого языка невозможен без учета эмодзи-коммуникации, которая существенным образом дополняет, расширяет возможности письменного текста, меняет его стилистику, повышая число невербальных компонентов. По мнению Ю. А. Иноземцевой современные “смайлы” уже являются самостоятельными речевыми действиями, выходящими за рамки дополнения к высказываниям, поясняющим эмоциональное состояние \cite{Inozemtseva2010,Inozemtseva2009}. Эмодзи служат своего рода суррогатом отсутствующего при дистанционном общении экспрессивного компонента эмоций, и чем дальше, тем больше вытесняют развернутые описания своих состояний у пользователей. Вероятно, 2020 год можно считать вехой, с которой начинается отсчет новой интернет-эры "--- времени, когда дистанционное общение начинает превалировать по количеству затраченного времени над реальным (очным). Именно поэтому важным является изучение цифрового языка и того, как он проявляет себя в различных социальных общностях.

Особое значение это имеет для оценки самочувствия молодежи, для которой цифровой язык становится важным инструментом “цифровой” социализации, как “опосредованный всеми доступными инфокоммуникационными технологиями процесс овладения и присвоения человеком социального опыта, приобретаемого в онлайн-контекстах, воспроизводства этого опыта в смешанной офлайн/онлайн реальности и формирующий его цифровую личность как часть реальной личности” \cite{Soldatova}.

В контексте задач нашего исследования, направленных на установление связи частоты встречаемости определенного вида пиктограмм-эмодзи в группах молодежи с разной степенью использования в сетевом общении обсценной лексики, отдельный интерес представляют уже имеющиеся в науке данные, отражающие созданные семантические классификации эмотиконов и эмодзи, а также факторы, оказывающие влияние на частоту их использования в сети. В первом случае, нас привлекла классификация, предложенная К. И. Белоусовым и И. А. Обуховой, материалом для которой выступили тексты 274 пользователей студенческого возраста с использованием эмотиконов и эмодзи, анализ которых позволил не только выявить наиболее типичные символы, но и установить их связь с конкретными эмоциональными состояниями. Авторами было выделено 15 категорий (рубрик), из которых однозначно положительной направленностью обладают всего четыре – одобрение, благодарность, любовь и радость. При этом, последние две имеют явное доминирование по частоте реплик в корпусе. Достаточно представленными оказались также печаль, сарказм (ирония), намек \cite{BelousovObukhova}. Последний, обозначаемый авторами термином «подмигивание», на наш взгляд, с трудом можно ставить в один ряд с другими эмоциональными состояниями в этой классификации (удивление, безразличие, недоумение, страх, вина, смущение, задумчивость, злость). В предложенном списке встречаются все классы «эмомаркеров»: положительной (позитивной) направленности, отрицательной направленности, а также нейтральные \cite{Mozgovaya}. К последним, например, можно отнести удивление, безразличие, задумчивость. Однако, это отнесение весьма условно, так как всегда важен контекст использования эмодзи, из-за которого ему может быть присвоено иное значение, вплоть до противоположного.

Перспективными направлениями исследований можно считать ориентированные на установление связи частоты использования тех или иных эмоциональных символов с социально-демографическими и личностными особенностями учащейся молодежи. Есть данные, подтверждающие, что эмотиконы и эмодзи, отражающие «радость», используются чаще других независимо от специфики выборки \cite{BelousovObukhova}. Помимо оценки самочувствия молодежи как, в целом, благополучного (доминирование позитивных эмомаркеров), данный результат также можно объяснить в русле так называемого «этикета эмодзи» (негласных правил употребления эмодзи в тексте). Помимо расположения символа в тексте, например, вставка после логически выделяемого сегмента, в конце предложения или иллюстрации последовательности событий, в сообщении чаще рекомендуется использовать символы, выражающие одобрение, симпатию, поддержку, а также использовать их в ситуациях межличностного, а не делового общения \cite{Kosmarskaya}. Последнее подтверждается и результатами исследований роли эмодзи в онлайн-сообществах. В сообществах с повышенной эмоциональной вовлеченностью (например, фанатские группы), эмодзи чаще используют как эмоциональный катализатор, не влияющий на снижение интереса аудитории к постам. Тогда как в образовательных сообществах есть тенденция к снижению комментариев к постам, содержащим эмодзи, то есть, с увеличением количества использования эмодзи снижается интерес к сообщениям сообщества. Эмодзи становятся объектом, отвлекающим внимание от смысла сообщения, а не привлекающим и удерживающим его \cite{ShapovalovaGusarovaDobrenko}. Полагаем, это может иметь и другое объяснение "--- закономерное снижение числа эмодзи при росте объема вербальных комментариев. Это особенно актуально при обсуждении деловых вопросов, к числу которых относится и образовательный контент. Всё это соответствует негласным правилам этикета употребления эмодзи в деловом онлайн-общении. Кроме того, есть данные, что эмоциональная нагрузка интеллектуального контента отличается спецификой фонетики "--- эмоции чаще отражаются в «умножении» гласных, их растягивании, повторе (дааааа..., вообщеееее... ), а также многократном использовании для передачи экспрессии восклицательного знака. Что касается графической фиксации состояния на страницах «интеллектуальных» групп сети ВКонтакте, то пользователи чаще прибегают к использованию эмотиконов (комбинации, сформированных из знаков препинания "--- скобок, двоеточия, знака «равно» и других графических символов), предпочитая их системным смайлам-картинкам \cite{Krylova}.

Есть результаты, доказывающие связь между полом и частотой использования конкретных эмотиконов и эмодзи. Например, мужчины не используют в текстах символы задумчивости, злости, вины и смущения, тогда как женщины обращаются ко всем эмоциональным категориям без исключения \cite{BelousovObukhova}. Это согласуется с утверждениями о большей эмоциональности женщин и принятии обществом проявлений ими любых эмоций в любых коммуникативных формах. Установлено, что женская речь в интернет-общении более эмоциональна, экспрессивна, оценочна, у мужчин же оценочная лексика выражена слабее на фоне доминирования отрицательных коннотаций над положительными. Смайликами мужчины пользуются реже, чем женщины, которые чаще используют экспрессивный синтаксис и пиктограммы \cite{NovikovaVerentsova}. Выявлена также гендерно обусловленная связь между использованием тех или иных эмоциональных символов и самооценкой: у женщин повышенная самооценка положительно коррелирует с обращением к ироничным эмодзи и отрицательно с символами флирта и любви, тогда как у мужчин при повышенной самооценке реже встречается использование эмомаркеров радости \cite{BelousovObukhova}.

Анализ текстов позволяет обнаружить в интернет-пространстве целые «эмотивные поля» "--- «совокупность вербальных и невербальных информационных единиц, выполняющих в определенной коммуникативной ситуации функцию манифестации человеческих эмоций» \cite[c.~125]{10}. Тип поля определяется по доминирующим эмомаркерам: положительный, отрицательный, бивалентный (одновременное присутствие положительных и отрицательных эмомаркеров) и амбивалентный (невозможность однозначного определения и интерпретации полярности из-за двойственности и неопределенности контекста). Принципиально важно, что сообщение может включать в себя несколько пересекающихся эмотивных полей (каждое имеет ядро, маркеры, периферию, силовые линии, границы) \cite{Mozgovaya}.

\paragraph{Организация исследования.} Объектом настоящего исследования выступили комментарии в нескольких популярных у молодежи сообществах социальной сети «ВКонтакте», предметом исследования "--- использование пиктограмм для отражения эмоций и чувств в интернет-общении. Целью исследования было сравнение предпочитаемых пиктограмм-эмодзи, которыми выражают свои эмоции и чувства пользователи, состоящие в сообществах с разной частотой использования обсценной лексики. Задачи исследования включали:

\begin{enumerate}
	\item Поиск популярных у молодежи сообществ социальной сети «ВКонтакте» разной тематики с открытыми комментариями;
	\item Анализ выбранных сообществ с точки зрения частоты использования в них обсценной лексики и составление рейтинга;
	\item Анализ выбранных сообществ с точки зрения частоты используемых в них пиктограмм-эмодзи;
	\item Сравнение рейтингов предпочитаемых пиктограмм-эмодзи в сообществах, полярных по частоте использования обсценной лексики.
\end{enumerate}

Методами исследования выступили опрос и контент-анализ. Для реализации первого метода была составлена анкета, в которой респондентам были заданы вопросы об особенностях их интернет-предпочтений, наиболее привлекательных для них интернет-сообществах в социальных сетях.

Исследование проводилось в два этапа. На первом из них выборка включала 854 человека (большая часть "--- обучающиеся), из них 505 девушек (59,1\%), 349 юношей (40,9\%). Средний возраст респондентов 14,52 года. Большая часть участников находилась в возрастном диапазоне от 11 до 17 дет. Основная часть респондентов "--- жители г. Челябинска. По данным опроса были выделены 128 наиболее популярных сообществ социальной сети «ВКонтакте» различной тематики (юмор, кино, познавательный контент, эстетика, бытовые советы, спорт, политика и многое другое).

Второй этап предполагал сбор и контент-анализ сообщений, представленных в исследуемых интернет-сообществах по двум направлениям: 1) изучение частоты использования обсценной лексики, 2) предпочитаемые пиктограммы-эмодзи, которые пользователи размещают в своих комментариях в сообществе. Сбор данных (data crawling) о частотности пиктограмм в текстах перечисленных выше сообществ был произведен с помощью поискового робота, основанного на использовании программного интерфейса API ВКонтакте. Далее группы были ранжированы по критерию частоты употребления в них обсценной лексики, выделены группы с высокой и низкой частотой использования ненормативных выражений. Также для каждого из сообществ был составлен рейтинг пиктограмм-эмодзи с учетом частоты их использования. Были выделены полярные с точки зрения частоты употребления нецензурной лексики сообщества, а в них "--- наиболее популярные пиктограммы-эмодзи. При формировании групп сообществ (с низкой и высокой частотой употребления обсценной лексики) осуществлялся подсчет средней частоты употребления той или иной пиктограммы в группе.

Гипотезой исследования выступило предположение, что группы с разной частотой использования обсценной лексики отличаются эмоциональной тональностью и используемыми пиктограммами-эмодзи для выражения эмоций и чувств. 

\paragraph{Результаты исследования.} Для определения частоты употребления обсценной лексики был применен список обсценных слов и выражений из 7500 единиц, составленный О. Дайховской \newline [https://github.com/odaykhovskaya/obscene\_words\_ru]. Данный перечень был несколько модифицирован и дополнен, исходя из частоты встречаемости лексем в текстах социальной сети ВКонтакте. Выбранные сообщества были проверены, а затем ранжированы по частоте встречаемости в них обсценной лексики. Отобранные сообщества представлены в таблице~\cref{tab:vkGroupFeatures}.

\begin{table} [htbp]%
	\centering
	\caption{Характеристики исследуемых интернет-сообществ.}%
	\label{tab:vkGroupFeatures}% label всегда желательно идти после caption
	\renewcommand{\arraystretch}{1.6}%% Увеличение расстояния между рядами, для улучшения восприятия.
	\def\tabularxcolumn#1{m{#1}}
	\begin{adjustbox}{width=1\textwidth}
	\small
	\begin{tabularx}{\textwidth}{@{}>{\raggedright}X >{\centering}m{5.5cm} >{\centering}m{3.0cm} >{\centering}m{3.0cm} >{\centering\arraybackslash}m{3.0cm}@{}}% Вертикальные полосы не используются принципиально, как и лишние горизонтальные (допускается по ГОСТ 2.105 пункт 4.4.5) % @{} позволяет прижиматься к краям
		\toprule     %%% верхняя линейка
		 № & Название группы & Количество всех комментариев & Количество комментариев с обсценной лексикой & Частота использования обсценной лексики, \% \\
		\midrule %%% тонкий разделитель. Отделяет названия столбцов. Обязателен по ГОСТ 2.105 пункт 4.4.5
		1 & «Рифмы и Панчи» & 1542579 & 282601 & 18,32 \\ 
		2 & «Fast Food Music» & 79870 & 12339 & 15,45 \\
		3 & «Некультурные мемы про субкультуры» & 13288 & 1814 & 13,65 \\ 
		4 & «CMH» & 6327 & 713 & 11,27 \\ 
		5 & «DK» & 9718 & 1094 & 11,26 \\
		6 & «ТВОЕ ВПШ» & 656679 & 73348 & 11,17 \\ 
		7 & «Фиксики» & 8863 & 20 & 0,23 \\
		8 & «Большая перемена» & 65919 & 54 & 0,08 \\
		9 & «Radio Province | Province Special» & 13868 & 11 & 0,08 \\
		10 & «The Vyshka» & 247 & 0 & 0,00 \\
		11 & «Летняя школа для будущих 11-классников | Умскул» & 8786 & 0 & 0,00 \\
		12 & «Citrus Fitness | Цитрус Фитнес» & 85 & 0 & 0,00 \\ 
		13 & «Православные Челябинска» & 507 & 0 & 0,00 \\
		14 & «Подростково-молодежный клуб "Z"» & 105 & 0 & 0,00 \\
		 \bottomrule %%% нижняя линейка
		 \multicolumn{5}{@{}p{\textwidth}}{%
			            \vspace*{-4ex}% этим подтягиваем повыше
			            \hspace*{2.5em}% абзацный отступ - требование ГОСТ 2.105
			            Примечание "---  Частота использования обсценной лексики вычислялась как процент комментариев с нецензурными выражениями в конкретном сообществе от общего числа комментариев в этом сообществе.
			            }
		            \\
	\end{tabularx}%
	\end{adjustbox}
\end{table}

Первые шесть сообществ в представленной таблице были отнесены к подгруппе с высокой частотой использования обсценной лексики, остальные восемь "--- к подгруппе с низкой частотой использования данного вида лексики. Далее в изучаемых сообществах были проанализированы наиболее популярные используемые их подписчиками пиктограммы. В целом для каждой группы сообществ (с низкой и высокой частотой употребления обсценной лексики) был составлен перечень из 40 самых часто употребляемых эмодзи, в таблице~\cref{tab:emojiFrequency} представлены (по убыванию) частоты встречаемости 20 наиболее предпочитаемых эмодзи в сообществах с различной частотой использования обсценной лексики.

\begin{table}[ht]%
	\centering
	\caption{Частота использования в сравниваемых сообществах эмодзи, характеризующих различные виды эмоций и чувств (\%).}%
	\label{tab:emojiFrequency}% label всегда желательно идти после caption
	\begin{adjustbox}{width=1\textwidth}
		\small
		\begin{tabular}{ c  c  c  c  c  c  c }% Вертикальные полосы не используются принципиально, как и лишние горизонтальные (допускается по ГОСТ 2.105 пункт 4.4.5) % @{} позволяет прижиматься к краям
			\toprule
			& \multicolumn{3}{c}{\makecell{Сообщества без обсценной лексики}} & \multicolumn{3}{c}{\makecell{Сообщества с обсценной лексикой}}\\
			\cline{2-7}
			№ & Эмодзи & Вид эмоций и чувств & Сред.Ч & Эмодзи & Вид эмоций и чувств & Сред.Ч \\
			\hline
			1 & \redheart & Сближающие & 14,04 & \facewithtearsofjoy & Гедонические & 5,44 \\
			2 & \foldedhands & Сближающие & 11,99 & \rollingonthefloorlaughing & Гедонические & 4,80 \\ 
			3 & \thumbsup & Сближающие & 6,33 & \smilingfacewithhearteyes & Сближающие & 3,63 \\ 
			4 & \smilingfacewithhearteyes & Сближающие & 2,62 & \redheart & Сближающие & 3,58 \\
			5 & \facewithtearsofjoy & Гедонические & 1,68 & \loudlycryingface & Меланхолические & 2,36 \\
			6 & \smilingfacewithsmilingeyes & Гедонические & 1,68 & \thumbsup & Сближающие & 2,28 \\
			7 & \smilingfacewithopenhands & Сближающие & 1,54 & \smilingfacewithsunglasses & Гедонические & 1,81 \\
			8 & \rollingonthefloorlaughing & Гедонические & 1,50 & \poutingface & Удаляющие & 1,79 \\ 
			9 & \winkingface & Сближающие & 1,33 & \pensiveface & Меланхолические & 1,63 \\
			10 & \flexingbiceps & Сближающие & 0,98 & \flushedface & Астенические & 1,49 \\
			11 & \beamingfacewithsmilingeyes & Гедонические & 0,85 & \moyai & Удаляющие & 1,47 \\
			12 & \clappinghands & Сближающие & 0,82 & \pleadingface & Меланхолические & 1,37 \\
			13 & \grinningfacewithsmilingeyes & Гедонические & 0,77 & \thinkingface & Гедонические & 1,18 \\
			14 & \loudlycryingface & Меланхолические & 0,75 & \smilingfacewithhearts & Сближающие & 1,15 \\
			15 & \grinningfacewithsweat & Гедонические & 0,70 & \smilingfacewithhorns & Удаляющие & 1,12 \\
			16 & \faceblowingakiss & Сближающие & 0,70 & \clownface & Гедонические & 1,08 \\
			17 & \partyingface & Гедонические & 0,70 & \beamingfacewithsmilingeyes & Гедонические & 0,95 \\
			18 & \smirkingface & Гедонические & 0,61 & \facesavoringfood & Сближающие & 0,88 \\
			19 & \slightlysmilingface & Гедонические & 0,60 & \catwithtearsofjoy & Гедонические & 0,85 \\
			20 & \thinkingface & Гедонические & 0,54 & \personfacepalming & Удаляющие & 0,79 \\
			\bottomrule
			\multicolumn{7}{@{}p{\textwidth}}{%
%				\vspace*{-4ex}% этим подтягиваем повыше
				\hspace*{2.5em}% абзацный отступ - требование ГОСТ 2.105
				Обозначения: «Сред.Ч» "--- средняя частоты встречаемости эмодзи в группе сообществ.
			}\\
		\end{tabular}%
	\end{adjustbox}
\end{table}

Анализируя содержание используемых в сообществах с низкой частотой употребления обсценной лексики пиктограмм, следует отметить, что в них чаще отображены позитивные эмоции (восторг, радость, одобрение, надежда, удовольствие), при этом положительные эмоции имеют, как правило, более высокий рейтинг, чем отрицательные "--- 13 первых в данном перечне состояний могут считаться позитивными. В группе сообществ с высокой частотой использования обсценной лексики эмодзи негативных эмоций (в рейтинге частотности) появляются не на 14-й позиции, как в первой группе сообществ, а уже на 5-й позиции. Встречаемость пиктограмм, отражающих пять видов (кластеров) эмоций, представлена на диаграмме частот эмодзи в рейтингах групп с разной частотой употребления обсценной лексики (см. рис.~\cref{fig:emojiFrequency}).

\begin{figure}[ht]
	\centerfloat{
		\includegraphics[scale=0.7]{emojiFrequency}
	}
	\caption{Сравнительная гистограмма университетов России и Великобритании.}\label{fig:emojiFrequency}
\end{figure}

Подавляющее большинство эмодзи в обеих исследуемых группах сообществ связаны с отображением позитивных эмоций, однако в группах с редким использованием обсценной лексики чаще встречается описание гедонических эмоций "--- радости, восторга, удовольствия. В негативном регистре несколько чаще встречаются пиктограммы, отображающие удаляющие (напр., гнев, отвращение, стыд) и меланхолические (напр., печаль, тоска) эмоции и чувства, причем оба этих кластера чувств чаще встречаются в группе сообществ с частым употреблением брани.

Для более детальной оценки различий в исследуемых группах сообществ по предпочтению пиктограмм было произведено сравнение средней частоты и ранга каждой из совпадающих пиктограмм в группах сообществ с разной частотой употребления обсценизмов. Были отобраны пиктограммы, которые употреблялись как в сообществах с частым употреблением нецензурной лексики, так и в сообществах, где подобная лексика практически не представлена. Из изначальных рейтингов, включающих 40 позиций, 30 пиктограмм совпадало в обеих группах сообществ. Сравнение рангов частотности этих 30 наиболее популярных пиктограмм представлено в таблице~\cref{tab:emojiDifference}. В ней пиктограммы перечислены по убыванию разности рангов, которые эмодзи занимают в группах с различной встречаемостью нецензурных выражений. В начале списка (разность рангов характеризуется положительным числом) стоят эмодзи, которые чаще встречаются в сообществах с редким употреблением бранных слов, в конце списка (разность рангов отрицательная) стоят пиктограммы, которые имеют более высокий ранг в группах с частым употреблением обсценной лексики. Другими словами, в верхних строках таблицы изображены пиктограммы с преобладанием частоты встречаемости в сообществах без обсценной лексики в сравнении с частотой встречаемости в сообществах с обсценной лексикой, а в нижних строках наоборот "--- с преобладанием по частоте встречаемости в сообществах с обсценной лексикой в сравнении с сообществами без обсценной лексики.

\begin{table}[ht]%
	\centering
	\caption{Различия в предпочитаемых эмодзи в сообществах с разной частотой использования обсценной лексики.}%
	\label{tab:emojiDifference}% label всегда желательно идти после caption
		\begin{adjustbox}{width=1\textwidth}
		\small
		\begin{tabular}{ c  c  c  c  c  c }% Вертикальные полосы не используются принципиально, как и лишние горизонтальные (допускается по ГОСТ 2.105 пункт 4.4.5) % @{} позволяет прижиматься к краям
			\toprule
			Эмодзи & \multicolumn{2}{c}{\makecell{Ранги эмодзи в\\группах}} & \makecell{Разность\\рангов} & \makecell{Словесная\\трактовка в Юникоде} & \makecell{Номер в\\Юникоде} \\
			\cline{2-3}
			& \makecell{без\\обсценной\\лексики} & \makecell{с\\обсценной\\лексикой} & & &  \\
			\hline
			\flexingbiceps & 9 & 28 & 19 & Бицепс & U+1F4AA \\
			\winkingface & 8 & 27 & 19 & \makecell{Подмигивающее лицо} & U+1F609 \\
			\foldedhands & 2 & 20 & 18 & \makecell{Человек, сложивший руки} & U+1F64F \\
			\grinningfacewithsweat & 13 & 26 & 13 & \makecell{Улыбающееся лицо в холодном\\поту с открытым ртом} & U+1F605 \\
			\smilingfacewithsmilingeyes & 6 & 18 & 12 & \makecell{Улыбающееся лицо со смеющимися глазами} & U+1F60A  \\
			\facescreaminginfear & 18 & 30 & 12 & \makecell{Лицо, кричащее от страха} & U+1F631  \\
			\smirkingface & 14 & 25 & 11 & \makecell{Усмехающееся лицо} & U+1F60F  \\
			\grinningfacewithsmilingeyes & 11 & 21 & 10 & \makecell{Улыбающееся лицо с открытым\\ртом и смеющимися глазами} & U+1F604  \\
			\beamingfacewithsmilingeyes & 10 & 15 & 5 & \makecell{Ухмыляющееся лицо со смеющимися глазами} & U+1F601  \\
			\grinningfacewithbigeyes & 16 & 22 & 6 & \makecell{Улыбающееся лицо с открытым ртом} & U+1F603 \\
			\redheart & 1 & 4 & 3 & \makecell{Закрашенное жирное сердце} & U+2764  \\
			\thumbsup & 3 & 6 & 3 & \makecell{Жест "хорошо" (большой палец поднят)} & U+1F44D  \\
			\smilingcatwithhearteyes & 21 & 24 & 3 & \makecell{Улыбающийся кот с глазами-сердечками} & U+1F63B  \\
			\personshrugging & 27 & 29 & 2 & \makecell{Пожимание плечами} & U+1F937 \\
			\cryingface & 20 & 19 & -1 & Плачущее лицо & U+1F622  \\
			\personfacepalming & 17 & 16 & -1 & \makecell{«Фейспалм» (Человек поднес\\руку к лицу, закрыв его)} & U+1F926  \\
			\thinkingface & 15 & 13 & -2 & \makecell{Задумчивое лицо} & U+1F914  \\
			\facewithtearsofjoy & 5 & 1 & -4 & \makecell{Лицо со слезами радости} & U+1F602  \\
			\rollingonthefloorlaughing & 7 & 2 & -5 & \makecell{Катается по полу от смеха} & U+1F923 \\
			\loudlycryingface & 12 & 5 & -7 & \makecell{Лицо громко плачет} & U+1F62D  \\
			\grinningsquintingface & 24 & 17 & -7 & \makecell{Улыбающееся лицо с открытым\\ртом и плотно закрытыми глазами} & U+1F606  \\
			\droolingface & 30 & 23 & -7 & \makecell{«Слюнки текут» (Лицо с закрытыми\\глазами и подтекающей слюной)} & U+1F924 \\  
			\smilingfacewithhearts & 22 & 14 & -8 & \makecell{Улыбающееся лицо с улыбающимися\\глазами и тремя сердечками} & U+1F970  \\
			\smilingfacewithhearteyes & 22 & 13 & -9 & \makecell{Улыбающееся лицо с глазами-сердечками} & U+1F60D  \\
			\smilingfacewithsunglasses & 19 & 7 & -12 & \makecell{Улыбающееся лицо в солнечных очках} & U+1F60E  \\
			\pleadingface & 25 & 12 & -13 & \makecell{Лицо с умоляющими глазами} & U+1F97A  \\
			\pensiveface & 23 & 9 & -14 & \makecell{Задумчивое лицо} & U+1F614  \\
			\flushedface & 26 & 10 & -16 &\makecell{Покрасневшее лицо} & U+1F633  \\
			\moyai & 28 & 11 & -17 & \makecell{«Мояи» (Каменное лицо)} & U+1F5FF  \\
			\poutingface & 29 & 8 & -21 & \makecell{Лицо, надувшее губы} & U+1F621 \\
			\bottomrule
		\end{tabular}%
	\end{adjustbox}
\end{table}

Первые десять пиктограмм (верхние позиции таблицы) характерны для сообществ с низкой частотой употребления нецензурных выражений, они обозначают:
\begin{itemize}
	\item одобрение "--- на пиктограммах, которые в Юникоде обозначены, как «Бицепс», «Подмигивающее лицо»;
	\item радость "--- «Улыбающееся лицо со смеющимися глазами»;
	\item разные градации веселья "--- «Улыбающееся лицо с открытым ртом и смеющимися глазами», «Ухмыляющееся лицо со смеющимися глазами», «Улыбающееся лицо с открытым ртом»;
	\item просьба или благодарность "--- «Человек, сложивший руки»;
	\item смущение и удовольствие: «Улыбающееся лицо в холодном поту с открытым ртом».
\end{itemize}

Заметно, что большая часть этих пиктограмм отражает гедонические и сближающие коммуникативные эмоции и чувства. Трудно с определенностью трактовать эмодзи «Усмехающееся лицо», поскольку ухмылка или насмешка могут быть как проявлением высокомерного отношения к собеседнику, так и неприятия чего-либо нежелательного "--- некоторой защитной реакцией на нерадостные новости, события или слова коммуниканта. В этом ряду одна пиктограмма может быть отнесена к символам негативных переживаний "--- «Лицо, кричащее от страха». Однако необходимо учитывать немалую условность пиктографических изображений. Поскольку мало кто из пользователей обращается к словесным интерпретациям смайликов, то вполне можно предположить, что данное изображение многие связывают с реакцией не страха, а сильного удивления. 

В средней части таблицы находятся пиктограммы, не обладающие выраженной специфичностью в плане стилистики коммуникации в сообществах.

В нижних строках таблицы представлены эмодзи, характерные для интернет-сообществ с частым использованием обсценной лексики. Обращает на себя внимание наличие существенно большего числа пиктограмм, выражающих негативные эмоции:
\begin{itemize}
	\item гнев "--- «Лицо, надувшее губы»; 
	\item страх "--- «Покрасневшее лицо»; 
	\item печаль "--- «Задумчивое лицо», «Лицо с умоляющими глазами»; 
	\item непонимание ситуации, затруднения в её оценке "--- «Мояи» (Каменное лицо).
\end{itemize}

В данной подвыборке встречаются также пиктограммы, отображающие позитивные эмоции:
\begin{itemize}
	\item уверенность в себе: «Улыбающееся лицо в солнечных очках»; 
	\item симпатия и любовь: «Улыбающееся лицо с глазами-сердечками», «Улыбающееся лицо с улыбающимися глазами и тремя сердечками»;
	\item удовольствие (возможно, связанное с удовлетворением потребности в пище: «Слюнки текут»); 
	\item веселье: «Улыбающееся лицо с открытым ртом и плотно закрытыми глазами».
\end{itemize}

Можно утверждать, что в группах с частым использованием обсценной лексики тональность смещена в сторону негативных эмоций и чувств. Другое отличие этих сообществ заключается в большем разнообразии пиктограмм, имеющих повышенную частотность использования "--- в сообществах без обсценной лексики заметно больше сходства в предпочитаемых пиктограммах. Данное различие свидетельствует, что в группах с редким использованием обсценной речи люди в большей степени настроены на выражение одобрения и симпатии, поддержки по отношению друг к другу. Рейтинг отрицательных эмоций в группах с часто используемой обсценной лексикой несколько выше, что может говорить о большей свободе в выражении негативных эмоциональных реакций, о повышенном настрое в данных группах на конфликт, меньших ограничениях, связанных с правилами этикета внутри подобных сообществ. 

Следующим шагом в анализе информации о частоте используемых пиктограмм выступило сопоставление эмодзи, попавших в рейтинг популярных, но встречающихся лишь в одном виде сообществ (см. табл.~\cref{tab:emojiEmotion}). 

\begin{table}[ht]%
	\centering
	\caption{Отличительные пиктограммы для сообществ с редким и частым употреблением нецензурной брани (средняя частота употребления, \% от общего числа эмодзи в сообществе).}%
	\label{tab:emojiEmotion}% label всегда желательно идти после caption
%	\begin{adjustbox}{width=1\textwidth}
%		\small
		\begin{tabular}{ c  c  c  }% Вертикальные полосы не используются принципиально, как и лишние горизонтальные (допускается по ГОСТ 2.105 пункт 4.4.5) % @{} позволяет прижиматься к краям
			\toprule
			Виды эмоций и чувств & \multicolumn{2}{c}{\makecell{Частота употребления эмодзи в группах}} \\
			\cline{2-3}
			& \makecell{с редким употреблением\\обсценной лексики} & \makecell{с частым употреблением\\обсценной лексики} \\
			\hline
			Гедонические & \makecell{\partyingface (0,7), \smilingface (0,6), \starstruck (0,43),\\\slightlysmilingface (0,42), \relievedface (0,31)} & \clownface (1,08), \catwithtearsofjoy (0,86) \\
			Сближающие & \smilingfacewithopenhands (1,54), \clappinghands (0,82), \faceblowingakiss (0,70) & \facesavoringfood (0,88) \\
			Удаляющие & \catwithwrysmile (0,07) & \smilingfacewithhorns (1,12), \facewithsteamfromnose (0,45), \facewithrollingeyes (0,56) \\
			Меланхолические & \sadbutrelievedface (0,06) & \nauseatedface (0,59), \wearyface (0,5) \\
			Астенические &  -- & \neutralface (0,58), \anxiousfacewithsweat (0,28) \\
			\bottomrule
		\end{tabular}%
%	\end{adjustbox}
\end{table}

Распределение отличительных (специфичных по использованию в сравниваемых сообществах) пиктограмм согласуется с обнаруженной тенденцией относительно отрицательных эмоций: их отображение чаще встречается в сообществах, в которых часто используется обсценная лексика.

\paragraph{Обсуждение результатов.} Полученные результаты показывают, что в сообществах, где обсценная лексика встречается часто, более высокий рейтинг имеют гедонические, удаляющие, меланхолические и астенические эмоции. В сообществах, где люди обходятся без использования нецензурной брани при общении преобладают эмодзи, описывающие сближающие эмоции.

Косвенно полученные результаты могут свидетельствовать, что в группах с редким использованием обсценной речи люди в большей степени настроены на позитивную и доверительную коммуникацию. В группах с нецензурной лексикой среди позитивных превалируют гедонические эмоции, вероятно, потому что многие из этих сообществ изобилуют юмористическим контентом, реакцией на который служит выражение экспрессивного компонента радости, а именно смеха, улыбки. Рейтинг отрицательных эмоций в группах с часто используемой обсценной лексикой несколько выше, что может говорить о более свободном выражении эмоций различного рода, возможностью снять внутреннее напряжение, в том числе за счет конфликтного взаимодействия с другими пользователями.

В группах без нецензурной брани больше количество сближающих эмоций и их содержание четко формирует позитивный посыл при общении с другим человеком, одобрение его действий, поступков (раскрытые для объятий руки, аплодисменты, поцелуй). Сближающая эмоция, обозначенная эмодзи \facesavoringfood, чаще встречающаяся в группах с частой обсценной бранью, хоть и обозначает чаще всего симпатию к собеседнику, но может иметь и двойственное значение. Показанный язык носит характер поддразнивания, раззадоривания, а в сочетании с комментариями негативного рода подобная пиктограмма и вовсе может трактоваться как издевка. В случае с гедоническими эмоциями в группах без нецензурной брани они выражают радость, но их экспрессия, как правило, менее ярко выражена. Гедонические эмоции в группах с частым употреблением обсценной лексики носят более яркий, выразительный характер "--- обозначение смеха.

\paragraph{Заключение.} Проведенное исследование позволяет сделать ряд выводов относительно особенностей использования эмодзи:
\begin{enumerate}
	\item Частота использования обсценной лексики выступает маркером эмоционального фона молодых участников сетевой коммуникации.
	\item Существуют сходства и различия при использовании эмодзи в интернет-сообществах с разной частотой использования обсценной лексики.
	\item Независимо от частоты обсценной лексики, пиктограммы, отражающие позитивные эмоции (восторг, радость, одобрение, надежду, удовольствие), встречаются чаще, чем символы отрицательных переживаний, при этом в иерархии их рейтинг выше в сообществе с низкой частотой использования обсценной лексики.
	\item Существуют различия в частоте использования эмодзи, характеризующих различные виды эмоций и чувств: в группах с частым употреблением брани чаще встречаются изображения удаляющих (гнев, отвращение, стыд) и меланхолических (печаль, тоска) переживаний, тогда как в группах с редким использованием обсценной лексики чаще встречаются символы гедонических эмоций (радость, восторг, удовольствие);
	\item Общая оценка эмоционального контекста коммуникации в сообществах разного типа показывает, что гедонические, удаляющие, меланхолические и астенические эмоции чаще проявляют члены сообществ с частым использованием обсценной лексики, тогда как с редким – ориентированы на проявление эмоций гедонических и сближающих.
\end{enumerate}

Подводя общий итог, следует отметить, что, несмотря на определенное сходство в перечнях популярных эмодзи, используемых для общения в группах с редким и частым употреблением нецензурной лексики, более детальный анализ рейтингов, специфических для каждой группы пиктограмм, позволяет выявить особенности в использовании экспрессивных средств сетевой невербальной коммуникации в сообществах. В группах с частым употреблением обсценной лексики границы в общении определены не так жестко, они реже табуируют негативное отношение собеседников друг к другу, позволяют (а может, и провоцируют) конфликтные ситуации, стычки между подписчиками, насмешливое и придирчиво-высокомерное отношение к другим людям. В сообществах с редким употреблением обсценизмов поощряют позитивные эмоции, выражение одобрения и поддержки в отношении других подписчиков, в них более выражен уважительный и доброжелательный тон.

%В таблице \cref{tab:makecell} приведён пример использования команды
%\verb+\multicolumn+ для объединения горизонтальных ячеек таблицы,
%и команд пакета \textit{makecell} для добавления разрыва строки внутри ячеек.
%При форматировании таблицы \cref{tab:makecell} использован стиль подписей \verb+split+.
%Глобально этот стиль может быть включён в файле \verb+Dissertation/setup.tex+ для диссертации и в
%файле \verb+Synopsis/setup.tex+ для автореферата.
%Однако такое оформление не~соответствует ГОСТ.
%
%\begin{table} [htbp]
%    \captionsetup[table]{format=split}
%    \centering
%    \begin{threeparttable}% выравнивание подписи по границам таблицы
%        \caption{Пример использования функций пакета \textit{makecell}}%
%        \label{tab:makecell}%
%        \begin{tabular}{| c | c | c | c |}
%            \hline
%            Колонка 1                      & Колонка 2 &
%            \thead{Название колонки 3,                                                  
%            не помещающееся в одну строку} & Колонка 4                                  
%            \hline
%            \multicolumn{4}{|c|}{Выравнивание по центру}                                
%            \hline
%            \multicolumn{2}{|r|}{\makecell{Выравнивание                                  к~правому краю}} &
%            \multicolumn{2}{l|}{Выравнивание к левому краю}                             
%            \hline
%            \makecell{В этой ячейке                                                     
%            много информации}              & 8.72      & 8.55                   & 8.44  
%            \cline{3-4}
%            А в этой мало                  & 8.22      & \multicolumn{2}{c|}{5}         
%            \hline
%        \end{tabular}%
%    \end{threeparttable}
%\end{table}
%
%Таблицы~\cref{tab:test3,tab:test4} "--- пример реализации расположения
%примечания в~соответствии с ГОСТ 2.105. Каждый вариант со своими достоинствами
%и~недостатками. Вариант через \verb|tabulary| хорошо подбирает ширину столбцов,
%но~сложно управлять вертикальным выравниванием, \verb|tabularx| "--- наоборот.
%\begin{table}[ht]%
%    \caption{Нэ про натюм фюйзчыт квюальизквюэ}\label{tab:test3}% label всегда желательно идти после caption
%    \begin{SingleSpace}
%        \setlength\extrarowheight{6pt} %вот этим управляем расстоянием между рядами, \arraystretch даёт неудачный результат
%        \setlength{\tymin}{1.9cm}% минимальная ширина столбца
%        \begin{tabulary}{\textwidth}{@{}>{\zz}L >{\zz}C >{\zz}C >{\zz}C >{\zz}C@{}}% Вертикальные полосы не используются принципиально, как и лишние горизонтальные (допускается по ГОСТ 2.105 пункт 4.4.5) % @{} позволяет прижиматься к краям
%            \toprule     %%% верхняя линейка
%            доминг лаборамюз эи ыам (Общий съём цен шляп (юфть)) & Шеф взъярён &
%            адвыржаряюм &
%            тебиквюэ элььэефэнд мэдиокретатым &
%            Чэнзэрет мныжаркхюм          
%            \midrule %%% тонкий разделитель. Отделяет названия столбцов. Обязателен по ГОСТ 2.105 пункт 4.4.5
%            Эй, жлоб! Где туз? Прячь юных съёмщиц в~шкаф Плюш изъят. Бьём чуждый цен хвощ! &
%            \({\approx}\) &
%            \({\approx}\) &
%            \({\approx}\) &
%            \( + \)  
%            Эх, чужак! Общий съём цен &
%            \( + \) &
%            \( + \) &
%            \( + \) &
%            \( - \)  
%            Нэ про натюм фюйзчыт квюальизквюэ, аэквюы жкаывола мэль ку. Ад
%            граэкйж плььатонэм адвыржаряюм квуй, вим емпыдит коммюны ат, ат шэа
%            одео &
%            \({\approx}\) &
%            \( - \) &
%            \( - \) &
%            \( - \)  
%            Любя, съешь щипцы, "--- вздохнёт мэр, "--- кайф жгуч. &
%            \( - \) &
%            \( + \) &
%            \( + \) &
%            \({\approx}\)  
%            Нэ про натюм фюйзчыт квюальизквюэ, аэквюы жкаывола мэль ку. Ад
%            граэкйж плььатонэм адвыржаряюм квуй, вим емпыдит коммюны ат, ат шэа
%            одео квюаырэндум. Вёртюты ажжынтиор эффикеэнди эож нэ. &
%            \( + \) &
%            \( - \) &
%            \({\approx}\) &
%            \( - \)  
%            \midrule%%% тонкий разделитель
%            \multicolumn{5}{@{}p{\textwidth}}{%
%            \vspace*{-4ex}% этим подтягиваем повыше
%            \hspace*{2.5em}% абзацный отступ - требование ГОСТ 2.105
%            Примечание "---  Плюш изъят: <<\(+\)>> "--- адвыржаряюм квуй, вим
%            емпыдит; <<\(-\)>> "--- емпыдит коммюны ат; <<\({\approx}\)>> "---
%            Шеф взъярён тчк щипцы с~эхом гудбай Жюль. Эй, жлоб! Где туз?
%            Прячь юных съёмщиц в~шкаф. Экс-граф?
%            }
%             
%            \bottomrule %%% нижняя линейка
%        \end{tabulary}%
%    \end{SingleSpace}
%\end{table}
%
%Если таблица~\cref{tab:test3} не помещается на той же странице, всё
%её~содержимое переносится на~следующую, ближайшую, а~этот текст идёт перед ней.
%\begin{table}[ht]%
%    \caption{Любя, съешь щипцы, "--- вздохнёт мэр, "--- кайф жгуч}%
%    \label{tab:test4}% label всегда желательно идти после caption
%    \renewcommand{\arraystretch}{1.6}%% Увеличение расстояния между рядами, для улучшения восприятия.
%    \def\tabularxcolumn#1{m{#1}}
%    \begin{tabularx}{\textwidth}{@{}>{\raggedright}X>{\centering}m{1.9cm} >{\centering}m{1.9cm} >{\centering}m{1.9cm} >{\centering\arraybackslash}m{1.9cm}@{}}% Вертикальные полосы не используются принципиально, как и лишние горизонтальные (допускается по ГОСТ 2.105 пункт 4.4.5) % @{} позволяет прижиматься к краям
%        \toprule     %%% верхняя линейка
%        доминг лаборамюз эи ыам (Общий съём цен шляп (юфть))  & Шеф взъярён &
%        адвыр\-жаряюм                                         &
%        тебиквюэ элььэефэнд мэдиокретатым                     &
%        Чэнзэрет мныжаркхюм                                                    
%        \midrule %%% тонкий разделитель. Отделяет названия столбцов. Обязателен по ГОСТ 2.105 пункт 4.4.5
%        Эй, жлоб! Где туз? Прячь юных съёмщиц в~шкаф Плюш изъят.
%        Бьём чуждый цен хвощ!                                 &
%        \({\approx}\)                                         &
%        \({\approx}\)                                         &
%        \({\approx}\)                                         &
%        \( + \)                                                                
%        Эх, чужак! Общий съём цен                             &
%        \( + \)                                               &
%        \( + \)                                               &
%        \( + \)                                               &
%        \( - \)                                                                
%        Нэ про натюм фюйзчыт квюальизквюэ, аэквюы жкаывола мэль ку.
%        Ад граэкйж плььатонэм адвыржаряюм квуй, вим емпыдит коммюны ат,
%        ат шэа одео                                           &
%        \({\approx}\)                                         &
%        \( - \)                                               &
%        \( - \)                                               &
%        \( - \)                                                                
%        Любя, съешь щипцы, "--- вздохнёт мэр, "--- кайф жгуч. &
%        \( - \)                                               &
%        \( + \)                                               &
%        \( + \)                                               &
%        \({\approx}\)                                                          
%        Нэ про натюм фюйзчыт квюальизквюэ, аэквюы жкаывола мэль ку. Ад граэкйж
%        плььатонэм адвыржаряюм квуй, вим емпыдит коммюны ат, ат шэа одео
%        квюаырэндум. Вёртюты ажжынтиор эффикеэнди эож нэ.     &
%        \( + \)                                               &
%        \( - \)                                               &
%        \({\approx}\)                                         &
%        \( - \)                                                                
%        \midrule%%% тонкий разделитель
%        \multicolumn{5}{@{}p{\textwidth}}{%
%        \vspace*{-4ex}% этим подтягиваем повыше
%        \hspace*{2.5em}% абзацный отступ - требование ГОСТ 2.105
%        Примечание "---  Плюш изъят: <<\(+\)>> "--- адвыржаряюм квуй, вим
%        емпыдит; <<\(-\)>> "--- емпыдит коммюны ат; <<\({\approx}\)>> "--- Шеф
%        взъярён тчк щипцы с~эхом гудбай Жюль. Эй, жлоб! Где туз? Прячь юных
%        съёмщиц в~шкаф. Экс-граф?
%        }
%         
%        \bottomrule %%% нижняя линейка
%    \end{tabularx}%
%\end{table}

\section{Поисковый робот (кроллер)}\label{sec:ch3/formatted-numbers}

%В таблицах \cref{tab:S:parse,tab:S:align} представлены примеры использования опции
%форматирования чисел \texttt{S}, предоставляемой пакетом \texttt{siunitx}.
%
%\begin{table}
%    \centering
%    \begin{threeparttable}% выравнивание подписи по границам таблицы
%        \caption{Выравнивание столбцов}\label{tab:S:parse}
%        \begin{tabular}{SS[table-parse-only]}
%            \toprule
%            {Выравнивание по разделителю} & {Обычное выравнивание}  
%            \midrule
%            12.345                        & 12.345                  
%            6,78                          & 6,78                    
%            -88.8(9)                      & -88.8(9)                
%            4.5e3                         & 4.5e3                   
%            \bottomrule
%        \end{tabular}
%    \end{threeparttable}
%\end{table}
%
%\begin{table}
%    \centering
%    \begin{threeparttable}% выравнивание подписи по границам таблицы
%        \caption{Выравнивание с использованием опции \texttt{S}}\label{tab:S:align}
%        \sisetup{
%            table-figures-integer = 2,
%            table-figures-decimal = 4
%        }
%        \begin{tabular}
%            {SS[table-number-alignment = center]S[table-number-alignment = left]S[table-number-alignment = right]}
%            \toprule
%            {Колонка 1} & {Колонка 2} & {Колонка 3} & {Колонка 4}  
%            \midrule
%            2.3456      & 2.3456      & 2.3456      & 2.3456       
%            34.2345     & 34.2345     & 34.2345     & 34.2345      
%            56.7835     & 56.7835     & 56.7835     & 56.7835      
%            90.473      & 90.473      & 90.473      & 90.473       
%            \bottomrule
%        \end{tabular}
%    \end{threeparttable}
%\end{table}
%
%\section{Виртуальный сетевой агент}\label{sec:ch3/sect2}
%% Не все (xe|lua)latex совместимые шрифты умеют работать с русским тире "---
%
%Некоторый текст.
%
%%\section{Параграф с подпараграфами}\label{sec:ch3/sect3}
%
%%\subsection{Подпараграф \cyrdash{} один}\label{subsec:ch3/sect3/sub1}
%
%Некоторый текст.
%
%%\subsection{Подпараграф \cyrdash{} два}\label{subsec:ch3/sect3/sub2}
%
%Некоторый текст.

\FloatBarrier
           % Глава 3
\chapter{Конструирование виртуальных мультиагентных подсетей}\label{ch:ch4}

\section{Конструирование Web-ресурсов}\label{sec:ch4/sect1}

\section{Типы виртуальных сетевых агентов}\label{sec:ch4/sect2}

\section{Оптимальное конструирование сборщика (цепочка роботов)}\label{sec:ch4/sect3}

\section{Реализация программного комплекса}\label{sec:ch4/sect3}

\FloatBarrier

           % Глава 4
\chapter{Интеллектуальных анализ данных Web-ресурсов}\label{ch:ch5}

\section{Анализ Web-графов}\label{sec:ch5/sect1}

\subsection{Beyond left and right: Real-world political polarization in twitter discussions on inter-ethnic conflicts}\label{subsec:ch5/sec1/sub1}

\subsubsection{1. Introduction}

Today, social polarization is believed to be growing both along traditional and newer lines along which schisms form \cite{DucaSaving}, of which political ones are, arguably, the sharpest. Despite the ever-increasing body of knowledge on political attitudes and alignments online, we still lack understanding of how political divisions show up in issue-oriented discussions and whether there is a cross-country pattern.

Despite all the well-described representation distortions \cite{Daniels}, the content of social media is still used for predicting consumer and/or electoral choices \cite{ColleoniRozzaArvidsson}, and the studies of political polarization on social media, including Twitter, are growing in popularity \cite{Barbera}. However, there are several shortcomings in today’s studies of political polarization in user-generated content.

Thus, in most cases, audience polarization is studied by examining purely political issues or events, while social conflicts of race, gender or religious origins with both evident and idiosyncratic polarization and politicisation \cite{McCrightDunlap} are rarely studied. Due to context and language differences, multi-country studies are also rare, especially where both established democracies and countries beyond the Euro-Atlantics are included, as, for most observers, these remain politically incomparable. However, conditions other than political regimes may create grounds for cross-cultural juxtapositions \cite{BodrunovaLitvinenkoBlekanov,BodrunovaBlekanovMaksimov}.

Another conceptual limitation is that, even in the most advanced studies, the detection of users’ political affiliations or ideologies is done via proxies, most often via structural network factors, such as: friendship affiliations; patterns of following \cite{BarberaJostNagler,Rivero}; or content sharing \cite{ColleoniRozzaArvidsson}, which could be misleading. Addressing this gap, newer works show that group polarization in social media may be studied by looking at user texts, including complex referrals to specific phenomena that matter for group identity \cite{Evolvi}. We argue that the analysis of political divisions needs to unite both structural and content aspects \cite{Bodrunova}.

In order to bridge these gaps in previous studies, we look at Twitter discussions regarding inter-ethnic clashes; they have similar conflict triggers and structure of social groups involved into conflict \cite{BodrunovaLitvinenkoBlekanov2017}. Whilst avoiding making straightfor- ward comparisons, we explore users’ political polariza- tion and suggest a mixed method to detect it across three cases in different political regimes: the USA; Germany; and Russia. By the UN estimates of 2013–2017, these countries have recently been the three most attractive countries to migrants in the world \cite{UN2013,UN2017} and have all witnessed violent inter-ethnic clashes that became global trending topics on Twitter.

This article, thus, is organized as follows: In Section 1, we review the approaches of assessing user polarization on social media and the conflicts under our scrutiny. In Section 2, we formulate the research questions and describe our methodology. In Section 3, we provide the results; in Section 4, we interpret and discuss them.

\subsubsection{2. Political Polarization on Twitter: The Current State of Research}

\paragraph{2.1. Political Polarization Studies and the Current Research Gaps}

Throughout recent years, mixed evidence has persisted in social media studies on whether users go online to agree or to argue \cite{YardiBoyd}. Research into echo chambers \cite{ColleoniRozzaArvidsson,Sunstein2002} has shown that user homophily, both structurally and semantically, may prevent the formation of online ‘opinion crossroads’, as there is ‘evidence of persistent ideological sorting in online communication networks’ \cite[p.~2]{Barbera}. However, a range of works point to the opposite effects in Twitter communication, with weaker ties responsible for the diversification of the consumption of political information \cite{Barbera} as well as different platform features on Twitter leading simultaneously to echo chambers and inter-community communication \cite{ConoverRatkiewiczFrancisco}. Thus, evidence suggests more research is needed to assess the patterns of users’ political clusterization on social networks.

Until today, most Twitter polarization studies are bound to the one-country-one-case strategy -- with a few notable exceptions \cite{Barbera,BarberaJostNagler}. Another problem arises from today’s understanding of online political polarization \cite{BramsonGrimSinger} as a situation when ‘a social or political group is divided into two opposing sub-groups having conflicting and contrasting positions, goals and viewpoints, with few individuals remaining neutral or holding an intermediate position’ \cite[p.~215]{CalaisGuerraMeiraJrCardie} \cite{Isenberg,Sunstein2002}.

Empirical evidence suggests that, if a heterogeneous group containing users with two opposing views has a non-zero cross-view retweet rate, it will end up as two polarized communities \cite{ConoverRatkiewiczFrancisco}. Following this logic, the studies of political polarization result in pre- defined binary descriptions of polarized communities -- see \cite{MoralesBorondoLosada} for Venezuela; 
\cite{AgathangelouKatakisRori} for Greece; or \cite{WeberGarimellaBatayneh} for Egypt.

However, for studies beyond the two-party electoral process, it seems useful to remember that polarization is an individual case of clusterization along schismatic lines, disregarding the number of resulting clusters \cite{EstebanRay}. In social conflicts, conflicting groups are not necessarily structured along binary political party divisions. The classic work of \cite{TajfelTurnerAustin} shows how social identity (including ethnic identity) divides in- and out-groups, while a later normative model of dissent in social groups \cite{Packer} implies that, in inter-ethnic conflicts, the majority may divide into pro- minority and anti-minority clusters if the anti-minority at- titude is perceived as harmful to the collective \cite[p.~5]{PackerChasteen}. Also, the very political spectra may be highly multi-dimensional, as The Manifesto Project (https://manifesto-project.wzb.eu/) or Polity Project (http://www.systemicpeace.org/polity/polity4.htm) suggest. Thus, we consider polarization more as multi-polar fragmentation of divergent clusters, of which bipolar clusterization is just an option.

Non-bipolar clusterization seems to be especially probable for ‘issue’ or ‘ad hoc’ publics \cite{BrunsBurgess,Papacharissi} that emerge on social networks. This claim is supported by research on the topic and issue-based discussions \cite{ElgesemSteskalDiakopoulos}. In single case studies, user polar- ization has been studied in regard to abortion, same-sex marriage, gun control, and climate change \cite{Elgesem,CalaisGuerraMeiraJrCardie,YardiBoyd}, with varying degrees and directions of polarization detected. Moreover, there is a clear difference in polarization patterns between political and non-political issues \cite{BarberaJostNagler}. But the evidence of differences in polarization patterns is still scarce in academic literature.

The biggest challenge in today’s polarization studies is that instead of taking into account the actual content of user posts, detection of users’ political affiliations is conducted via proxies. Of those, the most interesting results come from assessing structural network factors such as friendship affiliations \cite{BarberaRivero}, retweeting patterns \cite{CalaisGuerraVelosoMeiraJr}, patterns of political following \cite{BarberaJostNagler,Rivero} or content-sharing patterns \cite{AdamicGlance,BakshyMessingAdamic,ColleoniRozzaArvidsson,Elgesem}. However, using proxies may be misleading \cite{AdamicGlance} and even express analysis of actual tweets shows the extreme diversity of political views, both in the form of direct expression and in opinionated content.

However, if not proxies, then what? Several studies suggest that group polarization in social media may be examined by analysing complex user referrals to phenomena that matter for their identity and group alignment \cite{Evolvi}, as it is how the attitudes are expressed in natural language. In the simplest possible terms, one would take user attitudes (positive, negative, and neutral) towards particular objects for such referrals.

Thus, we will try to construct group divisions from the actual Tweet content by coding user referrals towards political players and then defining which of these attitudes divide the users most, and for how many clusters.

\paragraph{2.2. Lexicon-Based Approaches to the Analysis of User Polarization}

The area of research closest to our idea of bringing content into polarization studies is a lexicon-based analysis of Twitter data. In recent years, the field has experienced explosive growth, predominantly based on the analysis of sentiments. Without delving fully into these methodologies, we will simply note that the possibility of use of vocabulary-based approaches for polarization assessment tasks \cite{HillmanTrier} is usually based on combining lexical and structural analysis. Several researchers went beyond so-called ‘naïve’ sentiment and have tried to link affect \cite{StieglitzDangXuan} or appraisal \cite{DangXuanStieglitzWladarsch} in user texts, types of lexical units \cite{SperiosuSudanUpadhyay}, and structural elements of Twitter discussions, like graphs of following or speed and volume of Tweet dissemination. Using these and other works, one could conclude that a sentiment-based approach to detecting left and right differences would imply developing a ‘negative’ (say, leftist) + neutral + ‘positive’ (say, rightist) lexicon and applying them to the discussion bulk. However, the problem that we have run into with this approach is the following:

\begin{enumerate}
	\item The users expressed not ‘left views’ or ‘right views’ but attitudes (with their lexical markers) towards politicians, institutions, social groups, or events (‘actors’);
	\item A given user would express attitudes towards not just one but many actors of different political stances;
	\item The same user could express recognized- as-rightist attitudes towards one actor and recognized-as-leftist attitudes towards another actor of comparable significance (e.g., immigrants and nationalists);
	\item The same user could express negative views on both leftist and rightist actors (say, Barak Obama and the KKK in the USA).
\end{enumerate}

In case 3, the user’s preferences, as measured by one-dimensional positive/negative sentiment analysis, would create a zero-sum, and assigning the bias would not be possible. In case 4, an at least two-dimensional measurement of the political spectrum is needed. Taking this into consideration, we have further developed our research questions and the exploratory research design based on user sentiment, but not on pre-defined target-independent lexicons. Instead, to better capture user attitudes, we will use expert coding of Tweets, standardising the coding process with the help of the idea of ‘complex user referrals’ by Evolvi (2017) \cite{Evolvi}.

\paragraph{2.3. The Research Cases}

As stated above, we have studied three intergroup conflicts of ethnic or racial origins in the three leading immigration recipient countries: the USA, Germany, and Russia. Direct comparisons of ad hoc discussions \cite{BrunsBurgess} are currently viewed with some doubt in academic literature. Without developing a strictly com- parative research design, we have argued elsewhere \cite{BodrunovaLitvinenkoBlekanov2017,BodrunovaBlekanovMaksimov} that the conflicts we picked for the analysis are similar enough as research cases. They share a range of attributes: a violent interpersonal trigger, outbursts of public discussion across media platforms (becoming trending topics on Twitter), social polarization along the inter-ethnic or inter-race chasms, street action, and involvement of federal authorities. In addition, they were chosen because they were the first in a line of similar conflictual cases and, at least partly, set the communicative patterns for later discussions.

The cases are described as follows:
\begin{enumerate}
	\item A violent uprising against immigrants from Central Asia in the district of Biryulyovo, Moscow, Russia, in September 2013. After immigrant Orkhan Zeinalov, allegedly, killed local youngster Egor Scherbakov, the Biryolyovo residents destroyed the local warehouse and a trade centre around which hundreds of illegal immigrants had been dwelling. Several non-violent ‘people’s gatherings’ followed.
	
	\item Ferguson riots, Missouri, USA, in August 2014. There, unarmed African American teenager, Mike Brown was shot to death by white police officer Darren Wilson. The killing, as well as the defensive behaviour of the local police department, spurred several waves of street protests and peaceful support actions, including crowds at Mike Brown’s funeral ceremony.
	
	\item Mass harassment and rape of females on New Year’s Eve of 2016 in Cologne, Germany. Over 1,000 women reported being harassed during the celebrations on the city’s main square, allegedly, by re-settlers from North Africa and the Middle East. After that, demonstrations in protest were organized by radical political actors (PEGIDA movement and the party ‘Alternative for Germany’).
\end{enumerate}

\subsubsection{3. Research Questions and Methodology}

\paragraph{3.1. The Outline of the Research Design}

\paragraph{3.1.1. Research Questions} From what was said above, we have formulated the fol- lowing research questions:

\paragraph{RQ1. How, if at all, do the users cluster within the discussions, based on their attitudes to the major conflict actors? Does binary clusterization best describe user grouping?}

\paragraph{RQ2. Can the clusters be described as left or right in relation to the respective national political spectra? If not, then how could these clusters be described?}

\paragraph{RQ3. Are there similarities in the cluster structure of the discussions?}

\paragraph{3.1.2. The Research Design}

The way the RQs were formulated demanded an exploratory research design. To answer the research questions, we had to see which user groups emerged among the influencers and what discourses they conveyed.

Our concept for detecting user polarization was that political grouping within a discussion was constructed via a multi-dimensional combination of attitudes towards political actors (as defined above). These major political actors needed to be deduced from the discussions themselves. Then, the attitudes towards these players would be decrypted by expert coders for the key users, or influencers \cite{BodrunovaBlekanovMaksimov}, usually the bearers of the spectrum of attitudes.

The data received after coding would undergo clusterization, with each user belonging to one non-fuzzy cluster. Tweets by the users in the detected clusters would provide the word frequency vocabularies, which, after expert assessment, would turn into clusterization vocabularies. The latter would then be applied to all the users in the discussions, to see which users get into clusters and which discourses form there.

This approach, even if simple enough and reliant on expert intrusion, allows us to take into account the nature of the users’ political discussion, as well as the lack of linearity of their political positioning. We consider this crucial for studies on conflict discussions, as it may allow the inclusion of conflict-invoked (e.g., pro- or anti- minority), actor-oriented (e.g., authorities), and traditional political divisions (e.g., left and right and centre and radical). At this exploratory stage, our method does not imply machine learning or supervised approaches to data classification; we use big datasets at this stage of data collection only.

\paragraph{3.2. The Research Procedures}

\paragraph{3.2.1. Data Collection and Pre-Processing}

As this work is part of a bigger research project, our methods of data collection were described in detail previously \cite{BodrunovaLitvinenkoBlekanov2017}. Here, we briefly describe the steps we followed.

We used trendinalia.com to detect the initial discussion keywords and snowball reading to amplify this collection, thus forming the vocabularies for crawling. Trendinalia.com is a web service that allows daily monitoring of both world, regional, and national Twitter trending hashtags and words with no hashtag on an hourly basis, with the possibility of backdating; it has worked best in terms of detecting the trending topics, compared to the over ten other websites we had tried since 2013.

Using an API-independent Twitter crawler \cite{BlekanovSergeevMartynenko}, we collected the content of the discussions. All publicly available Tweets and the data on user interactions (likes, retweets, comments) were collected by a two-step procedure. Step one in- cluded the users who posted under the hashtags. Step two detected a wider community of likers, retweeters and commenters. On the discussion graph, only the step one nodes have been visualized.

Due to reasons regarding feasibility and sample comparability, collection periods differed. Thus, for Russia and Germany, the download period was 30 days after the trigger event. For the USA, we had to select the two weeks following the shooting, with Mike Brown’s funeral as the central event. The user samples included:

\begin{itemize}
	\item For Biryulyovo -- Step1: over 3,700 users; Step2: over 12,000 users;
	\item For Ferguson -- Step1: over 70,000 users; Step2: over 210,000 users;
	\item For Cologne -- Step1: over 12,000 users; Step2: over 99,000 users.
\end{itemize}

\paragraph{3.2.2. Data Analysis}

To answer RQ1, we needed to cluster key users by their political views, define the cluster vocabularies, and apply these vocabularies to the rest of the users, in order to see the discussion clusters and interpret their discourses.

As we expected the influencers to be bearers of the polarizing discourse, for each case, we defined the group of influencers based on nine parameters: the number of Tweets, likes, retweets, comments, in-degree, out- degree, degree, betweenness, and pagerank centralities. After these experiments, using various thresholds, the top 50 users were chosen as the cutting line for each parameter. As many users were repeated in the top due to several specifications, the duplicates were deleted. After elimination of influencers with low numbers of Tweets and bot-like influencers (with a percentage of repeated Tweets over 50\%) 156 users for Germany, 105 users for Russia, and 105 the USA were left. But, for the USA sample (which was several times bigger in the number of Tweets), the number of users was reduced by half, to 52 users. Their respective tweet collections for reading and coding included 13,359, 3,012, and 9,540 tweets.

To define user attitudes towards political actors, we developed scales for coding and coded the users (not their Tweets). The coders were experts in inter-ethnic conflict and, additionally, academic native-speakers, and the level of inter-coder reliability as measured by Cohen’s kappa reached at least 0.68 for any two sub-samples. The coders used the scale from \(-2\) to 2 to assess the attitude of each user to the following groups that had been identified as attitude triggers by reading the tweets before coding: 1) the minority (immigrants or African Americans); 2) ‘radical right’ or ‘radical white’ -- nationalists in Russia, PEGIDA and AFD in Germany, and the Tea Party and the KKK (as a label for radical whites) in the USA; 3) the incumbent country leaders -- Vladimir Putin and Dmitry Medvedev in Russia, Angela Merkel in Germany, and Barak Obama and Hillary Clinton in the USA; 4) local authorities and police forces grouped together as the ‘oppressive and responsible’ actors. Attitudes towards liberal opposition in Russia and to Republicans and Democrats were also coded, as we found them salient in the Tweets.

Based on this coding, the influences were clustered to form groups with similar combinations of attitudes. After clustering, the Tweets of each group were merged, and fully divergent frequency vocabularies of their discourses were formed with the use of expert vocabulary ‘cleaning’. Then, we applied the vocabularies to all the users in the discussions, to see how the discourses distribute within the discussion structure. We also measured whether these discourses formed distinct nebulae; but even if they did not, we assessed which users belonged to these discourses and interpreted the semantics of their speech qualitatively. What mattered for our analysis was whether the influencers formed distinct groups; all the following steps were the consequence.

In detail, the research steps were the following.

\begin{enumerate}
	\item Based on our coding, the influences of each case were grouped with the help of a k-means clustering algorithm with sorted the distances. With the number of clusters and the number of variables being diverse, the best solutions were finally chosen based on Silhouette metric S, within- and between-cluster square sums, examinations of variable means in each cluster (see Figures~\cref{fig:userAttitudesRussia} to~\cref{fig:userAttitudesUSA} for Russia, Germany, and the USA, respectively), and expert reading of tweets in each cluster. All the three influencer groups clustered well; Germany clustered best, Russia followed, and the USA the least, but all the solutions were sufficient by Silhouette from 2 to 10 possible clusters. To identify the best solutions, other aforementioned metrics were used. Those solutions were:
	\begin{itemize}
		\item For Russia: 4 variables (attitude to liberal opposition excluded), 3 clusters of 49, 36, and 20 users, \(S > 0.4\); 
		
		\item For Germany: 4 variables, 3 clusters of 99, 48, and 9 users, \(S > 0.5\);
		
		\item For the USA: 5 variables (attitude to Democrats excluded), 4 clusters of 15, 15, 12, and 10 users, \(S > 0.2\).
	\end{itemize}
	\item For each cluster, word frequency vocabularies were formed by merging the Tweets, ranging the words by frequency, and expert reading. After reading, only the words unique or highly characteristic for each cluster (for example, two mentions in one cluster and 160 in another would result into eliminating the word in the first one and leaving it in the other). If the difference between the numbers of mentions in any two clusters of the case was smaller than ten times, the word was eliminated in all the clusters.
	
	\item We applied the thesaurus to the rest of the users in each case; we wanted to identify the users who use the words from the divergent thesauri. As a result, we have received three types of users in each case: 1) the users who belonged to clusters 1 to 3 or 4; 2) the ‘overlappers’ who used the language of more than one cluster; 3) the users who did not use the discourses (mostly due to a low number of their Tweets). To ensure a higher quality of marking users, rather than using individual words from the thesauri, two-word combinations were used.
	
	\item Based on this information, we constructed the graphs of discourse distribution, with users as apexes and user interactions (comments and/or retweets) as edges, and calculated the indices for user centralities. We assessed who were the most influential discourse bearers and what they spoke about. We used Gephi algorithms OpenOrd and Force Atlas 2 for graph construction (see \cite{MartinBrownKlavans}, on OpenOrd), as the former favours centripetal graphing and the latter better shows visual homophily (see Figures~\cref{fig:openOrdRussia} to~\cref{fig:openOrdUSA}). To see whether the groups bearing the discourses were tighter than inter-group connections, we calculated the mean number of in-group and inter-group edges.
\end{enumerate}

\begin{figure}[ht]
	\centerfloat{
		\includegraphics[scale=0.5]{userAttitudesRussia}
	}
	\caption{Mean values of user attitudes to the selected political actors in attitude-based clusters for Russia.}\label{fig:userAttitudesRussia}
\end{figure}

\begin{figure}[ht]
	\centerfloat{
		\includegraphics[scale=0.5]{userAttitudesGermany}
	}
	\caption{Mean values of user attitudes to the selected political actors in attitude-based clusters for Germany.}\label{fig:userAttitudesGermany}
\end{figure}

\begin{figure}[ht]
	\centerfloat{
		\includegraphics[scale=0.52]{userAttitudesUSA}
	}
	\caption{Mean values of user attitudes to the selected political actors in attitude-based clusters for the USA.}\label{fig:userAttitudesUSA}
\end{figure}

\begin{figure}[ht]
	\centerfloat{
		\hfill
		\subcaptionbox[List-of-Figures entry]{\label{fig:openOrdRussia-1}}{%
			\includegraphics[width=0.419\linewidth]{openOrdRussia1}}
		\subcaptionbox{\label{fig:openOrdRussia-2}}{%
			\includegraphics[width=0.4\linewidth]{openOrdRussia2}}
		\hfill
	}
	\legend{Blue: Cluster 1, ‘anti-establishment nationalists’; red: Cluster 2, ‘news disseminators’; green: Cluster 3, ‘angry citizens’; black: ‘overlappers’; grey: non-clustered users.}
	\caption{Communication within and between discursive groups of users in the discussions, with users as vertices and interactions (retweets and comments) as edges; reconstructed by OpenOrd and Force Atlas 2 algorithms for Russia.}\label{fig:openOrdRussia}
\end{figure}

\begin{figure}[ht]
	\centerfloat{
		\hfill
		\subcaptionbox[List-of-Figures entry]{\label{fig:openOrdGermany-1}}{%
			\includegraphics[width=0.419\linewidth]{openOrdGermany1}}
		\subcaptionbox{\label{fig:openOrdGermany-2}}{%
			\includegraphics[width=0.4\linewidth]{openOrdGermany2}}
		\hfill
	}
	\legend{Blue: Cluster 1, ‘nationalists’; red: Cluster 2, ‘news disseminators’; green: Cluster 3, ‘anti-nationalists’; black: ‘overlappers’; grey: non-clustered users.}
	\caption{Communication within and between discursive groups of users in the discussions, with users as vertices and interactions (retweets and comments) as edges; reconstructed by OpenOrd and Force Atlas 2 algorithms for Germany.}\label{fig:openOrdGermany}
\end{figure}

\begin{figure}[ht]
	\centerfloat{
		\hfill
		\subcaptionbox[List-of-Figures entry]{\label{fig:openOrdUSA-1}}{%
			\includegraphics[width=0.419\linewidth]{openOrdUSA1}}
		\subcaptionbox{\label{fig:openOrdUSA-2}}{%
			\includegraphics[width=0.4\linewidth]{openOrdUSA2}}
		\hfill
	}
	\legend{Blue: Cluster 1, ‘politicized observers’; red: Cluster 2, ‘media-oriented users’; green: Cluster 3, ‘human rights activists’; purple: Cluster 4, ‘whites’ blamers’; black: ‘overlappers’; grey: non-clustered users.}
	\caption{Communication within and between discursive groups of users in the discussions, with users as vertices and interactions (retweets and comments) as edges; reconstructed by OpenOrd and Force Atlas 2 algorithms for the USA.}\label{fig:openOrdUSA}
\end{figure}

To answer RQ2 about the left or right nature of the clusters, we partly recoded our coding data and corrected the graphs of means (Figures~\cref{fig:userAttitudesRussia} to~\cref{fig:userAttitudesUSA}) accordingly. Recoding was needed to re-interpret attitudes for and against a given actor as pro-left or pro-right. E.g., the influencers expressed attitudes towards political leaders (Obama, Merkel, and Putin), coded \(-2\) to 2. But, for the respective political spectra, Obama is leftist, while Merkel and Putin \cite{BluhmVarga} represent the rightist spectrum side. To ‘normalize’ the user attitudes, we recoded all the pro-left views as \(-1\) to \(-2\), and all pro-right views as 1 to 2 (see Table~\cref{tab:variableLeftRightNormalization}). By doing this, we could show on the graphs of means whether the clusters (and how many of them) were pro-left, pro-right, or mixed -- see Figures~\cref{fig:recordedDataRussia} to~\cref{fig:recordedDataUSA} for Russia, Germany, and the USA, respectively.

\begin{table}[ht]%
	\centering
	\caption{Recoding of variables for their left-right normalization.}%
	\label{tab:variableLeftRightNormalization}% label всегда желательно идти после caption
		\begin{adjustbox}{width=1\textwidth}
				\small
		\begin{tabular}{ c  c  c  c  c  c  c  c }% Вертикальные полосы не используются принципиально, как и лишние горизонтальные (допускается по ГОСТ 2.105 пункт 4.4.5) % @{} позволяет прижиматься к краям
			\toprule
			Country & Minority & President & Police-Authorities & Nationalists & Opposition & Democrats & Republicans\\
			\hline
			Russia & Recoded & Not & Not & Not & Recoded & -- & --\\
			Germany & Recoded & Not & Not & Not & -- & -- & -- \\
			USA & Recoded & Recoded & Not & Not & -- & Recoded & Not \\
			\bottomrule
		\end{tabular}%
			\end{adjustbox}
\end{table}

\begin{figure}[ht]
	\centerfloat{
		\includegraphics[scale=0.5]{recordedDataRussia}
	}
	\caption{Mean values for the recoded data on user attitudes towards the selected political actors for Russia.}\label{fig:recordedDataRussia}
\end{figure}

\begin{figure}[ht]
	\centerfloat{
		\includegraphics[scale=0.5]{recordedDataGermany}
	}
	\caption{Mean values for the recoded data on user attitudes towards the selected political actors for Germany.}\label{fig:recordedDataGermany}
\end{figure}

\begin{figure}[ht]
	\centerfloat{
		\includegraphics[scale=0.5]{recordedDataUSA}
	}
	\caption{Mean values for the recoded data on user attitudes towards the selected political actors for the USA.}\label{fig:recordedDataUSA}
\end{figure}

To answer RQ3, we qualitatively assessed the results for RQ1 and RQ2.

\subsubsection{4. Results}

Our results show that the discourses identified by coding influencers cover a substantial part of the discourse in all the cases: for Russia, the thesauri covered 31,5\%, in Germany, 63,4\% and, in the USA, 73,5\% of the users. This shows that influencers’ talk reflects the discourse of ‘ordinary users’ to different extents in each country, but everywhere we were able to detect the discourses that were important for the overall discussion.

As the figures suggest, in all the three cases, group structure was not binary; moreover, binary solutions for each country would hide important discourses that actually constituted the discussions. Neither did the group divisions correspond to the minority/pro-minority majority/anti-minority majority scheme. Instead, the clusters may be described as follows:

For Russia, the clusters include: ‘news disseminators’; ‘anti-establishment nationalists’; and ‘angry citizens’. The first group was mostly neutral but formed a substantial part of the political discussions by supplying (posting or retweeting) news at each stage of the conflict. The second cluster was clearly anti-immigrant and nationalistic but differed from European nationalism. Within the discussion, there was also an evident divide between the nationalist groups who supported the current establishment and those who actively opposed it. The former saw the incumbent leadership as the flesh of the 1990s’ elites who ‘had stolen the country’; such users, therefore, blamed the national policymakers for supporting the post-Soviet immigration. The second type of nationalism -- the pro-establishment one -- showed up in the third cluster of ‘angry citizens’. This cluster united anti-institutionalists who were raising voices against \textit{bespredel} (‘the absence of limits’ and rules of the game), but in differing ways. This diverse group included pro-Putin nationalists who were ready to fight with the Moscow riot police, liberal oppositional media and public figures who criticized the policymakers, and ‘tired citizens’ who negatively treated the immigrants, and the country leaders, and the local authorities, and the nationalists. Unlike in the ‘news disseminators’ cluster, the close-to-zero means for these variables here were the result of pro- and anti-establishment views compensating each other while the users united against police (see Figure~\cref{fig:userAttitudesRussia}).

For Germany, the clusters include: ‘news disseminators’; ‘nationalists’; and ‘anti-nationalists’. Discursively, the biggest group of ‘nationalists’ unites two similar sub-groups, one with slightly more aggressive tendencies towards small liberal-oriented parties and activist movements (like Antifa), and the other more critical of the national government. The anti-nationalist group is, however, also salient, making the German picture one-dimensional in terms of political divisions (pro- and anti-minority), even if the dimension is not political-party but issue-based. Also, the overlappers play a significant role here, as they visually stand in between the two opposing clusters, thus creating bridges for public dialogue (see Figure~\cref{fig:openOrdGermany}).

For the USA, the clusters include: ‘media-oriented users’; ‘human rights activists’; ‘politicized observers’; and ‘whites’ blamers’ (see Figure~\cref{fig:userAttitudesUSA}). Within the influencers, the clusters were similar in volume, but, on the big graph, the last two groups were relatively small- scale, while the first two dominated the graph. Just as in Russia, the media-oriented discourse was a part of the political discussion, but the three other groups were not neutral, especially ‘whites’ blamers’ and ‘human rights activists.’ The former actively blamed ‘the white dominance’ and called for action against oppression. Interestingly, the hashtag \#blacklivesmatter was less important for this group than for the media-oriented discourse. However, blaming hashtags and words like ‘murderer,’ ‘republikkklan,’ or ‘kkkop,’ and calls for action (like ‘\#arrestdarrenwilson,’ ‘\#boycottgofundme,’ or ‘\#donotshopmonday’), were prominent. The other group, very different from ‘whites’ haters,’ and linked the case to human rights issues like abortion (\#prolife), gender inequality (\#womeninequalityday), morality (\#moralmonday), and others. The group itself, as one can see even from the hashtags, was polar in itself in terms of left and right divisions on human rights. For this group, positioning on Mike Brown’s death was different, expressed mostly by ‘don’t shoot’ hashtags. ‘Politicized observers’ abstained from taking clear sides, but discussed the Ferguson events in terms of its influence upon the political process in America. Interestingly, the cluster that mostly reposted media, was the most pro-Wilson, as media, evidently, tried to remain balanced; they also reported police press conferences that were modestly defensive towards Darren Wilson.

Then, we looked at how the discourses we described spread inside the graph. Our task was not to calculate the level of homophily and prove user clustering for all the discussions; the goal was to see how the discourses actually spread and whether they spread in a similar way -- and they did not. For Russia and the USA, the discourses mixed, but if in Russia we saw inter-cluster talk, in the USA overlappers took almost all the space in the graph centre. And in Germany, the graph was clearly structurally divided. This was also proved by the mean in- and inter-cluster weighted number of edges: in Russia, the inter-cluster links took over (216 vs. 323.5, respectively), while in Germany (4392.75 vs. 2890.25) and the USA (21114.4 vs. 3755.2) in-group connections were stronger.

Thus, the attitude-based grouping was different in each of the three cases. Also, it was far from clear left-right identifications. In order to show it, we have recoded the variables as stated above, making pro- left views negative (\(-1\) to \(-2\)) and pro-right views positive (1 to 2). We considered anti-minority, anti-Obama/Clinton, pro-Putin/Medvedev, pro-Merkel, pro-police, pro-nationalist, anti-opposition (in Russia), anti-Democrat, and pro-Republican (in the USA) views pro-right, while the opposite was marked pro-left. See the full recoding scheme in Table~\cref{tab:variableLeftRightNormalization}.

The resulting graphs of means are quite telling (see Figures~\cref{fig:recordedDataRussia} to~\cref{fig:recordedDataUSA}). Both in Russia and Germany, the leaders representing rightist sides of the spectra have actually taken pro-migration stance, and this has made right-wing users who support nationalist movements and speak against immigrants, move left and be against the incumbent leaders, as well as against the local authorities and police for ‘not protecting’ the host communities. But the other clusters in the two countries quite strongly differ from each other. While in Germany issue-based leftism is clearly seen, the other Russian cluster of ‘angry citizens’ diverges into three discourses that combine clearly rightist, pro-establishment nationalism; liberal, anti-establishment oppositional speakers; and politicised citizens. These politicised citizens, paradoxically for external observers, do not support any of the existing political factions, due to their impotence in resolving local problems. Thus, at least two nationalist discourses were detected by us for Russia -- while in the USA there are two very different left-wing clusters, one clearly left, supportive of either Obama or Clinton and based on human rights’ discourse, and another that was sharply anti-white, even blaming Obama for not being protective enough, which, in our rough coding, made the cluster stick out to anti-Obama views on the rightist side of Figure~\cref{fig:recordedDataUSA} (in effect, being extreme left). The cluster of ‘politicized observers’, interestingly, is reminiscent of the ‘tired citizens’ in Russia, as they are, on average, only slightly pro-African-American and, more strongly, anti- leader, anti-police, and anti-majority.

Another crucial observation is that, while the divisions in the discussion clearly stem from local political contexts, they are quite far from expectations determined by the systemic political features of the countries. Thus, in the majoritarian USA where one would expect two-sided polarization, the clusters were, in fact, numerous and the discussion was based on overlappers. It was rather coalitional Germany that showed polarization. And in Russia, just one side of the spectrum was present in the discussion. Thus, it is not only the local political markets but also the nature of the issue and issue-based divisions that shape political clustering
.
Overall conclusions are thus the following: The discursive schisms do exist in issue-based discussions, but they do not fall into binary categories according to majoritarian political divisions, and; they only partially fall into the three-side divisions expected by the nature of the issue. Instead, local political spectra may provoke the formation of, for example, two leftist or two rightist clusters. Only Germany has demonstrated the expected divisions between anti- and pro-minority majority, while the minority remained highly under-represented at all, like in Russia -- and unlike in America.

The similarities can also be traced, but not in terms of left and right divisions. First, in all the discussions, a politically neutral news-based cluster played a significant structural role. Second, all three discussions revealed harsh anti-institutionalism, including that from the users who, in conventional logic, were expected to support the incumbents. Third, Germany and Russia were similar in how nationalist clusters were against the conservative governments, and Russia and the USA were similar in how the ‘tired citizens’ were politicised against all the political sides.

\subsubsection{5. Conclusion}

In our article, we have combined content analysis of social media with cluster analysis and graph construction. Our method has revealed greater complexity of politicised discourse within ad hoc Twitter discussions on inter-ethnic conflicts. Thus, we have found that there may be several clusters of leftist or rightist views even if the number of clusters is minimal, and users may combine formally leftist and rightist views if positions of political actors or the nature of the issue demand it. The groups we have detected differ highly in their conceptualisation from the traditional left and right divisions and left or right labels cannot be attached to individual users based on their preferences, like pro- and anti-minority stances or treatments of country leaders or parties. We have also shown that, on the graphs, the discourses intertwine quite intensely if we do not force the graphs to artificially diverge according to users’ political views.

Our research provides new input for rethinking the political divisions that form online, on what grounds they form, and how to detect them. The local political contexts, as well as the nature of the issues under scrutiny, are major factors to be taken into account. In our article, the ‘issue publics’ provide clues on how political opinion is veering away from traditional left and right divisions, and Twitter communication is more complicated than the imaginary cocooned talk in echo chambers, especially for issues beyond elections and direct policing.

Limitations of our method stem from the subjectivity of coding and from the low number of coded influencers, but these may be partially overcome by automatisation of coding collections and the increase of the number of coded users thanks to automatisation. Our method may be applied to detect hidden issue-oriented polarization beyond one-dimensional left-right political spectra.

\subsection{Please Follow Us: Media roles in Twitter discussions in the United States, Germany, France, and Russia}\label{subsec:ch5/sec1/sub2}

\subsection{Multi-dimensional echo chambers: Language and sentiment structure of Twitter discussions on the Charlie Hebdo case}\label{subsec:ch5/sec1/sub3}

\subsubsection{1. Introduction}

Public discussions on social networks potentially have trans-border and multilingual nature. This comes true in heated conflictual discussions that reach global trending topics. Such discussions are expected to demonstrate ‘civilizational clashes’ \cite{AnKwakMejova}.

Being part of the global public sphere, since the 1990s, such discussions were expected by many observers to be more horizontal, all-involving, and democratically efficient \cite{Fuchs} than the traditional mass-mediated discussions \cite{McQuail}. But, with time, criticism towards the democratic quality of discussions in social media arose, with many works discovering the patterns of echo chambering and discourse polarization in social networks \cite{Sunstein2001,Sunstein2002,BarberaJostNagler,BastosMerceaBaronchelli,ColleoniRozzaArvidsson,ConoverRatkiewiczFrancisco}, which lowered the capacities of inter-group discussions and, thus, just formed an additional line of social segregation.

Object-oriented hashtagged discussions have been thoroughly studied in the 2010s, including those on political and social conflicts. But there is still scarce knowledge on whether affective hashtags \cite{Papacharissi} that convey emotions -- either of solidarity with or of anger towards a particular social group -- work in terms of user clusterization. Also, there is no clear understanding of comparative democratic quality of emotionally ‘positive’ and ‘negative’ hashtags in terms of echo chambering.

In this paper, we address these gaps by analyzing the Twitter discussion on the \textit{Charlie Hebdo} massacre of 2015. In the discussion upon the mass killings, the Twittershpere has created \#jesuischarlie and \#jenesuispascharlie -- two emotionally differing discussion clusters with, allegedly, opposite sentiments towards the journal’s ethics and freedom of speech; the hashtags soon became ‘role models’ for online solidarity towards the victims of terrorist attacks and anthropogenic disasters.

To analyze the echo chambering patterns in the two discussions, we have focused upon two levels of echo chambers. We were wondering whether echo chambers formed on the level of a hashtag (based on language use) and within a particular language (based on user sentiment of French-speaking users).
The remainder of the paper is organized as follows. Section 2 reviews the literature on echo chambering in social media. Section 3 presents our methodology and the conduct of the research. Section 4 presents our results and discusses them.

\subsubsection{2. User Groupings on Twitter and the Efficacy of Public Sphere}

\paragraph{2.1 Social Media and the Public Sphere: Echo Chambers vs. Opinion Crossroads}
Public sphere as a spatial metaphor for a complex of discussions and procedures with a public status and decision-making goals \cite{Kleinstuber} has been amplified by the appearance of social media in the 2000s. By 1990s, it had been established in the academic literature that mediatized public sphere with traditional media playing the role of information hubs was uneven and hardly efficient in terms of access to opinion expression, as whole social groups remained under-represented, and newsmakers privileged in comparison to \textit{vox populi}. With the appearance of social media, hopes arose that the new communicative milieus would foster horizontalization of communication and provide for democratization and higher political participation \cite{Fuchs}. Also, hopes for better understanding and resolution of non-political inter-group conflicts existed.

But with time, these hopes fainted, as offline disparities seemed to reproduce online, including political interests, race, gender, and other inequalities \cite{Daniels}; emotion and affect proved to rule the discourse \cite{Papacharissi}, with publics even in most democratically developed countries moving from diverse in opinion to dissonant and disconnected \cite{Pfetsch}. With the development of social network analysis (SNA) and its application to social media research, the question of the efficacy of the public discussions in social media \cite{BrunsHighfield} became linked to network and structural features of the discussions, such as the influencer status \cite{BodrunovaBlekanovMaksimov,BodrunovaLitvinenkoBlekanov2016} and clusterization of users also known as user polarization \cite{BarberaJostNagler,BastosMerceaBaronchelli} and echo chambering \cite{ColleoniRozzaArvidsson,ConoverRatkiewiczFrancisco}. Some evidence was also gathered on discussion sphericules forming on the global scale just as well as nationally \cite{CammaertsAudenhove}.

\paragraph{2.2 Why the Twitter Discussions Fragment: Linguistic Properties of Speech as Catalyzers of Echo Chambering}
In early studies of social networks and its users, authors interested in testing the ability of networks to pull together users from distant locations and weak ties linked geo- graphical distance with factors like residence of users and their language profile \cite{TakhteyevGruzdWellman}. This is why Twitter that enabled the (arguably) quickest possible information spread across locations and languages became a major attractor of scholarly attention \cite{LotanGraeffAnanny,HongConvertinoChi}. But despite the global reach of the platform, several studies have found that people were still connected locally on Twitter \cite{CammaertsAudenhove,YardiBoyd}.

Along with locality and residence, linguistic factors, arguably, play a major role in user grouping on the global scale. Thus, the language(s) used by the discussion participants is the first natural barrier that is expected to make users group together and communicate within their language-based echo chambers \cite{ChenTuZheng}, both on Twitter on the whole and within particular hashtags \cite{BastosPuschmannTravitzki}.

Other factors have also been discussed as the catalyzers of user grouping on Twitter. Among those, political attitudes lead the research agenda \cite{BarberaJostNagler,BastosMerceaBaronchelli}. Here, several ways to detect user clusterization exist. Of them, use of network or semantic proxies like friendship ties \cite{BarberaRivero}, patterns of following \cite{Rivero} and retweeting \cite{CalaisGuerraMeiraJrCardie}, content sharing \cite{ColleoniRozzaArvidsson,BakshyMessingAdamic} etc. is till today the most prominent; another is automated analysis of user sentiment, either in general or toward an issue/actor in question (object-oriented) \cite{ConoverGoncalvesRatkiewicz}.

But all these studies depict user groupings within a single dimension; our idea is to try and trace user groupings multi-dimensionally -- both on the level of a hashtag (based on language use) and within a language nebula (based on user sentiment).

\paragraph{2.3 The \textit{Charlie Hebdo} Case: Emotional Hashtags and User Groupings}
To search for multi-dimensional echo chambers, we have chosen the case of the \textit{Charlie Hebdo} massacre of 2015. Here, we could hypothesize that the existence of emotionally opposite hashtags (\#jesuischarlie and \#jenesuispascharlie) already creates enclaves within the general discussion on the case. Then, within the hashtags, several clusters based on language structure may exist. Then, on the third level, we will look whether within the language clusters sub-clusters of sentiment form. In general, our idea is to see how exactly the language clusters correspond to the sentiment clusters in each language and whether ‘positive’ users within one language are linked to such in another language, while ‘negative’ users also group across languages in a similar way. But here we present only preliminary results that check if the user clusters may at all be detected based on language and on sentiment within a language.

\paragraph{2.4 Research Hypotheses}

Thus, our hypotheses are the following:
\begin{itemize}
	\item H1a. Non-random user groups will be detected for both hashtags, as based on
	language use.
	\item H1b. \#jesuischarie will not differ from \#jenesuispascharlie in their language
	structure, as both hashtags have reached global trending topics and are expected to show ‘civilizational clashes’.
	\item H2a. Non-random user groups will be detected within one language (French), as based on positive and negative sentiment.
	\item H2b. \#jesuischarlie will differ from \#jenesuispascharlie in the grouping based on user sentiment, due to the emotional opposition of the hashtags themselves.
	\item H3. Multi-layer echo chambering may be detected in trans-border hashtagged discussions of global reach.
\end{itemize}

\subsubsection{3. Data Collection and Conduct of Research}

\paragraph{3.1 Data Collection and the Datasets}
Using a web crawler developed especially for the Twitter data collection, we gathered all the tweets published openly under the hashtags \#jesuischarlie and \#jenesuispascharlie (by separate crawls) within January 7 to 9, 2015, as these three days covered the active conflict (from the killings in the editorial office to the assailants’ death) when the users provided virtually millions of tweets for collection.

The collected datasets included: for \#jesuischarlie: 420,080 tweets; 266,904 tweeters; 719,503 users who interacted with the posted tweets (by likes, retweets, or comments); for \#jenesuispascharlie: 7,698 tweets; 5,466 tweeters; 17,872 users who posted and interacted with the posted tweets (by likes, retweets, or comments).

These full datasets were later used to reconstruct the overall web graphs for the two hashtags. But the datasets were very different in size, and be able to color them with language markers, we needed to sample the users for coding having in mind the volume difference of the datasets.

\paragraph{3.2 Language Analysis: Sampling, Coding, and Graph Reconstruction} 
To answer H1a and H1b, we have coded the users for their language use and then applied these data to the datasets for web graph reconstruction.

After eliminating bots and bot-like users (those who posted over 60\% of doubled tweets) as well as hashtag-only tweets, we have followed the strategy developed by the research group for previous Twitter studies \cite{Authors2016a,Authors2016b,Authors2018}, namely uniting random sampling with detection of influential users (influencers) for taking them into account. Then, we have coded all the influencers (disregarding the number of tweets they posted; for \#jesuiuscharlie, 402 users, for \#jenesuispascharlie, 85 users) and ‘ordinary users’ sampled in the feasible and comparable way. For \#jenesuispascharlie that was substantially smaller, all the users with 3 and more tweets were coded (339 users); for \#jesuischarlie, the ‘ordinary users’ with 5 tweets or more were taken into account (9,090 users), and of them, each second was coded (4500 users).

All the sampled underwent expert reading and were coded manually marking the number of tweets in language 1, language 2, and other languages; thus, users posting on one, two, and three or more languages were defined. The languages were identified for each user; in case of rare languages, Yandex language identifier was employed.

To reconstruct the graphs, we use Gephi API algorithms openly available online. Of the available algorithms, two were chosen: Hu \cite{Hu} and OpenOrd \cite{MartinBrownKlavans}; here, YifanHu-based graphs are presented, as the OpenOrd graphs require more space for presentation. We colored both the nodes (users) and the edges (connections between users). To prove that the visual nebulae are not artifacts of subjective viewership, we calculated the percentage of edges between and inside language groups, eliminating the ‘loops’ of self-commenting/liking by the users.

\paragraph{3.3 Sentiment Analysis: Sampling, Vocabulary Building, and Graph Reconstruction}After we have proved that the French-speaking users show non-random grouping in both cases (see below), we have taken them for sentiment analysis. The number of users for \#jesuischarlie included 1291 user; for \#jenesuispascharlie, 117 users.

Our strategy for French-language sentiment analysis was the following. We have united three sources in our vocabulary: the existing French dictionary with sentiment marking, machine translation from an additional Wordnet vocabulary, and the case-based vocabulary created from the collected tweets and manually marked for positive, negative, and neutral sentiment regardless of the case-specific meanings.

This vocabulary was applied to each tweet of the abovementioned French-speaking users; for each tweet, the sentiment was calculated. Then, the thresholds were defined: positive and negative were the users with positive(+neutral) and negative(+neutral) tweets, respectively; neutral were those with neutral tweets only; mixed were those with positive + negative(+neutral) tweets.

Then, we also checked the groupings with by calculating the percentage of edges both between and inside language groups, eliminating the ‘loops’ of self-commenting/liking by the users.

\subsubsection{4. Results and Discussion}

Our results are described below with regard to the hypotheses stated above.

\textit{H1a/H1b.} To assess the user groupings in both hashtagged discussions, we have reconstructed the web graphs for the coded users (see Fig.~\cref{fig:yifanHuGraphs-2} for \#jesuischarlie and Fig.~\cref{fig:yifanHuGraphs-2} for \#jenesuispascharlie, respectively). What we see on the graphs are three nebulae for~\cref{fig:yifanHuGraphs-1}: French, English, and other European, and two for~\cref{fig:yifanHuGraphs-2}: French and Engish. But the results of calculations of percentage of edges between and inside groups tell that the actual grouping is slightly different from what we see with unaided eyes. For \#jesuischarlie, the nebulae with density higher than the inter-group ones are French, English, and French/English (52.1\%, 16.7\%, and 16.9\%, respectively, against 6.17\% for inter-group edges) and not other European. For \#jenesuispascharlie, the graph is much denser (26\%), but still the same three clusters show up, with 26.2\%, 22.4\%, and 18.8\%, respectively; in both cases, other language clusters are virtually non-existent and do not mount to 3\%. As we stated in our earlier investigations \cite{Authors2018}, we have not seen a sign of ‘civilizational clashes’ in any of the hashtags.

\begin{figure}[ht]
	\centerfloat{
		\hfill
		\subcaptionbox[List-of-Figures entry]{\label{fig:yifanHuGraphs-1}}{%
			\includegraphics[width=0.419\linewidth]{yifanHuGraphs1}}
		\subcaptionbox{\label{fig:yifanHuGraphs-2}}{%
			\includegraphics[width=0.382\linewidth]{yifanHuGraphs2}}
		\hfill
	}
	\caption{The YifanHu graphs (fragments) for language distribution in \cref{fig:yifanHuGraphs-1} \#jesuischarlie and \cref{fig:yifanHuGraphs-2} \#jenesuispascharlie. Red: French; blue: English; lilac: French/English; green: other European. (Color figure online)}\label{fig:yifanHuGraphs-12}
\end{figure}

Thus, H1a is proven; H1b is proven too but not due to ‘civilizational clashes’.

\textit{H2a/H2b.} To see the user grouping and sentiment cleavages within the French-speaking parts of the discussions, we have reconstructed the web graphs for them (see Fig.~\cref{fig:yifanHuGraphs-3} for \#jesuischarlie and Fig.~\cref{fig:yifanHuGraphs-4} for \#jenesuispascharlie, respectively).

\begin{figure}[ht]
	\centerfloat{
		\hfill
		\subcaptionbox[List-of-Figures entry]{\label{fig:yifanHuGraphs-3}}{%
			\includegraphics[width=0.419\linewidth]{yifanHuGraphs3}}
		\subcaptionbox{\label{fig:yifanHuGraphs-4}}{%
			\includegraphics[width=0.4\linewidth]{yifanHuGraphs4}}
		\hfill
	}
	\caption{The YifanHu graphs (fragments) for language distribution in \cref{fig:yifanHuGraphs-3} \#jesuischarlie and \cref{fig:yifanHuGraphs-4} \#jenesuispascharlie. Red: French; blue: English; lilac: French/English; green: other European. (Color figure online)}\label{fig:yifanHuGraphs-34}
\end{figure}

Here, H2a should be rejected for \#jesuischarlie and partly supported for the second hashtag. Both in the graph and in the edge percentage calculations, it is only users with mixed sentiment who form a group (38.33\% against 56.05\% for the inter-group connections) in \#jesuischarlie. But for \#jenesuispascharlie, both mixed and negative user groups seem to have a potential for grouping (19.5\% and 15\% against 62.1\% for inter-group connections). Thus, H2b is supported: the cases do differ.

Even with the cases of such different sizes, H3 is supported: thanks to the negative nebula in \#jenesuispascharlie, we can state that echo chambers are able to form on at least two levels of the trans-border conflictual discussions of global (or, more precisely, macro-regional) reach. This adds to our understanding of the nature of public discussions in social media, even if lowers hopes for all-encompassing public spheres.

\subsection{Power laws in \textit{Ad Hoc} conflictual discussions on twitter. Digital Transformation and Global Society}\label{subsec:ch5/sec1/sub5}

\subsubsection{1. Introduction}

Public discussions online, and, of these, the ones on social networking sites, have become a growing area of scholarly attention, as they have been perceived as a manifestation of the public sphere, a crucial condition for efficient democratic deliberation \cite{Habermas}. Thus, understanding how networked discussions form and evolve is necessary for elaboration of proper criteria for their efficiency evaluation.

One of the issues recently raised by social science and communication scholars is that of comparability of the online discussions formed by the so-called issue publics \cite{Dahlgren}\cite[p.~422]{Habermas}. Such discussions form quickly (or even burst out) around events or burning issues. Due to their \textit{ad hoc} nature \cite{BrunsBurgess}, they may dissolve just as quickly, involve various actors, and are affective \cite{Papacharissi} and, thus, are shaped by emotions rather than by rational argumentation. The question remains whether we have grounds for comparing such discussions, as the differences in discussion substance may lead to critical differences in the discussion structure and connectivity, which would make, e.g., cross-cultural and cross-language comparisons of same-topic discussions impossible. Also, conclusions made for one discussion in terms of actor roles, dynamics, other constitutive parameters would not allow for predicting them for other discussions if the network structures are not assessed and recognized as similar.

This question has not yet been properly addressed in the social network analysis literature, as structural similarities of networked discussions, despite all the attention given to them, remain understudied in comparative perspective, especially in the view of the public sphere theory. The latter has elaborated its own vision of how to assess the efficacy of public discussions. Linking the two research areas for elaborating the SNA parameters for such assessment of the discussion networks would address one of the existing gaps in social network studies. For example, a number of more recent works have juxtaposed calculated and \textit{ad hoc} publics \cite{BrunsBurgess2015,LynnRosatiNair} in search of networking patterns of both types of discussions, but these studies were single-case, and we still lack the knowledge whether structural patterns vary across cases, cultures, and types of vocabularies used for data collection.

Another input for social network analysis from the public sphere theory would be addressing the difference between users who are key for the discussion outburst and random discussion participants. Traditional view would imply various manifestations of power distance between the two user groups, and one needs to know whether online discussions show stable patterns of differences between key and random users.

We address these research gaps by collecting data and analyzing web graph structures for five discussion outbursts about inter-ethnic (inter-national and inter-race) conflicts in the USA, Germany, France, and Russia of the 2010 s, reviewing altogether six discussion cases, as we split one into neutral and affective parts. We collect data from Twitter based on single hashtags and keyword conglomerates and show that there are repeatable structural patterns of the discussions across these cases.

The remainder of the paper is organized as follows. Section 2 discusses today’s literature on various aspects of discussion structure for comparative tasks. In Sect. 3, we describe the cases and formulate the hypotheses for their comparison. Section 4 shows and discusses the discovered results. In conclusion, we discuss applicability of our findings for further studies of online conflictual discussions.

\subsubsection{2. Comparing Discussion Structure on Social Networks: A Literature Review}

Ad hoc \textit{discussions: a plea for grounds for comparison}. Public discussions on social networking sites have drawn scholarly attention to their various aspects, including the connection between the deliberative power of ‘ordinary citizens’ and the platform and network features that might either empower or disempower user groups in their opinion expression, shape the discussion outcomes, and impact the respective decision-making, as well as inspire political mobilization via the ‘logic of connective action’ \cite{BennettSegerberg}. Thus, the discussions in online networks have been thoroughly studied in their substantial aspects. But one of the key issues in this research area is the principal possibility of comparisons between discussions on similar topics happening in various parts of the world and in varying times.

So far, comparative studies of cross-country social network discussions have been rare \cite{BodrunovaLitvinenkoBlekanov2017}; one of the reasons for that is the scholarly argument of low (or, rather, unknown) comparability of the discussions that are formed by the issue publics. This type of the discussion raises, arguably, the biggest amount of doubt ‘whether Habermas is on Twitter’ \cite{BrunsHighfeld2016}, as, from case to case, current research demonstrates varying results on echo chamber formation \cite{BastosMerceaBaronchelli,ColleoniRozzaArvidsson,Sunstein2001,YardiBoyd}, cross-group discussion potential \cite{Barbera,BarberaJostNagler}, and the roles of influencers detected by various means \cite{BastosRaimundoTravitzki,BodrunovaLitvinenkoBlekanov2017,DuboisGaffney}.

The reason for low comparability, as scholars argue, lies in the very nature of the discussions, as they are all \textit{ad hoc} -- that is, case-specific and random in formation. Such discussions have the outburst nature, may have varying patterns of dissolution, involve case-specific actors, and based on affect \cite{Papacharissi} -- that is, are shaped by emotions rather than by rational argumentation, which may add to the non-comparability of the discussions. Thus, the major research question is -- are \textit{ad hoc} discussions comparable in topical and sociological terms also comparable in the discussion structures? Or, in other words, are there structural features characteristic of such discussions, that would, ideally, distinguish them from other discussion types on social networks and on the Web on the whole? And what could be the markers for such \textit{ad hoc} discussion outbursts?

Today, in Twitter studies, there is scarce but growing evidence that \textit{ad hoc} discussions differ in their nature from other discussion types. Thus, several papers underline the differences between calculated and \textit{ad hoc} publics \cite{BrunsBurgess2015,LynnRosatiNair}, but no network parameters were used to prove the differences between these discussion types. Also, these studies were based on single cases, and we still lack the knowledge whether structural patterns for \textit{ad hoc} discussions vary or repeat across cases, cultures, and types of vocabularies used for data collection. This creates a focus for our enquiry.

\paragraph{Degree centralities as a proxy for discussion comparability and a potential discussion quality metric.} Degree centralities (in-degree, out-degree, and degree accumulating both of these) have been long ago recognized as the key structural metrics of user relations and random graph assessment \cite{AlbertBarabasi,BoccalettiLatoraMoreno,MisloveMarconGummadi}; degree distributions are recognized as a key variable describing users’ interest toward each other \cite{EdigerJiangRiedy}. At the same time, from the normative viewpoint relevant in social sciences, degree centralities are an important metric of user in-network influence \cite{BodrunovaLitvinenkoBlekanov2016}, as well as the metric important in deliberative terms: the more users are reached by the same user (or reach a given user, which is the same for non-directed graphs but not for directed graphs), the more probable is the chance for cross-opinion discussion. More importantly, it can also characterize general user involvement into the discussion in comparison with other discussions; in terms of user involvement, degree centralities are, arguably, the most telling. A range of more topic-specific research papers have also argued that degree-based metrics are useful for network-based studies of social conflict and dangerous networks. Of those, one work \cite{KarthinkaGeethaBose} has shown the importance of relational measures (as the authors noted, ‘who is related to whom’) in detecting the key attackers in the terrorist network.

There are, of course, several criticisms about degree distributions, as they gener- alize the network structure without taking into account the role of influencers (for the reviews on detecting and comparing influencer structures, see \cite{BodrunovaBlekanovMaksimov,BodrunovaLitvinenkoBlekanov2017,BodrunovaLitvinenkoBlekanov2016}), or ‘hub users’. Their roles, in various works, are assessed in different, if not directly opposite, ways. Thus, one stream of works insists on their disproportionately big role in information distribution (see \cite{BastosRaimundoTravitzki,DuboisGaffney}, and many others), while another line shows that such hubs may be inefficient due to their overload and incapability of transmitting information due to that \cite{HarriganAchananuparpLim}. But our goal is to test whether the discussions on the whole may be distinguished from Twitter on the whole.

As major research works in the area note \cite{BarabasiAlbert,BroderKumarMaghoul,HubermanAdamic}, due to preferential connectivity in real-world networks, a certain power law (just as its absence) in degree distribution may become the characteristic that describes specific network structures. Thus, networks expressing power-law-like degree distributions are known as scale-free \cite{AlbertBarabasi,EdigerJiangRiedy}, and ‘scale-free tweet mention graphs would imply that a few Twitter users and ‘mentioners’ are responsible for a disproportionately high fraction of a community’s discourse’ \cite[p.~587]{EdigerJiangRiedy}. Thus, degree distribution may serve as a proxy for, e.g., core vs. periphery assessment in terms of aggregate influencer power: it can tell whether the network is dominated by a small number of users \cite{YeWu}. Thus, for Twitter, it would allow for: (1) differentiating one network type from another; (2) proving that the discussions are similar in their core vs. periphery relations, which suits our research goals. Moreover, our previous studies \cite{BlekanovSergeevMaksimov2017} have shown that power law exponents, indeed, vary for different types of web structures; e.g. for university websites, the average exponent is 1.8 instead of the expected 2.1.

Early works on World Wide Web topology mentioned above, as well as other research papers (see, e.g., \cite{FaloustosFaloustosFaloustos}), all stated that power law distributions are characteristic for the Web on the whole. But another, later line of research papers has argued that the Web, by just evolving in time, has blurred the initial power-law-like degree distribution \cite{ChenChangGovindan,MeuselVignaLehmberg}. Smaller-scale research, though, tells that degree distributions are still valid for smaller real-world samples.

Today, more and more criticism is raised about the explanatory potential and the very existence of scale-free networks -- that is, of the power law in degree distributions -- in large human-based datasets \cite{BroidoClauset}. We, thus, want to test whether the certain power laws show up in ad hoc discussions, as distinguished from Twitter on the whole, cultivated long-term discussions, and random talk.

\paragraph{Power laws on Twitter: lack of comparative studies on the nature of the discussions.} So far, degree distributions have been studied for either the whole Twittersphere or its geographical segments; and, in these studies, the evidence on whether all the discussions on Twitter manifest power-law degree distributions is mixed.

Thus, cumulative degree distributions \cite{Newman} for Twitter on the whole studied a decade ago \cite{JavaSongFinin} showed that the slopes cin and cout [were] both approximately \(-2.4\). The authors have interpreted these figures as similar to those for the Web of those times (\(-2.1\) for in-degree, cf. \cite{DonatoLauraLeonardi}) and blogosphere (\(-2,38\), as based on Blogpulse conference dataset of 2003). With the latter claim, one could agree, but with the former one we would not, as our research shows that differences of 0.3 in exponent values may be characteristic of graph origin and/or discussion type \cite{BlekanovSergeevMaksimov2017}. Another study \cite{WengLimJiang} based on the data dated mostly from April 2008 to April 2009 and featuring the most followed Twitter users from Singapore has also demonstrated power-law-like degree distributions of the user networks. A Twitter-large study \cite{WelchSchonfeldHe} has shown a very different exponent value of \(-1.6\). But authors \cite{KwakLeePark} have crawled the whole Twittersphere and have not discovered any power law in link distribution on the global level, as other authors underline \cite{HansenArvidssonNielsen}.

Several studies have focused on country-based segments of Twitter. Thus, one study \cite{MyersSharmaGupta} has dealt with the Twitter follow graphs. The researchers have shown that in-degree and degree on Twitter are best fit by power law, while out-degree is best fit by log-normal distribution. Besides analyzing the entire Twitter, it also did country- based segment studies for Brazil, Japan, and the United States; very little variance was found between the countries in terms of in-degree, out-degree, and degree distribution \cite[p.~494]{MyersSharmaGupta}. Another study has dealt with the follow graphs of 10 country-bound Twitter segments around the world \cite{PobleteGarciaMendoza}. It has also demonstrated power-law in-degree and out-degree distributions, with power law coefficient ranging 5.91 to 9.51 and 8.12 to 13.62, respectively. But one also needs to note that user influencer status is linked much more to the number and activity levels of active followers (who retweet and/or mention a given user) than to the number of followers \cite{ChaHaddadiBenevenuto}; thus, one needs to look at the actual discussion graphs rather than at the graphs of following (much less linked to the real-world issue-based discussions) to more precisely determine the influential users, as well as to define the discussion type.

Even a fewer number of works have examined the degree distributions in hash- tagged discussions. We can name one work \cite{ZhouBandariKong} on Iranian elections; the discussion there also showed power law distributions, with in-degree being \(-2.85\) and out-degree being \(-2.42\), while the retweet-based network had in-degree of \(-1.94\). Another group of authors \cite{ConoverRatkiewiczFrancisco} have also observed that retweet- and mention-based networks virtually did not differ in their scale-free topology.

From the review above, one can conclude that scholarly evidence for power laws in degree distributions is greater than that of the opposite. But, despite their potential, degree distributions have not been tested for \textit{ad hoc} discussion outbursts in comparative perspective.

Our idea is to look how degree distributions work if we step-by-step eliminate the users with low degree index, starting from isolate users \((D = 0\)), and look at the cumulative degree distributions for each case.

\subsubsection{3. The Research Methodology}

\paragraph{The research questions.} Most of the available research proves that power laws are characteristic for Twitter ‘calm’ discussions, be it the whole Twitter or its parts, either hashtag-based or limited by region. But the exponent values of degree distribution vary highly, not allowing for any particular expectation. Thus, we ask: Will degree distribution of all the \textit{ad hoc} discussion be fit by a power law? Will the exponent values be similar, thus indicating that the discussions are similar? Will the exponent values differ from \(\lvert2.1\rvert\)? Will the exponent values be similar enough across world regions, hashtag- only / keyword conglomerates, and neutral / affective hashtags?

The research hypotheses that emerge of these research questions look as follows:

\begin{itemize}
	\item H1. All the discussions under scrutiny will be characterized by power law in degree distribution.
	
	\item H2. All the discussions will diverge from the \(\lvert2.1\rvert\) exponent value to the same direction and on comparable percentage. This includes hashtag vs. keyword conglomerates (H3).
	
	\item H4. Neutral and affective (expressing emotions of either sympathy and compassion or negation and hatred) hashtags will be comparable in their power law exponents.
\end{itemize}

\paragraph{The cases under scrutiny and their substantial comparability.} The cases of our attention all have the same set of features that is to ensure that the cases are comparable in sociological terms. Thus, all the cases have a violent trigger (a killing or rape/harassment) of inter-ethnic nature; the discussions are outburst -- that is, the number of users involved has one sharp peak and then slows down gradually; media report the communities to split into the minority, pro-minority majority, and anti-minority majority groups; there is peaceful protest or mass commemoration of the victims in the aftermath of the conflict; there is direct involvement of authorities of several levels into conflict resolution; and all the conflicts provoke a discussion on Twitter that gets to national (sometimes also to global) Twitter trending topics.

We are looking at the following cases (in chronological order):
\begin{itemize}
	\item A killing of a Russian Muscovite by an Uzbek immigrant and the subsequent anti- immigrant clashes in the Moscow district of Biryuliovo, Russia (2013);
	\item A killing of an African American teenager by a white police officer and the sub- sequent city riots in Ferguson, the USA (2014);
	\item The attack to the editorial office of Charlie Hebdo and the subsequent peaceful demonstrations in Paris, France (2015);
	\item The mass harassment of German females by male re-settlers from Middle East and North Africa in the New Year Eve in Cologne, Germany (2015--2016), and the subsequent protest meetings by PEGIDA and the ‘Alternative for Germany’ party;
	\item A bus attack at one of the Christmas markets in Berlin (2016).
\end{itemize}

In each case, the discussion data were collected and the respective web graph reconstructed; in the case of \textit{Charlie Hebdo}, two graphs (one for a neutral hashtag and one for a compassion hashtag) were reconstructed.

\paragraph{Data collection.} To collect the discussion bulk, we have created a specialized web crawler with adjustable modules \cite{BlekanovSergeevMartynenko}. It was done especially to overcome the well-known Twitter API limitations, like the ones on the number of requests to server and on the number of tweets available for download. It also bypasses the popularity algorithm and allows for human-like backfolding.

For bigger-scale discussions (Ferguson and Charlie Hebdo), one hashtag per graph was used. For smaller-scale discussions, snowballing reading of 1.000+ random tweets containing the primary hashtag/keyword was performed, thus bringing on keyword collections.

\begin{table} [htbp]%
	\centering
	\caption{The datasets collected.}%
	\label{tab:collectedDatasets}% label всегда желательно идти после caption
	\renewcommand{\arraystretch}{1.5}%% Увеличение расстояния между рядами, для улучшения восприятия.
	\begin{SingleSpace}
		\begin{tabulary}{\textwidth}{@{}>{\zz}L >{\zz}C >{\zz}C@{}} %Вертикальные полосы не используются принципиально, как и лишние горизонтальные (допускается по ГОСТ 2.105 пункт 4.4.5) % @{} позволяет прижиматься к краям
			\toprule     %%% верхняя линейка
			The case & The number of nodes & The number of edges \\
			\midrule %%% тонкий разделитель. Отделяет названия столбцов. Обязателен по ГОСТ 2.105 пункт 4.4.5
			Biryuliovo, 2013 & 11429 & 20106 \\
			Ferguson, 2014 & 169677 & 334050 \\
			\textit{Charlie Hebdo}, 2015 (neutral) & 952615 & 1782863 \\
			\textit{Charlie Hebdo}, 2015 (affective) & 719503 & 981131 \\
			Cologne, 2015--2016 & 40117 & 98508 \\
			Berlin, 2016 & 194937 & 298562 \\
			\bottomrule %%% нижняя линейка
		\end{tabulary}%
	\end{SingleSpace}
\end{table}

The initial datasets collected are represented in Table~\cref{tab:collectedDatasets}. Altogether, tweets by over 2 mln users were included into the research. Edges between the users were formed if any type of substantial interaction emerged between them (we counted retweets, mentions, and comments as such).

\paragraph{Web graph reconstruction and analytics.} We have reconstructed the graphs for each case. The graphs were non-directed, as we were interested in the aforementioned deliberative aspects of the discussions. We have used the OpenOrd algorithm for the graph reconstruction. This methods is based on calculating the maximum distance (in \%) between the two nodes and on a particular number of iterations (in all our cases, the algorithm converged the proposed 800 iterations to 750). This algorithm was chosen, as it brings to the discussion core (that is, to the center of the graph) the influential users defined by a wide set of variables, including absolute-figure ones (the number of tweets, likes, retweets, mentions, and comments) and centrality metrics (including degree, betweenness, parerank, and eigenvector centralities).

Then, the degree distribution exponents were calculated for each case, using eleven steps of isolate user elimination (for users with \(D = 0\) to 10 in the initial graph).

The results and answers to our hypotheses are presented below.

\subsubsection{4. The Research Methodology}

As stated above, we have calculated degree distribution exponents for each case. In order to clearly highlight the differences with the expected exponent value of \(\lvert2.1\rvert\) based on previous Web studies, we have also calculated the exponents for the eliminated users. The exponent values for both active users and eliminated users, as well as their absolute deviations from \(\lvert2.1\rvert\), are presented in Table 2. The degree distribution graphs are presented in Figs.~\cref{fig:biryuliovoDegreeDistribution},~\cref{fig:fergusonDegreeDistribution},~\cref{fig:charlieNDegreeDistribution},~\cref{fig:charlieADegreeDistribution},~\cref{fig:cologneDegreeDistribution} and~\cref{fig:berlinDegreeDistribution} for the respective cases, as ordered in Table~\cref{tab:caseExponentValues}.

\begin{table} [htbp]%
	\centering
	\caption{Exponent values for the cases and their divergence from the expects exponent value.}%
	\label{tab:caseExponentValues}% label всегда желательно идти после caption
	\renewcommand{\arraystretch}{1.5}%% Увеличение расстояния между рядами, для улучшения восприятия.
	\begin{SingleSpace}
		\begin{tabulary}{\textwidth}{@{}>{\zz}L >{\zz}C >{\zz}C >{\zz}C >{\zz}C@{}} %Вертикальные полосы не используются принципиально, как и лишние горизонтальные (допускается по ГОСТ 2.105 пункт 4.4.5) % @{} позволяет прижиматься к краям
			\toprule     %%% верхняя линейка
			The case & Exponent value, active users &  Absolute deviation, active users & Exponent value, eliminated users & Absolute deviation, eliminated users \\
			\midrule %%% тонкий разделитель. Отделяет названия столбцов. Обязателен по ГОСТ 2.105 пункт 4.4.5
			Biryuliovo, 2013 &  \(\lvert1.56\rvert \) & 0.54 & \(\lvert2.04\rvert\) & 0.06 \\
			Ferguson, 2014 & \(\lvert1.39\rvert\) & 0.71 & \(\lvert2.16\rvert\) &\( -0.06\)      \\
			\textit{Charlie Hebdo}, 2015 (neutral) & \(\lvert 1.59 \rvert\) & 0.51 & \(\lvert2.01\rvert\) & 0.09 \\
			\textit{Charlie Hebdo}, 2015 (affective) & \(\lvert1.80\rvert\) & 0.3 & \(\lvert2.45\rvert\) & \(-0.35\) \\
			Cologne, 2015--2016 & \(\lvert1.28\rvert\) & 0.82 & \(\lvert1.91\rvert\) & 0.19 \\
			Berlin, 2016 & \(\lvert1.63\rvert\) & 0.47 & \(\lvert2.34\rvert\) &  \(-0.24\) \\
			\bottomrule %%% нижняя линейка
		\end{tabulary}%
	\end{SingleSpace}
\end{table}

\begin{figure}[ht]
	\centerfloat{
		\hfill
		\subcaptionbox[List-of-Figures entry]{\label{fig:biryuliovoDegreeDistribution-1}}{%
			\includegraphics[width=0.5\linewidth]{biryuliovoDegreeDistribution1}}
		\subcaptionbox{\label{fig:biryuliovoDegreeDistribution-2}}{%
			\includegraphics[width=0.56\linewidth]{biryuliovoDegreeDistribution2}}
		\hfill
	}
	\caption{Degree distributions for the Biryuliovo case: (а) active users; (б) eliminated users. \textit{Source:} authors.}\label{fig:biryuliovoDegreeDistribution}
\end{figure}

\begin{figure}[ht]
	\centerfloat{
		\hfill
		\subcaptionbox[List-of-Figures entry]{\label{fig:fergusonDegreeDistribution-1}}{%
			\includegraphics[width=0.509\linewidth]{fergusonDegreeDistribution1}}
		\subcaptionbox{\label{fig:fergusonDegreeDistribution-2}}{%
			\includegraphics[width=0.55\linewidth]{fergusonDegreeDistribution2}}
		\hfill
	}
	\caption{Degree distributions for the Ferguson case: (а) active users; (б) eliminated users. \textit{Source:} authors.}\label{fig:fergusonDegreeDistribution}
\end{figure}

\begin{figure}[ht]
	\centerfloat{
		\hfill
		\subcaptionbox[List-of-Figures entry]{\label{fig:charlieNDegreeDistribution-1}}{%
			\includegraphics[width=0.56\linewidth]{charlieNDegreeDistribution1}}
		\subcaptionbox{\label{fig:charlieNDegreeDistribution-2}}{%
			\includegraphics[width=0.5\linewidth]{charlieNDegreeDistribution2}}
		\hfill
	}
	\caption{Degree distributions for the \textit{Charlie Hebdo} neutral case: (а) active users; (б) eliminated users. \textit{Source:} authors.}\label{fig:charlieNDegreeDistribution}
\end{figure}

\begin{figure}[ht]
	\centerfloat{
		\hfill
		\subcaptionbox[List-of-Figures entry]{\label{fig:charlieADegreeDistribution-1}}{%
			\includegraphics[width=0.5\linewidth]{charlieADegreeDistribution1}}
		\subcaptionbox{\label{fig:charlieADegreeDistribution-2}}{%
			\includegraphics[width=0.55\linewidth]{charlieADegreeDistribution2}}
		\hfill
	}
	\caption{Degree distributions for the \textit{Charlie Hebdo} emotional case: (а) active users; (б) eliminated users. \textit{Source:} authors.}\label{fig:charlieADegreeDistribution}
\end{figure}

\begin{figure}[ht]
	\centerfloat{
		\hfill
		\subcaptionbox[List-of-Figures entry]{\label{fig:cologneDegreeDistribution-1}}{%
			\includegraphics[width=0.51\linewidth]{cologneDegreeDistribution1}}
		\subcaptionbox{\label{fig:cologneDegreeDistribution-2}}{%
			\includegraphics[width=0.55\linewidth]{cologneDegreeDistribution2}}
		\hfill
	}
	\caption{Degree distributions for the Cologne case: (а) active users; (б) eliminated users. \textit{Source:} authors.}\label{fig:cologneDegreeDistribution}
\end{figure}

\begin{figure}[ht]
	\centerfloat{
		\hfill
		\subcaptionbox[List-of-Figures entry]{\label{fig:berlinDegreeDistribution-1}}{%
			\includegraphics[width=0.51\linewidth]{berlinDegreeDistribution1}}
		\subcaptionbox{\label{fig:berlinDegreeDistribution-2}}{%
			\includegraphics[width=0.55\linewidth]{berlinDegreeDistribution2}}
		\hfill
	}
	\caption{Degree distributions for the Berlin case: (а) active users; (б) eliminated users. \textit{Source:} authors.}\label{fig:berlinDegreeDistribution}
\end{figure}

\textit{H1.} All the cased that we have studied do, indeed, demonstrate power laws in degree distribution, and this is true for both active users (see Figs.~\cref{fig:biryuliovoDegreeDistribution-1},~\cref{fig:fergusonDegreeDistribution-1},~\cref{fig:charlieNDegreeDistribution-1},~\cref{fig:charlieADegreeDistribution-1},~\cref{fig:cologneDegreeDistribution-1} and~\cref{fig:berlinDegreeDistribution-1}) and eliminated users (see Figs.~\cref{fig:biryuliovoDegreeDistribution-2},~\cref{fig:fergusonDegreeDistribution-2},~\cref{fig:charlieNDegreeDistribution-2},~\cref{fig:charlieADegreeDistribution-2},~\cref{fig:cologneDegreeDistribution-2} and~\cref{fig:berlinDegreeDistribution-2}). Thus, power law is indicative for \textit{ad hoc} discussion outbursts across cultures and cases; H1 is supported.

\textit{H2.} All the cases do, indeed, diverge from the \(\lvert2.1\rvert\) expected exponent value to the same direction (see Table~\cref{tab:caseExponentValues}, column 2, exponent values for active users): the received exponent values are all equal or below 1.8, which indicates the lower power distance between the users. At the same time, the exponent values for eliminated users fluctuate around \(\lvert2.1\rvert\) to \(\lvert2.4\rvert\), the figures indicative for the Web degree distributions in various studies mentioned above. Thus, power laws with exponent values definitely below those discovered for the World Wide Web are indicative for \textit{ad hoc} discussions, which makes them, at least in terms of core vs. periphery relations, similar and comparable. The second part of the hypothesis is, though, not clearly supported: deviation from \(\lvert2.1\rvert\) substantially varies in percentage (from almost 40\% for the Cologne case to 14.3\% for \#jesuischarlie), and thus we cannot indicate any figure more precise than the fluctuation between \(\lvert1.28\rvert\) and \(\lvert1.80\rvert\) for ad hoc discussions; due to this, H2 is only partly proven. But we also need to state that, for three of six discussions, the exponent values were between \(\lvert1.56\rvert\) and \(\lvert1.63\rvert\), which corresponds to earlier results in \cite{WelchSchonfeldHe}.

\textit{H3.} For three of the six discussions (the Ferguson case and both Charlie Hebdo cases) the data were collected based on single hashtags (\#ferguson, \#charliehebdo, and \#jesuischarlie, respectively), while other cases were collected by keyword conglomerates ranging from 6 keywords (for the Biryuliovo case) to over a dozen keywords (for both German cases). As we see from Table~\cref{tab:caseExponentValues} and Figs.~\cref{fig:biryuliovoDegreeDistribution-1},~\cref{fig:fergusonDegreeDistribution-1},~\cref{fig:charlieNDegreeDistribution-1},~\cref{fig:charlieADegreeDistribution-1},~\cref{fig:cologneDegreeDistribution-1} and~\cref{fig:berlinDegreeDistribution-1}, there is no difference between single-hashtag and keyword-conglomerate data collection. H3 is supported, but this, paradoxically, might add not only to the evidence that \textit{ad hoc} discussions have similar patterns of degree distribution but also to the evidence that Twitter as a platform fosters the power law degree distributions in any type of discussion. To answer this, more research is needed.

\textit{H4.} As seen from Table~\cref{tab:caseExponentValues} and  Figs.~\cref{fig:biryuliovoDegreeDistribution-1},~\cref{fig:fergusonDegreeDistribution-1},~\cref{fig:charlieNDegreeDistribution-1},~\cref{fig:charlieADegreeDistribution-1},~\cref{fig:cologneDegreeDistribution-1} and~\cref{fig:berlinDegreeDistribution-1}, both neutral and affective hashtags are subjected to power law for both active and eliminated users, and exponent values deviate from \(\lvert2.1\rvert\) to the same direction. H4 is supported. We just need to mention that the compassion hashtag \#jesuischarlie has shown the highest exponent values, thus demonstrating bigger gaps between the influential and ‘peripheral’ users, and this might create room for further comparative investigations.

\paragraph{Discussion.} In search for proof of comparability of \textit{ad hoc} online discussions, we have applied our idea of evaluating degree distributions to six datasets of five comparable Twitter discussions on inter-ethnic conflicts in four countries that happened in the 2010 s. What we have discovered is the following.

First, we have shown that all the discussions we have observed are subjected to power law in degree distributions. Second, we have shown that exponent values for degree distributions in the discussion graphs diverge from the figures indicated for the Web on the whole in previous research. Moreover, they diverge in the same direction and to varying but, to a certain degree, also comparable percentage. Taken together, these findings indicate that, at least in terms of influence, interest, and/or power distribution in such discussions, they are comparable across countries and years, as well as across vocabulary types and neural/affective hashtags. Thus, exponent values may serve as indicators for the type of an online discussion. We have also shown that the peripheral part of the \textit{ad hoc} discussions was always closer to \(\lvert2.1\rvert\) to \(\lvert2.4\rvert\) exponent values discovered earlier for the Web and some social networks, which may be a sign that the \textit{ad hoc} discussions do differ from the ‘average’ network structure of the Web.

But we have also discovered that the exponent values for the discussions were fluctuating around \(\lvert1.6\rvert\) indicated earlier for Twitter on the whole; also, variance in value divergence from the expected figures was too big to state that a particular array of meanings may be indicative for ad hoc discussions and could become their structural marker. This needs further investigation, which might imply experimental design.

Last but not least, we have seen that \textit{ad hoc} discussions, if judged by the exponent values, show the patterns of lower power distance between core and periphery. This may be due exactly to their spontaneous nature, as institutional actors who shape and frame the offline discussions compete with ordinary users, crisis witnesses, and grassroots leaders. This may broaden our views upon discussion outbursts on social media.

\section{Тематическое моделирование контента}\label{sec:ch5/sect2}

\subsection{Topic modeling of conflict ad hoc discussions in social networks}\label{subsec:ch5/sec2/sub1}

\subsubsection{Introduction}

As a result of different media platforms achieving a steady user growth in a recent years more and more people begin to use different social networks as the main source of news on economical, political and social events. In particular, ad hoc discussions emerged which can be defined as a debate about a specific problem. In most cases such discussions appear in the case of controversial events and involve large number of participants.

Presence of such user activity raises the problem of analyzing large volumes of this type of data which has become one of the most important problems in many data analysis tasks, including topic modeling. Topic modeling algorithm in this case is an algorithm that, given the number of topics and the list of user messages can output two distributions: topics over documents and of words over topics. Such algorithm can help in the understanding of different points of view and highlight the main arguments. This can be useful in many ways one of which is the case in which the number of documents is big and we want to know their general content without reading them all. Another use case is data prepossessing, reducing the dimensions of data to use in other analysis tasks such semantic analysis. However directly applying traditional topic models like LDA and PLSA to short texts can be problematic primarily due to the sparsity of data given the specificity of short texts. In this paper we are studying the usage of different models on a large scale data which can effectively infer hidden topics in big discussions taking these features into account.

\subsubsection{Prior Work}

Early studies of the problem of topic modeling on short texts mainly concerned the use of external knowledge to improve the representation of text data. For example, Phan and others \cite{HoriguchiPhanNguyen} used the modeling of those short texts based on the traditional topic model, the effectiveness of which was tested on a large-scale data set for the purpose of short texts classification. In these works, it is assumed that the use of data derived from long texts could help improving the model for short texts. However, these methods are effective only when the auxiliary data are closely related to the original data. In the case of texts, obtained from social networks this task is impossible due to most of user messages being self-contained. Another assumption is based on the use of different aggregation methods. In the case of data, obtained from the Twitter, user messages can be aggregated by authors, publication time and hashtags. The \cite{WrayLexingRishabh} shows that the best performing method is based on hashtag based aggregation, however it can’t be applied then the texts themselves are collected using a series of hashtags (which is one of better ways for collecting data on big events). Author based aggregation can also be unreliable since the majority of users will have very few messages on any given topic. In this paper we are focusing on models, relying on statistical information about the data.

\subsubsection{Topic Models}

\paragraph{LDA.} Latent Dirichlet Allocation is a three-level hierarchical Bayesian model in which each element of the collection is modeled as a finite distribution over the set of topics. Each topic, in turn, is modeled as an infinite mixture by the set of topic probabilities \cite{MichaelJohnDavidAndrew}.

Lda assumes the following generative process:
\begin{itemize}
	\item Choose \(\theta_i \sim \textit{Dir}(\alpha)\)
	\item Choose \(\phi_i \sim \textit{Dir}(\beta)\)
	\item For every word position \(i, j\):
	\begin{itemize}
		\item Choose a topic \(z_{i, j} \sim \textit{Multinomial}(\theta_i)\)
		\item Choose a word \(w_{i, j} \sim \textit{Multinomial}({\phi_z}_{i,j})\)
	\end{itemize}
\end{itemize}
The model parameters \(\alpha\) and \(\beta\) are typically chosen sparse for better performance on short texts. In this paper, LDA is used as a baseline model, the effectiveness of which for standard texts has been proven both theoretically and in many experimental results.

\paragraph{BTM.} Biterm Topic Model performs topic modeling task by modeling a set of biterms (unordered word pair cooccurring in a short context). The main idea is that if two words co-occur more frequently, they are more likely to belong to a same topic \cite{YanyanJiafengXueqi}.

Btm assumes the following generative process:
\begin{itemize}
	\item Draw \(\theta_i \sim \textit{Dir}(\alpha)\)
	\item For each topic \(k\):
	\begin{itemize}
		\item draw \(\phi_k \sim \textit{Multinomial}(\beta)\)
	\end{itemize}
	\item For each biterm \(b_i\):
	\begin{itemize}
		\item draw \(z_i \sim \textit{Multinomial}(\theta)\)
		\item draw \(w_{i,1},w_{i,2} \sim \textit{Multinomial}({\phi_z}_i)\)
	\end{itemize}
\end{itemize}

As BTM does not model documents explicitly, we must provide a way to infer the topics in a document, i.e., evaluating the topic posterior. Using the chain rule the following equation was obtained:
\begin{equation}
	\label{eqn:29}
	P(z \mid b) = \sum P(z \mid b_i) P(b_i \mid d)
\end{equation}

Where \(P (z \mid b_i)\) can be obtained using via Bayes’ formula based on the parameters learned in BTM and \(P(b_i \mid d)\) can be calculated using empirical distribution of words in a document.

\paragraph{WNTM.} Word Network Topic Model’s idea is based on the following observations. When the texts are short, the word-document space is very sparse, but the word-word space still contains a large number of non-zero elements. Since the topic distribution for each doc- ument can not be recognized accurately in short or unbalanced texts, instead WNTM uses the topic distribution for each word \cite{KeYuanJichang}. Therefore, WNTM studies the distribution by topics for words, rather than those for documents. Studying the topics of the word, rather than those of the document make WNTM less sensitive to the length of the document. In addition, a network of words can be built with any type of text, which makes the WNTM model simple and universal in real applications, unlike other models, such as the mixture of unigrams \cite{ThrunMitchellNigam} and BTM.

The generative process of the model is in many respects similar to that of the LDA, but due to the use of a different distribution it has its own features:
\begin{itemize}
	\item For every latent word group \(z\) choose \(\phi \sim Dir(\beta)\)
	\item Choose \(\vartheta_i \sim \textit{Dir}(\alpha)\) distribution of a latent word group for adjacent word list \(L_i\) for word \(w_i\)
	\item For every word \(w_j \in L_i\):
	\begin{itemize}
		\item Choose a latent word group \(z_j \sim \vartheta_i\) 
		\item Choose an adjacent word \(w_j \sim {\phi_z}_j\)
	\end{itemize}
\end{itemize}

\begin{figure}[ht]
	\centerfloat{
		\includegraphics[scale=1.0]{wntmGeneration}
	}
	\caption{Generation of word network for WNTM model \cite{KeYuanJichang}.}\label{fig:wntmGeneration}
\end{figure}

Similarly to BTM this model can’t be directly applied to get topic distributions over documents. To get the topics in the document, we assume that the proportions of the words generated by the document is equal to the proportions of the document topics, that is:
\begin{equation}
	\label{eqn:30}
	P(z \mid d) = \sum P(z \mid w_i) P(w_i \mid d)
\end{equation}
Where \(P(z \mid w_i\)) equal \(\vartheta_{i,z}\), obtained in the generative process of WNTM and \(P(w_i \mid d)\) obtained as an empirical distribution of words in documents.

\subsubsection{Experiment}

The experiment was conducted to evaluate the quality of these three models, on three sets of ad hoc discussions. The following coherence measures were used to test the effectiveness of each model:

\begin{itemize}
	\item UMass \cite{MimnoWallachTalley} measures how often a common word of each topic is in average a good predictor for a less common word.
	\item NPMI \cite{StevesonAletras} is a normalized version of pointwise mutual information.
\end{itemize}

\paragraph{Data sets.} In this work models were tested on data, collected on three ad hoc discussions from Twitter social network: Riots in Biryulevo (Russia), October 2013 \cite{BodrunovaLitvinenkoBlekanov}, Ferguson unrest (USA), August 2014 \cite{SmoliarovaBlekanovBodrunova} Charlie Hebdo shooting (France), January 2015 \cite{SmoliarovaBlekanovLitvinenko}. The data was crawled based on hashtags in user messages.

\textit{Byrulevo. Riots in Biryulevo}
\begin{itemize}
	\item Total number of user messages: 10215
	\item Total number of users participated in the discussion: 11429
	\item Surveyed time period: 1.10.2013 - 31.10.2013
	\item Number of users who published tweets in the period under consideration: 3574
\end{itemize}

\textit{Ferguson. Ferguson unrest}
\begin{itemize}
	\item Total number of user messages: 193812
	\item Total number of users participated in the discussion: 169677
	\item Surveyed time period: 22.08.2014 - 31.08.2014
	\item Number of users who published tweets in the period under consideration: 70018
\end{itemize}

Charlie Hebdo. Charlie Hebdo shooting
\begin{itemize}
	\item Total number of user messages: 505069
	\item Total number of users participated in the discussion: 952615
	\item Surveyed time period: 07.01.2015 - 10.01.2015
	\item Number of users who published tweets in the period under consideration: 238491
\end{itemize}


\begin{table}[ht]%
	\centering
	\caption{Topics for Byrulevo data set.}%
	\label{tab:byrulevoTopics}% label всегда желательно идти после caption
	%	\begin{adjustbox}{width=1\textwidth}
		%		\small
		\begin{tabular}{ c  c  c  c }% Вертикальные полосы не используются принципиально, как и лишние горизонтальные (допускается по ГОСТ 2.105 пункт 4.4.5) % @{} позволяет прижиматься к краям
			\toprule
			Topic 1 & Topic 2 & Topic 3 & Topic 4 \\
			\hline
			\multicolumn{4}{c}{\makecell{LDA}} \\
			migrant & warehouse & broadcast & riot  \\
			Zeynalov & work & live & Manezhka \\
			police & man & moscow & moscow \\
			murder & boutique & photo & migrant \\
			Sherbakov & moscow & find & block \\
			\hline
			\multicolumn{4}{c}{\makecell{WNTM}} \\
			Moscow & news & Moscow & police \\
			event & Sherbakov & OMON & authorities \\
			Russia & migrant & Zeynalov & killer \\
			riot & murder & arrest & russian \\
			mayhem & killer & Sherbakov & meetings \\
			\hline
			\multicolumn{4}{c}{\makecell{BTM}} \\
			citizen & OMON & Sherbakov &  russian \\
			police & warehouse & Zeynalov & government \\
			local & police & killer & riot \\
			riot & arrest & arrest & migrant \\
			Moscow & killer & moscow & news \\
			\bottomrule
		\end{tabular}%
		%	\end{adjustbox}
\end{table}

\begin{table}[ht]%
	\centering
	\caption{Topics for Charlie Hebdo data set.}%
	\label{tab:charlieTopics}% label всегда желательно идти после caption
	%	\begin{adjustbox}{width=1\textwidth}
		%		\small
		\begin{tabular}{ c  c  c  c }% Вертикальные полосы не используются принципиально, как и лишние горизонтальные (допускается по ГОСТ 2.105 пункт 4.4.5) % @{} позволяет прижиматься к краям
			\toprule
			Topic 1 & Topic 2 & Topic 3 & Topic 4 \\
			\hline
			\multicolumn{4}{c}{\makecell{LDA}} \\
			policia & die & police & islam \\
			Paris & satire & shooting & religion \\
			terroristas & cartoonist & attack & youngest  \\
			sospechosos & frankreich & suspects & local \\
			ataque & attentater & update & extrimists \\
			\hline
			\multicolumn{4}{c}{\makecell{WNTM}} \\
			attack & suspects & cartoonists & media\\
			french & police & support & cartoons  \\
			today & two & editor & toxic \\
			terror & attack & respond & image \\
			killed & hostage & journalism & caricatures \\
			\hline
			\multicolumn{4}{c}{\makecell{BTM}} \\
			police & french & victims & muslims\\
			suspects & shooting & solidarity & islam \\
			hostage & gunman & attack & must \\
			killed & dead & France & say\\
			breaking & killed & jesuischarlie & religion \\
			\bottomrule
		\end{tabular}%
		%	\end{adjustbox}
\end{table}

\begin{table}[ht]%
	\centering
	\caption{Topics for Ferguson data set.}%
	\label{tab:fergusonTopics}% label всегда желательно идти после caption
	%	\begin{adjustbox}{width=1\textwidth}
		%		\small
		\begin{tabular}{ c  c  c  c }% Вертикальные полосы не используются принципиально, как и лишние горизонтальные (допускается по ГОСТ 2.105 пункт 4.4.5) % @{} позволяет прижиматься к краям
			\toprule
			Topic 1 & Topic 2 & Topic 3 & Topic 4 \\
			\hline
			\multicolumn{4}{c}{\makecell{LDA}} \\
			police & MikeBrown & militarization & Miami  \\
			life & black & police & overtown \\
			surrender & brown & law & America \\
			must & racism & reason & vote  \\
			dissa & justice & end & jail \\
			\hline
			\multicolumn{4}{c}{\makecell{WNTM}} \\
			MikeBrown & CNN & must & movement  \\
			amp & cops & surrender & speak \\
			police & Times & police & join  \\
			black & black & dissa & support \\
			people & shooting & see & now\\
			\hline
			\multicolumn{4}{c}{\makecell{BTM}} \\
			must & join & community & pd\\
			police & movement & Miami & look \\
			surrender & now & support & closer  \\
			dise & speak & lot & MikebBrown  \\
			click & die & overtown & msnbc \\
			\bottomrule
		\end{tabular}%
		%	\end{adjustbox}
\end{table}

\paragraph{Results.} Coherence measures were calculated and plotted to estimate the models’ effectiveness on different number of topics. The results (Figures~\cref{fig:charlieCoherence,fig:byrulevoCoherence,fig:fergusonCoherence}) allow us to talk about the approximate number of topics that is optimal for this task, but due to imperfection of quality indicators it requires manual clarification.

\begin{figure}[ht]
	\centerfloat{
		\includegraphics[scale=1.0]{charlieCoherence}
	}
	\caption{Coherence scores on Charlie Hebdo data set.}\label{fig:charlieCoherence}
\end{figure}

\begin{figure}[ht]
	\centerfloat{
		\includegraphics[scale=1.0]{byrulevoCoherence}
	}
	\caption{Coherence scores on Byrulevo data set.}\label{fig:byrulevoCoherence}
\end{figure}

\begin{figure}[ht]
	\centerfloat{
		\includegraphics[scale=1.0]{fergusonCoherence}
	}
	\caption{Coherence scores on Ferguson data set.}\label{fig:fergusonCoherence}
\end{figure}

For each data set and topic model a few most coherent topics (topics focused on specific words which have the higher probability to occur in a specific topics) were chosen and can be seen in Tables~\cref{tab:byrulevoTopics,tab:charlieTopics,tab:fergusonTopics}.

As a result of the analysis of the given conflict ad hoc discussions, Biterm Topic Model came out to be the best performing and most stable method based on all coherence measures, baseline Lda model not specialized for working with short text has shown to not be suitable without the additional preprocessing of data mainly due to the data sparsity. However topics identified using this model can still be used to analyze key moments of discussions, participants’ arguments, and as a basis for other data analysis tasks.

Future work involves continuing the work on peer review, the first part of which was done for the Byrulevo data set by students and researchers from the journalism department of SPBU university. It has shown that the best results were obtained by the WNTM model. Comparing results of the models to the manual evaluation will allow us to measure not only the effectiveness of each model but also to compare coherence measures and find which one correlates better with human perception. The different prepocessing methods can also be used to improve the final result, tf-idf can be used instead of regular bag words approach to omit less frequent words more effectively.

\subsection{The ideal topic: interdependence of topic interpretability and other quality features in topic modelling for short texts}\label{subsec:ch5/sec2/sub3}

\subsubsection{1. Introduction}

The studies of large text corpora via automated detection of topicality are a growing area of social research. Topic modelling is a probabilistic tool for discovering non-evident topics in a collection of texts. After the potential of topic modelling for Twitter has been proven \cite{RamageDumaisLiebling}, it has become a separate and growing area within topic detection research. Several works have shown that topic modelling of Twitter data is suitable for both detection of hidden discussion topics and data pre-processing aimed at reducing the topical dimensionality of the studied corpora \cite{BlekanovTarasovMaksimov,KoltsovaKoltcov}, as well as for a range of comparative tasks. In both capacities, it has already been applied to practical tasks like drug use analysis \cite{JonnagaddalaJueDai} or assessment of popular behavior during natural disasters \cite{LigutomOrioRamacho,MacedaLlovidoPalaoag}.

But, despite the promise, topic modelling is highly problematic when applied to short texts \cite{MazaruraDeWaalKanifer}, especially of the oral-written nature \cite{Lutovinova} often met online, of which Twitter is exemplary. Short texts, including Twitter posts (tweets), are today a popular form of public speech online, but their nature creates substantial obstacles to modelling topicality in Twitter corpora of virtually any size.

Topic modelling algorithms for the English language on Twitter have been explored widely, even if not extensively. Several models beyond LDA with Gibbs sampling, such as vector models or pLSA, as well as extensions for LDA, have been proposed. In most cases, classic unsupervised LDA performed worse than extended LDA or other models like biterm topic modelling (BTM). But, in some cases, LDA is used unquestioned, with results interpretable enough; this is why we will compare LDA which is in wide use with two other models described below.

Till today, comparative works on Twitter topic models for various languages beyond English remain very rare, mostly focusing on German and Spanish. Comparative studies of topic modelling for the francophone segment of Twitter are truly scarce, and, for the Russian-language Twitter, except for the earlier pilot works by our working group \cite{RamageDumaisLiebling,BlekanovTarasovMaksimov,SmoliarovaBodrunovaYakunin}, are almost non-extant. For a rare comparative work on five languages including French and Russian. This work, similar to ours in its goals, tests three models of topic detection, including the authors’ own development, and tests them by human coders and objective measures. The paper reports the comparative results on topic coherence and the so-called ‘utility’ of topic clusters (not topics!) without explaining how utility relates to human interpretability (whether this actually is interpretability or not). The authors also claim that their own model performs better than both LDA and BTM, but they do not reveal its full details for that would allow for comparison of the algorithms. Moreover, the results are not presented in full but, e.g. for all the languages beyond English, only on LDA and the authors’ model. Due to these reasons, we will look to French and Russian to once again compare the quality of topic detection.

Another gap in the current research on Twitter topics detection is that, except for e-health and natural disasters, issue- or event-oriented discussions are practically not studied. The methodological papers, including \cite{Sridhar}, do a lot to avoid limitations to overall Twitter topicality and work with large data that contain all possible topics. But, within a case-oriented discussion, detection of topics must be seriously complicated by the fact that the inner sub-topics share very similar lexicons.

To partly cover these gaps, we have taken under scrutiny three conflictual cases in three different countries, as they have provoked the respective Twitter discussion of a large scale but at the same time were focused enough for our experiment. In our earlier research, we have applied unsupervised LDA, as well as WNTM and BTM, to these three cases (see below), with BTM showing the best results by automated quality assessment metrics, namely Umass and NPMI \cite{BlekanovTarasovMaksimov}. But here we need to re-discuss our conclusions and juxtapose them to other, more subjective quality metrics such as topic saliency and human interpretation. Also, we try to see whether topic saliency and topic interpretability are linked to each other; to our best knowledge, this idea has not yet been explored by the academic community.

The remainder of the paper is organized as follows. In Sect. 2, we discuss the topic modelling algorithms that we use and the results we have received by applying the automated quality metrics and human coding, with the conclusion to reject both baselines. In Sect. 3, we present our ideas on the two other measures of topic detection quality, namely topic saliency and topic robustness. Section 4 serves for describing the cases and formulating the hypotheses. Section 5 provides the description of methods. In Sect. 6, we present the results and discuss them. To conclude, we formulate our idea of the ‘ideal topic’.

\subsubsection{2. Human and Automated Evaluation of Quality: Rejection of Both Baselines}

\paragraph{2.1 Automated Quality Assessment}

For automated quality assessment, we have used the coherence metrics that have gained popularity in the recent academic research, namely Umass and normalized PMI (NPMI). As we have shown in our earlier pre-tests, by both metrics, we have seen BTM to be performing better than the two other algorithms \cite{BlekanovTarasovMaksimov}.

But, after preliminary evaluation by human coders, we have seen that there are problems in topic interpretation that are not captured by the topic coherence metrics. Thus, the problems that we have discovered may be summarized as follows: 1) seemingly low number of interpretable topics; 2) a relatively big number of topics too similar to each other; 3) one-retweet-based topics gathered around one popular tweet rather than around a real topic in the discussion; 4) chained topics that unite several micro-themes but are non-clear to the human coders.

\paragraph{2.2 Human Coding and Its Artifacts}

Interpretability is understood here as an ability of a human coder to say why the topic descriptors (the most relevant words that describe the topic) stay together in a given topic. This is a traditional and ultimate metric of topic quality, as human interpretability is the ultimate goal of the whole modelling procedure. But this metric is non-feasible if the number of topics in multiple runs of the algorithms is overwhelming.

Interpretability is really hard to measure in an objective way. Moreover, below we show that human coding is critically dependent on three types of knowledge: that on the case, that on method, and that on the dataset. The difference in preparation to coding critically alters the coding results.

In Table~\cref{tab:coderTrainingArtifacts}, we show the coding results for three pairs of coders: unexperienced coders who have no detailed knowledge on the case, unexperienced coders who have received detailed instructions on the case and a lot of background knowledge, and members of our working group who have knowledge on the case, the dataset, and the algorithms.

\begin{table}[ht]%
	\centering
	\caption{Human coding of the Russian dataset, 100 topics: the artifacts of coder training.}%
	\label{tab:coderTrainingArtifacts}% label всегда желательно идти после caption
	%	\begin{adjustbox}{width=1\textwidth}
		%		\small
		\begin{tabular}{ c  c  c  c  c  c  c }% Вертикальные полосы не используются принципиально, как и лишние горизонтальные (допускается по ГОСТ 2.105 пункт 4.4.5) % @{} позволяет прижиматься к краям
			\toprule
			 & \multicolumn{2}{c}{\makecell{Unexperienced}} & \multicolumn{2}{c}{\makecell{Trained}} & \multicolumn{2}{c}{\makecell{Experienced}}\\
			\cline{2-7}
			& Coder 1 & Coder 2 & Coder 1 & Coder 2 & Coder 1 & Coder 2\\
			\hline
			LDA & 33\% & 11\% & 65\% & 94\% & 53\% & 63\% \\
			WNTM & 12\% & 11\% & 75\% & 92\% & 76\% & 73\%\\
			BTM & 22\% & 8\% & 88\% & 86\% & 76\% & 83\%\\
			\bottomrule
		\end{tabular}%
		%	\end{adjustbox}
\end{table}

As we see from Table~\cref{tab:coderTrainingArtifacts}, the coding result is, indeed, highly skewed by the coder background knowledge. Also, experienced coders from the working group show the tendency to give out a more balanced coding output, as they are more strict on the method performance and do not imply the inner grammar relations to the topics, filtering interpretability by the knowledge on the internal discourses in the dataset.

Thus, we see both automated dataset-level and human-based topic-level metrics to be unreliable in terms of finding ‘real topics’. This is why we suggest two other metrics and test their interconnectedness with human interpretability, to be able in future to substitute human coding by measuring these metrics.

\subsubsection{3. Topic-Level Quality Assessment: In Search of Independence Form Human Eye}

Thus, we have decided to introduce topic-level metrics of quality of topic modelling and test whether human interpretability will be related to them; we do it in an attempt to find topic-level metrics to substitute human coding without the necessity to assess the quality of the algorithm. At the same time, we still consider the well-known \textit{tf-idf} metric relevant for topic quality detection and will test our metrics against it in future.

\paragraph{3.1 Topic Saliency}

The idea of topic saliency lies in the fact that, for various time slots within a discussion, different sub-themes may emerge as the leading ones. Saliency is calculated in the following way. Each document in the dataset is assigned a topic distribution; then, based on the time of publication, the topic distributions in the tweets are summarized, and, for each topic, the saliency is calculated as the sum of the tweets belonging (with a certain probability) to the particular topic. Since the intensity of the discussion may vary in time, the overall topic saliency for all topics may vary (from 0 to 1, in our case). We have used a 24-h step and measured the topic saliency for each day in aggregate. This strategy constitutes one of the limitations for our study, as the French dataset encompasses four days only, due to its overwhelming volume (see below); in future, for the 4-day French dataset, we plan to change the step and conduct an hourly analysis.

To our best knowledge, topic saliency has not yet been discussed in the literature as a possible quality metric for topic modelling. But, by commonsense logic, topic saliency should be related to topic interpretability: the more salient a topic is, the bigger it is within the dataset and must be easier to interpret. And, if so, topic saliency may in future be seen a quality metric -- a proxy for understanding topics by human coders.

\paragraph{3.2 Topic Robustness}

Topic robustness shows whether the topic is composed of a relatively large number of relevant texts/words. We will look at the relevance levels of topic descriptors (top words). The relevance is, basically, the extent to which a particular word belongs to a given topic, depending on how many times it is met in the topic-relevant tweets; the relevance is measured 0 to 1. There is no clear agreement in today’s research what relevance is to be considered high; this metric is dataset-dependent and thus needs to be assessed for each topic depending on the highest levels for a given run of topic modelling. Usual values for word relevance range from 0.01 to 0.1 for short-text small datasets to 0.001 to 0.05 for longer-text larger datasets.

Topic robustness, thus, combines two internal parameters: the word relevance to a topic and the number of relevant words in a topic. Robust topics are those who have many highly relevant words; it means they attract a lot of similar texts. We will provide the exact measurements for topic robustness in Sect. 5. But topics with lower robustness must not be viewed as ‘bad’: robustness highly depends on the nature of the discussion.
We will check the relationships between these metrics, to see how one can describe an ideal topic: a robust, understandable, salient one -- and assess which datasets (bigger/smaller, monolingual/multilingual) provide for a bigger quantity of ideal topics.

\paragraph{4. The Cases Under Scrutiny and the Research Hypotheses}

Here, we describe the cases we work with, which will help formulate the hypotheses. For this research, we work with the datasets from Twitter on the three cases:

\begin{itemize}
	\item Anti-immigrant riots in the Moscow district of Biryulevo, Russia, 2013: number of users who published tweets -- 3574, total number of user messages -- 10215;
	\item The Ferguson unrest, USA, 2014: number of users who published tweets -- 70018, total number of messages -- 193812;
	\item The Charlie Hebdo shooting, France, 2015: number of users who published tweets -- 238491, total number of messages -- 505069.
\end{itemize}

The datasets are highly uneven in volume, and thus we had conducted pre-tests of the number of topics that would provide the most interpretable results and, at the same time, remain comparable. We ran the topic modelling for all the three cases with the BTM algorithm for the following number of topics set arbitrarily: for Biryulevo, 10, 50, and 100; for Ferguson, 50, 100, and 200; for Charlie Hebdo, 100, 200, and 400 topics. We have pre-tested their interpretability for small numbers of topics for each run (10 to 20), and we have seen that, for 100 topics, the results may be comparable. This is why we will further on use the 100-topic runs for comparing the quality assessment results.

Having in mind the aforementioned metrics for quality assessment and the number of topics in each run, we have set the following hypotheses:
\begin{itemize}
	\item H1. Human-based interpretability will correspond to the volume of the dataset: the bigger the dataset, the more interpretable the topics are.
	\item H2. Human-based interpretability will be higher for monolingual discussions (Russia, the USA) than for multilingual discussions (France).
	\item H3. More interpretable topics will have higher saliency in all the datasets. 
	\item H4. More interpretable topics will show higher robustness in all the datasets. 
	\item H5. More robust topics will have higher saliency in all the datasets.
\end{itemize}

We note that H1 and H2 are mutually exclusive.

\subsubsection{5. The Research Methods}

Here, we will in short describe the methods we use to test the hypotheses. As tour data collection has been well-described in our previous works, we will only provide the description for the methods directly used in this paper.

For H1 and H2, we have worked with native speakers as human coders. Unlike in other research, we do not test the inter-rater reliability for the coders and then let them code separate segments of the task. Instead, two coders were asked to code the topic descriptors independently, and then, their inter-coder reliability was checked. We used this to see how divergent the opinions of the coders could be and whether the number of interpretable topics corresponds to earlier studies of longer texts, e.g. Russian blogs \cite{KoltsovaKoltcov}. Thus, six coders were instructed to code the topics as interpretable or non-interpretable based on the abovestated assumption of comprehensibility of why the topic descriptors stay together in a particular topic. Then, the number of interpretable topics for each coder and the meta-coding results were calculated. Each topic was assigned an interpretability index of 0 (non-interpretable for both coders), 1 (interpretable for one coder), or 2 (interpretable for both coders).

In case of France, multi-lingual coders had to work (those who could recognize English, French, Spanish, Italian, German, and, with additional help, Russian). Finding such coders constitutes a separate problem in assessing the quality of multilingual global- scale discussions. We consider the procedures of assignment of meaning comparable for Russia, the USA, and France, as interpreting the topic descriptors requires simple word recognition procedures and basic knowledge of language(s), which makes recognition of a top word the unit of interpretation, disregarding the language belonging of a word.

For H3, topic saliency was automatically calculated based on the modelling results, and the saliency thresholds were defined based on the overall saliency picture for a given case. But, as the saliency measurements have shown, topics could be considered salient if they reached circa 30\% of the overall saliency of topics in a given day. Then each topic was assigned a saliency index of 0 (the topic has never reached over 30\% of the overall saliency for a particular time slot), 1 (the topic has at least once reached over 50\% of saliency for a particular time slot), or 2 (the topic stably reached over 50\% of saliency in at least 30\% of the overall time span for the case). Then, Spearman’s rho was used to see the dependencies between topic interpretability and topic saliency.

For H4 and H5, we have calculated the topic robustness for each topic. First, we have calculated the word relevance and have established 0.02 as the high relevance threshold. Then, we have introduced the robustness score: if a topic has no top words with the relevance reaching 0.02, it is non-robust (0); if there are up to 5 words of the relevance of 0.02 or higher, the topic is acceptably robust (1); if there are more than 5 words with the relevance of 0.02 or higher, the topic is fully robust (2).

\subsubsection{6. Results and Discussion}

\textit{H1.} We have presupposed that the datasets with a bigger amount of tweets will provide for better topic extraction, if the number of topics is the same. But H1 has to be rejected, as we have found that the number of highly interpretable topics was almost the same for all the datasets: 45 of 99, for Russia; 40 of 96, for the USA; and 41 of 100, for France (topics on non-Russian and non-English were eliminated from the first two cases). Thus, we do not observe any growth of interpretability for bigger tweet collections.

\textit{H2.} The same goes for H2: we do not observe higher interpretability for monolingual cases, given that the coders understand the languages of the multilingual discussion.

Rejection of H1 and H2 is telling, as it provides input for understanding the nature of topic interpretability on Twitter using BTM. Only circa 40--45\% of the topics get surely interpreted, which needs to be addressed in the future research. In the French case, most non-interpretable topics were composed by tweets in different languages, first and foremost French, English, and Spanish. The algorithm that puts the tweets in different languages into one topic is still to be analysed, to prevent the mixing of languages in future. Another problem discovered both in the ‘international’ topics and in the francophone one is the over-abundance of pronouns, service words and particles as top words. In the biterm-based approaches, eliminating them from the pre-processed dataset would significantly change the clustering results; but maybe a decision here is simple -- they should not be allowed to the lists of topic descriptors, thus leaving more space to the meaningful top words.

\textit{H3--H5.} For the results for H3 to H5, see Table~\cref{tab:topicQualityMetrics}.

\begin{table}[ht]%
	\centering
	\caption{Human coding of the Russian dataset, 100 topics: the artifacts of coder training.}%
	\label{tab:topicQualityMetrics}% label всегда желательно идти после caption
		\begin{adjustbox}{width=1\textwidth}
				\small
		\begin{tabular}{ c  c  c  c  c  c  c  c  c  c }% Вертикальные полосы не используются принципиально, как и лишние горизонтальные (допускается по ГОСТ 2.105 пункт 4.4.5) % @{} позволяет прижиматься к краям
			\toprule
			& \multicolumn{3}{c}{\makecell{Russia}} & \multicolumn{3}{c}{\makecell{The USA}} & \multicolumn{3}{c}{\makecell{France}}\\
			\cline{2-10}
			& Interpr. & Saliency & Rob. & Interpr. & Saliency & Rob. & Interpr. & Saliency & Rob. \\
			\hline
			Interpretability & -- & \(-0,080\) & 0,064 & -- & 0,226* & 0,192 & -- & 0,337*** & 0,396*** \\
			Saliency & \(-0,080\) & -- & 0,261** & 0,226* & -- & \(-0,086\) & 0,337*** & -- & 0,345*** \\
			Robustness & 0,064 & 0,261** & -- & 0,192 & \(-0,086\) & -- & 0,396*** & 0,345*** & -- \\
			\hline
			\multicolumn{10}{c}{\makecell{Note. * -- \(p \le 0,05\); ** -- \(p \le 0,01\); *** -- \(p \le 0,001\).}}\\
			\bottomrule
		\end{tabular}%
			\end{adjustbox}
\end{table}

Table~\cref{tab:topicQualityMetrics} does not provide for any systematic picture of the interdependence of the three topic quality metrics. Our hypotheses were formulated the way that they demanded a clear picture, which is not always the case in practice. In the way they are formulated, they have to be rejected; but important conclusions can be drawn.

Thus, more interpretable topics are more salient in the US and French cases. This might have to do with the dataset volume. And the French case, despite its multilingualism, shows that the three aspects of the topic detection quality may be interdependent.

The correlations for the French case are stronger than for the monolingual cases, and this definitely demands future research on how multilingual discussions are constructed in terms of topicality and why the multilingual discussions perform better than monolingual ones, which is quite counter-intuitive.

To conclude, we underline the following. Putting our results against previous research, we need to state that interpretability for Twitter is significantly lower than for longer-text datasets (cf. \cite{KoltsovaKoltcov}). Also, looking at topic/dataset characteristics like robustness and saliency may provide for better understanding of the nature of an ideal topic. In all 300 topics we have assessed, only 7 were all robust, salient, and well-interpretable.

\subsection{Topics in the Russian Twitter and relations between their interpretability and sentiment}\label{subsec:ch5/sec2/sub4}

\subsubsection{1. Introduction}

After the principal possibility of topic modelling on such short and noisy texts as tweets has been demonstrated over ten years ago \cite{RamageDumaisLiebling}, it has become a focus of scholarly attention as a method of detection of hidden substance of the tweet corpora, tracking user behavior, and pre-processing to reducing the dimensionality of the corpora. For English, an extensive number of studies have been made, and several algorithms have been proposed. For other languages, though, there have been much fewer experiments, especially of comparative thought.

Despite, by the nature of topic modelling, short and noisy natural-language texts are virtually the worst object for topic detection, and despite the diminishing relevance of Twitter in the media ecosystems of many countries, the attractiveness of Twitter as a source of datasets for modelling remains high -- again, mostly for the English-language cases.

Russian social media have so far received moderate attention in terms of experiments with topic modelling, and the Russian Twitter has not been its primary focus. In particular, topic interpretability has not received, to our viewpoint, enough scholarly attention. Despite topic quality is quite an expansive research area (for a review, see \cite{MavrinFilchenkovKoltcov}), only several studies have addressed the divergence between the objective and subjective topic quality assessment, especially rarely for Russian \cite{BodrunovaKoltcovKoltsova}. The causes of low interpretability of topics for short texts have virtually escaped proper discussion, and no studies have been dedicated to linking topic semantics and their interpretability.

Our work aims at partly covering the gaps described above. We model the topics for a dataset of a conflictual discussion of as early as 2013, as at that moment, according to our previous studies \cite{BodrunovaLitvinenkoBlekanov2017}, the Russian Twitter was not as bot-occupied as it has been since 2016 \cite{StukalSanovichBonneau}. We have chosen to work with conflicts online, as they provide for emotionally loaded texts and link the discussions to a variety of wider issues of user interest.

Earlier, we have tested three topic modelling algorithms, namely unsupervised LDA, BTM, and WNTM, on three raw-data Twitter datasets in English, French, and Russian, including the aforementioned dataset \cite{BlekanovTarasovMaksimov}. But, as it often happens in topic modelling studies, the results must be tested also by human coders. Below, we will describe in short what and why we have done in terms of algorithm development and what we have tested additionally to provide more freedom to human coders.

The remainder of the paper is organized as follows. In Section 2, we review the literature on topic modelling and Twitter studies in Russia and beyond. In Section 3, we formulate the hypotheses, tell of the case under scrutiny, our previous topic modelling experience with the three algorithms, and suggest data amplification for human coders. In Section 4, we describe the method and the research procedures. Section 5 provides the results and discusses them for both the process of human coding and the results based on sentiment analysis. The concluding remarks reformulate our findings against prior knowledge and provide the guidelines for further research.

\subsubsection{2. Topic Modelling for Russian Social Media: A Literature Review}

To our best knowledge, there has so far been no extensive review of how topic modeling has developed for the Russian language. Not aspiring for a truly representative one, we anyway need to review the field to show how scarce attention has so far been given to the relations between the interpretability of modelling results and the topics’ inherent features. This, at least partly, is also true for topic modelling studies for other languages.

For Russian, topic modelling studies may be divided into methodological (that develop, compare, and extend models as well as evaluate their quality), applied (that apply topic modelling to extract the meanings from datasets), and relational (that relate topic modelling results to other features of the datasets or external factors).

\paragraph{A. Methodological studies of topic modelling} 
There are several groups within Russia who have been focusing on various topic modelling algorithms. Thus, in a sequence of influential works, Koltsova and colleagues develop LDA \cite{KoltcovKoltsovaNikolenko2014} and a range of extensions to it. The latter include semi-supervised interval LDA, or ISLDA \cite{BodrunovaKoltcovKoltsova}, granulated LDA \cite{KoltcovNikolenkoKoltsova0516} and LDA with local density regularization \cite{KoltcovNikolenkoKoltsova0916}. In relation to this work, Koltsov and colleagues have optimized Gibbs sampling \cite{KoltcovNikolenkoKoltsova16} and the model itself applying Rényi and Tsallis entropies \cite{Koltcov}. This group of authors has linked the use of topic modelling to qualitative studies of social media and beyond \cite{NikolenkoKoltcovKoltsova} and applied their instruments to several corpora of longer and shorter texts (see below). In these works, Russian is used more as an example of ‘a language as such’, just as English is used in topic modeling, often without discussing inherent linguistic or contextual limitations. Similarly, the works by Vorontsov and colleagues (e.g. \cite{VorontsovFreiApishev1015,VorontsovPotapenko,VorontsovFreiApishev0415}) have been influential in pLSA and its modifications based on non- Bayesian regularization. Among the rest, they have tested two algorithms for the Russian-language short texts, namely biterm topic modelling (BTM) and word network topic model (WNTM) \cite[p.~191]{KochedykovApishevGolitsyn} which we have also employed for our Twitter study (see \cite{BlekanovTarasovMaksimov} and below). The two research groups have collaborated on additive topic models \cite{ApishevKoltcovKoltsova1016,ApishevKoltcovKoltsova16} and have published important methodological papers in Russian (which we do not review here due to the scarcity of space). We will also omit from our review the works that suggest text clustering instruments alternative to topic modelling.

Several other authors have amplified this corpus of works by adding automated labeling to Russian-language topics \cite{MirzagitovaMitrofanova}, showing the possibility of single term extraction by topic modelling \cite{BolshakovaLoukachevitchNokel}, and multi-model tests on optimization of the number of topics \cite{KrasnovSen}.

\paragraph{B. Use of topic modelling for content interpretation: applied and relational works}
Content-exploring research has scrutinized both social media and text corpora beyond them. Thus, topic modelling has been employed to map the agenda of the Russian Livejournal of 2013 \cite{KoltsovaKoltcov} and ethnic contents of the Russian blogs \cite{Nagornyy}, including detection of most hated ethnicities \cite{BodrunovaKoltsovaKoltcov}.

Beyond the social networking realm, LDA has been applied to TV topics \cite{KoltsovaPashakhin}, Russian prose \cite{SedovaMitrofanova}, a corpus of musicological texts \cite{Mitrofanova}, newspaper news texts on climate change \cite{BoussalisCoanPoberezhskaya}, and, along with BTM, to Q\&A queries \cite{VolskeBraslavskiHagen}, with reasonable success.

The works that, above, we have called ‘relational’ are almost exclusively written by Koltsova and colleagues. They have shown how activity in political blogs in 2011-2012 correlated with the politicians’ ratings \cite{KoltsovaShcherbak}; the topicality of Livejournal top blogs corresponded to the structure of the Livejournal co-commenting communities \cite{KoltsovaKoltcovNikolenko}; and the linkage between news topics on TV and user feedback \cite{KoltsovPashakhinDokuka}, among other works.

\paragraph{C. Topic modelling for the Russian Twitter}
There are three methodological works except ours that explore topic modelling for the Russian Twitter \cite{MimnoWallachNaradowsky,Sridhar,GutierrezShutovaLichtenstein}. They all use Russian datasets for developing multilingual modelling tools. The first two do not discuss individual results for any single language, and the third only observes one difference in description of sports between Russian- and English-language Twitter. Both works assess the quality of the models but lack the discussion on human interpretability of the topics.

Similarly, only a small handful of works applies topic modelling to Russian Twitter to detect substantial meanings or discussion features. Thus, one work \cite{ChewTurnley} has shown the divergence between Russian- and English-language ‘master narratives’ on Russian cyber-operations.

Our works appear to the only continuous effort (since 2013) to apply automated text analysis to the Russian Twitter. Thus, we have tested three topic models \cite{BlekanovTarasovMaksimov} (see the details below) and have also applied BTM to detect the dynamics of topicality in conflictual discussions \cite{SmoliarovaBodrunovaYakunin}.

But, so far, there has been no discussion on whether topic models work well for the Russian Twitter in terms of human interpretability and what would be the features that could help in raising it.

\paragraph{D. Quality assessment and interpretability of the Russian-language topics}
An extensive study of nine automated metrics juxtaposed to the human-coding baseline was performed in \cite{Nikolenko}. Based on word2vec approaches suggested earlier, the author shows that normalized PMI (NPMI) outperforms PMI as well as other conventional metrics like tf-idf, and that vector-based metrics work better than all others. But, for short texts, the question remains how wor2vec metrics work with pooled short texts (what is most often done for tweets), this is why we use NPMI, as well as Umass \cite{BodrunovaKoltcovKoltsova}. Another important attempt to introduce a quality metric for longer texts was made in \cite{MavrinFilchenkovKoltcov}.

But none of these works has primarily focused on the issues in human (non-)interpretability of the topics, mostly seeing human coding as a baseline -- perhaps because, for longer texts, when interpretability was at stake, the models performed well enough. Thus, the authors \cite{KoltsovaKoltcov} have shown that, for long texts like Livejournal posts, circa two-thirds of the topics are interpretable after LDA has been applied. They have also identified three types of uninterpretable topics: ‘language’ (other than Russian), ‘style’ (writing styles, including offensive language), and ‘noise’ (uninterpretable texts / combinations of texts) \cite[p.~218]{KoltsovaKoltcov}.

In our pilot studies, though, we have seen that topics for Twitter are less interpretable, and we explore the coders’ experience below. The general features of Russian as a highly inflected language are well-known \cite{Whittaker}; for such languages, several solutions have been suggested \cite{MaucecKacicHorvat}, but it is not the nature of Russian alone that seems to be causing lower topic interpretability in the case of Russian Twitter.

Also, no studies have been done to understand the features of the topics that could raise the topic interpretability. In other languages, first and foremost English, relations between sentiment and topicality of discussions in social media have been explored in a range of works (for a short review, see \cite{XiangZhou}). But, so far, as soon as we know, there have been no works that would explore topic sentiment in relation to human interpretability. Our experiments with sentiment on Twitter in three languages \cite{BlekanovKukarkinMaksimov,BodrunovaBlekanovKukarkin2018}, including Russian \cite{BodrunovaLitvinenkoBlekanov2017,NigmatullinaBodrunova}, show that sentiment of tweets is tightly linked to topicality and user status as influencer. Our intuition is that more interpretable topics are more sentiment-loaded, in particular -- negativity-loaded.

\subsubsection{3. Our Hypotheses and the Case Under Scrutiny}

\paragraph{A. The research hypotheses}
Previously, we have applied three topic modelling algorithms to three datasets from Twitter, all featuring nation- or world-scale conflictual discussions \cite{BlekanovTarasovMaksimov}. We have tested the results with two topic coherence metrics: Umass (as suggested by \cite{MimnoWallachTalley}) and NPMI (as suggested by \cite{StevesonAletras}). There, we have shown, in particular, that, for all the three cases that: 1) BTM outperforms LDA and WNTM for our data; 2) coherence metrics show that topics are extracted well; 3) that the best number of topics for our datasets is 100.

But, as we have stated above, later pilot tests of the models on human coders have demonstrated that the model runs labeled as good by automated metrics do not perform that well when read by humans. Thus, for this study, we turn the testing logic upside down and question the ‘human baseline’ approach, to discuss the shortcomings of the algorithms noted by the human coders. For this, we amplify our coding interface a word-to-topic relevance metric \(\lambda\) suggested by \cite{SievertShirley} and check its application upon a small number of topics with human coders. ‘Playing’ with this function allows to lower the most relevant (and, thus, most frequent) terms among the topic descriptors and bring more specific terms up.

We also link human interpretability to sentiment load of the topic descriptors (30 words).

Thus, we hypothesize that:
\begin{itemize}
	\item H1. For all the three algorithms, lowering \(\lambda\) leads to higher interpretability.
	\item H2. With the best relevance level, BTM outperforms LDA and WNTM in human coding, just as it does by objective metrics.
	\item H3. More interpretable topics are more negativity-loaded.
	\item H4. The correlation between human interpretability and sentiment will be higher for the model that performs best by automated quality metrics.
	\item H5. The correlation between human interpretability and sentiment will be higher for the model that performs best by human coders’ assessment.
\end{itemize}

The two last hypotheses become mutually exclusive if H2 is not proven.

\paragraph{B. The case under scrutiny} 
The case that we work with is a month-long discussion that dates back to October 2013 and describes a conflict between locals and re-settlers from Central Asia in the Moscow district of Biryulevo. The discussion unfolded after a killing of a Russian Egor Scherbakov by an Uzbek re-settler Orham Zeinalov and opened up a nationwide discussion on migration issues like illegal dwelling of immigrants, visas for immigrant workers, cultural differences, and ‘radical white’ Russia-for-Russians opposition to re-settling. The Twitter part of this discussion has immediately reached national trending topics and has sparked participation of local authorities, police special troops, diasporas, international commentators, and NGOs.

The collected corpus of tweets is not very large; it covers October 1 to 31, 2013, and, after cleaning of spam, tweets in languages other than Russian, and links-only tweets, includes 10215 tweets from 3574 users who published them.

It suits our goals by several reasons. First, its dating back to 2013 ensures that it is almost bot-free. Second, it has high conflicting potential, and its tweets are highly emotionally loaded. Third, half of our coders have previously worked with this dataset and could link the issues they experienced in coding to the nature of the discussion already studied in other aspects.

\subsubsection{4. The Research Methods and Conduct of the Study}

For this part of the study, our methods look simple, but we still consider this part of our research crucial for further development. We will describe the methods and the conduct of the research together, to be able to comment on methods use.

\begin{enumerate}
	\item To select the proper meaning of word relevance for further assessment, we have tested the relevance metric \(\lambda\) for the three algorithms on a minor number of topics (30 for each algorithm).
	\item We have hand-coded the 100-topic runs for the three algorithms. The coders were not familiar with the dataset and were just explained the task and provided some knowledge on the details of the case; the trainers belonged to the group that was familiar with the data. Then, for each topic, we have calculated the interpretability score \(K\) of 0 (both coders could not interpret it), 1 (one coder managed to interpret it), or 2 (both coders interpreted it) and checked the Cohen’s Kappa for inter-rater reliability.
	\item We applied our sentiment assigning algorithm \cite{BlekanovKukarkinMaksimov} to the topic descriptors and hand-corrected it. What we were focusing upon was negativity, very characteristic for this discussion. This is why we coded the sentiment in the following way: positive, neutral, mixed, and undefined terms = 0, negative terms (by dictionary) = 1, hate speech including harsh pejoratives and obscene lexicon = 2. For each topic, we have calculated the descriptor negativity score \((N)\), dividing the negativity score by 30 (the number of word descriptors assessed).
	\item We calculated Spearman’s rho for each algorithm juxtaposing the interpretability score and the negativity indices for the topics, to see if there is any correlation between them.
\end{enumerate}

We need to state that this stage of research is preliminary, and our work in progress involves identical tasks for the datasets on similar inter-ethnic conflicts in English, German, and French.

\subsubsection{5. The Results and Their Discussion}

\paragraph{A. H1: term relevance and interpretability}
H1. We have pre-tested \(\lambda\) for our data, to be able to choose its proper meaning for further coding. The results are shown in Table~\cref{tab:threeAlgorithmRelevance}. We have conducted the tests with \(\lambda\) changing from 1 to 0.4 with the step of 0.1, but show only the results for \(\lambda = 1\) and \(\lambda = 0.5\), as 0.5 has behaved as a threshold changing the word descriptors picture the most. Codings for various meanings \(\lambda\) were made with significant time intervals (5 to 8 days), to avoid results skewing.

\begin{table}[ht]%
	\centering
	\caption{Pre-testing the relevance metric \(\lambda\) for the three algorithms, \% of 30 topics.}%
	\label{tab:threeAlgorithmRelevance}% label всегда желательно идти после caption
	%	\begin{adjustbox}{width=1\textwidth}
		%		\small
		\begin{tabular}{ c  c  c  c  c }% Вертикальные полосы не используются принципиально, как и лишние горизонтальные (допускается по ГОСТ 2.105 пункт 4.4.5) % @{} позволяет прижиматься к краям
			\toprule
			\makecell[c]{Model} & \multicolumn{2}{c}{\makecell{\(\lambda = 1\)}} & \multicolumn{2}{c}{\makecell{\(\lambda = 0.5\)}} \\
			\cline{2-5}
			 & Coder 1 & Coder 2 & Coder 1 & Coder 2 \\
			 \hline
			 LDA & 53\% & 63\% & 40\% & 33\% \\
			 WNTM & 76\% & 73\% & 70\% & 67\%\\
			 BTM & 76\% & 83\% & 69\% & 75\% \\
			\bottomrule
		\end{tabular}%
		%	\end{adjustbox}
\end{table}

As Table~\cref{tab:threeAlgorithmRelevance} shows, H1 proves wrong for all the models. The general line in interpretability is that it either gradually declines with lowering  \(\lambda\) or keeps steady and then sharply drops at circa 0.55 to 0.45. This is perfectly logical if we consider the idea of word relevance itself, but it does not correspond to how the instrument is intended to work. In most cases, well-interpretable topics remained in place; only rarely, one or another coder found some dubious topics more interpretable than before, but mostly, the results were negative, For some good topics, unfortunately, the most frequent words remained in place, while the words \#10 to \#20 changed, completely altering the meaning of the topic to the worse in terms of interpretability. We need more experiments with bigger number of topics and coders, to detect the optimal  \(\lambda\) that would help avoid the shifts of meaning and at the same time eliminate the frequent terms. So far, we will use \(\lambda = 1\) for further coding, despite a relatively big number of frequent terms in the top of the topic descriptors. For a number of topics in WNTM, though, we have seen significant improvement in interpretability where the term relevance change worked exactly as expected.

\paragraph{B. H2: interpretability of the topics with \(\lambda = 1\)}
The aggregated results of the human coding of the topics are represented in Table~\cref{tab:topicHumanInterpretability}.

\begin{table}[ht]%
	\centering
	\caption{Human interpretability of the twitter topics, \% of 100 topics.}%
	\label{tab:topicHumanInterpretability}% label всегда желательно идти после caption
	%	\begin{adjustbox}{width=1\textwidth}
		%		\small
		\begin{tabular}{ c  c  c  c  c }% Вертикальные полосы не используются принципиально, как и лишние горизонтальные (допускается по ГОСТ 2.105 пункт 4.4.5) % @{} позволяет прижиматься к краям
			\toprule
			Model & Coder1 & Coder2 & Metacoding \((K = 2)\) &  Cohen’s Kappa \\
			\hline
			LDA & 66\% & 52\% & 46\% & 60\%; 0,109 \\
			WNTM & 66\% & 56\% & 34\% & 64\%; 0,277 \\
			BTM & 68\% & 53\% & 43\% & 72\%; 0,437 \\
			\bottomrule
		\end{tabular}%
		%	\end{adjustbox}
\end{table}

As we see from Table~\cref{tab:topicHumanInterpretability}, the difference between the coders in terms of the number of the interpretable topics was not that big (14\% of topics maximum). Coder 1 has shown the levels of interpretability similar to those reported for longer texts \cite{KoltsovaKoltcov}, while Coder 2 could not reach 60\% with any of the algorithms.

For the number of well-interpretable topics, H2 must be rejected as well, as LDA outperformed both WNTM and BTM in \(K = 2\). But, at the same time, BTM shows much higher inter-coder agreement, and this makes us conclude that, for H2, we have mixed evidence. For BTM and WNTM, there were more topics that were filtered out by both coders (14 for LDA; 27 for WNTM; 33 for BTM, 5 of which by language). This makes us pose a question whether high number of well-interpretable topics must always be the goal of topic modelling, especially for short texts. Instead, a smaller number of ‘ideal’ (robust, salient, and well- interpretable) topics might be sought for.

\paragraph{C. Qualitative topic assessment: interpretability issues}
Before going any further, we feel necessary to describe the difficulties that the coders have encountered. Some of them fall into the categories of topic fallacies described previously for LDA \cite{BoydGraberMimnoNewman}. Among them, there are general words, identical topics and mixed/chained topics that involve several chains of semantically related terms.

But, for Twitter, these problems have a specific face. Thus, for an issue-based discussion composed of short texts, even well-interpretable topics are problematic. They seem to be mostly based on the corpus of frequent terms of the discussion, like ‘migrant’, ‘situation’, ‘power’, ‘killer’, ‘detained’ etc., and some additional topic descriptors attracted to them by the model. The frequent lexicon is responsible for the coder’s feeling of comprehensibility, while the attracted lexicon contributes to the new meaning. But altogether this creates the sense of high similarity of the topics; the coders have complained on similarity of multiple topics rather than on non-comprehension. Also, the frequent terms blur the meaning of topics. Here, we see a possible decision in defining topic clusters (sub-themes) instead of individual topics, based on topic similarity metrics such as Jaccard or Kullback-Leibler indices.

Another issue is that, often enough, the topics seem to be formed by one popular (highly retweeted) tweet. Such tweets, allegedly, drag the users’ pools to each other and start the topics, further attracting to the topics the tweets of the respective users, rather than the topically similar tweets. Here, the problem might lie in the pooling technique, as we have used all tweets by a given user as ‘one text’. To resolve this issue, pooling vs. non-pooling must be tested.

And yet another problem lies in the case specificity. Working with a particular event- or issue-based discussion demands a certain level of contextual knowledge, without which one cannot assign meanings to the topic descriptors. Thus, during pre-tests, the members of our working group could interpret more topics than the coders, as the nature of the topics was details-dependent. This poses a wider problem of how the research community should treat interpretability of a topic in general. For longer texts from non-case-specific corpora, topics appear to be more distinguishable from each other, as they, indeed, relate to very different themes in users’ online speech. But defining sub-topics within a highly thematized corpus demands a different understanding of topic interpretability, coherence, and distance. Also, we recommend that, for coding thematic discussions, the coders become familiar with the theme and its context, and not by preliminary reading but from other sources.

\paragraph{D. Negativity in topics and its relation to interpretability}
To test H3, we have conducted Spearman’s test between topic interpretability score \((K)\) and topic negativity score \((N)\). For all the three algorithms, H3 is proven, with relatively high Spearman’s correlation values: for LDA, 0,375**; for WNTM, 0,383**; for BTM, 0,494. This proves H3 and demands further research on whether negatively loaded topics show higher interpretability for datasets of other nature -- non-thematic, non-conflictual, platform-diverse.

We also see that the Spearman’s value is the highest for BTM which proves our H4. It also proves H5 if we consider inter-coder reliability, but H5 is rejected if we look at \(K\) for each algorithm.

\subsubsection{6. Conclusion}

This paper describes the work in progress upon a dataset from the Russian Twitter; we amplify earlier topic modelling results and objective topic quality measurement with a discussion on human interpretation of the modelling results. Our results contribute to the field of the topic modelling studies in three ways. We, first, show that term relevance tool works only partly the way that could help raise interpretability. Second, we see that automated quality assessment does not always correspond with how coders assess the topics, and designate a range of problematic issues in coding. Third, we show that negativity in topic descriptors might be related to interpretability of the topics for human coders, which is not captured by automated assessment of topic detection quality.

More generally, we can cautiously claim that the difficulties with Russian as a highly inflected language are not what cause the problems in topic detection and interpretation. Rather, it is the algorithmic imperfection, dataset preparation technique, and context dependence that altogether cause lower interpretability of short texts that, in addition, are noisy and ‘oral-written’ \cite{Lutovinova}.

We, at least partly, support what was earlier reported on the number of interpretable topics in Russian-language datasets but pose the question of what is more important for Twitter -- a higher number of interpretable (but identical) topics or a better-defined (but smaller in quantity) topics in the run, as this question seems crucial for Twitter studies.

Our research will continue in the areas designated above, including stretching our method to three other languages, experimenting with term relevance, and exploring the nature of impact of negativity upon topic interpretation.

\subsection{Detecting pivotal points in social conflicts via topic modeling of twitter content}\label{subsec:ch5/sec2/sub5}

\subsubsection{1. Theoretical Framework}

Twitter is believed to be actively used in protest activities for mobilization \cite{TheocharisLoweVanDeth,CarenGaby}. Researchers tried to study as many factors as possible to understand why the street protests start in one case and do not start others; why their dynamic may be different (see review in \cite{KoltsovaSelivanova}). Majority of researchers agree upon the fact that social media platforms, including Twitter, are actively used during social unrests or street actions as a main source of information \cite{AnduizaGallegoCantijoch,TufekciWilson}. Social media help activists organize street actions and mobilize potential supporters \cite{EarlKimport,FlesherFominaya}. But within all the corpus of the existing studies, our knowledge on the relations between street mobilization during social conflicts and digital technologies remains contradictory, as studies of causal relation between mobilization results and \textit{ad hoc} discussions on Twitter did not reach agreement \cite{Mercea}.

For instance, authors \cite{BastosMerceaCharperntier} have studied the role of Twitter and Facebook in mobilization of protest movements and have shown the causality between communication in social networks and protest activities. Their results are based on the Granger causality tests for three cases (the Indignados in Spain, the Occupy, and the Vinegar protests in Brazil) and prove that protest-related social media activity on Twitter and Facebook enables forecasting protest-related onsite activity. But the opposite results have been proven by another group of authors \cite{TheocharisLoweVanDeth}. They have shown a low level of interrelation between offline political activity and Twitter discussions for the same events.

We aim at suggesting a qualitative way to look at the linkages between discussion topicality and intensity, on one hand, and people’s street actions, on the other hand, for the cases where causality tests are either impossible to perform due to scarcity of data or provide mixed results. For this, we first show that the Granger causality test does not provide evidence for the online-offline linkages and then employ topic modeling, time series reconstructions of online and offline activities, and qualitative assessment of topics to formalize the pattern of online-offline linkages in time.

The paper is organized as follows. In Sect. 2, we describe the case under scrutiny and data collection. Section 3 shows the result of the Granger causality test as the premise for further research. Section 4 describes our method of topic modeling and the results of topic assessment against the timing of major political events within the case. Section 5 discusses and generalizes the results.

\subsubsection{2. Case Description}

New Year’s Eve sexual assaults in Germany were chosen for the case study due to several reasons. The event was reported massively on Twitter but did not reach the key quality news media in Germany. This fact triggered a significant discussion about accountability of German quality news media. Although local newspapers published the news already on January 1, national media have not paid attention to the event until January 5. Press service of Cologne police has published the official statement on the 2nd of January and this one-day delay has also provoked accusations of concealment against the elites. The absence of coverage of the event on the public service broadcasting and in quality newspapers was explained by the editorial houses first in terms of lack of human resources during the last days of the Christmas holidays; involvement of the rest of journalists into coverage of Munich terrorist attack; banality of criminal activities on the New Year’s Eve. The major German news agency DPA evaluated the importance of the event as very low and this decision could also influence the level of attention given by weekend editors to the event. While quality news media were blamed in a significant failure in detecting an important social problem, it remains understudied whether the causality between online and onsite protest activity can be detected in this case.

The data from Twitter was gathered during the period from 1.01.2016 to 31.01.2016 via a set of hashtags. The sample includes more than 99000 users (the core consists of 12000 users) and more than 64000 tweets. The specially developed web-crawler was used to download the data. We have described our data collection method extensively in our previous works \cite{BodrunovaBlekanovMaksimov,BodrunovaLitvinenkoBlekanov}, and here we will not stop on this.

\subsubsection{3. Granger Test}

The data about political activities in Germany during January 2016 was collected from the media reports and publications \cite{KoopmansRucht}. The level of the protest mobilization was estimated accordingly to the number of participants of a political activity (see Table~\cref{tab:actionsAndStatementsGermany}): not only anti-migrant activists were calculated, instead the feminist human rights movement activists and counter-PEGIDA unions were also included into the estimation, as suggested by \cite{BastosMerceaCharperntier} and \cite{Castells}. In case of contradictory data from mass media, the average between the data in media publications was calculated \cite{BastosMerceaCharperntier}.

\begin{table}[ht]%
	\centering
	\caption{Street actions and official statements of German authorities during January 2016.}%
	\label{tab:actionsAndStatementsGermany}% label всегда желательно идти после caption
	\begin{adjustbox}{width=0.9\textwidth}
		\small
		\begin{tabular}{ l  l }% Вертикальные полосы не используются принципиально, как и лишние горизонтальные (допускается по ГОСТ 2.105 пункт 4.4.5) % @{} позволяет прижиматься к краям
			\toprule
			Date & Political activity (estimated number of partaking people -- in brackets) \\
			\hline
			02.01.2016 & \makecell[l]{Press service of Cologne police has published the official statement mentioning\\ sexual assaults} \\
			04.01.2016 & \makecell[l]{Demonstrations of PEGIDA supporters in Dresden (400) and Leipzig (370)\\ No clashes with police, nobody has been arrested\\ The Head of Cologne’s police Wolfgang Albers, North Rhine-Westphalia’s\\ Interior Minister Ralf Jäger name the suspects as North African refugees} \\
			05.01.2016 & \makecell[l]{Demonstration of women in Cologne (400) \\ No clashes with police, nobody has been arrested \\ The mayor of Cologne Genriette Recker promises to provide more resources for\\ the police in the future but states that she still doesn’t have evidences about\\ refugees involved into the assaults} \\
			06.01.2016 & \makecell[l]{A public statement made by the press secretary of Angela Merkel about the\\ negotiations between the chancellor and the mayor of Cologne Genriette\\ Recker. Merkel expressed her ‘outrage’ with the violence and demanded a\\ ‘harsh response from the state’} \\
			07.01.2016 & \makecell[l]{Major quality newspapers -- \textit{FAZ}, \textit{Sueddeutsche Zeitung}, \textit{Tagesspiegel}, \textit{Die}\\ \textit{Welt} -- publish official responses about the reasons of silence of their outlets in \\first days of January} \\
			08.01.2016 & \makecell[l]{North Rhine-Westphalia’s Interior Minister Ralf Jäger dismissed the Head of\\ Cologne’s police Wolfgang Albers} \\
			09.01.2016 & \makecell[l]{Protests in Cologne: demonstration of women (1000), PEGIDA supporters\\ (1700), counter-protest demonstration ‘Cologne against the right-wing’ (1300)\\ Clashes with police; seizures of participants} \\
			11.01.2016 & \makecell[l]{Demonstrations in Leipzig: Anti-LEGIDA -- 2500, LEGIDA supporters -- 3000\\ Nobody was arrested\\North Rhine-Westphalia’s Interior Minister Ralf Jäger presents a report to an\\ extraordinary committee of the Internal Affairs Committee of North Rhine-\\Westphalia’s state parliament and confirms that majority of suspects are people\\ with migrant background \\ Right extremists attack group of refugees. Minister of Justice Heiko Maas\\ condemned the attack}\\
			12.01.2016 & \makecell[l]{Demonstration in Cologne: PEGIDA supporters (2000)\\ No clashes with police, nobody was arrested}\\
			13.01.2016 & \makecell[l]{Demonstration in Erfurt: PEGIDA supporters (2200)\\ No clashes with police, nobody was arrested} \\
			14.01.2016 & \makecell[l]{Minister President of North Rhine-Westphalia Hannelore Kraft presents the\\ package of actions to be taken for improving security and integration} \\
			18.01.2016 & \makecell[l]{Demonstrations in Dresden: PEGIDA supporters (4000), NOPEGIDA (500)\\ No clashes with police, nobody was arrested} \\
			20.01.2016 & \makecell[l]{Demonstrations in Dresden: PEGIDA supporters (700), NOPEGIDA (900)\\ No clashes with police, nobody was arrested} \\
			23.01.2016 & \makecell[l]{Demonstration in Stolberg PEGIDA supporters (1600)\\No clashes with police, nobody has been arrested} \\
			25.01.2016 & \makecell[l]{Demonstration of PEGIDA supporters in Dresden and in Cologne (4000) \\ No clashes with police, nobody was arrested} \\
			\bottomrule
		\end{tabular}%
	\end{adjustbox}
\end{table}

\paragraph{3.1 Casual Relations Between Online and Offline Activities in the Case Study} 
To test whether Twitter discussion on the New Year’s Eve’s sexual assaults in Germany is causally related with political activities in Germany during January 2016, we applied the Granger causality test to measure if the time lags of level of the protest mobilization (number of participants in the days of street actions in the dataset) relates to the distribution of the level of Twitter activity (number of tweets for each day of the studied period) over time. The data about daily users’ activity in Twitter and street protests is visualized as the time series graph (Fig.~\cref{fig:onlineDiscussionTimeSeries}).

\begin{figure}[ht]
	\centerfloat{
		\includegraphics[scale=1.0]{onlineDiscussionTimeSeries}
	}
	\caption{Time series for online discussion and offline events.}\label{fig:onlineDiscussionTimeSeries}
\end{figure}

Granger causality test can statistically determine when one variable precedes and predictively explains another variable \cite{GroshekCloughGroshek} in case when the influence of one 1-dimensional time series parameters \((xt)\) on another time series parameters \((yt)\). The Granger causality analysis are carried out by evaluating of two types of regression:
\[ 
yt = \alpha0 + \alpha1yt - 1 + \ldots + \alpha p yt - p + \beta1xt - 1 + \ldots + \beta1xt - p + \epsilon t
\] and
\[ 
xt = \alpha0 + \alpha1xt - 1 + \ldots + \alpha p xt - p + \beta1yt - 1 + \ldots + \beta1yt - p + \upsilon t
\]
to prove the null hypothesis \(\beta1 = \ldots = \beta p = 0\).

The regressions were carried out in a special program for statistical calculation EViews (Quantitative Micro Software, Version 10). The time lag is set at 2, the level of significance, \(\alpha\), is set at 5\%, We compared two time series -- data with number of participants and data about number of tweets. The null hypothesis about the absence of influence of online activities on the offline activities has been confirmed (Probability = 0,4696), as well as null hypothesis about absence of the opposite causal relation (Probability = 0,8315).

Hence, the idea of causal relationship between street demonstrations and Twitter communication about the event is rejected. Political activities in German cities triggered by New Year’s Eve’s sexual assaults in January 2016 and discussion about them on Twitter demonstrated an asynchronous character and do not depend on each other in the time dynamic.

\subsubsection{4. Detecting Pivotal Points via Topic Modeling}

\paragraph{4.1 The Method}
On the second step of our work with the case, we have tried an alternative methodology to detect patterns of interaction between online and offline activities. Our idea was to trace the change of topicality during the immediate aftermath of the conflict and to see whether the patterns repeat across several cases. In this paper, we look at just one case but we already see a particular pattern in interaction of topicality and protest activity.

To define the pattern, topic modeling was performed on the tweets to model the topic distribution for each day from the period of the study. The algorithm Biterm Topic Model has been chosen due to time restrictions and computational capabilities. BTM has been developed for short texts: this is an efficient algorithm that allows getting coherence topics for short text collections \cite{YanyanJiafengXueqi}.

The BTM algorithm assigns the topics to each document; then, the grouping of the messages and the corresponding topic distribution is performed, whereas the time of publication is used as a criterion for the grouping procedure. After several pre-tests with the number of topics from 20 to 200, 50 topics were chosen to perform the modeling procedure. The list of the topics with distribution coefficients \(p(z)\) for each day of the period has been visualized in a form of a heat map (see Fig.~\cref{fig:topicDistributionJan2016}).

\begin{figure}[ht]
	\centerfloat{
		\includegraphics[scale=1.0]{topicDistributionJan2016}
	}
	\caption{Distribution of topics during January 2016 (axis \(x\)) according to their coefficients \(p(z)\) (axis \(y\)).}\label{fig:topicDistributionJan2016}
\end{figure}

To analyze the semantic structure of the discussion, we formed the list of the most salient topics. If the weight of a topic (p) on a particular date exceeded 0,047 (more than 50\% in the daily topic distribution), it indicated that the actuality of this topic has increased sharply. 26 topics that were at particular points more salient than 0,047 are unevenly distributed in the course of January 2016, with most salient topics visible in the first three days of the conflict (five topics on January 1 and 2, six at January 3). In the discussion below, the topics are listed 1 to 50 (instead of 0 to 49 on the heat map).

\subsubsection{5. Results}

As visualized on Fig.~\cref{fig:mostSalientTopicsNYE}, four topics were salient on more than five days during the period of study. Their semantic structure is listed in the Table~\cref{tab:salientTopicsSemanticStructure}. Every topic includes nominations of place, actors, and hashtags employed by the users to frame the issue.

\begin{figure}[ht]
	\centerfloat{
		\includegraphics[scale=1.0]{mostSalientTopicsNYE}
	}
	\caption{The most salient topics in the Twitter discussion about New Year’s Eve’s sexual assaults (axis \(y =\) weight of a topic).}\label{fig:mostSalientTopicsNYE}
\end{figure}

\begin{table} [htbp]%
	\centering
	\caption{Semantic structure of the regularly salient topics.}%
	\label{tab:salientTopicsSemanticStructure}% label всегда желательно идти после caption
	\renewcommand{\arraystretch}{1.5}%% Увеличение расстояния между рядами, для улучшения восприятия.
	\begin{SingleSpace}
		\begin{tabulary}{\textwidth}{@{}>{\zz}L >{\zz}C >{\zz}C@{}} %Вертикальные полосы не используются принципиально, как и лишние горизонтальные (допускается по ГОСТ 2.105 пункт 4.4.5) % @{} позволяет прижиматься к краям
			\toprule     %%% верхняя линейка
			Topic N & Top words & Dates of saliency \\
			\midrule %%% тонкий разделитель. Отделяет названия столбцов. Обязателен по ГОСТ 2.105 пункт 4.4.5
			2 &  Ausnahmslos, gewalt, koelnhbf, rassismus, sexualisierte, frauen, sexismus, immer, sexuelle, aufschrei & January 1, 2, 4, 7, 22 \\
			33 &Koelnhbf, ausnahmslos, koeln, muslime, rapefugees, versuch, neue & January 1, 2, 4, 7, 18, 22, 23\\
			30 & Koelnhbf, koeln, polizei, medien, silvesternacht, amp, mueller & January 11 to 17 \\
			38 & Koelnhbf, ausnahmslos, merkel, amp, pegida, schon, mal, abmerkeln, brief & January 19, 20, 21, 23, 25, 27 \\
			\bottomrule %%% нижняя линейка
		\end{tabulary}%
	\end{SingleSpace}
\end{table}

Juxtaposition of the most salient topics clearly shows that there are four periods on the time-series graph when different topics dominated with a varying degree of sal- iency: January 1 to 7; January 8 to 15; January 16 to 23; January 24 to 31 (see Fig.~\cref{fig:topicDistributionPeriods}).

\begin{figure}[ht]
	\centerfloat{
		\includegraphics[scale=1.0]{topicDistributionPeriods}
	}
	\caption{Periods in topic distribution.}\label{fig:topicDistributionPeriods}
\end{figure}

In the first period (January 1 to 7, 2016), nominations among the top words define direct participants (topics N2 and N33): women (‘frauen’) against Muslims and ‘rapefugees’. Hashtags \#ausnahmslos and \#aufshrei indicate a feminist voice in the Twitter discussion: the hashtag \#aufschrei was introduced first in 2013 by feminism activist Anne Wizorek. This period was the most ‘buzzy’, according to the number of salient topics.

In the second period (January 8 to 15, 2016), topic N30 is the most salient. It demonstrates the public debate on the role of police and German media during the conflict: both actors were heavily criticized on Twitter. The second most salient topic during one day was the topic N20 that discusses more extensively political authorities who participated in the debate: Minister President of North Rhine-Westphalia Hannelore Kraft, North Rhine-Westphalia’s Interior Minister Ralf Jäger, and North Rhine-Westphalia’s parliamentarian from FDP Christian Lindner.

The third period (January 16 to 21, 2016), as the heat map shows, is a ‘gap’ period between topic N30 that is becoming less salient from its maximum on January 11, and topic N38 that increases in salience tremendously since 24\(^{th}\) of January. According to our time series analysis (see Table~\cref{tab:actionsAndStatementsGermany}), both PEGIDA and NOPEGIDA demonstrations have taken place but the number of tweets reduced significantly from more than 2500 on January 14 to less than 500 on January 23.

Topic N38 dominates the discussion in the last period (January 22 to 31, 2016). It reflects the activity of the PEGIDA movement and its supporters who use the hashtags \#abmerkeln and \#merkelmussweg -- both urge Angel Merkel to leave her Chancellor post.

We tried to trace manually whether the topics relate to the events. Concerning street actions, PEGIDA movement was mentioned only in 3 topics of 50 in our model (N21, N28, N38). Topic N21 goes beyond the saliency threshold only once, on January 2, while the first street action of PEGIDA supporters happened on January 4. Topic N28 passes the salience threshold only once, on January 31, after a wave of street actions. Topic N38 has been most salient before the street protests on January 18 as well as on the days when street actions took place (January 9, 20, 23, and 25). It reaches the highest level of salience (weight of a topic (p) exceeds 0,08) on January 24 and 25 in comparison to \(p = 0,048\) on January 23. In both cases, on January 18 and 25, the biggest number of participants took part in PEGIDA street actions (4000, as estimated by German media reports).

On January 14, Minister President of North Rhine-Westphalia Hannelore Kraft presented the package of actions to be taken for improving security and integration. Salient topics do not mention Kraft and do not relate directly neither to N11 (most salient, \(p = 0,089\), mentions the Head of Cologne’s police Albers fired on January 4 and North Rhine-Westphalia’s Interior Minister Ralf Jäger) nor to N25 (mentions political parties CDU/CSU, AFD, SPD and FDP), nor to N30 (mentions feminist hashtags \#aufschrei and \#ausnahmlos as well as popular hashtag related to the mayor Reker \#einearmlaenge). On the next day, the only topic N30 exceeded the level of \(p > 0,047\). In this case, we do not observe any connection between a political act of authorities and Twitter debate.

\subsubsection{6. Conclusions}

Scholars have tested a wide range of factors influencing political mobilization and, in particular, mass participation in a street action. Internet penetration, online user behavior and participation in online discussions were studied among other parameters. We argue that online behavior cannot be seen is a sole, dominant predictor for spillovers of protest to the street, as scholars that focused on interrelation between online and offline political activities have revealed controversial evidences on the causal relations between them.

In this paper, we presented the pilot study of relations between online and offline political activities triggered by New Year’s Eve’s sexual assaults in Germany 2016. The Granger causality test conducted in a way suggested by authors \cite{BastosMerceaCharperntier}, has failed, partly because the case study is characterized with the lack of parameters and data to compare the level of commitment among Internet users and street actions participants. In such cases, alternative ways to assess the relations between online and offline activities might be needed.

To prove the inter-relation between online discussion and offline political activities, new criteria and parameters are needed to be operationalized. To reveal these parameters for further research, we analyzed the data with the help of topic modelling that has visualized how the topics develop through the online debate.

Based on the results of the topic modeling, we juxtaposed the development of the most salient topics and the timing of the offline political activities in January 2016. Four periods were revealed, including the ‘gap’ period where topic saliency was the lowest. While in two first periods authorities and their (dis)activities were the main focus of the discussion, the fourth period is dominated by the topic that reflects far-right discourse (N38).

Several topics have gone beyond the saliency threshold (weight of a topic \((p)\) exceeding 0,047) before, on, and after the days of street actions; the most visible topics have also reached the highest level of saliency (weight of a topic \((p)\) exceeding 0,08) before the offline events. These intersections should be studied in a more nuanced way, \textit{i.a.} hour by hour. Hence, topic modeling has provided evidence that visibility of political activists precedes the street actions they organize. Moreover, the rapid change of the level of saliency that we observed for the topic N38 took place before the day the biggest number of participants took part in PEGIDA street actions in two cities simultaneously on January 25. If the peak of the topic in an online discussion precedes the mass participation in an offline street action, one might argue that the platform participated actively in forming of protest consensus. Rapid appearance of a peak within a topic, just like we observed for topic N38, might be a predictor for the forthcoming successful mobilization.

We assume that the Granger causality test might provide more accurate results in cases like New Year’s Eve’s sexual assaults in Germany 2016 if the level of street activity is examined against the level of salience of the dominant topics. Significant changes in the level of salience (rapid appearance of a peak with the topic distribution, as observed for topic N38) might be a predictor for the forthcoming successful mobilization. For more conclusions, we will in future compare several case studies of the same nature, to see whether the patterns that we have discovered for Cologne repeats for other countries and cases within Germany.

\section{Методы анализа тональности контента}\label{sec:ch5/sect3}

\subsection{Sentiment analysis for ad hoc discussions using multilingual knowledge-based approach}\label{subsec:ch5/sec3/sub1}

\subsubsection{1. Introduction}

Recently, the number of active users of social networks has increased many times, for example according to the website statista.com: Twitter has 313 million active users per month, 100 million users per day and 500 million tweets are published daily; Facebook has 1.79 billion active users per month, 1.18 billion per day. Such activity also influences people’s reaction to various social, political and economic events in the world. As a result, this means that such events are covered and distributed in social networks several times faster than in other media. Various discussions that arise as a result of any major events have an unpredictable structure of development and a number of specific properties, such as the peculiarity of user interaction among themselves, their relationship to the event, etc. Such discussions are called ad hoc discussions. For example, on Twitter and other social networks on major world events, there are long and branched discussions, the collection and analysis of which is a non-trivial task.

A separate task in the analysis of social networks is to determine the sentiment of the posts of users that were published in the context of any discussion, in particular, in the ad hoc discussion. Sentiment analysis \cite{Liu} is a class of methods that combines algorithms of computer linguistics, semantic and lexical analysis, as well as machine learning. One of the serious problems of analyzing the tonality of texts is the separation of a useful part of the text, that is, one that contains really useful information, and noise -- information that either does not affect the positivity or negativity of the statement, or complicates the precise definition. Another problem that the authors are also trying to solve is the analysis of short texts. So, for example, social network Twitter seriously differs from many other platforms with its limitations on the length of the post, which in 2018 is only 280 characters. This suggests that the messages in in most cases can contain only some essential information, or be spam and noise, detection and screening of which is another task. In addition, another distinctive feature of Twitter is the opportunity to conduct large and branched discussions on certain topics: whether it’s politics, economics, sports, etc., which can be easily monitored using hashtags.

In this paper, taking into account the specifics of short messages on Twitter, the authors set the task to develop a knowledge-based method for sentiment analysis of user messages in two major ad hoc discussions (Ferguson unrest (USA) and Biryuliovo bashings (Russia)). One of the problems that we are trying to solve is the lack of a decent russian lexicon. In fact, there is not one complete and accessible. There have been attempts to create such lexicons, but their effectiveness has not been verified to date. So we create a dictionary of russian affective words and compare its effectiveness with english lexicons.

\subsubsection{2. Current Research in Sentiment Analysis}

Researchers have studied almost all the main aspects of the problem. In their paper \cite{SchullerKnaup}, B. Schuller and T. Knaup propose an approach based on N-Grams and support vector machines \cite{ScholkopfCristianini}, which is very popular among researchers nowadays. In addition, the authors propose an approach that does not require pre-tagged data, and operates on the basis of knowledge-based resources, as well as linguistic analysis. They describe the use of online dictionaries containing both information about words and phrases, and the connections between them. Also, the article considers the possibility of using information about parts of speech of words to improve the accuracy of the analysis results. Also, K. Yessenov \cite{MisailovicYessenov} presented an analysis of the tone of comments to the reviews on films. The authors empirically studied the effectiveness of machine learning algorithms in the classification of text messages to their semantic value. The researchers used comments on the articles on the pop- ular Digg.com resource and investigated them for positivity and negativity. In \cite{RavindranSuchdevKotkar} the authors consider the possibility of combining knowledge-based methods and machine learning algorithms for analyzing users’ messages on the Twitter social network, which differ in their length. In the context of their study, they create a feature vector that contains such information as the number of positive and negative words, emoji and hashtags, and then classify data using machine learning algorithms (in particular, the support vector machine \cite{ScholkopfCristianini}). In addition, the issues of data preprocessing, use of metadata of each tweet and ways of its representation for further classification are considered.

Previous works also described methods for detecting influencers and ordinary users of ad hoc discussions. Often, users with fewer views and likes are more influential than others. This is shown in the articles \cite{BodrunovaLitvinenkoBlekanov2016,BodrunovaLitvinenkoBlekanov2017}, in which users and their connections were presented in the form of graphs and evaluated by various measures of centralities, such as pagerank and betweenness, as well as an expert analysis of the sentiment of users.

There are not many works devoted to the analysis of tonality in Russian. In \cite{BelyatskayaZagibalovCaroll} authors describe comparable corpora of book reviews in English and Russian and present experiments with a sentiment classification system. In \cite{ChetviorkinLoukachevitch} researchers consider a new approach for domain-specific sentiment lexicon extraction in Russian. They propose a set of statistical features and algorithm combination that can discriminate sentiment words in a specific domain. The extraction model is trained in the movie domain and then utilized to other domains.

\subsubsection{3. Proposed Sentiment Analysis Method}

The knowledge-based method developed by the authors includes the following steps:
\begin{enumerate}
	\item Construction of affective dictionaries.
	\item Data preprocessing.
	\item Classification.
\end{enumerate}

\paragraph{3.1 Construction of affective dictionaries} 
To analyze the sentiment of users’ messages in the study of major ad hoc discussions in the social network Twitter authors used a knowledge-based approach based on the construction and application of tagged dictionaries of positive and negative words. These dictionaries were constructed for English, French, Italian, Russian and Spanish languages according to the following principle. The most complete tagged list of English positive and negative words from open sources was selected. In our case, a dictionary that was constructed during the work of the University of Illinois at Chicago staff was used and is being expanded since then \cite{EnglishOpinionWords}. This dictionary contains approximately 6785 words. Initially there were 2005 positive words and 4780 negative words. Further, the procedure of extending the resulting dictionary was applied by adding synonyms from the WordNet database to it as follows. The authors developed a software component that matches each word from the already existing dictionary with a list of synonyms. As a result of the procedure, the dictionary consisted of 5787 and 11660 positive and negative words, respectively. This component checked the extended dictionary for repetitions, and, in case of occurrence of any identical words, deleted them. Using the machine translation procedure, new dictionaries were constructed in the appropriate natural languages. Since the results of the machine translation the repetitions could be formed again, the resulting dictionaries passed the procedure of removing duplicates. As a result of the performed actions, the dictionaries of Russian positive and negative words now consist of 3832 and 3890 words, respectively.

\paragraph{3.2 Preprocessing}
In order to improve the results and improve the accuracy of solving the classification problem, it is first necessary to bring the tweets to a normal view. This process is called normalization and, in addition to bringing all the symbols to the lowercase, includes following three steps.

\textit{3.2.1 Stop-words removal.} First of all, when processing the case, it is necessary to delete all noise words or stop-words. Stop-words are words that are often found in the language and carry almost no semantic meaning. For example, in English, such words can be prepositions (in, on, at, etc.), conjunction (and, but, for, etc.) and other function parts of speech. In addition, numbers, dates and times are also words that are removed from the tweet. In the context of processing tweets, we also remove all the mentions, that are, links to other Twitter users that start with the \@ symbol, as well as all possible links to third-party resources and sites that are not carrying any practically useful information and basically are noise. Since we do not have a tool capable of evaluating the significance of punctuation marks, we also expose them.

\textit{3.2.2 Tokenization.} The next step is to tokenize the text. This means we have to separate the text into significant units (tokens) for further processing. In the case of tweets, the most optimal size of token is one word. Therefore, after the separation, the tweets will be represented by a list of words that remained after the process of stop-words removal.

\textit{3.2.3 Stemming.} The third and final step is stemming \cite{SavoyDolamic}. Stemming is the process of extracting the basis of a word, which in this case may not coincide with the root of the word. In tasks of text classification, stemming is necessary first of all in order to reduce the dimensionality of the problem and increase the productivity of the classifier. There are many stemming algorithms for different languages. The most common is the stemmer of Martin Porter, developed originally only for the English language. But later the main idea of the algorithm was used to create the project "Snowball" and stemmers for different languages. Snowball is a small programming language designed for writing stemming algorithms. In the context of this study, we also subjected the dictionaries constructed in the previous paragraph to the stemming, since in the process of classifying it is necessary to define the polarity of a word, which is fully expressed in its basis.

\paragraph{3.3 Classification}
Our classification algorithm is fed to the input received after a preprocessed list of words that could affect the classification. Each word of this list is checked for the presence in dictionaries of positive and negative words constructed in the previous section. Ultimately, the number of positive and negative words in the list is counted. If the number of positive tokens exceeds zero and the number of negative tokens is zero, then the corresponding tweet is considered as positive, if the number of negative tokens exceeds zero, and the number of positive tokens is zero, then tweet is negative. In all other cases: if these numbers are equal or in some tweet there are both positive and negative tokens, then the tweet is marked as neutral. This approach shows different results on different test cases. This is our first approach.

We also offer a more obvious version of this classifier. In this version, the algorithm is also fed with a list of words obtained after preprocessing, and the number of occurrences of positive and negative words is also considered. However, those tweets whose word lists contain more positive than negative elements are considered positive, tweets that contain more negative than positive words are getting negative tag. Finally, tweet is marked as neutral if numbers of positive and negative words are equal. In this approach, you can see that in most cases tweets will be classified as positive or negative and only a small part of them will be neutral. Because of the fact that in this paper we consider the ad hoc discussion, we propose some adjustment of this algorithm: depending on the topic and features of the current discussions, it is possible to introduce a measure that would allow more precise definition of neutrality of some tweet. In other words, the number of tweets that would be labeled as neutral in the original version of the algorithm would be greater. The proposed metric \(m\) is introduced as follows:
\begin{equation}
	\label{eqn:31}
	m = \frac{p - n}{N}
\end{equation}
where \(p\) is number of positive words in tweet, \(n\) -- of negative words and \(N\) is total number of words in tweet. After the introduction of this metric, you can add to the algorithm some coefficient (let’s say \(c\)), which would be the threshold value for neutral tweets. In other words, if the number of positive and negative words differs slightly, then the tweet is considered neutral. This approach allows you to customize the work of the algorithm, depending on the characteristics of a certain ad hoc debate.

\subsubsection{4. Experiment}
An experiment was conducted to assess the quality of the developed method in three modifications. There were two qualitative measures used: accuracy score and f1 score \cite{Powers}. As a test data there were chosen two real large ad hoc discussions \cite{SmoliarovaBlekanovBodrunova} in Russian and English.

\paragraph{4.1 Test cases description and tagging}
In this paper authors have the opportunity to analyze the data of ad hoc discussions in the social network Twitter of two well-known events: Ferguson unrest (USA) in August 2014 and Biryuliovo bashings (Russia) in October 2013. For each of the events above, there was compiled a list of popular hashtags that describe it:
\begin{enumerate}
	\item Biryuliovo: \#бирюлево, \#мигранты, \#овощебаза, \#зейналов.
	\item Ferguson: \#ferguson.
\end{enumerate}
As a result of the collecting and processing the set of hashtags, we got the following data:
\begin{enumerate}
	\item Biryuliovo: Number of users who have published their tweets in the period since 1st of October 2013 till Thirty-first of October 2013 is 3574. Number of published posts is 10215.
	\item Ferguson: Number of users who have published their tweets in the period since twenty-second of August 2014 till Thirty- first of August 2014 is 70018. Number of published posts is 193812.
\end{enumerate}
On the Figures~\cref{fig:fergusonPublicationActivity} and~\cref{fig:biryulievoPublicationActivity}, one can observe the publication activity of users of discussions about Ferguson and Biryulyovo, respectively. It can be seen that the ad hoc discussions have a different structure: only one "splash" of activity is visible in the discussion of Biryulyovo and the discussion of the situation in Ferguson that has not ceased throughout the week.

\begin{figure}[ht]
	\centerfloat{
		\includegraphics[scale=0.5]{fergusonPublicationActivity}
	}
	\caption{Periods in topic distribution.}\label{fig:fergusonPublicationActivity}
\end{figure}

\begin{figure}[ht]
	\centerfloat{
		\includegraphics[scale=0.5]{biryulievoPublicationActivity}
	}
	\caption{Periods in topic distribution.}\label{fig:biryulievoPublicationActivity}
\end{figure}

To calculate the qualitaty measures, the authors performed the tagging of the above-described data sets as follows:

\begin{enumerate}
	\item For each ad hoc discussion, lists of top-100 users were compiled according to the number of published messages, as far as betweenness centrality and as far as pagerank centrality \cite{BodrunovaLitvinenkoBlekanov}.
	\item Aggregation of three top-100 lists and deletion of duplicate users. As a result, 213 users were obtained in the Biryuliovo case and 200 in the Ferguson case.
	\item By the aggregated list of users in the Biryuliovo case there were selected all user messages, the total number of which was 3885. In the Ferguson case there were collected just fifteen random messages by each user in order to keep size of the collection slightly equal to the size of Biryuliovo collection.
\end{enumerate}

\paragraph{4.2 Results}
In the course of carrying out the work described above, the problem of the analysis of the ad hoc discussions in the social network Twitter was considered. The authors presented a multi-lingual approach, which is based on the construction and use of affective dictionaries for Russian language and compare it efficiency to expanded existing English lexicon. To assess the quality of the three modifications of this approach, measures accuracy score and f1 score were used on manually annotated data. These results are presented in the Table~\cref{tab:fergusonAccuracyF1}. Results for Biryuliovo case are represented in the Table~\cref{tab:biryliovoAccuracyF1}. These tables show that our best approach gives 0.7 f1 score which is the average result in English \cite{KingHopkins,RavichandranRao}. The results for the Russian dictionaries are slightly lower than the average English. We assume that this is due to the imperfection of machine translation and the difference in the structure of Russian and English.

\begin{table} [htbp]%
	\centering
	\caption{Accuracy and F1 score for Ferguson.}%
	\label{tab:fergusonAccuracyF1}% label всегда желательно идти после caption
	\renewcommand{\arraystretch}{1.5}%% Увеличение расстояния между рядами, для улучшения восприятия.
	\begin{SingleSpace}
		\begin{tabulary}{\textwidth}{@{}>{\zz}L >{\zz}C >{\zz}C >{\zz}C@{}} %Вертикальные полосы не используются принципиально, как и лишние горизонтальные (допускается по ГОСТ 2.105 пункт 4.4.5) % @{} позволяет прижиматься к краям
			\toprule     %%% верхняя линейка
			& First version & Second version & With \(m\) metricy \\
			\midrule %%% тонкий разделитель. Отделяет названия столбцов. Обязателен по ГОСТ 2.105 пункт 4.4.5
			accuracy score  & 0.65 & 0.49 & 0.62\\
			f1 score & 0.70 & 0.52 & 0.66 \\
			\bottomrule %%% нижняя линейка
		\end{tabulary}%
	\end{SingleSpace}
\end{table}

\begin{table} [htbp]%
	\centering
	\caption{Accuracy and F1 score for Biryliovo.}%
	\label{tab:biryliovoAccuracyF1}% label всегда желательно идти после caption
	\renewcommand{\arraystretch}{1.5}%% Увеличение расстояния между рядами, для улучшения восприятия.
	\begin{SingleSpace}
		\begin{tabulary}{\textwidth}{@{}>{\zz}L >{\zz}C >{\zz}C >{\zz}C@{}} %Вертикальные полосы не используются принципиально, как и лишние горизонтальные (допускается по ГОСТ 2.105 пункт 4.4.5) % @{} позволяет прижиматься к краям
			\toprule     %%% верхняя линейка
			& First version & Second version & With \(m\) metricy \\
			\midrule %%% тонкий разделитель. Отделяет названия столбцов. Обязателен по ГОСТ 2.105 пункт 4.4.5
			accuracy score  & 0.57 & 0.41 & 0.55\\
			f1 score & 0.65 & 0.46 & 0.6\\
			\bottomrule %%% нижняя линейка
		\end{tabulary}%
	\end{SingleSpace}
\end{table}

\subsubsection{5. Conclusion}

In this article, authors proposed a multi-lingual approach to the analysis of ad hoc discussions and, in particular, the sentiment analysis of messages in such discussions. The presented algorithms showed quite good results, which were represented in the relevant tables, taking into account the specifics of the given discussions. In the future work, the authors plan to develop this approach and compare it with the work of such classification algorithms as the support vector machines.

\subsection{When emotions grow: cross-cultural differences in the role of emotions in the dynamics of conflictual discussions on social media}\label{subsec:ch5/sec3/sub2}

\subsubsection{1. Introduction}

The spread of affective content on social media, as well as user grouping based on affect \cite{Papacharissi}, has been a focus of scholarly attention due to multiple reasons. Emotional content is seen mostly as a threat to the democratic quality of public discussions \cite{Cortese}, as it undermines its integrity and rationality \cite{BadjatiyaGuptaGupta}, as well as its civility and constructive character \cite{BurnapWilliams}, especially when one deals with hate speech and toxic expressions \cite{ParkFung,GeorgakopoulosTasoulisVrahatis}.

Despite the evident necessity to learn how emotions work within the discussion, we lack evidence on what roles various particular emotions play in the dynamics of discussions on social media. Detection of particular emotions has been mostly linked to political emotionality and political functions of emotions \cite{Lyman,TicentoCloughHalley,HeaneyFlam,WahlJorgensen}, as well as their linkage to other features of content like fake news \cite{VosoughiRoyAral}.

The substantially negative role of hate speech is practically non-questioned today; but several works still draw attention to controversial relations between freedom of speech and hate speech \cite{Dorsett,Cammaerts} and to the possible use of offensive language in positive sense in discriminated communities like LGBTQ \cite{WaseemDavidsonWarmsley}. Another potentially positive role of emotional content (negative just as well as positive) might lie in spurring the discussions, providing them flame and substance and not letting them ebb.

In this paper, we explore just two aspects of this issue, namely -- the growth of emotional load in conflictual discussions and the relations between the number of overall users in the discussion and the angry and compassionate ones.

For this, we use the data from three highly conflictual and emotionally polar discussions, more or less territorially localized, despite potential global participation on Twitter. Use of three discussions from different communicative cultures will allow us generalize on the nature of the relations between the overall growth of the discussion and its highly emotional content.

To explore the discussion bulk, we collected the data via Twitter crawling, then coded the emotions with the help of native speakers, then taught the machine to recognize the emotions, and then detected the users who used emotional speech and reconstructed the hourly graphs of the discussions, as well as time series. We have also performed Granger test to spot the relations between the number of emotional and neutral users in the discussions.

The remainder of the paper is organized as follows. Section 2 discusses our idea of the emotions as discussion drivers. Section 3 poses the hypotheses, and Sect. 4 describes the data collection procedures, sampling, and the method of detecting emotions. Section 5 provides and discusses the results.

\subsubsection{2. Emotions Aggregated: What Can They Tell?}

Emotional contagion theory \cite{HatfieldBensmanThornton} adapted for social media suggests that diffusion of emotions happens on individual level, via direct one-time contact with emotionalized content \cite{CovielloSohnKramer}. Individual emotions, though, play a key role in group behavior online, as they form the so-called ad hoc discussions \cite{BrunsBurgess} of affective publics \cite{Papacharissi} and, thus, matter more on the aggregate level and might work differently from what is expected, especially in fast-growing conflictual discourse online. Other theories, like theories of social influence or social learning \cite{Young}, suggest multiple, hierarchical, and/or topically-restricted contacts of users with emotional content; whether these limitations are also true for how the emotions grow and spread collectively, still remains unclear.

The idea of affective agenda \cite{ColemanWu} implies that the dynamics of an emotional discussion might be dependent on mood changes or spillovers from group to group, language to language, or topic to topic. Also, the relations between various emotions might be a com- petitive one, since they work as the polarizing agents within the discussion: anger may overcome fear, and denial might beat compassion. A shockingly strong example of this is, of course, the case of \textit{Charlie Hebdo} which polarized the discussion participants to the extent of creating formally opposing hashtags (‘jesuischarlie’ and ‘jenesuispascharlie’, representing compassion and its rejection, anger, and discontent).

The change of affective agenda or mood within the discussion might lead to its growth (or even outburst) or ebbing, and this question remains under-explored. We suggest that emotions of different stance (positive/negative) may spur/slow down the discussions in various ways.

\subsubsection{3. The Research Hypotheses}

In this paper, we have posed three hypotheses:

\begin{itemize}
	\item H1. In all the three conflictual discussions under our scrutiny, the dynamics of emotions will go ahead of the general discussion dynamics: the number of emotional users will grow and diminish faster than the number of emotional tweets.
	\item H2. In all the three discussions, the dynamics of the number of involved users will depend on the dynamics of the number of emotional users.
	\item H3. The dynamics of compassion and anger will differ in all the three discussions and will not correlate within the cases.
\end{itemize}

We will also qualitatively describe the picture we see on the web graphs and on time series graphs, to contextualize and reflect more on the received data. By analyzing the co-dynamics of the overall discussions and these two emotions we can conclude whether the pattern of the spread of emotions and its link with the discussion dynamics is the same in various language segments of Twitter.

\subsubsection{4. Data Collection and Methods of Research}

\paragraph{4.1 The Cases Under Scrutiny}

We analyze the spread of two polar emotions -- anger and compassion -- in three Twitter discussions on inter-ethnic conflicts, namely:

\begin{itemize}
	\item The Ferguson unrest, USA, 2014: number of users who published tweets -- 70018, total number of messages -- 193812;
	\item The Charlie Hebdo shooting, France, 2015: number of users who published tweets -- 238491, total number of messages -- 505069;
	\item The mass harassment of women by the re-settlers from the Arab world in Cologne, Germany, in the New Year night of 2015-2016; number of users who published tweets -- 99981, total number of messages -- 215253.
\end{itemize}

The overall number of tweets collected based on the relevant hashtags and keywords is actually much bigger. The data we use were collected by our Twitter crawler in the aftermath of the conflicts and include altogether over 2,5 M tweets. The samples described above resulted from the fact that we selected the mono-language parts of the discussions only -- German for Cologne, French for Charlie Hebdo, and English for Ferguson. We have also cleaned spam, non-relevant tweets, and hashtag-only tweets.

\paragraph{4.2 Emotion Detection}

To be able to detect emotions, we have used an enhanced SVM-based approach with the use of Numpy and scikit-learn libraries. As we have been working with multi-lingual sentiment analysis for several years, we have modified the approach with two features. First, the experiments have shown that elimination of stop words from tweets actually diminishes the quality of the classifier. Thus, the stop words have been returned to the initial dataset, which has raised the quality by 4-5\%. Second, we have experimentally found the streaming hyper-parameters and adjusted them: the regularization parameter was 0.9, and we have also used the linear nuclear function and the OVR function. We have also spotted that, in several cases, the machine tended to over-learn; thus, we have balanced the samples to reach better results.

Tweets from real-world discussions, despite we have used the well-preprocessed datasets, remain one of the ‘noisiest’ types of textual data; this has cast some negative impact upon the quality metrics of emotion detection. This is why we have used quite large collections for hand coding. We have selected the tweets for hand coding by the procedure described in our earlier works \cite{BodrunovaLitvinenkoBlekanov2017,BodrunovaBlekanovSmoliarova}, by selecting influential users by nine parameters -- four absolute ones (the number of posts, retweets, likes, and comments) and five graph-dependent ones (in-degree, out-degree, degree, betweenness and pagerank centralities). The top users’ tweets constitute the samples for coding. Altogether, we have coded: for Germany -- 13,620 tweets, for the USA -- 7200 tweets, and for France -- 7433 tweets.

We have received the following best levels for the quality metrics (see Table~\cref{tab:threeCasesQualityMetrics}).

\begin{table}[ht]%
	\centering
	\caption{Quality metrics for emotion detection for the three cases.}%
	\label{tab:threeCasesQualityMetrics}% label всегда желательно идти после caption
	\begin{adjustbox}{width=1\textwidth}
		\small
		\begin{tabular}{ c  c  c  c  c  c  c  }% Вертикальные полосы не используются принципиально, как и лишние горизонтальные (допускается по ГОСТ 2.105 пункт 4.4.5) % @{} позволяет прижиматься к краям
			\toprule
			& \multicolumn{2}{c}{\makecell{Germany}} & \multicolumn{2}{c}{\makecell{The USA}} & \multicolumn{2}{c}{\makecell{France}} \\
			& Accuracy & F-measure & Accuracy & F-measure & Accuracy & F-measure \\
			\hline
			Anger & 0.7 & 0.7 & 0.77 & 0.77 & 0.78 & 0.77 \\
			Compassion & 0.75 & 0.74 & 0.8 & 0.8 & 0.84 & 0.84 \\
			\bottomrule
		\end{tabular}%
	\end{adjustbox}
\end{table}

\paragraph{4.3 Data Analysis and Visualization}

We have marked the tweet datasets by emotion; then we have aggregated the tweets per user and have calculated hourly numbers of neutral, angry, and compassionate users, according on whether they have expressed these emotions within a given hour. Thus, we have received he hourly numbers of users who were emotional, either compassionate or angry. Of course, in a small number of cases, a user could express both compassion and anger within one hour; but this mattered only within the first couple of hours, while in the course of the discussion the number of such users became negligibly small.

To visualize the results of the emotion detection, we have constructed the hourly web graphs for three days using the YifanHu algorithm from the Gephi library, as well as the time series graphs with the hourly lag. For Cologne and \textit{Charie Hebdo}, we have taken the first three days of after the trigger events; for Ferguson, we did not have such data and, thus, the three days were taken from the middle of the discussion.

To assess the inter-emotional relations (in terms of the number of users), we have used the Spearmans’ Rho.

To test the hypotheses, we have performed the Granger test with the hourly data and have also qualitatively assessed the graphs \cite{Wessa}. The results are presented below.

\subsubsection{5. Results and Discussion}

\textit{H1-H2.} When assessing the cases, we have seen that the basic pattern of the discussions is different from the expected: in general, the lines indicating the number of neutral users are generally followed by those of anger and compassion (see Fig.~\cref{fig:timeSeriesGraphs}). But, if we look at details and the results of the Granger test, the three start to diverge.

\begin{figure}[ht]
	\centerfloat{
		\includegraphics[scale=0.5]{timeSeriesGraphs}
	}
	\caption{a.Time series graphs, the Cologne case. b. Time series graphs, the Ferguson case. c. Time series graphs, the Charlie Hebdo case. (Color figure online)}\label{fig:timeSeriesGraphs}
\end{figure}

First of all, the three cases differ in whether the growth of the number of angry or compassionate users affects the rest of the speech in the case (see Table~\cref{tab:grangerTestResults}). In some sense, Cologne and Ferguson ‘mirror’ each other in terms of when exactly the emotional users spur the number of the neutral users. We can see from the Table~\cref{tab:grangerTestResults} that it is compassion, not anger, that plays a bigger role in predicting the number of neutral users -- that is, the very growth or ebbing of the discussions. As we see from the time series graphs, the regression shows the dependency when the discussion starts to grow. This happens with the first two days of Cologne where the journalists were reluctant to cover the mass harassment story as they were not sure how to report on the nationality of the harassers. It also happens during an outburst of emotions in the Ferguson case, on Day 3. Thus, the role of emotional content in the early-stage discussions needs to be well researched upon.

\begin{table}[ht]%
	\centering
	\caption{Quality metrics for emotion detection for the three cases.}%
	\label{tab:grangerTestResults}% label всегда желательно идти после caption
	\begin{adjustbox}{width=1\textwidth}
		\small
		\begin{tabular}{ c  c  c  c }% Вертикальные полосы не используются принципиально, как и лишние горизонтальные (допускается по ГОСТ 2.105 пункт 4.4.5) % @{} позволяет прижиматься к краям
			\toprule
			& Days 1-3 & Days 2-3 & Day 3 \\
			\hline
			Cologne anger & No No & No No & \colorbox{yellow}{11,17** 13,3**}\\
			Cologne compassion & \colorbox{green}{13,1*** No} & \colorbox{green}{8,71** No} & \colorbox{yellow}{5* 9,8**} \\
			Ferguson anger & \colorbox{violet}{No 4,13*} &	No No  & 4,76* No\\
			Ferguson compassion & \colorbox{yellow}{5,53* 8,1**} & \colorbox{yellow}{4,68* 6,52*} & \colorbox{green}{5,89* No} \\
			Charlie Hebdo anger & No No & No No & No No \\
			Chare Hebdo compassion & No No & \colorbox{violet}{No 10,2**} & \colorbox{violet}{No 5,65*} \\
			\hline
			\multicolumn{4}{l}{\makecell{Note. Green: the test shows the dependency of the overall discussion upon these users and \\authors. Yellow: the test shows that the numbers of the neutral users and emotional users \\correlate due to a third factor.\\Violet: The linkage between emotional and neutral users shows the dependency of the\\ other direction: the higher the number of neutral tweets, the higher the number of the\\ angry and/or compassionate tweets. That is, when the discussion is boiling itself, \\the number of angry users also grows without any coordinated effort.}} \\
			\bottomrule
		\end{tabular}%
	\end{adjustbox}
\end{table}

In the American case, there is also a smaller trend covered by a more general one. There are several moments in the course of the discussion when anger starts to diminish earlier than the number of the neutral tweets, and the blue line follows.

Another important conclusion is that the emotions do not grow/diminish within the discussions but have intense hourly dynamics. The ‘jumps’ of the number and connectivity of users who express anger and compassion (see Fig.~\cref{fig:fastDynamcs}, the example of Cologne) need to be further explored within the aforementioned theories of emotional contagion and/or social learning.

\begin{figure}[ht]
	\centerfloat{
		\includegraphics[scale=0.5]{fastDynamcs}
	}
	\caption{a. Fast dynamics of growth of the number of angry users within 3h (a ‘jump’ of anger), Cologne, January 4, 2016. b. Fast dynamics of growth of the number of compassionate users within 3h (a ‘jump’ of compassion), Cologne, January 4, 2016.}\label{fig:fastDynamcs}
\end{figure} 

Thus, H1 and H2, in how they were formulated, needs to be rejected: the number of the neutral users does not grow depending of the spread of emotions between emotional users. But what is also important is that the dynamics of the emotions is intensive even within the hourly perspective and that the moments when the discussion grows might depend on compassion, rather than anger. The emotions matter in the very beginning of the discussions or in the moments of discussion outbursts.

\textit{H3.} Ourthirdhypothesiswascreatedtoberejected,but,again,thisisnotwhathappens in the datasets. Thus, the correlations between the two hourly spreads of people are: for Cologne -- 0.950****, for Ferguson -- 0,988***, for Charlie Hebdo -- 0,937**. This adds to the evidence that emotional users do not lead the way and do not cause excessive discussion around them but appear within the discussion the same way as neutral users. But, when on an emotional peak, such users might create discussion outbursts.

We have tested the dependence of the discussion dynamics on the number of emotional (angry and compassionate) users of Twitter. Within the three cases, there were no similar patterns of such dependence, which might be explained by socio-cultural differences or the structure of the sample. We have also detected a special role of emotional tweets in the beginnings of the discussions and the discussion outbursts; we suggest that, in future, not users but the tweets themselves need to become the unit of analysis.

\section{Методы суммаризации}\label{sec:ch5/sect4}

\FloatBarrier

           % Глава 5
\chapter{Цифровой двойник сетевых физических процессов}\label{ch:ch6}

\section{Цепь поставок}\label{sec:ch6/sect1}

\section{Визуализация}\label{sec:ch6/sect2}

\section{Геопозиция}\label{sec:ch6/sect3}

\section{Сеть маршрутов распространения физических ресурсов на основе данных (обработанных подсетью виртуальных сетевых агентов)}\label{sec:ch6/sect3}

\FloatBarrier

           % Глава 6
\chapter{Программный комплекс}\label{ch:ch7}

%\section{Конструирование Web-ресурсов}\label{sec:ch4/sect1}

%\section{Типы виртуальных сетевых агентов}\label{sec:ch4/sect2}

%\section{Оптимальное конструирование сборщика (цепочка роботов)}\label{sec:ch4/sect3}

%\section{Реализация программного комплекса}\label{sec:ch4/sect3}

\clearpage
           % Глава 7
\include{Dissertation/conclusion}      % Заключение
%\nomenclature{DirLouv}{Directed Louvain}
\nomenclature{GKM}{Generalized K-means}
\nomenclature{LeidMod}{Leiden modularity}
\nomenclature{LeidCPM}{Leiden CPM}
\nomenclature{GANXiSd}{SLPA disjoint}
\nomenclature{GANXiSo}{SLPA overlapping}
\nomenclature{DirMod}{Directed modularity}
\nomenclature{ClusCoe}{Clustering coefficient}
\nomenclature{NEDind}{NEDindex}
\nomenclature{Conduc}{Conduction}
\nomenclature{Contrac}{Contraction}
\nomenclature{Expans}{Expansion}
\nomenclature{Bir}{the Biryulevo case}
\nomenclature{Col}{the Cologne case}
\nomenclature{Fer}{the Ferguson case}
\nomenclature{ChE}{the Charlie Hebdo case}

\printnomenclature[3.5cm] % Значение ширины столбца с обозначениями стоит подбирать вручную        % Список сокращений и условных обозначений
%\include{Dissertation/dictionary}      % Словарь терминов
\include{Dissertation/references}      % Список литературы
\include{Dissertation/lists}           % Списки таблиц и изображений (иллюстративный материал)

\setcounter{totalchapter}{\value{chapter}} % Подсчёт количества глав

%%% Настройки для приложений
%\appendix
% Оформление заголовков приложений ближе к ГОСТ:
%\setlength{\midchapskip}{20pt}
%\renewcommand*{\afterchapternum}{\par\nobreak\vskip \midchapskip}
%\renewcommand\thechapter{\Asbuk{chapter}} % Чтобы приложения русскими буквами нумеровались

%nclude{Dissertation/appendix}        % Приложения

%\setcounter{totalappendix}{\value{chapter}} % Подсчёт количества приложений

\end{document}
